% -------------------------------------------------------------
% NOTE ON THE DETAILED AND SHORT VERSIONS:
% -------------------------------------------------------------
% This paper comes in two versions, a detailed and a short one.
% The short version should be more than sufficient for any
% reasonable use; the detailed one was written purely to
% convince the author of its correctness.
% To switch between the two versions, find the line containing
% "\newenvironment{noncompile}{}{}" in this LaTeX file.
% Look at the two lines right beneath this line.
% To compile the detailed version, they should be as follows:
%   \includecomment{verlong}
%   \excludecomment{vershort}
% To compile the short version, they should be as follows:
%   \excludecomment{verlong}
%   \includecomment{vershort}
% As a rule, the line
%   \excludecomment{noncompile}
% should stay as it is.
% -------------------------------------------------------------
% NOTES ON SOME HACKS USED IN THIS FILE:
% -------------------------------------------------------------
% One of my pet peeves with amsthm is its use of italics in the theorem and
% proposition environments; this makes math and text indistinguishable in said
% enviroments. To avoid this, I redefine the enviroments to use the standard
% font and to use a hanging indent, along with a bold vertical bar to its
% left, to distinguish these environments from surrounding text. (Along with
% the advantage of distinguishing math from text, this also allows nesting
% several such environments inside each other, like a definition inside a
% remark. I'm not sure how good of an idea this is, though. There are also
% downsides related to the hanging indentation, such as footnotes out of it
% being painful to do right.) This is done starting from the line
%   \theoremstyle{definition}
% and until the line
%   {\end{leftbar}\end{exmp}}

\documentclass[numbers=enddot,12pt,final,onecolumn,notitlepage]{scrartcl}%
%\usepackage[headsepline,footsepline,manualmark]{scrlayer-scrpage}
\usepackage[all,cmtip]{xy}
\usepackage{amsfonts}
\usepackage{amssymb}
\usepackage{framed}
\usepackage{amsmath}
\usepackage{comment}
\usepackage{needspace}
\usepackage{color}
\usepackage[breaklinks]{hyperref}
\usepackage[sc]{mathpazo}
\usepackage[T1]{fontenc}
\usepackage{amsthm}
%TCIDATA{OutputFilter=latex2.dll}
%TCIDATA{Version=5.50.0.2960}
%TCIDATA{LastRevised=Monday, April 11, 2016 10:09:14}
%TCIDATA{SuppressPackageManagement}
%TCIDATA{<META NAME="GraphicsSave" CONTENT="32">}
%TCIDATA{<META NAME="SaveForMode" CONTENT="1">}
%TCIDATA{BibliographyScheme=Manual}
%BeginMSIPreambleData
\providecommand{\U}[1]{\protect\rule{.1in}{.1in}}
%EndMSIPreambleData
\theoremstyle{definition}
\newtheorem{theo}{Theorem}[section]
\newenvironment{theorem}[1][]
{\begin{theo}[#1]\begin{leftbar}}
{\end{leftbar}\end{theo}}
\newtheorem{lem}[theo]{Lemma}
\newenvironment{lemma}[1][]
{\begin{lem}[#1]\begin{leftbar}}
{\end{leftbar}\end{lem}}
\newtheorem{prop}[theo]{Proposition}
\newenvironment{proposition}[1][]
{\begin{prop}[#1]\begin{leftbar}}
{\end{leftbar}\end{prop}}
\newtheorem{defi}[theo]{Definition}
\newenvironment{definition}[1][]
{\begin{defi}[#1]\begin{leftbar}}
{\end{leftbar}\end{defi}}
\newtheorem{remk}[theo]{Remark}
\newenvironment{remark}[1][]
{\begin{remk}[#1]\begin{leftbar}}
{\end{leftbar}\end{remk}}
\newtheorem{coro}[theo]{Corollary}
\newenvironment{corollary}[1][]
{\begin{coro}[#1]\begin{leftbar}}
{\end{leftbar}\end{coro}}
\newtheorem{conv}[theo]{Convention}
\newenvironment{condition}[1][]
{\begin{conv}[#1]\begin{leftbar}}
{\end{leftbar}\end{conv}}
\newtheorem{quest}[theo]{TODO}
\newenvironment{todo}[1][]
{\begin{quest}[#1]\begin{leftbar}}
{\end{leftbar}\end{quest}}
\newtheorem{warn}[theo]{Warning}
\newenvironment{conclusion}[1][]
{\begin{warn}[#1]\begin{leftbar}}
{\end{leftbar}\end{warn}}
\newtheorem{conj}[theo]{Conjecture}
\newenvironment{conjecture}[1][]
{\begin{conj}[#1]\begin{leftbar}}
{\end{leftbar}\end{conj}}
\newtheorem{exmp}[theo]{Example}
\newenvironment{example}[1][]
{\begin{exmp}[#1]\begin{leftbar}}
{\end{leftbar}\end{exmp}}
\newenvironment{statement}{\begin{quote}}{\end{quote}}
\iffalse
\newenvironment{proof}[1][Proof]{\noindent\textbf{#1.} }{\ \rule{0.5em}{0.5em}}
\fi
\newenvironment{verlong}{}{}
\newenvironment{vershort}{}{}
\newenvironment{noncompile}{}{}
\newenvironment{obsolete}{}{}
\excludecomment{verlong}
\includecomment{vershort}
\excludecomment{noncompile}
\excludecomment{obsolete}
\newcommand\arxiv[1]{\href{http://www.arxiv.org/abs/#1}{\texttt{arXiv:#1}}}
\let\sumnonlimits\sum
\let\prodnonlimits\prod
\let\bigcapnonlimits\bigcap
\renewcommand{\sum}{\sumnonlimits\limits}
\renewcommand{\prod}{\prodnonlimits\limits}
\renewcommand{\bigcap}{\bigcapnonlimits\limits}
\setlength\textheight{22.5cm}
\setlength\textwidth{15cm}
%\ihead{Note on NBC sets and chromatic polynomial}
%\ohead{\today}
\begin{document}

\title{A note on non-broken-circuit sets and the chromatic polynomial}
\author{Darij Grinberg}
\date{version 1.1, \today}
\maketitle
\tableofcontents

\section*{***}

The purpose of this note is to demonstrate several generalizations of
Whitney's theorem \cite{BlaSag86} -- a classical formula for the chromatic
polynomial of a graph. The directions in which we generalize this formula are
the following:

\begin{itemize}
\item Instead of summing over the sets which contain no broken circuits, we
can sum over the sets which are \textquotedblleft$\mathfrak{K}$%
-free\textquotedblright\ (i.e., contain no element of $\mathfrak{K}$ as a
subset), where $\mathfrak{K}$ is some fixed set of broken circuits (in
particular, $\mathfrak{K}$ can be $\varnothing$, yielding another well-known
formula for the chromatic polynomial).

\item Even more generally, instead of summing over $\mathfrak{K}$-free
subsets, we can make a weighted sum over all subsets, where the weight depends
on the broken circuits contained in the subset.

\item Analogous (and more general) results hold for chromatic symmetric functions.

\item Analogous (and more general) results hold for matroids instead of graphs.
\end{itemize}

Note that, to my knowledge, the last two generalizations cannot be combined:
Unlike graphs, matroids do not seem to have a well-defined notion of a
chromatic symmetric function.

We shall explore these generalizations in the note below. We shall also use
them to prove an apparently new formula for the chromatic polynomial of a
graph obtained from a transitive digraph by forgetting the orientations of the
edges (Proposition \ref{prop.digraph.2pf-chrom}). This latter formula was
suggested to me as a conjecture by Alexander Postnikov, during a discussion on
hyperplane arrangements on a space with a bilinear form; it is this formula
which gave rise to this whole note. The subject of hyperplane arrangements,
however, will not be breached here.

\subsection*{Acknowledgments}

I thank Alexander Postnikov and Richard Stanley for discussions on hyperplane
arrangements that led to the results in this note.

\section{Definitions and a main result}

\subsection{Graphs and colorings}

This note will be concerned with finite graphs. While some results of this
note can be generalized to matroids, we shall not discuss this generalization
here. Let us start with the definition of a graph that we shall be using:

\begin{definition}
\label{def.graph}\textbf{(a)} If $V$ is any set, then $\dbinom{V}{2}$ will
denote the set of all $2$-element subsets of $V$. In other words, if $V$ is
any set, then we set%
\begin{align*}
\dbinom{V}{2}  &  =\left\{  S\in\mathcal{P}\left(  V\right)  \ \mid
\ \left\vert S\right\vert =2\right\} \\
&  =\left\{  \left\{  s,t\right\}  \ \mid\ s\in V,\ t\in V,\ s\neq t\right\}
\end{align*}
(where $\mathcal{P}\left(  V\right)  $ denotes the powerset of $V$).

\textbf{(b)} A \textit{graph} means a pair $\left(  V,E\right)  $, where $V$
is a set, and where $E$ is a subset of $\dbinom{V}{2}$. A graph $\left(
V,E\right)  $ is said to be \textit{finite} if the set $V$ is finite. If
$G=\left(  V,E\right)  $ is a graph, then the elements of $V$ are called the
\textit{vertices} of the graph $G$, while the elements of $E$ are called the
\textit{edges} of the graph $G$. If $e$ is an edge of a graph $G$, then the
two elements of $e$ are called the \textit{endpoints} of the edge $e$. If
$e=\left\{  s,t\right\}  $ is an edge of a graph $G$, then we say that the
edge $e$ \textit{connects the vertices }$s$ \textit{and }$t$ of $G$.
\end{definition}

Comparing our definition of a graph with some of the other definitions used in
the literature, we thus observe that our graphs are undirected (i.e., their
edges are sets, not pairs), loopless (i.e., the two endpoints of an edge must
always be distinct), edge-unlabelled (i.e., their edges are just $2$-element
sets of vertices, rather than objects with \textquotedblleft their own
identity\textquotedblright), and do not have multiple edges (or, more
precisely, there is no notion of several edges connecting two vertices, since
the edges form a set, nor a multiset, and do not have labels).

\begin{definition}
\label{def.coloring}Let $G=\left(  V,E\right)  $ be a graph. Let $X$ be a set.

\textbf{(a)} An $X$\textit{-coloring} of $G$ is defined to mean a map
$V\rightarrow X$.

\textbf{(b)} An $X$-coloring $f$ of $G$ is said to be \textit{proper} if every
edge $\left\{  s,t\right\}  \in E$ satisfies $f\left(  s\right)  \neq f\left(
t\right)  $.
\end{definition}

\subsection{Symmetric functions}

We shall now briefly introduce the notion of symmetric functions. We shall not
use any nontrivial results about symmetric functions; we will merely need some
notations.\footnote{For an introduction to symmetric functions, see any of
\cite[Chapter 7]{Stanley-EC2}, \cite[Chapter 9]{Martin15} and \cite[Chapter
2]{Reiner} (and a variety of other texts).}

In the following, $\mathbb{N}$ means the set $\left\{  0,1,2,\ldots\right\}
$. Also, $\mathbb{N}_{+}$ shall mean the set $\left\{  1,2,3,\ldots\right\}  $.

A \textit{partition} will mean a sequence $\left(  \lambda_{1},\lambda
_{2},\lambda_{3},\ldots\right)  \in\mathbb{N}^{\infty}$ of nonnegative
integers such that $\lambda_{1}\geq\lambda_{2}\geq\lambda_{3}\geq\cdots$ and
such that all sufficiently high integers $i\geq1$ satisfy $\lambda_{i}=0$. If
$\lambda=\left(  \lambda_{1},\lambda_{2},\lambda_{3},\ldots\right)  $ is a
partition, and if a positive integer $n$ is such that all integers $i\geq n$
satisfy $\lambda_{i}=0$, then we shall identify the partition $\lambda$ with
the finite sequence $\left(  \lambda_{1},\lambda_{2},\ldots,\lambda
_{n-1}\right)  $. Thus, for example, the sequences $\left(  3,1\right)  $ and
$\left(  3,1,0\right)  $ and the partition $\left(  3,1,0,0,0,\ldots\right)  $
are all identified. Every weakly decreasing finite list of positive integers
thus is identified with a unique partition.

Let $\mathbf{k}$ be a commutative ring with unity. We shall keep $\mathbf{k}$
fixed throughout the paper. The reader will not be missing out on anything if
she assumes that $\mathbf{k}=\mathbb{Z}$.

We consider the $\mathbf{k}$-algebra $\mathbf{k}\left[  \left[  x_{1}%
,x_{2},x_{3},\ldots\right]  \right]  $ of (commutative) power series in
countably many distinct indeterminates $x_{1},x_{2},x_{3},\ldots$ over
$\mathbf{k}$. It is a topological $\mathbf{k}$-algebra\footnote{See
\cite[Section 2.6]{Reiner} or \cite[\S 2]{Grinbe16} for the definition of its
topology. This topology makes sure that a sequence $\left(  P_{n}\right)
_{n\in\mathbb{N}}$ of power series converges to some power series $P$ if and
only if, for every monomial $\mathfrak{m}$, all sufficiently high
$n\in\mathbb{N}$ satisfy%
\[
\left(  \text{the }\mathfrak{m}\text{-coefficient of }P_{n}\right)  =\left(
\text{the }\mathfrak{m}\text{-coefficient of }P\right)
\]
(where the meaning of \textquotedblleft sufficiently high\textquotedblright%
\ can depend on the $\mathfrak{m}$).}. A power series $P\in\mathbf{k}\left[
\left[  x_{1},x_{2},x_{3},\ldots\right]  \right]  $ is said to be
\textit{bounded-degree} if there exists an $N\in\mathbb{N}$ such that every
monomial of degree $>N$ appears with coefficient $0$ in $P$. A power series
$P\in\mathbf{k}\left[  \left[  x_{1},x_{2},x_{3},\ldots\right]  \right]  $ is
said to be \textit{symmetric} if and only if $P$ is invariant under any
permutation of the indeterminates. We let $\Lambda$ be the subset of
$\mathbf{k}\left[  \left[  x_{1},x_{2},x_{3},\ldots\right]  \right]  $
consisting of all symmetric bounded-degree power series $P\in\mathbf{k}\left[
\left[  x_{1},x_{2},x_{3},\ldots\right]  \right]  $. This subset $\Lambda$ is
a $\mathbf{k}$-subalgebra of $\mathbf{k}\left[  \left[  x_{1},x_{2}%
,x_{3},\ldots\right]  \right]  $, and is called the $\mathbf{k}$%
\textit{-algebra of symmetric functions} over $\mathbf{k}$.

We shall now define the few families of symmetric functions that we will be
concerned with in this note. The first are the \textit{power-sum symmetric
functions}:

\begin{definition}
\label{def.powersum}Let $n$ be a positive integer. We define a power series
$p_{n}\in\mathbf{k}\left[  \left[  x_{1},x_{2},x_{3},\ldots\right]  \right]  $
by%
\begin{equation}
p_{n}=x_{1}^{n}+x_{2}^{n}+x_{3}^{n}+\cdots=\sum_{j\geq1}x_{j}^{n}.
\label{eq.def.powersum.pn}%
\end{equation}
This power series $p_{n}$ lies in $\Lambda$, and is called the $n$\textit{-th
power-sum symmetric function}.

We also set $p_{0}=1\in\Lambda$. Thus, $p_{n}$ is defined not only for all
positive integers $n$, but also for all $n\in\mathbb{N}$.
\end{definition}

\begin{definition}
\label{def.powersum2}Let $\lambda=\left(  \lambda_{1},\lambda_{2},\lambda
_{3},\ldots\right)  $ be a partition. We define a power series $p_{\lambda}%
\in\mathbf{k}\left[  \left[  x_{1},x_{2},x_{3},\ldots\right]  \right]  $ by%
\[
p_{\lambda}=\prod_{i\geq1}p_{\lambda_{i}}.
\]
This is well-defined, because the infinite product $\prod_{i\geq1}%
p_{\lambda_{i}}$ converges (indeed, all but finitely many of its factors are
$1$ (because every sufficiently high integer $i$ satisfies $\lambda_{i}=0$ and
thus $p_{\lambda_{i}}=p_{0}=1$)).
\end{definition}

We notice that every partition $\lambda=\left(  \lambda_{1},\lambda_{2}%
,\ldots,\lambda_{k}\right)  $ (written as a finite list of nonnegative
integers) satisfies%
\begin{equation}
p_{\lambda}=p_{\lambda_{1}}p_{\lambda_{2}}\cdots p_{\lambda_{k}}.
\label{eq.def.powersum2.finite-expression}%
\end{equation}


\subsection{Chromatic symmetric functions}

The next symmetric functions we introduce are the actual subject of this note;
they are the \textit{chromatic symmetric functions} and originate in
\cite[Definition 2.1]{Stanley-chsym}:

\begin{definition}
\label{def.chromsym}Let $G=\left(  V,E\right)  $ be a finite graph. For every
$\mathbb{N}_{+}$-coloring $f:V\rightarrow\mathbb{N}_{+}$, we let
$\mathbf{x}_{f}$ denote the monomial $\prod_{v\in V}x_{f\left(  v\right)  }$
in the indeterminates $x_{1},x_{2},x_{3},\ldots$. We define a power series
$X_{G}\in\mathbf{k}\left[  \left[  x_{1},x_{2},x_{3},\ldots\right]  \right]  $
by%
\[
X_{G}=\sum_{\substack{f:V\rightarrow\mathbb{N}_{+}\text{ is a}\\\text{proper
}\mathbb{N}_{+}\text{-coloring of }G}}\mathbf{x}_{f}.
\]


This power series $X_{G}$ is called the \textit{chromatic symmetric function}
of $G$.
\end{definition}

We have $X_{G}\in\Lambda$ for every finite graph $G=\left(  V,E\right)  $;
this will follow from Theorem \ref{thm.chromsym.empty} further below (but is
also rather obvious).

We notice that $X_{G}$ is denoted by $\Psi\left[  G\right]  $ in
\cite[\S 7.3.3]{Reiner}.

\subsection{Connected components}

\begin{vershort}
We shall now briefly recall the notion of connected components of a graph.
\end{vershort}

\begin{verlong}
We shall now recall the notion of connected components of a graph. This notion
is a particular case of the notion of a quotient set by an equivalence
relation; thus, we shall briefly recall the latter first.

A \textit{binary relation} on a set $X$ means a subset of $X\times X$. If
$\sim$ is a binary relation on a set $X$, and if $x$ and $y$ are two elements
of $X$, then one says that $x$ \textit{is related to }$y$ \textit{by }$\sim$
if and only if $\left(  x,y\right)  $ belongs to the set $\sim$. An
\textit{equivalence relation} is a binary relation (on a set) that is
reflexive, symmetric and transitive.

\begin{definition}
\label{def.relquot}Let $X$ be a set. Let $\sim$ be an equivalence relation on
$X$. We shall write $\sim$ infix (that is, for any $x\in X$ and $y\in X$, we
shall abbreviate \textquotedblleft$\left(  x,y\right)  \in\sim$%
\textquotedblright\ as \textquotedblleft$x\sim y$\textquotedblright). For
every $x\in X$, we let $\left[  x\right]  _{\sim}$ be the subset $\left\{
y\in X\ \mid\ x\sim y\right\}  $ of $X$; this subset $\left[  x\right]
_{\sim}$ is called the $\sim$\textit{-equivalence class} of $x$ (or the
\textit{equivalence class of }$x$ \textit{with respect to the relation }$\sim
$). A $\sim$\textit{-equivalence class} means a subset of $X$ which is the
$\sim$-equivalence class of some $x\in X$. It is well-known that any two
$\sim$-equivalence classes are either identical or disjoint, and that $X$ is
the union of all $\sim$-equivalence classes.

We let $X/\left(  \sim\right)  $ denote the set of all $\sim$-equivalence
classes. This set $X/\left(  \sim\right)  $ is called the \textit{quotient of
the set }$X$ \textit{by the equivalence relation }$\sim$. We define a map
$\pi_{X}:X\rightarrow X/\left(  \sim\right)  $ by setting%
\[
\left(  \pi_{X}\left(  x\right)  =\left[  x\right]  _{\sim}%
\ \ \ \ \ \ \ \ \ \ \text{for every }x\in X\right)  .
\]
Thus, the map $\pi_{X}$ sends every $x\in X$ to its $\sim$-equivalence class.
\end{definition}

Let us introduce a few more notations:

\begin{definition}
\label{def.relquot.maps}Let $X$ and $Y$ be two sets.

\textbf{(a)} Then, $Y^{X}$ denotes the set of all maps $X\rightarrow Y$.

\textbf{(b)} Let $\sim$ be any binary relation on $X$. We shall write $\sim$
infix (that is, for any $x\in X$ and $y\in X$, we shall abbreviate
\textquotedblleft$\left(  x,y\right)  \in\sim$\textquotedblright\ as
\textquotedblleft$x\sim y$\textquotedblright). Then, we let $Y_{\sim}^{X}$
denote the subset%
\[
\left\{  g\in Y^{X}\ \mid\ g\left(  x\right)  =g\left(  y\right)  \text{ for
any }x\in X\text{ and }y\in X\text{ satisfying }x\sim y\right\}
\]
of $Y^{X}$.
\end{definition}

The map $\pi_{X}$ defined in Definition \ref{def.relquot} has an important
universal property:

\begin{proposition}
\label{prop.relquot.uniprop}Let $X$ be a set. Let $\sim$ be an equivalence
relation on $X$.

\textbf{(a)} The map $\pi_{X}:X\rightarrow X/\left(  \sim\right)  $ is
surjective and belongs to $\left(  X/\left(  \sim\right)  \right)  _{\sim}%
^{X}$.

\textbf{(b)} Let $Y$ be any set. Then, the map%
\[
Y^{X/\left(  \sim\right)  }\rightarrow Y_{\sim}^{X}%
,\ \ \ \ \ \ \ \ \ \ f\mapsto f\circ\pi_{X}%
\]
is a bijection.
\end{proposition}

\begin{proof}
[Proof of Proposition \ref{prop.relquot.uniprop}.]\textbf{(a)} The map
$\pi_{X}$ is surjective\footnote{\textit{Proof.} Let $u\in X/\left(
\sim\right)  $. Thus, $u$ is a $\sim$-equivalence class (since $X/\left(
\sim\right)  $ is the set of all $\sim$-equivalence classes). In other words,
there exists some $x\in X$ such that $u$ is the $\sim$-equivalence class of
$x$ (by the definition of a \textquotedblleft$\sim$-equivalence
class\textquotedblright). Consider this $x$.
\par
We know that $\left[  x\right]  _{\sim}$ is the $\sim$-equivalence class of
$x$. In other words, $\left[  x\right]  _{\sim}$ is $u$ (since $u$ is the the
$\sim$-equivalence class of $x$). Thus, $\left[  x\right]  _{\sim}=u$. Now,
the definition of $\pi_{X}$ yields $\pi_{X}\left(  x\right)  =\left[
x\right]  _{\sim}=u$. Thus, $u=\pi_{X}\left(  \underbrace{x}_{\in X}\right)
\in\pi_{X}\left(  X\right)  $.
\par
Now, let us forget that we fixed $u$. We thus have shown that $u\in\pi
_{X}\left(  X\right)  $ for every $u\in X/\left(  \sim\right)  $. In other
words, $X/\left(  \sim\right)  \subseteq\pi_{X}\left(  X\right)  $. In other
words, the map $\pi_{X}$ is surjective. Qed.}.

The definition of $\left(  X/\left(  \sim\right)  \right)  _{\sim}^{X}$ shows
that%
\begin{align}
&  \left(  X/\left(  \sim\right)  \right)  _{\sim}^{X}\nonumber\\
&  =\left\{  g\in\left(  X/\left(  \sim\right)  \right)  ^{X}\ \mid\ g\left(
x\right)  =g\left(  y\right)  \text{ for any }x\in X\text{ and }y\in X\text{
satisfying }x\sim y\right\}  \label{pf.prop.relquot.uniprop.a.def}%
\end{align}


Now, we claim that
\begin{equation}
\pi_{X}\left(  x\right)  =\pi_{X}\left(  y\right)
\ \ \ \ \ \ \ \ \ \ \text{for any }x\in X\text{ and }y\in X\text{ satisfying
}x\sim y. \label{pf.prop.relquot.uniprop.a.eq}%
\end{equation}
\textit{Proof of (\ref{pf.prop.relquot.uniprop.a.eq}):} Let $x\in X$ and $y\in
X$ be any two elements satisfying $x\sim y$. We shall show that $\pi
_{X}\left(  x\right)  =\pi_{X}\left(  y\right)  $.

Recall that $\pi_{X}\left(  x\right)  =\left[  x\right]  _{\sim}$. Thus,
$\pi_{X}\left(  x\right)  $ is the $\sim$-equivalence class of $x$ (since
$\left[  x\right]  _{\sim}$ is the $\sim$-equivalence class of $x$). The same
argument (applied to $y$ instead of $x$) shows that $\pi_{X}\left(  y\right)
$ is the $\sim$-equivalence class of $y$.

We have $x\sim y$. Thus, the elements $x$ and $y$ of $X$ lie in the same
$\sim$-equivalence class. In other words, the $\sim$-equivalence class of $x$
is the $\sim$-equivalence class of $y$. In other words, $\pi_{X}\left(
x\right)  $ is $\pi_{X}\left(  y\right)  $ (since $\pi_{X}\left(  x\right)  $
is the $\sim$-equivalence class of $x$, and since $\pi_{X}\left(  y\right)  $
is the $\sim$-equivalence class of $y$). In other words, $\pi_{X}\left(
x\right)  =\pi_{X}\left(  y\right)  $. This proves
(\ref{pf.prop.relquot.uniprop.a.eq}).

Now, $\pi_{X}$ is a map $g\in\left(  X/\left(  \sim\right)  \right)  ^{X}$
which satisfies $g\left(  x\right)  =g\left(  y\right)  $ for any $x\in X$ and
$y\in X$ satisfying $x\sim y$ (according to
(\ref{pf.prop.relquot.uniprop.a.eq})). In other words,%
\begin{align*}
\pi_{X}  &  \in\left\{  g\in\left(  X/\left(  \sim\right)  \right)  ^{X}%
\ \mid\ g\left(  x\right)  =g\left(  y\right)  \text{ for any }x\in X\text{
and }y\in X\text{ satisfying }x\sim y\right\} \\
&  =\left(  X/\left(  \sim\right)  \right)  _{\sim}^{X}%
\ \ \ \ \ \ \ \ \ \ \left(  \text{by (\ref{pf.prop.relquot.uniprop.a.def}%
)}\right)  .
\end{align*}
This completes the proof of Proposition \ref{prop.relquot.uniprop}
\textbf{(a)}.

\textbf{(b)} We have%
\begin{equation}
f\circ\pi_{X}\in Y_{\sim}^{X}\ \ \ \ \ \ \ \ \ \ \text{for every }f\in
Y^{X/\left(  \sim\right)  } \label{pf.prop.relquot.uniprop.wd}%
\end{equation}
\footnote{\textit{Proof of (\ref{pf.prop.relquot.uniprop.wd}):} Let $f\in
Y^{X/\left(  \sim\right)  }$. Then, clearly, $f\circ\pi_{X}\in Y^{X}$ (since
$\pi_{X}\in\left(  X/\left(  \sim\right)  \right)  ^{X}$). Now, let $x\in X$
and $y\in X$ be any two elements satisfying $x\sim y$. We shall show that
$\left(  f\circ\pi_{X}\right)  \left(  x\right)  =\left(  f\circ\pi
_{X}\right)  \left(  y\right)  $.
\par
From (\ref{pf.prop.relquot.uniprop.a.eq}), we obtain $\pi_{X}\left(  x\right)
=\pi_{X}\left(  y\right)  $. Now, $\left(  f\circ\pi_{X}\right)  \left(
x\right)  =f\left(  \underbrace{\pi_{X}\left(  x\right)  }_{=\pi_{X}\left(
y\right)  }\right)  =f\left(  \pi_{X}\left(  y\right)  \right)  =\left(
f\circ\pi_{X}\right)  \left(  y\right)  $.
\par
Now, let us forget that we fixed $x$ and $y$. We thus have shown that $\left(
f\circ\pi_{X}\right)  \left(  x\right)  =\left(  f\circ\pi_{X}\right)  \left(
y\right)  $ for any $x\in X$ and $y\in X$ satisfying $x\sim y$. Thus,
$f\circ\pi_{X}$ is a map $g\in Y^{X}$ which satisfies $g\left(  x\right)
=g\left(  y\right)  $ for any $x\in X$ and $y\in X$ satisfying $x\sim y$. In
other words,%
\begin{align*}
f\circ\pi_{X}  &  \in\left\{  g\in Y^{X}\ \mid\ g\left(  x\right)  =g\left(
y\right)  \text{ for any }x\in X\text{ and }y\in X\text{ satisfying }x\sim
y\right\} \\
&  =Y_{\sim}^{X}%
\end{align*}
(because $Y_{\sim}^{X}$ was defined to be $\left\{  g\in Y^{X}\ \mid\ g\left(
x\right)  =g\left(  y\right)  \text{ for any }x\in X\text{ and }y\in X\text{
satisfying }x\sim y\right\}  $). This proves (\ref{pf.prop.relquot.uniprop.wd}%
).}. Hence, we can define a map $\Phi:Y^{X/\left(  \sim\right)  }\rightarrow
Y_{\sim}^{X}$ by%
\begin{equation}
\left(  \Phi\left(  f\right)  =f\circ\pi_{X}\ \ \ \ \ \ \ \ \ \ \text{for
every }f\in Y^{X/\left(  \sim\right)  }\right)  .
\label{pf.prop.relquot.uniprop.Phi}%
\end{equation}
Consider this map $\Phi$. We shall now show that the map $\Phi$ is a bijection.

The map $\Phi$ is injective\footnote{\textit{Proof.} Let $f$ and $g$ be two
elements of $Y^{X/\left(  \sim\right)  }$ such that $\Phi\left(  f\right)
=\Phi\left(  g\right)  $. We shall show that $f=g$.
\par
Let $u\in X/\left(  \sim\right)  $. Thus, $u\in X/\left(  \sim\right)
=\pi_{X}\left(  X\right)  $ (since the map $\pi_{X}$ is surjective). In other
words, there exists some $x\in X$ such that $u=\pi_{X}\left(  x\right)  $.
Consider this $x$.
\par
Now, the definition of $\Phi$ yields $\Phi\left(  f\right)  =f\circ\pi_{X}$.
Hence, $\left(  \underbrace{\Phi\left(  f\right)  }_{=f\circ\pi_{X}}\right)
\left(  x\right)  =\left(  f\circ\pi_{X}\right)  \left(  x\right)  =f\left(
\underbrace{\pi_{X}\left(  x\right)  }_{=u}\right)  =f\left(  u\right)  $. The
same argument (but applied to $g$ instead of $f$) yields $\left(  \Phi\left(
g\right)  \right)  \left(  x\right)  =g\left(  u\right)  $. Now, $f\left(
u\right)  =\left(  \underbrace{\Phi\left(  f\right)  }_{=\Phi\left(  g\right)
}\right)  \left(  x\right)  =\left(  \Phi\left(  g\right)  \right)  \left(
x\right)  =g\left(  u\right)  $.
\par
Let us now forget that we fixed $u$. We thus have shown that $f\left(
u\right)  =g\left(  u\right)  $ for every $u\in X/\left(  \sim\right)  $. In
other words, $f=g$.
\par
Now, let us forget that we fixed $f$ and $g$. We thus have proven that if $f$
and $g$ are two elements of $Y^{X/\left(  \sim\right)  }$ such that
$\Phi\left(  f\right)  =\Phi\left(  g\right)  $, then $f=g$. In other words,
the map $\Phi$ is injective. Qed.}. We shall now show that the map $\Phi$ is surjective.

Indeed, let $h\in Y_{\sim}^{X}$.

We shall now define a map $f\in Y^{X/\left(  \sim\right)  }$ as follows: Let
$u\in X/\left(  \sim\right)  $. Thus, $u\in X/\left(  \sim\right)  =\pi
_{X}\left(  X\right)  $ (since the map $\pi_{X}$ is surjective). Thus, there
exists some $x\in X$ satisfying $u=\pi_{X}\left(  x\right)  $. Pick such an
$x$. Then, $h\left(  x\right)  $ is independent on the choice of
$x$\ \ \ \ \footnote{\textit{Proof.} Let $x_{1}$ and $x_{2}$ be any two $x\in
X$ satisfying $u=\pi_{X}\left(  x\right)  $. We shall show that $h\left(
x_{1}\right)  =h\left(  x_{2}\right)  $.
\par
Indeed, $x_{1}$ is an $x\in X$ satisfying $u=\pi_{X}\left(  x\right)  $. In
other words, $x_{1}$ is an element of $X$ and satisfies $u=\pi_{X}\left(
x_{1}\right)  $. But the definition of $\pi_{X}$ shows that $\pi_{X}\left(
x_{1}\right)  =\left[  x_{1}\right]  _{\sim}$. Thus, $u=\pi_{X}\left(
x_{1}\right)  =\left[  x_{1}\right]  _{\sim}$. In other words, $u$ is $\left[
x_{1}\right]  _{\sim}$. In other words, $u$ is the $\sim$-equivalence class of
$x_{1}$ (since $\left[  x_{1}\right]  _{\sim}$ is the $\sim$-equivalence class
of $x_{1}$ (by the definition of $\left[  x_{1}\right]  _{\sim}$).
\par
The same argument (applied to $x_{2}$ instead of $x_{1}$) shows that $u$ is
the $\sim$-equivalence class of $x_{2}$. Thus, the $\sim$-equivalence classes
of $x_{1}$ and $x_{2}$ are both $u$. Hence, these two $\sim$-equivalence
classes are identical. In other words, $x_{1}$ and $x_{2}$ belong to the same
$\sim$-equivalence class. In other words, $x_{1}\sim x_{2}$.
\par
But%
\[
h\in Y_{\sim}^{X}=\left\{  g\in Y^{X}\ \mid\ g\left(  x\right)  =g\left(
y\right)  \text{ for any }x\in X\text{ and }y\in X\text{ satisfying }x\sim
y\right\}
\]
(by the definition of $Y_{\sim}^{X}$). In other words, $h$ is a $g\in Y^{X}$
which satisfies $g\left(  x\right)  =g\left(  y\right)  $ for any $x\in X$ and
$y\in X$ satisfying $x\sim y$. In other words, $h$ is an element of $Y^{X}$
and satisfies%
\begin{equation}
h\left(  x\right)  =h\left(  y\right)  \text{ for any }x\in X\text{ and }y\in
X\text{ satisfying }x\sim y. \label{pf.prop.relquot.uniprop.b.surj.fn1.1}%
\end{equation}
\par
Now, we can apply (\ref{pf.prop.relquot.uniprop.b.surj.fn1.1}) to $x=x_{1}$
and $y=x_{2}$ (since $x_{1}\sim x_{2}$). As a result, we obtain $h\left(
x_{1}\right)  =h\left(  x_{2}\right)  $.
\par
Let us now forget that we fixed $x_{1}$ and $x_{2}$. We thus have shown that
if $x_{1}$ and $x_{2}$ are any two $x\in X$ satisfying $u=\pi_{X}\left(
x\right)  $, then $h\left(  x_{1}\right)  =h\left(  x_{2}\right)  $. In other
words, $h\left(  x\right)  $ is independent on the choice of $x$ (when $x$ is
chosen as an element of $X$ satisfying $u=\pi_{X}\left(  x\right)  $). Qed.}.
Hence, we can set $f\left(  u\right)  =h\left(  x\right)  $.

Thus, we have defined a map $f\in Y^{X/\left(  \sim\right)  }$. This map has
the following property: If $u\in X/\left(  \sim\right)  $, and if $x\in X$ is
such that $u=\pi_{X}\left(  x\right)  $, then%
\begin{equation}
f\left(  u\right)  =h\left(  x\right)  \label{pf.prop.relquot.uniprop.surj.1}%
\end{equation}
(because this is how $f\left(  u\right)  $ was defined).

Now, $f\circ\pi_{X}=h$\ \ \ \ \footnote{\textit{Proof.} Every $x\in X$
satisfies $\left(  f\circ\pi_{X}\right)  \left(  x\right)  =f\left(  \pi
_{X}\left(  x\right)  \right)  =h\left(  x\right)  $ (by
(\ref{pf.prop.relquot.uniprop.surj.1}), applied to $u=\pi_{X}\left(  x\right)
$). In other words, $f\circ\pi_{X}=h$, qed.}. The definition of $\Phi$ now
yields $\Phi\left(  f\right)  =f\circ\pi_{X}=h$. Hence, $h=\Phi\left(
\underbrace{f}_{\in Y^{X/\left(  \sim\right)  }}\right)  \in\Phi\left(
Y^{X/\left(  \sim\right)  }\right)  $.

Now, let us forget that we fixed $h$. We thus have proven that $h\in
\Phi\left(  Y^{X/\left(  \sim\right)  }\right)  $ for every $h\in Y_{\sim}%
^{X}$. In other words, the map $\Phi$ is surjective.

Thus, we know that the map $\Phi$ is injective and surjective. In other words,
$\Phi$ is bijective. In other words, $\Phi$ is a bijection. Since $\Phi$ is
the map
\[
Y^{X/\left(  \sim\right)  }\rightarrow Y_{\sim}^{X}%
,\ \ \ \ \ \ \ \ \ \ f\mapsto f\circ\pi_{X}%
\]
(because $\Phi$ is a map $Y^{X/\left(  \sim\right)  }\rightarrow Y_{\sim}^{X}$
and satisfies (\ref{pf.prop.relquot.uniprop.Phi})), this shows that the map
\[
Y^{X/\left(  \sim\right)  }\rightarrow Y_{\sim}^{X}%
,\ \ \ \ \ \ \ \ \ \ f\mapsto f\circ\pi_{X}%
\]
is a bijection. This proves Proposition \ref{prop.relquot.uniprop}
\textbf{(b)}.
\end{proof}

We shall now recall what connected components are:
\end{verlong}

\begin{definition}
\label{def.path}Let $G=\left(  V,E\right)  $ be a finite graph. Let $u$ and
$v$ be two elements of $V$ (that is, two vertices of $G$). A \textit{walk}
from $u$ to $v$ in $G$ will mean a sequence $\left(  w_{0},w_{1},\ldots
,w_{k}\right)  $ of elements of $V$ such that $w_{0}=u$ and $w_{k}=v$ and%
\[
\left(  \left\{  w_{i},w_{i+1}\right\}  \in E\ \ \ \ \ \ \ \ \ \ \text{for
every }i\in\left\{  0,1,\ldots,k-1\right\}  \right)  .
\]
We say that $u$ and $v$ are \textit{connected (in }$G$\textit{)} if there
exists a walk from $u$ to $v$ in $G$.
\end{definition}

\begin{definition}
\label{def.connectedness}Let $G=\left(  V,E\right)  $ be a graph.

\textbf{(a)} We define a binary relation $\sim_{G}$ (written infix) on the set
$V$ as follows: Given $u\in V$ and $v\in V$, we set $u\sim_{G}v$ if and only
if $u$ and $v$ are connected (in $G$). It is well-known that this relation
$\sim_{G}$ is an equivalence relation. The $\sim_{G}$-equivalence classes are
called the \textit{connected components} of $G$.

\textbf{(b)} Assume that the graph $G$ is finite. We let $\lambda\left(
G\right)  $ denote the list of the sizes of all connected components of $G$,
in weakly decreasing order. (Each connected component should contribute only
one entry to the list.) We view $\lambda\left(  G\right)  $ as a partition
(since $\lambda\left(  G\right)  $ is a weakly decreasing finite list of
positive integers).
\end{definition}

Now, we can state a formula for chromatic symmetric functions:

\begin{theorem}
\label{thm.chromsym.empty}Let $G=\left(  V,E\right)  $ be a finite graph.
Then,%
\[
X_{G}=\sum_{F\subseteq E}\left(  -1\right)  ^{\left\vert F\right\vert
}p_{\lambda\left(  V,F\right)  }.
\]
(Here, of course, the pair $\left(  V,F\right)  $ is regarded as a graph, and
the expression $\lambda\left(  V,F\right)  $ is understood according to
Definition \ref{def.connectedness} \textbf{(b)}.)
\end{theorem}

This theorem is not new; it appears, e.g., in \cite[Theorem 2.5]%
{Stanley-chsym}. We shall show a far-reaching generalization of it (Theorem
\ref{thm.chromsym.varis}) soon.

\subsection{Circuits and broken circuits}

Let us now define the notions of cycles and circuits of a graph:

\begin{definition}
\label{def.cycle}Let $G=\left(  V,E\right)  $ be a graph. A \textit{cycle} of
$G$ denotes a list $\left(  v_{1},v_{2},\ldots,v_{m+1}\right)  $ of elements
of $V$ with the following properties:

\begin{itemize}
\item We have $m>1$.

\item We have $v_{m+1}=v_{1}$.

\item The vertices $v_{1},v_{2},\ldots,v_{m}$ are pairwise distinct.

\item We have $\left\{  v_{i},v_{i+1}\right\}  \in E$ for every $i\in\left\{
1,2,\ldots,m\right\}  $.
\end{itemize}

If $\left(  v_{1},v_{2},\ldots,v_{m+1}\right)  $ is a cycle of $G$, then the
set $\left\{  \left\{  v_{1},v_{2}\right\}  ,\left\{  v_{2},v_{3}\right\}
,\ldots,\left\{  v_{m},v_{m+1}\right\}  \right\}  $ is called a
\textit{circuit} of $G$.
\end{definition}

\begin{definition}
\label{def.BC}Let $G=\left(  V,E\right)  $ be a graph. Let $X$ be a totally
ordered set. Let $\ell:E\rightarrow X$ be a function. We shall refer to $\ell$
as the \textit{labeling function}. For every edge $e$ of $G$, we shall refer
to $\ell\left(  e\right)  $ as the \textit{label} of $e$.

A \textit{broken circuit} of $G$ means a subset of $E$ having the form
$C\setminus\left\{  e\right\}  $, where $C$ is a circuit of $G$, and where $e$
is the unique edge in $C$ having maximum label (among the edges in $C$). Of
course, the notion of a broken circuit of $G$ depends on the function $\ell$;
however, we suppress the mention of $\ell$ in our notation, since we will not
consider situations where two different $\ell$'s coexist.
\end{definition}

Thus, if $G$ is a graph with a labeling function, then any circuit $C$ of $G$
gives rise to a broken circuit provided that among the edges in $C$, only one
attains the maximum label. (If more than one of the edges of $C$ attains the
maximum label, then $C$ does not give rise to a broken circuit.) Notice that
two different circuits may give rise to one and the same broken circuit.

\begin{theorem}
\label{thm.chromsym.varis}Let $G=\left(  V,E\right)  $ be a finite graph. Let
$X$ be a totally ordered set. Let $\ell:E\rightarrow X$ be a function. Let
$\mathfrak{K}$ be some set of broken circuits of $G$ (not necessarily
containing all of them). Let $a_{K}$ be an element of $\mathbf{k}$ for every
$K\in\mathfrak{K}$. Then,%
\[
X_{G}=\sum_{F\subseteq E}\left(  -1\right)  ^{\left\vert F\right\vert }\left(
\prod_{\substack{K\in\mathfrak{K};\\K\subseteq F}}a_{K}\right)  p_{\lambda
\left(  V,F\right)  }.
\]
(Here, of course, the pair $\left(  V,F\right)  $ is regarded as a graph, and
the expression $\lambda\left(  V,F\right)  $ is understood according to
Definition \ref{def.connectedness} \textbf{(b)}.)
\end{theorem}

Before we come to the proof of this result, let us explore some of its
particular cases. First, a definition is in order:

\begin{definition}
\label{def.K-free}Let $E$ be a set. Let $\mathfrak{K}$ be a subset of the
powerset of $E$ (that is, a set of subsets of $E$). A subset $F$ of $E$ is
said to be $\mathfrak{K}$\textit{-free} if $F$ contains no $K\in\mathfrak{K}$
as a subset. (For instance, if $\mathfrak{K}=\varnothing$, then every subset
$F$ of $E$ is $\mathfrak{K}$-free.)
\end{definition}

\begin{corollary}
\label{cor.chromsym.K-free}Let $G=\left(  V,E\right)  $ be a finite graph. Let
$X$ be a totally ordered set. Let $\ell:E\rightarrow X$ be a function. Let
$\mathfrak{K}$ be some set of broken circuits of $G$ (not necessarily
containing all of them). Then,%
\[
X_{G}=\sum_{\substack{F\subseteq E;\\F\text{ is }\mathfrak{K}\text{-free}%
}}\left(  -1\right)  ^{\left\vert F\right\vert }p_{\lambda\left(  V,F\right)
}.
\]

\end{corollary}

\begin{corollary}
\label{cor.chromsym.NBC}Let $G=\left(  V,E\right)  $ be a finite graph. Let
$X$ be a totally ordered set. Let $\ell:E\rightarrow X$ be a function. Then,%
\[
X_{G}=\sum_{\substack{F\subseteq E;\\F\text{ contains no broken}%
\\\text{circuit of }G\text{ as a subset}}}\left(  -1\right)  ^{\left\vert
F\right\vert }p_{\lambda\left(  V,F\right)  }.
\]

\end{corollary}

Corollary \ref{cor.chromsym.NBC} appears in \cite[Theorem 2.9]{Stanley-chsym},
at least in the particular case in which $\ell$ is supposed to be injective.

Let us now see how Theorem \ref{thm.chromsym.empty}, Corollary
\ref{cor.chromsym.K-free} and Corollary \ref{cor.chromsym.NBC} can be derived
from Theorem \ref{thm.chromsym.varis}:

\begin{vershort}
\begin{proof}
[Proof of Corollary \ref{cor.chromsym.K-free} using Theorem
\ref{thm.chromsym.varis}.]For every subset $F$ of $E$, we have%
\begin{equation}
\prod_{\substack{K\in\mathfrak{K};\\K\subseteq F}}0=%
\begin{cases}
1, & \text{if }F\text{ is }\mathfrak{K}\text{-free;}\\
0, & \text{if }F\text{ is not }\mathfrak{K}\text{-free}%
\end{cases}
\label{pf.cor.chromsym.K-free.short.prod0}%
\end{equation}
(because if $F$ is $\mathfrak{K}$-free, then the product $\prod
_{\substack{K\in\mathfrak{K};\\K\subseteq F}}0$ is empty and thus equals $1$;
otherwise, the product $\prod_{\substack{K\in\mathfrak{K};\\K\subseteq F}}0$
contains at least one factor and thus equals $0$). Now, Theorem
\ref{thm.chromsym.varis} (applied to $0$ instead of $a_{K}$) yields%
\begin{align*}
X_{G}  &  =\sum_{F\subseteq E}\left(  -1\right)  ^{\left\vert F\right\vert
}\underbrace{\left(  \prod_{\substack{K\in\mathfrak{K};\\K\subseteq
F}}0\right)  }_{\substack{=%
\begin{cases}
1, & \text{if }F\text{ is }\mathfrak{K}\text{-free;}\\
0, & \text{if }F\text{ is not }\mathfrak{K}\text{-free}%
\end{cases}
\\\text{(by (\ref{pf.cor.chromsym.K-free.short.prod0}))}}}p_{\lambda\left(
V,F\right)  }\\
&  =\sum_{F\subseteq E}\left(  -1\right)  ^{\left\vert F\right\vert }%
\begin{cases}
1, & \text{if }F\text{ is }\mathfrak{K}\text{-free;}\\
0, & \text{if }F\text{ is not }\mathfrak{K}\text{-free}%
\end{cases}
p_{\lambda\left(  V,F\right)  }=\sum_{\substack{F\subseteq E;\\F\text{ is
}\mathfrak{K}\text{-free}}}\left(  -1\right)  ^{\left\vert F\right\vert
}p_{\lambda\left(  V,F\right)  }.
\end{align*}
This proves Corollary \ref{cor.chromsym.K-free}.
\end{proof}
\end{vershort}

\begin{verlong}
\begin{proof}
[Proof of Corollary \ref{cor.chromsym.K-free} using Theorem
\ref{thm.chromsym.varis}.]We can apply Theorem \ref{thm.chromsym.varis} to $0$
instead of $a_{K}$. As a result, we obtain%
\begin{equation}
X_{G}=\sum_{F\subseteq E}\left(  -1\right)  ^{\left\vert F\right\vert }\left(
\prod_{\substack{K\in\mathfrak{K};\\K\subseteq F}}0\right)  p_{\lambda\left(
V,F\right)  }. \label{pf.cor.chromsym.K-free.0}%
\end{equation}
Now, if $F$ is any subset of $E$, then%
\begin{equation}
\prod_{\substack{K\in\mathfrak{K};\\K\subseteq F}}0=%
\begin{cases}
1, & \text{if }F\text{ is }\mathfrak{K}\text{-free;}\\
0, & \text{if }F\text{ is not }\mathfrak{K}\text{-free}%
\end{cases}
\label{pf.cor.chromsym.K-free.1}%
\end{equation}
\footnote{\textit{Proof of (\ref{pf.cor.chromsym.K-free.1}):} Let $F$ be any
subset of $E$. Recall that the set $F$ is $\mathfrak{K}$-free if and only if
$F$ contains no $K\in\mathfrak{K}$ as a subset (by the definition of
\textquotedblleft$\mathfrak{K}$-free\textquotedblright). Taking the
contrapositive of this equivalence statement, we obtain the following: The set
$F$ is \textbf{not} $\mathfrak{K}$-free if and only if $F$ contains some
$K\in\mathfrak{K}$ as a subset.
\par
We are in one of the following two cases:
\par
\textit{Case 1:} The set $F$ is $\mathfrak{K}$-free.
\par
\textit{Case 2:} The set $F$ is not $\mathfrak{K}$-free.
\par
Let us first consider Case 1. In this case, the set $F$ is $\mathfrak{K}%
$-free. In other words, $F$ contains no $K\in\mathfrak{K}$ as a subset
(because the set $F$ is $\mathfrak{K}$-free if and only if $F$ contains no
$K\in\mathfrak{K}$ as a subset (by the definition of \textquotedblleft%
$\mathfrak{K}$-free\textquotedblright)). In other words, there exists no
$K\in\mathfrak{K}$ such that $F$ contains $K$ as a subset. In other words,
there exists no $K\in\mathfrak{K}$ such that $K\subseteq F$. Hence, the
product $\prod_{\substack{K\in\mathfrak{K};\\K\subseteq F}}0$ is empty. Thus,
$\prod_{\substack{K\in\mathfrak{K};\\K\subseteq F}}0=\left(  \text{empty
product}\right)  =1$. Comparing this with
\[%
\begin{cases}
1, & \text{if }F\text{ is }\mathfrak{K}\text{-free;}\\
0, & \text{if }F\text{ is not }\mathfrak{K}\text{-free}%
\end{cases}
=1\ \ \ \ \ \ \ \ \ \ \left(  \text{since }F\text{ is }\mathfrak{K}%
\text{-free}\right)  ,
\]
we obtain $\prod_{\substack{K\in\mathfrak{K};\\K\subseteq F}}0=%
\begin{cases}
1, & \text{if }F\text{ is }\mathfrak{K}\text{-free;}\\
0, & \text{if }F\text{ is not }\mathfrak{K}\text{-free}%
\end{cases}
$. Thus, (\ref{pf.cor.chromsym.K-free.1}) is proven in Case 1.
\par
Let us now consider Case 2. In this case, the set $F$ is \textbf{not}
$\mathfrak{K}$-free. In other words, $F$ contains some $K\in\mathfrak{K}$ as a
subset (since the set $F$ is \textbf{not} $\mathfrak{K}$-free if and only if
$F$ contains some $K\in\mathfrak{K}$ as a subset). Let $L$ be such a $K$.
Thus, $L$ is an element of $\mathfrak{K}$, and the set $F$ contains $L$ as a
subset.
\par
Now, $L\subseteq F$ (since $F$ contains $L$ as a subset). Hence, the product
$\prod_{\substack{K\in\mathfrak{K};\\K\subseteq F}}0$ has at least one factor
(namely, the factor for $K=L$). This factor is clearly $0$. Therefore, the
whole product $\prod_{\substack{K\in\mathfrak{K};\\K\subseteq F}}0$ equals $0$
(because if a product contains a factor which is $0$, then the whole product
must equal $0$). In other words, $\prod_{\substack{K\in\mathfrak{K}%
;\\K\subseteq F}}0=0$. Comparing this with
\[%
\begin{cases}
1, & \text{if }F\text{ is }\mathfrak{K}\text{-free;}\\
0, & \text{if }F\text{ is not }\mathfrak{K}\text{-free}%
\end{cases}
=0\ \ \ \ \ \ \ \ \ \ \left(  \text{since }F\text{ is \textbf{not}
}\mathfrak{K}\text{-free}\right)  ,
\]
we obtain $\prod_{\substack{K\in\mathfrak{K};\\K\subseteq F}}0=%
\begin{cases}
1, & \text{if }F\text{ is }\mathfrak{K}\text{-free;}\\
0, & \text{if }F\text{ is not }\mathfrak{K}\text{-free}%
\end{cases}
$. Thus, (\ref{pf.cor.chromsym.K-free.1}) is proven in Case 2.
\par
We now have proven (\ref{pf.cor.chromsym.K-free.1}) in each of the two Cases 1
and 2. Since these two Cases cover all possibilities, this shows that
(\ref{pf.cor.chromsym.K-free.1}) always holds.}. Thus,
(\ref{pf.cor.chromsym.K-free.0}) becomes%
\begin{align*}
X_{G}  &  =\sum_{F\subseteq E}\left(  -1\right)  ^{\left\vert F\right\vert
}\underbrace{\left(  \prod_{\substack{K\in\mathfrak{K};\\K\subseteq
F}}0\right)  }_{\substack{=%
\begin{cases}
1, & \text{if }F\text{ is }\mathfrak{K}\text{-free;}\\
0, & \text{if }F\text{ is not }\mathfrak{K}\text{-free}%
\end{cases}
\\\text{(by (\ref{pf.cor.chromsym.K-free.1}))}}}p_{\lambda\left(  V,F\right)
}\\
&  =\sum_{F\subseteq E}\left(  -1\right)  ^{\left\vert F\right\vert }%
\begin{cases}
1, & \text{if }F\text{ is }\mathfrak{K}\text{-free;}\\
0, & \text{if }F\text{ is not }\mathfrak{K}\text{-free}%
\end{cases}
p_{\lambda\left(  V,F\right)  }\\
&  =\sum_{\substack{F\subseteq E;\\F\text{ is }\mathfrak{K}\text{-free}%
}}\left(  -1\right)  ^{\left\vert F\right\vert }\underbrace{%
\begin{cases}
1, & \text{if }F\text{ is }\mathfrak{K}\text{-free;}\\
0, & \text{if }F\text{ is not }\mathfrak{K}\text{-free}%
\end{cases}
}_{\substack{=1\\\text{(since }F\text{ is }\mathfrak{K}\text{-free)}%
}}p_{\lambda\left(  V,F\right)  }+\sum_{\substack{F\subseteq E;\\F\text{ is
not }\mathfrak{K}\text{-free}}}\left(  -1\right)  ^{\left\vert F\right\vert
}\underbrace{%
\begin{cases}
1, & \text{if }F\text{ is }\mathfrak{K}\text{-free;}\\
0, & \text{if }F\text{ is not }\mathfrak{K}\text{-free}%
\end{cases}
}_{\substack{=0\\\text{(since }F\text{ is not }\mathfrak{K}\text{-free)}%
}}p_{\lambda\left(  V,F\right)  }\\
&  =\sum_{\substack{F\subseteq E;\\F\text{ is }\mathfrak{K}\text{-free}%
}}\left(  -1\right)  ^{\left\vert F\right\vert }p_{\lambda\left(  V,F\right)
}+\underbrace{\sum_{\substack{F\subseteq E;\\F\text{ is not }\mathfrak{K}%
\text{-free}}}\left(  -1\right)  ^{\left\vert F\right\vert }0p_{\lambda\left(
V,F\right)  }}_{=0}\\
&  =\sum_{\substack{F\subseteq E;\\F\text{ is }\mathfrak{K}\text{-free}%
}}\left(  -1\right)  ^{\left\vert F\right\vert }p_{\lambda\left(  V,F\right)
}.
\end{align*}
This proves Corollary \ref{cor.chromsym.K-free}.
\end{proof}
\end{verlong}

\begin{vershort}
\begin{proof}
[Proof of Corollary \ref{cor.chromsym.NBC} using Corollary
\ref{cor.chromsym.K-free}.]Corollary \ref{cor.chromsym.NBC} follows from
Corollary \ref{cor.chromsym.K-free} when $\mathfrak{K}$ is set to be the set
of \textbf{all} broken circuits of $G$.
\end{proof}
\end{vershort}

\begin{verlong}
\begin{proof}
[Proof of Corollary \ref{cor.chromsym.NBC} using Corollary
\ref{cor.chromsym.K-free}.]Let $\mathfrak{K}$ be the set of all broken
circuits of $G$. Thus, the elements of $\mathfrak{K}$ are the broken circuits
of $G$.

Now, for every subset $F$ of $E$, we have the following equivalence of
statements:%
\begin{align*}
&  \ \left(  F\text{ is }\mathfrak{K}\text{-free}\right) \\
&  \Longleftrightarrow\ \left(  F\text{ contains no }K\in\mathfrak{K}\text{ as
a subset}\right) \\
&  \ \ \ \ \ \ \ \ \ \ \left(
\begin{array}
[c]{c}%
\text{because }F\text{ is }\mathfrak{K}\text{-free if and only if }F\text{
contains no }K\in\mathfrak{K}\text{ as a subset}\\
\text{(by the definition of \textquotedblleft}\mathfrak{K}%
\text{-free\textquotedblright)}%
\end{array}
\right) \\
&  \Longleftrightarrow\ \left(  F\text{ contains no element of }%
\mathfrak{K}\text{ as a subset}\right) \\
&  \Longleftrightarrow\ \left(  F\text{ contains no broken circuit of }G\text{
as a subset}\right)
\end{align*}
(since the elements of $\mathfrak{K}$ are the broken circuits of $G$). Hence,
$\sum_{\substack{F\subseteq E;\\F\text{ is }\mathfrak{K}\text{-free}}%
}=\sum_{\substack{F\subseteq E;\\F\text{ contains no broken}\\\text{circuit of
}G\text{ as a subset}}}$ (an equality between summation signs). Now, Corollary
\ref{cor.chromsym.K-free} yields%
\[
X_{G}=\underbrace{\sum_{\substack{F\subseteq E;\\F\text{ is }\mathfrak{K}%
\text{-free}}}}_{=\sum_{\substack{F\subseteq E;\\F\text{ contains no
broken}\\\text{circuit of }G\text{ as a subset}}}}\left(  -1\right)
^{\left\vert F\right\vert }p_{\lambda\left(  V,F\right)  }=\sum
_{\substack{F\subseteq E;\\F\text{ contains no broken}\\\text{circuit of
}G\text{ as a subset}}}\left(  -1\right)  ^{\left\vert F\right\vert
}p_{\lambda\left(  V,F\right)  }.
\]
This proves Corollary \ref{cor.chromsym.NBC}.
\end{proof}
\end{verlong}

\begin{vershort}
\begin{proof}
[Proof of Theorem \ref{thm.chromsym.empty} using Theorem
\ref{thm.chromsym.varis}.]Let $X$ be the totally ordered set $\left\{
1\right\}  $, and let $\ell:E\rightarrow X$ be the only possible map. Let
$\mathfrak{K}$ be the empty set. Clearly, $\mathfrak{K}$ is a set of broken
circuits of $G$. For every $F\subseteq E$, the product $\prod_{\substack{K\in
\mathfrak{K};\\K\subseteq F}}0$ is empty (since $K$ is the empty set), and
thus equals $1$. Now, Theorem \ref{thm.chromsym.varis} (applied to $0$ instead
of $a_{K}$) yields%
\[
X_{G}=\sum_{F\subseteq E}\left(  -1\right)  ^{\left\vert F\right\vert
}\underbrace{\left(  \prod_{\substack{K\in\mathfrak{K};\\K\subseteq
F}}0\right)  }_{=1}p_{\lambda\left(  V,F\right)  }=\sum_{F\subseteq E}\left(
-1\right)  ^{\left\vert F\right\vert }p_{\lambda\left(  V,F\right)  }.
\]
This proves Theorem \ref{thm.chromsym.empty}.
\end{proof}
\end{vershort}

\begin{verlong}
\begin{proof}
[Proof of Theorem \ref{thm.chromsym.empty} using Theorem
\ref{thm.chromsym.varis}.]Let $X$ be the totally ordered set $\left\{
1\right\}  $ (equipped with the only possible order on this set). Let
$\ell:E\rightarrow X$ be the function sending each $e\in E$ to $1\in X$. Let
$\mathfrak{K}$ be the empty set. Clearly, $\mathfrak{K}$ is a set of broken
circuits of $G$. Theorem \ref{thm.chromsym.varis} (applied to $0$ instead of
$a_{K}$) yields%
\begin{align*}
X_{G}  &  =\sum_{F\subseteq E}\left(  -1\right)  ^{\left\vert F\right\vert
}\underbrace{\left(  \prod_{\substack{K\in\mathfrak{K};\\K\subseteq
F}}0\right)  }_{\substack{=\left(  \text{empty product}\right)  \\\text{(since
}\mathfrak{K}\text{ is the empty set)}}}p_{\lambda\left(  V,F\right)  }\\
&  =\sum_{F\subseteq E}\left(  -1\right)  ^{\left\vert F\right\vert
}\underbrace{\left(  \text{empty product}\right)  }_{=1}p_{\lambda\left(
V,F\right)  }=\sum_{F\subseteq E}\left(  -1\right)  ^{\left\vert F\right\vert
}p_{\lambda\left(  V,F\right)  }.
\end{align*}
This proves Theorem \ref{thm.chromsym.empty}.
\end{proof}
\end{verlong}

\section{Proof of Theorem \ref{thm.chromsym.varis}}

We shall now prepare for the proof of Theorem \ref{thm.chromsym.varis} with
some notations and some lemmas. Our proof will imitate \cite[proof of
Whitney's theorem]{BlaSag86}.

\subsection{$\operatorname*{Eqs}f$ and basic lemmas}

\begin{definition}
\label{def.Eqs}Let $V$ and $X$ be two sets. Let $f:V\rightarrow X$ be a map.
We let $\operatorname*{Eqs}f$ denote the subset%
\[
\left\{  \left\{  s,t\right\}  \ \mid\ \left(  s,t\right)  \in V^{2},\ s\neq
t\text{ and }f\left(  s\right)  =f\left(  t\right)  \right\}
\]
of $\dbinom{V}{2}$. (This is well-defined, because any two elements $s$ and
$t$ of $V$ satisfying $s\neq t$ clearly satisfy $\left\{  s,t\right\}
\in\dbinom{V}{2}$.)
\end{definition}

We shall now state some first properties of this notion:

\begin{lemma}
\label{lem.Eqs.proper}Let $G=\left(  V,E\right)  $ be a graph. Let $X$ be a
set. Let $f:V\rightarrow X$ be a map. Then, the $X$-coloring $f$ of $G$ is
proper if and only if $E\cap\operatorname*{Eqs}f=\varnothing$.
\end{lemma}

\begin{vershort}
\begin{proof}
[Proof of Lemma \ref{lem.Eqs.proper}.]The set $E\cap\operatorname*{Eqs}f$ is
precisely the set of edges $\left\{  s,t\right\}  $ of $G$ satisfying
$f\left(  s\right)  =f\left(  t\right)  $; meanwhile, the $X$-coloring $f$ is
called proper if and only if no such edges exist. Thus, Lemma
\ref{lem.Eqs.proper} becomes obvious.
\end{proof}
\end{vershort}

\begin{verlong}
\begin{proof}
[Proof of Lemma \ref{lem.Eqs.proper}.]The definition of $\operatorname*{Eqs}f$
shows that%
\begin{align}
\operatorname*{Eqs}f  &  =\left\{  \left\{  s,t\right\}  \ \mid\ \left(
s,t\right)  \in V^{2},\ s\neq t\text{ and }f\left(  s\right)  =f\left(
t\right)  \right\} \nonumber\\
&  =\left\{  \left\{  x,y\right\}  \ \mid\ \left(  x,y\right)  \in
V^{2},\ x\neq y\text{ and }f\left(  x\right)  =f\left(  y\right)  \right\}
\label{pf.lem.Eqs.proper.0}%
\end{align}
(here, we renamed the index $\left(  s,t\right)  $ as $\left(  x,y\right)  $).

We shall first prove the logical implication%
\begin{equation}
\left(  \text{the }X\text{-coloring }f\text{ of }G\text{ is proper}\right)
\ \Longrightarrow\ \left(  E\cap\operatorname*{Eqs}f=\varnothing\right)  .
\label{pf.lem.Eqs.proper.im1}%
\end{equation}


\textit{Proof of (\ref{pf.lem.Eqs.proper.im1}):} Assume that the $X$-coloring
$f$ of $G$ is proper. We must show that $E\cap\operatorname*{Eqs}%
f=\varnothing$.

Recall that the $X$-coloring $f$ of $G$ is proper if and only if every edge
$\left\{  s,t\right\}  \in E$ satisfies $f\left(  s\right)  \neq f\left(
t\right)  $ (by the definition of \textquotedblleft proper\textquotedblright).
Thus,%
\begin{equation}
\text{every edge }\left\{  s,t\right\}  \in E\text{ satisfies }f\left(
s\right)  \neq f\left(  t\right)  . \label{pf.lem.Eqs.proper.im1.pf.1}%
\end{equation}


Now, let $e\in E\cap\operatorname*{Eqs}f$. Thus,%
\[
e\in E\cap\operatorname*{Eqs}f\subseteq\operatorname*{Eqs}f=\left\{  \left\{
s,t\right\}  \ \mid\ \left(  s,t\right)  \in V^{2},\ s\neq t\text{ and
}f\left(  s\right)  =f\left(  t\right)  \right\}
\]
(by the definition of $\operatorname*{Eqs}f$). In other words, $e=\left\{
s,t\right\}  $ for some $\left(  s,t\right)  \in V^{2}$ satisfying $s\neq t$
and $f\left(  s\right)  =f\left(  t\right)  $. Consider this $\left(
s,t\right)  $. We have $\left\{  s,t\right\}  =e\in E\cap\operatorname*{Eqs}%
f\subseteq E$ and therefore $f\left(  s\right)  \neq f\left(  t\right)  $ (by
(\ref{pf.lem.Eqs.proper.im1.pf.1})). This contradicts $f\left(  s\right)
=f\left(  t\right)  $.

Now, let us forget that we fixed $e$. We thus have found a contradiction for
every $e\in E\cap\operatorname*{Eqs}f$. Therefore, no $e\in E\cap
\operatorname*{Eqs}f$ exists. In other words, the set $E\cap
\operatorname*{Eqs}f$ is empty. In other words, $E\cap\operatorname*{Eqs}%
f=\varnothing$. Thus, the implication (\ref{pf.lem.Eqs.proper.im1}) is proven.

Now, we shall prove the implication%
\begin{equation}
\left(  E\cap\operatorname*{Eqs}f=\varnothing\right)  \ \Longrightarrow
\ \left(  \text{the }X\text{-coloring }f\text{ of }G\text{ is proper}\right)
. \label{pf.lem.Eqs.proper.im2}%
\end{equation}


\textit{Proof of (\ref{pf.lem.Eqs.proper.im2}):} Assume that $E\cap
\operatorname*{Eqs}f=\varnothing$. We have to show that the $X$-coloring $f$
of $G$ is proper.

Let $\left\{  s,t\right\}  \in E$ be an edge. We shall now show that $f\left(
s\right)  \neq f\left(  t\right)  $.

Indeed, assume the contrary. Thus, $f\left(  s\right)  =f\left(  t\right)  $.
Now, $\left\{  s,t\right\}  \in E\subseteq\dbinom{V}{2}$. In other words,
$\left\{  s,t\right\}  $ is a $2$-element subset of $V$ (since $\dbinom{V}{2}$
is the set of all $2$-element subsets of $V$). Thus, $\left\vert \left\{
s,t\right\}  \right\vert =2$, so that $s\neq t$. Also, $\left\{  s,t\right\}
$ is a subset of $V$; thus, $s\in V$ and $t\in V$. Hence, $\left(  s,t\right)
\in V^{2}$. So we know that $\left\{  s,t\right\}  $ has the form $\left\{
x,y\right\}  $ for some $\left(  x,y\right)  \in V^{2}$ satisfying $x\neq y$
and $f\left(  x\right)  =f\left(  y\right)  $ (namely, this $\left(
x,y\right)  $ is $\left(  s,t\right)  $). In other words,%
\[
\left\{  s,t\right\}  \in\left\{  \left\{  x,y\right\}  \ \mid\ \left(
x,y\right)  \in V^{2},\ x\neq y\text{ and }f\left(  x\right)  =f\left(
y\right)  \right\}  =\operatorname*{Eqs}f
\]
(by (\ref{pf.lem.Eqs.proper.0})). Combining this with $\left\{  s,t\right\}
\in E$, we obtain $\left\{  s,t\right\}  \in E\cap\operatorname*{Eqs}%
f=\varnothing$. Thus, the set $\varnothing$ has an element (namely, $\left\{
s,t\right\}  $). This contradicts the fact that the set $\varnothing$ is
empty. Thus, we have obtained a contradiction. This shows that our assumption
was wrong. Hence, $f\left(  s\right)  \neq f\left(  t\right)  $ is proven.

Let us now forget that we fixed $\left\{  s,t\right\}  $. We thus have shown
that every edge $\left\{  s,t\right\}  \in E$ satisfies $f\left(  s\right)
\neq f\left(  t\right)  $. Therefore, the $X$-coloring $f$ of $G$ is proper
(since the $X$-coloring $f$ of $G$ is proper if and only if every edge
$\left\{  s,t\right\}  \in E$ satisfies $f\left(  s\right)  \neq f\left(
t\right)  $ (by the definition of \textquotedblleft proper\textquotedblright%
)). This proves the implication (\ref{pf.lem.Eqs.proper.im2}).

Now we have proven the two implications (\ref{pf.lem.Eqs.proper.im1}) and
(\ref{pf.lem.Eqs.proper.im2}). Combining these two implications, we obtain the
equivalence%
\[
\left(  \text{the }X\text{-coloring }f\text{ of }G\text{ is proper}\right)
\ \Longleftrightarrow\ \left(  E\cap\operatorname*{Eqs}f=\varnothing\right)
.
\]
This proves Lemma \ref{lem.Eqs.proper}.
\end{proof}
\end{verlong}

\begin{lemma}
\label{lem.Eqs.circuit}Let $G=\left(  V,E\right)  $ be a graph. Let $X$ be a
set. Let $f:V\rightarrow X$ be a map. Let $C$ be a circuit of $G$. Let $e\in
C$ be such that $C\setminus\left\{  e\right\}  \subseteq\operatorname*{Eqs}f$.
Then, $e\in E\cap\operatorname*{Eqs}f$.
\end{lemma}

\begin{vershort}
\begin{proof}
[Proof of Lemma \ref{lem.Eqs.circuit}.]The set $C$ is a circuit of $G$. Hence,
we can write $C$ in the form%
\[
C=\left\{  \left\{  v_{1},v_{2}\right\}  ,\left\{  v_{2},v_{3}\right\}
,\ldots,\left\{  v_{m},v_{m+1}\right\}  \right\}
\]
for some cycle $\left(  v_{1},v_{2},\ldots,v_{m+1}\right)  $ of $G$. Consider
this cycle $\left(  v_{1},v_{2},\ldots,v_{m+1}\right)  $. According to the
definition of a \textquotedblleft cycle\textquotedblright, the cycle $\left(
v_{1},v_{2},\ldots,v_{m+1}\right)  $ is a list of elements of $V$ having the
following properties:

\begin{itemize}
\item We have $m>1$.

\item We have $v_{m+1}=v_{1}$.

\item The vertices $v_{1},v_{2},\ldots,v_{m}$ are pairwise distinct.

\item We have $\left\{  v_{i},v_{i+1}\right\}  \in E$ for every $i\in\left\{
1,2,\ldots,m\right\}  $.
\end{itemize}

Recall that $e\in C$. We can thus WLOG assume that $e=\left\{  v_{m}%
,v_{m+1}\right\}  $ (since otherwise, we can simply relabel the vertices along
the cycle $\left(  v_{1},v_{2},\ldots,v_{m+1}\right)  $). Assume this. Since
$\left\{  v_{m},v_{m+1}\right\}  =e$, we have%
\[
C\setminus\left\{  e\right\}  =\left\{  \left\{  v_{1},v_{2}\right\}
,\left\{  v_{2},v_{3}\right\}  ,\ldots,\left\{  v_{m-1},v_{m}\right\}
\right\}
\]
(since $v_{1},v_{2},\ldots,v_{m}$ are distinct, and since $m>1$ and
$v_{m+1}=v_{1}$). For every $i\in\left\{  1,2,\ldots,m-1\right\}  $, we have
$f\left(  v_{i}\right)  =f\left(  v_{i+1}\right)  $ (since
\[
\left\{  v_{i},v_{i+1}\right\}  \subseteq\left\{  \left\{  v_{1}%
,v_{2}\right\}  ,\left\{  v_{2},v_{3}\right\}  ,\ldots,\left\{  v_{m-1}%
,v_{m}\right\}  \right\}  =C\setminus\left\{  e\right\}  \subseteq
\operatorname*{Eqs}f
\]
). Hence, $f\left(  v_{1}\right)  =f\left(  v_{2}\right)  =\cdots=f\left(
v_{m}\right)  $, so that $f\left(  v_{m}\right)  =f\left(  \underbrace{v_{1}%
}_{=v_{m+1}}\right)  =f\left(  v_{m+1}\right)  $. Thus, $\left\{
v_{m},v_{m+1}\right\}  \in\operatorname*{Eqs}f$. Thus, $e=\left\{
v_{m},v_{m+1}\right\}  \in\operatorname*{Eqs}f$. Combined with $e\in E$, this
yields $e\in E\cap\operatorname*{Eqs}f$. This proves Lemma
\ref{lem.Eqs.circuit}.
\end{proof}
\end{vershort}

\begin{verlong}
\begin{proof}
[Proof of Lemma \ref{lem.Eqs.circuit}.]The set $C$ is a circuit of $G$. In
other words, the set $C$ has the form $\left\{  \left\{  v_{1},v_{2}\right\}
,\left\{  v_{2},v_{3}\right\}  ,\ldots,\left\{  v_{m},v_{m+1}\right\}
\right\}  $, where $\left(  v_{1},v_{2},\ldots,v_{m+1}\right)  $ is a cycle of
$G$ (by the definition of a \textquotedblleft circuit\textquotedblright).
Consider this cycle $\left(  v_{1},v_{2},\ldots,v_{m+1}\right)  $. We thus
have
\begin{equation}
C=\left\{  \left\{  v_{1},v_{2}\right\}  ,\left\{  v_{2},v_{3}\right\}
,\ldots,\left\{  v_{m},v_{m+1}\right\}  \right\}  .
\label{pf.lem.Eqs.circuit.1}%
\end{equation}


We have $e\in C=\left\{  \left\{  v_{1},v_{2}\right\}  ,\left\{  v_{2}%
,v_{3}\right\}  ,\ldots,\left\{  v_{m},v_{m+1}\right\}  \right\}  $. Thus,
$e=\left\{  v_{i},v_{i+1}\right\}  $ for some $i\in\left\{  1,2,\ldots
,m\right\}  $. Consider this $i$.

The list $\left(  v_{1},v_{2},\ldots,v_{m+1}\right)  $ is a cycle of $G$.
According to the definition of a \textquotedblleft cycle\textquotedblright,
this means that this list is a list of elements of $V$ satisfying the
following four properties:

\begin{itemize}
\item We have $m>1$.

\item We have $v_{m+1}=v_{1}$.

\item The vertices $v_{1},v_{2},\ldots,v_{m}$ are pairwise distinct.

\item We have $\left\{  v_{i},v_{i+1}\right\}  \in E$ for every $i\in\left\{
1,2,\ldots,m\right\}  $.
\end{itemize}

Thus, $\left(  v_{1},v_{2},\ldots,v_{m+1}\right)  $ is a list of elements of
$V$ satisfying the four properties that we have just mentioned. In particular,
$v_{m+1}=v_{1}$. Also,%
\begin{equation}
\left\{  v_{i},v_{i+1}\right\}  \in E\ \ \ \ \ \ \ \ \ \ \text{for every }%
i\in\left\{  1,2,\ldots,m\right\}  . \label{pf.lem.Eqs.circuit.inE}%
\end{equation}


Any two distinct elements $p$ and $q$ of $\left\{  1,2,\ldots,m\right\}  $
satisfy%
\begin{equation}
\left\{  v_{p},v_{p+1}\right\}  \neq\left\{  v_{q},v_{q+1}\right\}
\label{pf.lem.Eqs.circuit.distinct-edges}%
\end{equation}
\footnote{\textit{Proof of (\ref{pf.lem.Eqs.circuit.distinct-edges}):} Let $p$
and $q$ be two distinct elements of $\left\{  1,2,\ldots,m\right\}  $. We must
prove (\ref{pf.lem.Eqs.circuit.distinct-edges}).
\par
Assume the contrary. Thus, $\left\{  v_{p},v_{p+1}\right\}  =\left\{
v_{q},v_{q+1}\right\}  $.
\par
The vertices $v_{1},v_{2},\ldots,v_{m}$ are pairwise distinct. In other words,
any two distinct elements $a$ and $b$ of $\left\{  1,2,\ldots,m\right\}  $
satisfy $v_{a}\neq v_{b}$. Applying this to $a=p$ and $b=q$, we obtain
$v_{p}\neq v_{q}$. Combining $v_{p}\in\left\{  v_{p},v_{p+1}\right\}
=\left\{  v_{q},v_{q+1}\right\}  $ with $v_{p}\neq v_{q}$, we obtain $v_{p}%
\in\left\{  v_{q},v_{q+1}\right\}  \setminus\left\{  v_{q}\right\}
\subseteq\left\{  v_{q+1}\right\}  $. Thus, $v_{p}=v_{q+1}$. The same argument
(with $p$ and $q$ replaced by $q$ and $p$) yields $v_{q}=v_{p+1}$.
\par
We have $p\neq q$ (since $p$ and $q$ are distinct). Thus, we can WLOG assume
that $p<q$ (since otherwise, we can simply switch $p$ with $q$). Assume this.
From $q\in\left\{  1,2,\ldots,m\right\}  $, we obtain $q\leq m$, so that
$p<q\leq m$. Since $p$ and $m$ are integers, this shows that $p\leq m-1$.
Thus, $p+1\leq m$. Hence, $p+1\in\left\{  1,2,\ldots,m\right\}  $.
\par
Recall again that any two distinct elements $a$ and $b$ of $\left\{
1,2,\ldots,m\right\}  $ satisfy $v_{a}\neq v_{b}$. Applying this to $a=p+1$
and $b=q$, we obtain $v_{p+1}\neq v_{q}$ (since $p+1\in\left\{  1,2,\ldots
,m\right\}  $). Hence, $v_{p+1}\neq v_{q}=v_{p+1}$. This is absurd.
\par
We thus have found a contradiction. This contradiction proves that our
assumption was wrong. Hence, (\ref{pf.lem.Eqs.circuit.distinct-edges}) is
proven.}.

Now, we have%
\begin{equation}
f\left(  v_{j}\right)  =f\left(  v_{j+1}\right)  \ \ \ \ \ \ \ \ \ \ \text{for
every }j\in\left\{  1,2,\ldots,m\right\}  \setminus\left\{  i\right\}
\label{pf.lem.Eqs.circuit.2}%
\end{equation}
\footnote{\textit{Proof of (\ref{pf.lem.Eqs.circuit.2}):} Let $j\in\left\{
1,2,\ldots,m\right\}  \setminus\left\{  i\right\}  $. Thus, $j\in\left\{
1,2,\ldots,m\right\}  $ and $j\notin\left\{  i\right\}  $. From $j\notin%
\left\{  i\right\}  $, we obtain $j\neq i$. Therefore, the two elements $j$
and $i$ of $\left\{  1,2,\ldots,m\right\}  $ are distinct. Thus,
(\ref{pf.lem.Eqs.circuit.distinct-edges}) (applied to $p=j$ and $q=i$) shows
that $\left\{  v_{j},v_{j+1}\right\}  \neq\left\{  v_{i},v_{i+1}\right\}  =e$.
But from $j\in\left\{  1,2,\ldots,m\right\}  $, we obtain%
\begin{align*}
\left\{  v_{j},v_{j+1}\right\}   &  \in\left\{  \left\{  v_{k},v_{k+1}%
\right\}  \ \mid\ k\in\left\{  1,2,\ldots,m\right\}  \right\} \\
&  =\left\{  \left\{  v_{1},v_{2}\right\}  ,\left\{  v_{2},v_{3}\right\}
,\ldots,\left\{  v_{m},v_{m+1}\right\}  \right\}  =C
\end{align*}
(by (\ref{pf.lem.Eqs.circuit.1})). Combining this with $\left\{  v_{j}%
,v_{j+1}\right\}  \neq e$, we obtain%
\begin{align*}
\left\{  v_{j},v_{j+1}\right\}   &  \in C\setminus\left\{  e\right\}
\subseteq\operatorname*{Eqs}f\\
&  =\left\{  \left\{  s,t\right\}  \ \mid\ \left(  s,t\right)  \in
V^{2},\ s\neq t\text{ and }f\left(  s\right)  =f\left(  t\right)  \right\}  .
\end{align*}
In other words, $\left\{  v_{j},v_{j+1}\right\}  $ has the form $\left\{
s,t\right\}  $ for some $\left(  s,t\right)  \in V^{2}$ satisfying $s\neq t$
and $f\left(  s\right)  =f\left(  t\right)  $. Consider this $\left(
s,t\right)  $. Thus, $\left\{  v_{j},v_{j+1}\right\}  =\left\{  s,t\right\}
$.
\par
We have $f\left(  s\right)  =f\left(  t\right)  $. Therefore, set $g=f\left(
s\right)  =f\left(  t\right)  $. We have%
\[
f\left(  \left\{  s,t\right\}  \right)  =\left\{  \underbrace{f\left(
s\right)  }_{=g},\underbrace{f\left(  t\right)  }_{=g}\right\}  =\left\{
g,g\right\}  =\left\{  g\right\}  .
\]
Now, $v_{j}\in\left\{  v_{j},v_{j+1}\right\}  =\left\{  s,t\right\}  $, and
thus $f\left(  \underbrace{v_{j}}_{\in\left\{  s,t\right\}  }\right)  \in
f\left(  \left\{  s,t\right\}  \right)  =\left\{  g\right\}  $. In other
words, $f\left(  v_{j}\right)  =g$. Also, $v_{j+1}\in\left\{  v_{j}%
,v_{j+1}\right\}  =\left\{  s,t\right\}  $, and thus $f\left(
\underbrace{v_{j+1}}_{\in\left\{  s,t\right\}  }\right)  \in f\left(  \left\{
s,t\right\}  \right)  =\left\{  g\right\}  $. In other words, $f\left(
v_{j+1}\right)  =g$. Comparing this with $f\left(  v_{j}\right)  =g$, we
obtain $f\left(  v_{j}\right)  =f\left(  v_{j+1}\right)  $. This proves
(\ref{pf.lem.Eqs.circuit.2}).}. Hence,%
\begin{equation}
f\left(  v_{1}\right)  =f\left(  v_{i}\right)  \label{pf.lem.Eqs.circuit.3a}%
\end{equation}
\footnote{\textit{Proof of (\ref{pf.lem.Eqs.circuit.3a}):} Let $j\in\left\{
1,2,\ldots,i-1\right\}  $. Thus, $j\in\left\{  1,2,\ldots,i-1\right\}
\subseteq\left\{  1,2,\ldots,m\right\}  $. Combining this with $j\neq i$
(since $j<i$ (since $j\in\left\{  1,2,\ldots,i-1\right\}  $)), we obtain
$j\in\left\{  1,2,\ldots,m\right\}  \setminus\left\{  i\right\}  $. Hence,
$f\left(  v_{j}\right)  =f\left(  v_{j+1}\right)  $ (by
(\ref{pf.lem.Eqs.circuit.2})).
\par
Now, let us forget that we fixed $j$. We thus have proven that $f\left(
v_{j}\right)  =f\left(  v_{j+1}\right)  $ for every $j\in\left\{
1,2,\ldots,i-1\right\}  $. In other words, $f\left(  v_{1}\right)  =f\left(
v_{2}\right)  =\cdots=f\left(  v_{i}\right)  $. Hence, $f\left(  v_{1}\right)
=f\left(  v_{i}\right)  $. This proves (\ref{pf.lem.Eqs.circuit.3a}).}. Also,%
\begin{equation}
f\left(  v_{i+1}\right)  =f\left(  v_{m+1}\right)
\label{pf.lem.Eqs.circuit.3b}%
\end{equation}
\footnote{\textit{Proof of (\ref{pf.lem.Eqs.circuit.3b}):} Let $j\in\left\{
i+1,i+2,\ldots,m\right\}  $. Thus, $j\in\left\{  i+1,i+2,\ldots,m\right\}
\subseteq\left\{  1,2,\ldots,m\right\}  $. Combining this with $j\neq i$
(since $j>i$ (since $j\in\left\{  i+1,i+2,\ldots,m\right\}  $)), we obtain
$j\in\left\{  1,2,\ldots,m\right\}  \setminus\left\{  i\right\}  $. Hence,
$f\left(  v_{j}\right)  =f\left(  v_{j+1}\right)  $ (by
(\ref{pf.lem.Eqs.circuit.2})).
\par
Now, let us forget that we fixed $j$. We thus have proven that $f\left(
v_{j}\right)  =f\left(  v_{j+1}\right)  $ for every $j\in\left\{
i+1,i+2,\ldots,m\right\}  $. In other words, $f\left(  v_{i+1}\right)
=f\left(  v_{i+2}\right)  =\cdots=f\left(  v_{m+1}\right)  $. Hence, $f\left(
v_{i}\right)  =f\left(  v_{m+1}\right)  $. This proves
(\ref{pf.lem.Eqs.circuit.3b}).}. Now, (\ref{pf.lem.Eqs.circuit.3a}) yields%
\[
f\left(  v_{i}\right)  =f\left(  \underbrace{v_{1}}_{=v_{m+1}}\right)
=f\left(  v_{m+1}\right)  =f\left(  v_{i+1}\right)  .
\]


Moreover, $v_{i}$ and $v_{i+1}$ are elements of $V$ (since $\left(
v_{1},v_{2},\ldots,v_{m+1}\right)  $ is a list of elements of $V$). In other
words, $v_{i}\in V$ and $v_{i+1}\in V$. Hence, $\left(  v_{i},v_{i+1}\right)
\in V^{2}$.

Furthermore, $v_{i}\neq v_{i+1}$\ \ \ \ \footnote{\textit{Proof.} Assume the
contrary. Thus, $v_{i}=v_{i+1}$.
\par
Let us first assume (for the sake of contradiction) that $i=m$. Thus,
$v_{i}=v_{m}$. Also, from $i=m$, we obtain $v_{i+1}=v_{m+1}=v_{1}$. Hence,
$v_{m}=v_{i}=v_{i+1}=v_{1}$.
\par
The vertices $v_{1},v_{2},\ldots,v_{m}$ are pairwise distinct. In other words,
any two distinct elements $a$ and $b$ of $\left\{  1,2,\ldots,m\right\}  $
satisfy $v_{a}\neq v_{b}$. Applying this to $a=m$ and $b=1$, we obtain
$v_{m}\neq v_{1}$ (since $m$ and $1$ are distinct (since $m>1$)). This
contradicts $v_{m}=v_{1}$.
\par
This contradiction proves that our assumption (that $i=m$) was wrong. Hence,
we cannot have $i=m$. We thus have $i\neq m$. Combined with $i\in\left\{
1,2,\ldots,m\right\}  $, this yields $i\in\left\{  1,2,\ldots,m\right\}
\setminus\left\{  m\right\}  \subseteq\left\{  1,2,\ldots,m-1\right\}  $.
Thus, $i+1\in\left\{  2,3,\ldots,m\right\}  \subseteq\left\{  1,2,\ldots
,m\right\}  $.
\par
Now, recall that any two distinct elements $a$ and $b$ of $\left\{
1,2,\ldots,m\right\}  $ satisfy $v_{a}\neq v_{b}$. We can apply this to $a=i$
and $b=i+1$ (since $i+1\in\left\{  1,2,\ldots,m\right\}  $). Thus, we obtain
$v_{i}\neq v_{i+1}$. This contradicts $v_{i}=v_{i+1}$. This contradiction
shows that our assumption was wrong. Qed.}.

Now, the definition of $\operatorname*{Eqs}f$ shows that%
\begin{equation}
\operatorname*{Eqs}f=\left\{  \left\{  s,t\right\}  \ \mid\ \left(
s,t\right)  \in V^{2},\ s\neq t\text{ and }f\left(  s\right)  =f\left(
t\right)  \right\}  . \label{pf.lem.Eqs.circuit.Eqs-def}%
\end{equation}


But we have $\left(  v_{i},v_{i+1}\right)  \in V^{2}$, $v_{i}\neq v_{i+1}$ and
$f\left(  v_{i}\right)  =f\left(  v_{i+1}\right)  $. Hence, $\left\{
v_{i},v_{i+1}\right\}  $ has the form $\left\{  s,t\right\}  $ for some
$\left(  s,t\right)  \in V^{2}$ satisfying $s\neq t$ and $f\left(  s\right)
=f\left(  t\right)  $ (namely, for $\left(  s,t\right)  =\left(  v_{i}%
,v_{i+1}\right)  $). Thus,%
\[
\left\{  v_{i},v_{i+1}\right\}  \in\left\{  \left\{  s,t\right\}
\ \mid\ \left(  s,t\right)  \in V^{2},\ s\neq t\text{ and }f\left(  s\right)
=f\left(  t\right)  \right\}  =\operatorname*{Eqs}f
\]
(by (\ref{pf.lem.Eqs.circuit.Eqs-def})). Thus, $e=\left\{  v_{i}%
,v_{i+1}\right\}  \in\operatorname*{Eqs}f$.

But $e=\left\{  v_{i},v_{i+1}\right\}  \in E$ (by
(\ref{pf.lem.Eqs.circuit.inE})). Combining this with $e\in\operatorname*{Eqs}%
f$, we obtain $e\in E\cap\operatorname*{Eqs}f$. This proves Lemma
\ref{lem.Eqs.circuit}.
\end{proof}
\end{verlong}

\begin{verlong}
\begin{lemma}
\label{lem.Eqs.sum-aux}Let $\left(  V,B\right)  $ be a finite graph. Let
$\sim$ denote the equivalence relation $\sim_{\left(  V,B\right)  }$ (defined
as in Definition \ref{def.connectedness} \textbf{(a)}).

Let $Y$ be a set. A set $Y_{\sim}^{V}$ is defined (according to Definition
\ref{def.relquot.maps} \textbf{(b)}). Let $f:V\rightarrow Y$ be any map. Then,
we have the following logical equivalence of statements:%
\[
\left(  B\subseteq\operatorname*{Eqs}f\right)  \ \Longleftrightarrow\ \left(
f\in Y_{\sim}^{V}\right)  .
\]

\end{lemma}

\begin{proof}
[Proof of Lemma \ref{lem.Eqs.sum-aux}.]We have%
\begin{equation}
Y_{\sim}^{V}=\left\{  g\in Y^{V}\ \mid\ g\left(  x\right)  =g\left(  y\right)
\text{ for any }x\in V\text{ and }y\in V\text{ satisfying }x\sim y\right\}
\label{pf.lem.Eqs.sum.defN+Vsim}%
\end{equation}
(by the definition of $Y_{\sim}^{V}$).

The definition of $\operatorname*{Eqs}f$ shows that%
\begin{align}
\operatorname*{Eqs}f  &  =\left\{  \left\{  s,t\right\}  \ \mid\ \left(
s,t\right)  \in V^{2},\ s\neq t\text{ and }f\left(  s\right)  =f\left(
t\right)  \right\} \nonumber\\
&  =\left\{  \left\{  x,y\right\}  \ \mid\ \left(  x,y\right)  \in
V^{2},\ x\neq y\text{ and }f\left(  x\right)  =f\left(  y\right)  \right\}
\label{pf.lem.Eqs.sum.defEqs}%
\end{align}
(here, we renamed the index $\left(  s,t\right)  $ as $\left(  x,y\right)  $).

We shall now show the following logical implication:%
\begin{equation}
\left(  B\subseteq\operatorname*{Eqs}f\right)  \ \Longrightarrow\ \left(  f\in
Y_{\sim}^{V}\right)  \label{pf.lem.Eqs.sum.im1}%
\end{equation}


\textit{Proof of (\ref{pf.lem.Eqs.sum.im1}):} Assume that $B\subseteq
\operatorname*{Eqs}f$. We must show that $f\in Y_{\sim}^{V}$.

Let $x\in V$ and $y\in V$ be such that $x\sim y$. We have $x\sim y$. In other
words, $x\sim_{\left(  V,B\right)  }y$ (since $\sim$ is the equivalence
relation $\sim_{\left(  V,B\right)  }$). In other words, $x$ and $y$ are
connected in the graph $\left(  V,B\right)  $ (since $x\sim_{\left(
V,B\right)  }y$ holds if and only if $x$ and $y$ are connected in the graph
$\left(  V,B\right)  $ (by the definition of the relation $\sim_{\left(
V,B\right)  }$)). In other words, there exists a walk from $x$ to $y$ in
$\left(  V,B\right)  $ (since $x$ and $y$ are connected in $\left(
V,B\right)  $ if and only if there exists a walk from $x$ to $y$ in $\left(
V,B\right)  $ (by the definition of \textquotedblleft
connected\textquotedblright)). Let $\mathfrak{w}$ be this walk. Thus,
$\mathfrak{w}$ is a walk from $x$ to $y$ in $\left(  V,B\right)  $. In other
words, $\mathfrak{w}$ is a sequence $\left(  w_{0},w_{1},\ldots,w_{k}\right)
$ of elements of $V$ such that $w_{0}=x$ and $w_{k}=y$ and%
\begin{equation}
\left(  \left\{  w_{i},w_{i+1}\right\}  \in B\ \ \ \ \ \ \ \ \ \ \text{for
every }i\in\left\{  0,1,\ldots,k-1\right\}  \right)
\label{pf.lem.Eqs.sum.im1.pf.1}%
\end{equation}
(since a walk from $x$ to $y$ in $\left(  V,B\right)  $ is the same as a
sequence $\left(  w_{0},w_{1},\ldots,w_{k}\right)  $ of elements of $V$ such
that $w_{0}=x$ and $w_{k}=y$ and $\left(  \left\{  w_{i},w_{i+1}\right\}  \in
B\ \ \ \ \ \ \ \ \ \ \text{for every }i\in\left\{  0,1,\ldots,k-1\right\}
\right)  $ (by the definition of a \textquotedblleft walk\textquotedblright)).
Consider this sequence $\left(  w_{0},w_{1},\ldots,w_{k}\right)  $.

For every $i\in\left\{  0,1,\ldots,k-1\right\}  $, we have $f\left(
w_{i}\right)  =f\left(  w_{i+1}\right)  $\ \ \ \ \footnote{\textit{Proof.} Let
$i\in\left\{  0,1,\ldots,k-1\right\}  $. Thus, (\ref{pf.lem.Eqs.sum.im1.pf.1})
shows that
\[
\left\{  w_{i},w_{i+1}\right\}  \in B\subseteq\operatorname*{Eqs}f=\left\{
\left\{  s,t\right\}  \ \mid\ \left(  s,t\right)  \in V^{2},\ s\neq t\text{
and }f\left(  s\right)  =f\left(  t\right)  \right\}
\]
(by the definition of $\operatorname*{Eqs}f$). In other words, $\left\{
w_{i},w_{i+1}\right\}  =\left\{  s,t\right\}  $ for some $\left(  s,t\right)
\in V^{2}$ satisfying $s\neq t$ and $f\left(  s\right)  =f\left(  t\right)  $.
Consider this $\left(  s,t\right)  $.
\par
We have $f\left(  s\right)  =f\left(  t\right)  $. Therefore, set $g=f\left(
s\right)  =f\left(  t\right)  $. Then, $f\left(  \left\{  s,t\right\}
\right)  =\left\{  \underbrace{f\left(  s\right)  }_{=g},\underbrace{f\left(
t\right)  }_{=g}\right\}  =\left\{  g,g\right\}  =\left\{  g\right\}  $. Now,
$f\left(  \underbrace{w_{i}}_{\in\left\{  w_{i},w_{i+1}\right\}  =\left\{
s,t\right\}  }\right)  \in f\left(  \left\{  s,t\right\}  \right)  =\left\{
g\right\}  $, so that $f\left(  w_{i}\right)  =g$. Also, $f\left(
\underbrace{w_{i+1}}_{\in\left\{  w_{i},w_{i+1}\right\}  =\left\{
s,t\right\}  }\right)  \in f\left(  \left\{  s,t\right\}  \right)  =\left\{
g\right\}  $, so that $f\left(  w_{i+1}\right)  =g$. Hence, $f\left(
w_{i}\right)  =g=f\left(  w_{i+1}\right)  $, qed.}. In other words, $f\left(
w_{0}\right)  =f\left(  w_{1}\right)  =\cdots=f\left(  w_{k}\right)  $. Thus,
$f\left(  w_{0}\right)  =f\left(  w_{k}\right)  $. This rewrites as $f\left(
x\right)  =f\left(  y\right)  $ (since $w_{0}=x$ and $w_{k}=y$).

Now, let us forget that we fixed $x$ and $y$. We thus have shown that
$f\left(  x\right)  =f\left(  y\right)  $ for any $x\in V$ and $y\in V$
satisfying $x\sim y$. Hence, $f$ is a $g\in Y^{V}$ which satisfies $g\left(
x\right)  =g\left(  y\right)  $ for any $x\in V$ and $y\in V$ satisfying
$x\sim y$. In other words,%
\begin{align*}
f  &  \in\left\{  g\in Y^{V}\ \mid\ g\left(  x\right)  =g\left(  y\right)
\text{ for any }x\in V\text{ and }y\in V\text{ satisfying }x\sim y\right\} \\
&  =Y_{\sim}^{V}\ \ \ \ \ \ \ \ \ \ \left(  \text{by
(\ref{pf.lem.Eqs.sum.defN+Vsim})}\right)  .
\end{align*}
Thus, the implication (\ref{pf.lem.Eqs.sum.im1}) is proven.

Next, we shall show the following logical implication:%
\begin{equation}
\left(  f\in Y_{\sim}^{V}\right)  \ \Longrightarrow\ \left(  B\subseteq
\operatorname*{Eqs}f\right)  . \label{pf.lem.Eqs.sum.im2}%
\end{equation}


\textit{Proof of (\ref{pf.lem.Eqs.sum.im2}):} Assume that $f\in Y_{\sim}^{V}$.
We must show that $B\subseteq\operatorname*{Eqs}f$.

Since $\left(  V,B\right)  $ is a graph, we must have $B\subseteq\dbinom{V}%
{2}$.

Let $e\in B$. Then, $e\in B\subseteq\dbinom{V}{2}$. In other words, $e$ is a
$2$-element subset of $V$ (since $\dbinom{V}{2}$ is the set of all $2$-element
subsets of $V$). In other words, $e=\left\{  s,t\right\}  $ for two distinct
elements $s$ and $t$ of $V$. Consider these $s$ and $t$.

We have $\left\{  s,t\right\}  =e\in B$. Thus, $s\sim t$%
\ \ \ \ \footnote{\textit{Proof.} We have $s\in V$ and $t\in V$. Thus,
$\left(  s,t\right)  $ is a sequence of elements of $V$. Let us denote this
sequence $\left(  s,t\right)  $ by $\left(  p_{0},p_{1},\ldots,p_{\ell
}\right)  $. Thus, $\ell=1$, $p_{0}=s$ and $p_{1}=t$.
\par
We have $\ell=1$ and thus $p_{\ell}=p_{1}=t$.
\par
Let $i\in\left\{  0,1,\ldots,\ell-1\right\}  $. Thus, $i\geq0$ and
$i\leq\underbrace{\ell}_{=1}-1=1-1=0$. Combining $i\leq0$ and $i\geq0$, we
obtain $i=0$, so that $p_{i}=p_{0}=s$ and $p_{i+1}=p_{0+1}=p_{1}=t$. Now,
$\left\{  \underbrace{p_{i}}_{=s},\underbrace{p_{i+1}}_{=t}\right\}  =\left\{
s,t\right\}  \in B$.
\par
Now, let us forget that we fixed $i$. We thus have shown that $\left\{
p_{i},p_{i+1}\right\}  \in B$ for every $i\in\left\{  0,1,\ldots
,\ell-1\right\}  $. Thus, the sequence $\left(  p_{0},p_{1},\ldots,p_{\ell
}\right)  $ is a sequence of elements of $V$ satisfying $p_{0}=s$, $p_{\ell
}=t$ and $\left(  \left\{  p_{i},p_{i+1}\right\}  \in
B\ \ \ \ \ \ \ \ \ \ \text{for every }i\in\left\{  0,1,\ldots,\ell-1\right\}
\right)  $. In other words, the sequence $\left(  p_{0},p_{1},\ldots,p_{\ell
}\right)  $ is a sequence $\left(  w_{0},w_{1},\ldots,w_{k}\right)  $ of
elements of $V$ such that $w_{0}=s$ and $w_{k}=t$ and $\left(  \left\{
w_{i},w_{i+1}\right\}  \in B\ \ \ \ \ \ \ \ \ \ \text{for every }i\in\left\{
0,1,\ldots,k-1\right\}  \right)  $. In other words, the sequence $\left(
p_{0},p_{1},\ldots,p_{\ell}\right)  $ is walk from $s$ to $t$ in $\left(
V,B\right)  $ (since a walk from $s$ to $t$ in $\left(  V,B\right)  $ is the
same as a sequence $\left(  w_{0},w_{1},\ldots,w_{k}\right)  $ of elements of
$V$ such that $w_{0}=s$ and $w_{k}=t$ and $\left(  \left\{  w_{i}%
,w_{i+1}\right\}  \in B\ \ \ \ \ \ \ \ \ \ \text{for every }i\in\left\{
0,1,\ldots,k-1\right\}  \right)  $ (by the definition of a \textquotedblleft
walk\textquotedblright)). Hence, there exists a walk from $s$ to $t$ in
$\left(  V,B\right)  $. In other words, $s$ and $t$ are connected in the graph
$\left(  V,B\right)  $ (since $s$ and $t$ are connected in $\left(
V,B\right)  $ if and only if there exists a walk from $s$ to $t$ in $\left(
V,B\right)  $ (by the definition of \textquotedblleft
connected\textquotedblright)). In other words, $s\sim_{\left(  V,B\right)  }t$
(since $s\sim_{\left(  V,B\right)  }t$ holds if and only if $s$ and $t$ are
connected in the graph $\left(  V,B\right)  $ (by the definition of the
relation $\sim_{\left(  V,B\right)  }$)). In other words, $s\sim t$ (since
$\sim$ is the equivalence relation $\sim_{\left(  V,B\right)  }$). Qed.}.

But $f\in Y_{\sim}^{V}=\left\{  g\in Y^{V}\ \mid\ g\left(  x\right)  =g\left(
y\right)  \text{ for any }x\in V\text{ and }y\in V\text{ satisfying }x\sim
y\right\}  $ (by (\ref{pf.lem.Eqs.sum.defN+Vsim})). In other words, $f$ is a
$g\in Y^{V}$ such that $g\left(  x\right)  =g\left(  y\right)  $ for any $x\in
V$ and $y\in V$ satisfying $x\sim y$. In other words, $f$ is an element of
$Y^{V}$ and satisfies%
\begin{equation}
f\left(  x\right)  =f\left(  y\right)  \ \ \ \ \ \ \ \ \ \ \text{for any }x\in
V\text{ and }y\in V\text{ satisfying }x\sim y. \label{pf.lem.Eqs.sum.im2.pf.3}%
\end{equation}


Applying (\ref{pf.lem.Eqs.sum.im2.pf.3}) to $x=s$ and $y=t$, we obtain
$f\left(  s\right)  =f\left(  t\right)  $ (since $s\sim t$).

We have $s\in V$ and $t\in V$. Thus, $\left(  s,t\right)  \in V^{2}$. Also,
$s$ and $t$ are distinct; thus, $s\neq t$. Hence, $\left(  s,t\right)  $ is an
element of $V^{2}$ satisfying $s\neq t$ and $f\left(  s\right)  =f\left(
t\right)  $. In other words, $\left(  s,t\right)  $ is an $\left(  x,y\right)
\in V^{2}$ satisfying $x\neq y$ and $f\left(  x\right)  =f\left(  y\right)  $.
Therefore, $\left\{  s,t\right\}  $ has the form $\left\{  x,y\right\}  $ for
some $\left(  x,y\right)  \in V^{2}$ satisfying $x\neq y$ and $f\left(
x\right)  =f\left(  y\right)  $ (namely, for $\left(  x,y\right)  =\left(
s,t\right)  $). In other words,%
\[
\left\{  s,t\right\}  \in\left\{  \left\{  x,y\right\}  \ \mid\ \left(
x,y\right)  \in V^{2},\ x\neq y\text{ and }f\left(  x\right)  =f\left(
y\right)  \right\}  =\operatorname*{Eqs}f
\]
(by (\ref{pf.lem.Eqs.sum.defEqs})). Thus, $e=\left\{  s,t\right\}
\in\operatorname*{Eqs}f$.

Now, let us forget that we fixed $e$. We thus have shown that $e\in
\operatorname*{Eqs}f$ for every $e\in B$. In other words, $B\subseteq
\operatorname*{Eqs}f$. This proves the implication (\ref{pf.lem.Eqs.sum.im2}).

Now, we can combine the two implications (\ref{pf.lem.Eqs.sum.im1}) and
(\ref{pf.lem.Eqs.sum.im2}). As a result, we obtain the equivalence $\left(
B\subseteq\operatorname*{Eqs}f\right)  \ \Longleftrightarrow\ \left(  f\in
Y_{\sim}^{V}\right)  $. This proves Lemma \ref{lem.Eqs.sum-aux}.
\end{proof}

The next lemma is a fundamental fact about counting:

\begin{lemma}
\label{lem.function-count}Let $W$ be a finite set. Let $\left(  C_{1}%
,C_{2},\ldots,C_{k}\right)  $ be a list of all elements of $W$ which contains
each of these elements exactly once. Let $Y$ be any set. Then, the map
\begin{align*}
Y^{W}  &  \rightarrow Y^{k},\\
f  &  \mapsto\left(  f\left(  C_{1}\right)  ,f\left(  C_{2}\right)
,\ldots,f\left(  C_{k}\right)  \right)
\end{align*}
is a bijection.
\end{lemma}

\begin{proof}
[Proof of Lemma \ref{lem.function-count}.]Let $\Phi$ denote the map%
\begin{align*}
Y^{W}  &  \rightarrow Y^{k},\\
f  &  \mapsto\left(  f\left(  C_{1}\right)  ,f\left(  C_{2}\right)
,\ldots,f\left(  C_{k}\right)  \right)  .
\end{align*}
We shall show that the map $\Phi$ is a bijection.

First, we notice that the map $\Phi$ is injective\footnote{\textit{Proof.} Let
$f$ and $g$ be two elements of $Y^{W}$ such that $\Phi\left(  f\right)
=\Phi\left(  g\right)  $. We shall show that $f=g$.
\par
Let $w\in W$. Recall that $\left(  C_{1},C_{2},\ldots,C_{k}\right)  $ is a
list of all elements of $W$. Thus, $W=\left\{  C_{1},C_{2},\ldots
,C_{k}\right\}  $. Now, $w\in W=\left\{  C_{1},C_{2},\ldots,C_{k}\right\}  $.
Hence, there exists some $i\in\left\{  1,2,\ldots,k\right\}  $ such that
$w=C_{i}$. Consider this $i$.
\par
Now, the definition of $\Phi\left(  f\right)  $ yields $\Phi\left(  f\right)
=\left(  f\left(  C_{1}\right)  ,f\left(  C_{2}\right)  ,\ldots,f\left(
C_{k}\right)  \right)  $. Hence,%
\begin{align*}
&  \left(  \text{the }i\text{-th entry of }\underbrace{\Phi\left(  f\right)
}_{=\left(  f\left(  C_{1}\right)  ,f\left(  C_{2}\right)  ,\ldots,f\left(
C_{k}\right)  \right)  }\right) \\
&  =\left(  \text{the }i\text{-th entry of }\left(  f\left(  C_{1}\right)
,f\left(  C_{2}\right)  ,\ldots,f\left(  C_{k}\right)  \right)  \right)
=f\left(  \underbrace{C_{i}}_{=w}\right)  =f\left(  w\right)  .
\end{align*}
The same argument (applied to $g$ instead of $f$) yields%
\[
\left(  \text{the }i\text{-th entry of }\Phi\left(  g\right)  \right)
=g\left(  w\right)  .
\]
Hence, $f\left(  w\right)  =\left(  \text{the }i\text{-th entry of
}\underbrace{\Phi\left(  f\right)  }_{=\Phi\left(  g\right)  }\right)
=\left(  \text{the }i\text{-th entry of }\Phi\left(  g\right)  \right)
=g\left(  w\right)  $.
\par
Now, let us forget that we fixed $w$. We thus have proven that $f\left(
w\right)  =g\left(  w\right)  $ for every $w\in W$. In other words, $f=g$.
\par
Let us now forget that we fixed $f$ and $g$. We thus have shown that if $f$
and $g$ are two elements of $Y^{W}$ such that $\Phi\left(  f\right)
=\Phi\left(  g\right)  $, then $f=g$. In other words, the map $\Phi$ is
injective. Qed.}. Now, let $s\in Y^{k}$ be arbitrary. Let us write $s$ in the
form $\left(  y_{1},y_{2},\ldots,y_{k}\right)  $. Thus, $s=\left(  y_{1}%
,y_{2},\ldots,y_{k}\right)  $.

We now define a map $f:W\rightarrow Y$ as follows: Let $w\in W$. Then, there
exists a unique $i\in\left\{  1,2,\ldots,k\right\}  $ such that $w=C_{i}$
(since $\left(  C_{1},C_{2},\ldots,C_{k}\right)  $ is a list of all elements
of $W$ which contains each of these elements exactly once). Consider this $i$.
Then, we set $f\left(  w\right)  =y_{i}$.

Thus, we have defined a map $f:W\rightarrow Y$. It is clear that if $w\in W$,
and if $i\in\left\{  1,2,\ldots,k\right\}  $ is such that $w=C_{i}$, then%
\begin{equation}
f\left(  w\right)  =y_{i} \label{pf.lem.function-count.1}%
\end{equation}
(by the definition of $f\left(  w\right)  $).

Now, every $i\in\left\{  1,2,\ldots,k\right\}  $ satisfies $f\left(
C_{i}\right)  =y_{i}$ (by (\ref{pf.lem.function-count.1}), applied to
$w=C_{i}$). Hence, $\left(  f\left(  C_{1}\right)  ,f\left(  C_{2}\right)
,\ldots,f\left(  C_{k}\right)  \right)  =\left(  y_{1},y_{2},\ldots
,y_{k}\right)  $. Now, the definition of $\Phi$ yields $\Phi\left(  f\right)
=\left(  f\left(  C_{1}\right)  ,f\left(  C_{2}\right)  ,\ldots,f\left(
C_{k}\right)  \right)  =\left(  y_{1},y_{2},\ldots,y_{k}\right)  =s$. Thus,
$s=\Phi\left(  \underbrace{f}_{\in Y^{W}}\right)  \in\Phi\left(  Y^{W}\right)
$.

Now, let us forget that we fixed $s$. We thus have proven that $s\in
\Phi\left(  Y^{W}\right)  $ for every $s\in Y^{k}$. In other words,
$Y^{k}\subseteq\Phi\left(  Y^{W}\right)  $. In other words, the map $\Phi$ is surjective.

Since the map $\Phi$ is both injective and surjective, we see that the map
$\Phi$ is bijective. In other words, the map $\Phi$ is a bijection. In other
words, the map
\begin{align*}
Y^{W}  &  \rightarrow Y^{k},\\
f  &  \mapsto\left(  f\left(  C_{1}\right)  ,f\left(  C_{2}\right)
,\ldots,f\left(  C_{k}\right)  \right)
\end{align*}
is a bijection (since this map is $\Phi$). Lemma \ref{lem.function-count} is
thus proven.
\end{proof}
\end{verlong}

\begin{lemma}
\label{lem.Eqs.sum}Let $\left(  V,B\right)  $ be a finite graph. Then,%
\[
\sum_{\substack{f:V\rightarrow\mathbb{N}_{+};\\B\subseteq\operatorname*{Eqs}%
f}}\mathbf{x}_{f}=p_{\lambda\left(  V,B\right)  }%
\]
(Here, $\mathbf{x}_{f}$ is defined as in Definition \ref{def.chromsym}, and
the expression $\lambda\left(  V,B\right)  $ is understood according to
Definition \ref{def.connectedness} \textbf{(b)}.)
\end{lemma}

\begin{vershort}
\begin{proof}
[Proof of Lemma \ref{lem.Eqs.sum}.]Let $\left(  C_{1},C_{2},\ldots
,C_{k}\right)  $ be a list of all connected components of $\left(  V,B\right)
$, ordered such that $\left\vert C_{1}\right\vert \geq\left\vert
C_{2}\right\vert \geq\cdots\geq\left\vert C_{k}\right\vert $%
.\ \ \ \ \footnote{Every connected component of $\left(  V,B\right)  $ should
appear exactly once in this list.} Then, $\lambda\left(  V,B\right)  =\left(
\left\vert C_{1}\right\vert ,\left\vert C_{2}\right\vert ,\ldots,\left\vert
C_{k}\right\vert \right)  $ (by the definition of $\lambda\left(  V,B\right)
$). Hence, (\ref{eq.def.powersum2.finite-expression}) (applied to
$\lambda\left(  V,B\right)  $ and $\left\vert C_{i}\right\vert $ instead of
$\lambda$ and $\lambda_{i}$) shows that{}%
\begin{equation}
p_{\lambda\left(  V,B\right)  }=p_{\left\vert C_{1}\right\vert }p_{\left\vert
C_{2}\right\vert }\cdots p_{\left\vert C_{k}\right\vert }=\prod_{i=1}%
^{k}p_{\left\vert C_{i}\right\vert }. \label{pf.lem.Eqs.sum.short.p1}%
\end{equation}
But for every $i\in\left\{  1,2,\ldots,k\right\}  $, we have $p_{\left\vert
C_{i}\right\vert }=\sum_{s\in\mathbb{N}_{+}}x_{s}^{\left\vert C_{i}\right\vert
}$ (by the definition of $p_{\left\vert C_{i}\right\vert }$). Hence,
(\ref{pf.lem.Eqs.sum.short.p1}) becomes%
\begin{equation}
p_{\lambda\left(  V,B\right)  }=\prod_{i=1}^{k}\underbrace{p_{\left\vert
C_{i}\right\vert }}_{=\sum_{s\in\mathbb{N}_{+}}x_{s}^{\left\vert
C_{i}\right\vert }}=\prod_{i=1}^{k}\sum_{s\in\mathbb{N}_{+}}x_{s}^{\left\vert
C_{i}\right\vert }=\sum_{\left(  s_{1},s_{2},\ldots,s_{k}\right)  \in\left(
\mathbb{N}_{+}\right)  ^{k}}\prod_{i=1}^{k}x_{s_{i}}^{\left\vert
C_{i}\right\vert } \label{pf.lem.Eqs.sum.short.p2}%
\end{equation}
(by the product rule).

The list $\left(  C_{1},C_{2},\ldots,C_{k}\right)  $ contains all connected
components of $\left(  V,B\right)  $, each exactly once. Thus, $V=\bigsqcup
_{i=1}^{k}C_{i}$.

We now define a map%
\[
\Phi:\left(  \mathbb{N}_{+}\right)  ^{k}\rightarrow\left\{  f:V\rightarrow
\mathbb{N}_{+}\ \mid\ B\subseteq\operatorname*{Eqs}f\right\}
\]
as follows: Given any $\left(  s_{1},s_{2},\ldots,s_{k}\right)  \in\left(
\mathbb{N}_{+}\right)  ^{k}$, we let $\Phi\left(  s_{1},s_{2},\ldots
,s_{k}\right)  $ be the map $V\rightarrow\mathbb{N}_{+}$ which sends every
$v\in V$ to $s_{i}$, where $i\in\left\{  1,2,\ldots,k\right\}  $ is such that
$v\in C_{i}$. (This is well-defined, because for every $v\in V$, there exists
a unique $i\in\left\{  1,2,\ldots,k\right\}  $ such that $v\in C_{i}$; this
follows from $V=\bigsqcup_{i=1}^{k}C_{i}$.) This map $\Phi$ is well-defined,
because for every $\left(  s_{1},s_{2},\ldots,s_{k}\right)  \in\left(
\mathbb{N}_{+}\right)  ^{k}$, the map $\Phi\left(  s_{1},s_{2},\ldots
,s_{k}\right)  $ actually belongs to $\left\{  f:V\rightarrow\mathbb{N}%
_{+}\ \mid\ B\subseteq\operatorname*{Eqs}f\right\}  $%
\ \ \ \ \footnote{\textit{Proof.} We just need to check that $B\subseteq
\operatorname*{Eqs}\left(  \Phi\left(  s_{1},s_{2},\ldots,s_{k}\right)
\right)  $. But this is easy: For every $\left(  u,v\right)  \in B$, the
vertices $u$ and $v$ of $\left(  V,B\right)  $ lie in one and the same
connected component $C_{i}$ of the graph $\left(  V,B\right)  $, and thus (by
the definition of $\Phi\left(  s_{1},s_{2},\ldots,s_{k}\right)  $) the map
$\Phi\left(  s_{1},s_{2},\ldots,s_{k}\right)  $ sends both of them to $s_{i}$;
but this shows that $\left(  u,v\right)  \in\operatorname*{Eqs}\left(
\Phi\left(  s_{1},s_{2},\ldots,s_{k}\right)  \right)  $.}.

A moment's thought reveals that the map $\Phi$ is injective\footnote{In fact,
we can reconstruct $\left(  s_{1},s_{2},\ldots,s_{k}\right)  \in\left(
\mathbb{N}_{+}\right)  ^{k}$ from its image $\Phi\left(  s_{1},s_{2}%
,\ldots,s_{k}\right)  $, because each $s_{i}$ is the image of any element of
$C_{i}$ under $\Phi\left(  s_{1},s_{2},\ldots,s_{k}\right)  $ (and this allows
us to compute $s_{i}$, since $C_{i}$ is nonempty).}. Let us now show that the
map $\Phi$ is surjective.

In order to show this, we must prove that every map $f:V\rightarrow
\mathbb{N}_{+}$ satisfying $B\subseteq\operatorname*{Eqs}f$ has the form
$\Phi\left(  s_{1},s_{2},\ldots,s_{k}\right)  $ for some $\left(  s_{1}%
,s_{2},\ldots,s_{k}\right)  \in\left(  \mathbb{N}_{+}\right)  ^{k}$. So let us
fix a map $f:V\rightarrow\mathbb{N}_{+}$ satisfying $B\subseteq
\operatorname*{Eqs}f$. We must find some $\left(  s_{1},s_{2},\ldots
,s_{k}\right)  \in\left(  \mathbb{N}_{+}\right)  ^{k}$ such that
$f=\Phi\left(  s_{1},s_{2},\ldots,s_{k}\right)  $.

We have $B\subseteq\operatorname*{Eqs}f$. Thus, for every $\left\{
s,t\right\}  \in B$, we have $\left\{  s,t\right\}  \in B\subseteq
\operatorname*{Eqs}f$ and thus%
\begin{equation}
f\left(  s\right)  =f\left(  t\right)  . \label{pf.lem.Eqs.sum.short.surj.1}%
\end{equation}


Now, if $x$ and $y$ are two elements of $V$ lying in the same connected
component of $\left(  V,B\right)  $, then%
\begin{equation}
f\left(  x\right)  =f\left(  y\right)  \label{pf.lem.Eqs.sum.short.surj.2}%
\end{equation}
\footnote{\textit{Proof of (\ref{pf.lem.Eqs.sum.short.surj.2}):} Let $x$ and
$y$ be two elements of $V$ lying in the same connected component of $\left(
V,B\right)  $. Then, the vertices $x$ and $y$ are connected by a walk in the
graph $\left(  V,B\right)  $ (by the definition of a \textquotedblleft
connected component\textquotedblright). Let $\left(  v_{0},v_{1},\ldots
,v_{j}\right)  $ be this walk (regarded as a sequence of vertices); thus,
$v_{0}=x$ and $v_{j}=y$. For every $i\in\left\{  0,1,\ldots,j-1\right\}  $, we
have $\left\{  v_{i},v_{i+1}\right\}  \in B$ (since $\left(  v_{0}%
,v_{1},\ldots,v_{j}\right)  $ is a walk in the graph $\left(  V,B\right)  $)
and thus $f\left(  v_{i}\right)  =f\left(  v_{i+1}\right)  $ (by
(\ref{pf.lem.Eqs.sum.short.surj.1}), applied to $\left(  s,t\right)  =\left(
v_{i},v_{i+1}\right)  $). In other words, $f\left(  v_{0}\right)  =f\left(
v_{1}\right)  =\cdots=f\left(  v_{j}\right)  $. Hence, $f\left(  v_{0}\right)
=f\left(  v_{j}\right)  $, so that $f\left(  \underbrace{x}_{=v_{0}}\right)
=f\left(  v_{0}\right)  =f\left(  \underbrace{v_{j}}_{=y}\right)  =f\left(
y\right)  $, qed.}. In other words, the map $f$ is constant on each connected
component of $\left(  V,B\right)  $. Thus, the map $f$ is constant on $C_{i}$
for each $i\in\left\{  1,2,\ldots,k\right\}  $ (since $C_{i}$ is a connected
component of $\left(  V,B\right)  $). Hence, for each $i\in\left\{
1,2,\ldots,k\right\}  $, we can define a positive integer $s_{i}\in
\mathbb{N}_{+}$ to be the image of any element of $C_{i}$ under $f$ (this is
well-defined, because $f$ is constant on $C_{i}$ and thus the choice of the
element does not matter). Define $s_{i}\in\mathbb{N}_{+}$ for each
$i\in\left\{  1,2,\ldots,k\right\}  $ this way. Thus, we have defined a
$k$-tuple $\left(  s_{1},s_{2},\ldots,s_{k}\right)  \in\left(  \mathbb{N}%
_{+}\right)  ^{k}$. Now, $f=\Phi\left(  s_{1},s_{2},\ldots,s_{k}\right)  $
(this follows immediately by recalling the definitions of $\Phi$ and $s_{i}$).

Let us now forget that we fixed $f$. We thus have shown that for every map
$f:V\rightarrow\mathbb{N}_{+}$ satisfying $B\subseteq\operatorname*{Eqs}f$,
there exists some $\left(  s_{1},s_{2},\ldots,s_{k}\right)  \in\left(
\mathbb{N}_{+}\right)  ^{k}$ such that $f=\Phi\left(  s_{1},s_{2},\ldots
,s_{k}\right)  $. In other words, the map $\Phi$ is surjective. Since $\Phi$
is both injective and surjective, we conclude that $\Phi$ is a bijection.

Moreover, it is straightforward to see that every map $\left(  s_{1}%
,s_{2},\ldots,s_{k}\right)  \in\left(  \mathbb{N}_{+}\right)  ^{k}$ satisfies%
\begin{equation}
\mathbf{x}_{\Phi\left(  s_{1},s_{2},\ldots,s_{k}\right)  }=\prod_{i=1}%
^{k}x_{s_{i}}^{\left\vert C_{i}\right\vert } \label{pf.lem.Eqs.sum.short.4}%
\end{equation}
(by the definitions of $\mathbf{x}_{\Phi\left(  s_{1},s_{2},\ldots
,s_{k}\right)  }$ and of $\Phi$). Now,%
\begin{align*}
\sum_{\substack{f:V\rightarrow\mathbb{N}_{+};\\B\subseteq\operatorname*{Eqs}%
f}}\mathbf{x}_{f}  &  =\sum_{\left(  s_{1},s_{2},\ldots,s_{k}\right)
\in\left(  \mathbb{N}_{+}\right)  ^{k}}\underbrace{\mathbf{x}_{\Phi\left(
s_{1},s_{2},\ldots,s_{k}\right)  }}_{\substack{=\prod_{i=1}^{k}x_{s_{i}%
}^{\left\vert C_{i}\right\vert }\\\text{(by (\ref{pf.lem.Eqs.sum.short.4}))}%
}}\\
&  \ \ \ \ \ \ \ \ \ \ \left(
\begin{array}
[c]{c}%
\text{here, we have substituted }\Phi\left(  s_{1},s_{2},\ldots,s_{k}\right)
\text{ for }f\text{ in the sum,}\\
\text{since the map }\Phi:\left(  \mathbb{N}_{+}\right)  ^{k}\rightarrow
\left\{  f:V\rightarrow\mathbb{N}_{+}\ \mid\ B\subseteq\operatorname*{Eqs}%
f\right\} \\
\text{is a bijection}%
\end{array}
\right) \\
&  =\sum_{\left(  s_{1},s_{2},\ldots,s_{k}\right)  \in\left(  \mathbb{N}%
_{+}\right)  ^{k}}\prod_{i=1}^{k}x_{s_{i}}^{\left\vert C_{i}\right\vert
}=p_{\lambda\left(  V,B\right)  }\ \ \ \ \ \ \ \ \ \ \left(  \text{by
(\ref{pf.lem.Eqs.sum.short.p2})}\right)  .
\end{align*}
This proves Lemma \ref{lem.Eqs.sum}.
\end{proof}
\end{vershort}

\begin{verlong}
\begin{proof}
[Proof of Lemma \ref{lem.Eqs.sum}.]Let $\sim$ denote the equivalence relation
$\sim_{\left(  V,B\right)  }$ (defined as in Definition
\ref{def.connectedness} \textbf{(a)}). The connected components of $\left(
V,B\right)  $ are the $\sim_{\left(  V,B\right)  }$-equivalence classes
(because this is how the connected components of $\left(  V,B\right)  $ are
defined). In other words, the connected components of $\left(  V,B\right)  $
are the $\sim$-equivalence classes (since $\sim$ is the relation
$\sim_{\left(  V,B\right)  }$). In other words, the connected components of
$\left(  V,B\right)  $ are the elements of $V/\left(  \sim\right)  $ (since
the elements of $V/\left(  \sim\right)  $ are the $\sim$-equivalence classes
(by the definition of $V/\left(  \sim\right)  $)).

A set $\left(  \mathbb{N}_{+}\right)  _{\sim}^{V}$ is defined (according to
Definition \ref{def.relquot.maps} \textbf{(b)}).

Proposition \ref{prop.relquot.uniprop} \textbf{(b)} (applied to $X=V$ and
$Y=\mathbb{N}_{+}$) shows that the map%
\[
\left(  \mathbb{N}_{+}\right)  ^{V/\left(  \sim\right)  }\rightarrow\left(
\mathbb{N}_{+}\right)  _{\sim}^{V},\ \ \ \ \ \ \ \ \ \ f\mapsto f\circ\pi_{V}%
\]
is a bijection.

For every map $f:V\rightarrow\mathbb{N}_{+}$, we have the following
equivalence:%
\begin{equation}
\left(  B\subseteq\operatorname*{Eqs}f\right)  \ \Longleftrightarrow\ \left(
f\in\left(  \mathbb{N}_{+}\right)  _{\sim}^{V}\right)
\label{pf.lem.Eqs.sum.equivalence}%
\end{equation}
(according to Lemma \ref{lem.Eqs.sum-aux}, applied to $Y=\mathbb{N}_{+}$).
Thus, we have the following equality of summation signs:
\[
\sum_{\substack{f:V\rightarrow\mathbb{N}_{+};\\B\subseteq\operatorname*{Eqs}%
f}}=\sum_{\substack{f:V\rightarrow\mathbb{N}_{+};\\f\in\left(  \mathbb{N}%
_{+}\right)  _{\sim}^{V}}}=\sum_{\substack{f\in\left(  \mathbb{N}_{+}\right)
^{V};\\f\in\left(  \mathbb{N}_{+}\right)  _{\sim}^{V}}}=\sum_{f\in\left(
\mathbb{N}_{+}\right)  _{\sim}^{V}}%
\]
(since $\left(  \mathbb{N}_{+}\right)  _{\sim}^{V}$ is a subset of $\left(
\mathbb{N}_{+}\right)  ^{V}$). Hence,%
\begin{equation}
\underbrace{\sum_{\substack{f:V\rightarrow\mathbb{N}_{+};\\B\subseteq
\operatorname*{Eqs}f}}}_{=\sum_{f\in\left(  \mathbb{N}_{+}\right)  _{\sim}%
^{V}}}\mathbf{x}_{f}=\sum_{f\in\left(  \mathbb{N}_{+}\right)  _{\sim}^{V}%
}\mathbf{x}_{f}=\sum_{f\in\left(  \mathbb{N}_{+}\right)  ^{V/\left(
\sim\right)  }}\mathbf{x}_{f\circ\pi_{V}} \label{pf.lem.Eqs.sum.1}%
\end{equation}
(here, we have substituted $f\circ\pi_{V}$ for $f$ in the sum, since the map
$\left(  \mathbb{N}_{+}\right)  ^{V/\left(  \sim\right)  }\rightarrow\left(
\mathbb{N}_{+}\right)  _{\sim}^{V},\ f\mapsto f\circ\pi_{V}$ is a bijection).

Now, let $\left(  C_{1},C_{2},\ldots,C_{k}\right)  $ be a list of all
connected components of $\left(  V,B\right)  $, ordered such that $\left\vert
C_{1}\right\vert \geq\left\vert C_{2}\right\vert \geq\cdots\geq\left\vert
C_{k}\right\vert $.\ \ \ \ \footnote{Every connected component of $\left(
V,B\right)  $ should appear exactly once in this list.} Then, $\left(
\left\vert C_{1}\right\vert ,\left\vert C_{2}\right\vert ,\ldots,\left\vert
C_{k}\right\vert \right)  $ is the list of the sizes of all connected
components of $\left(  V,B\right)  $, in weakly decreasing order (since
$\left\vert C_{1}\right\vert \geq\left\vert C_{2}\right\vert \geq\cdots
\geq\left\vert C_{k}\right\vert $). In other words, $\left(  \left\vert
C_{1}\right\vert ,\left\vert C_{2}\right\vert ,\ldots,\left\vert
C_{k}\right\vert \right)  $ is $\lambda\left(  V,B\right)  $ (since
$\lambda\left(  V,B\right)  $ is the list of the sizes of all connected
components of $\left(  V,B\right)  $, in weakly decreasing order (by the
definition of $\lambda\left(  V,B\right)  $)). In other words, $\lambda\left(
V,B\right)  =\left(  \left\vert C_{1}\right\vert ,\left\vert C_{2}\right\vert
,\ldots,\left\vert C_{k}\right\vert \right)  $. Thus,
(\ref{eq.def.powersum2.finite-expression}) (applied to $\lambda\left(
V,B\right)  $ and $\left\vert C_{i}\right\vert $ instead of $\lambda$ and
$\lambda_{i}$) shows that{}%
\begin{equation}
p_{\lambda\left(  V,B\right)  }=p_{\left\vert C_{1}\right\vert }p_{\left\vert
C_{2}\right\vert }\cdots p_{\left\vert C_{k}\right\vert }=\prod_{i=1}%
^{k}p_{\left\vert C_{i}\right\vert }. \label{pf.lem.Eqs.sum.p1}%
\end{equation}


But for every $i\in\left\{  1,2,\ldots,k\right\}  $, we have $p_{\left\vert
C_{i}\right\vert }=\sum_{s\in\mathbb{N}_{+}}x_{s}^{\left\vert C_{i}\right\vert
}$\ \ \ \ \footnote{\textit{Proof.} Let $i\in\left\{  1,2,\ldots,k\right\}  $.
Then, $C_{i}$ is a connected component of $\left(  V,B\right)  $ (since
$\left(  C_{1},C_{2},\ldots,C_{k}\right)  $ is a list of all connected
components of $\left(  V,B\right)  $). Hence, $C_{i}$ is a nonempty subset of
$V$ (since every connected component of $\left(  V,B\right)  $ is a nonempty
subset of $V$). Hence, $\left\vert C_{i}\right\vert $ is a positive integer.
Thus, (\ref{eq.def.powersum.pn}) (applied to $n=\left\vert C_{i}\right\vert $)
shows that $p_{\left\vert C_{i}\right\vert }=\underbrace{\sum_{j\geq1}}%
_{=\sum_{j\in\mathbb{N}_{+}}}x_{j}^{\left\vert C_{i}\right\vert }=\sum
_{j\in\mathbb{N}_{+}}x_{j}^{\left\vert C_{i}\right\vert }=\sum_{s\in
\mathbb{N}_{+}}x_{s}^{\left\vert C_{i}\right\vert }$ (here, we have renamed
the summation index $j$ as $s$). Qed.}. Hence, (\ref{pf.lem.Eqs.sum.p1})
becomes%
\begin{equation}
p_{\lambda\left(  V,B\right)  }=\prod_{i=1}^{k}\underbrace{p_{\left\vert
C_{i}\right\vert }}_{=\sum_{s\in\mathbb{N}_{+}}x_{s}^{\left\vert
C_{i}\right\vert }}=\prod_{i=1}^{k}\sum_{s\in\mathbb{N}_{+}}x_{s}^{\left\vert
C_{i}\right\vert }=\sum_{\left(  s_{1},s_{2},\ldots,s_{k}\right)  \in\left(
\mathbb{N}_{+}\right)  ^{k}}\prod_{i=1}^{k}x_{s_{i}}^{\left\vert
C_{i}\right\vert } \label{pf.lem.Eqs.sum.p2}%
\end{equation}
(by the product rule).

Recall that $\left(  C_{1},C_{2},\ldots,C_{k}\right)  $ is a list of all
connected components of $\left(  V,B\right)  $. In other words, $\left(
C_{1},C_{2},\ldots,C_{k}\right)  $ is a list of all elements of $V/\left(
\sim\right)  $ (since the elements of $V/\left(  \sim\right)  $ are the
connected components of $\left(  V,B\right)  $). Moreover, every element of
$V/\left(  \sim\right)  $ appears exactly once in this list $\left(
C_{1},C_{2},\ldots,C_{k}\right)  $ (since the entries of the list $\left(
C_{1},C_{2},\ldots,C_{k}\right)  $ are pairwise distinct\footnote{since every
connected component of $\left(  V,B\right)  $ appears exactly once in this
list}). Thus, $\left(  C_{1},C_{2},\ldots,C_{k}\right)  $ is a list of all
elements of $V/\left(  \sim\right)  $, and contains each of these elements
exactly once. Hence, the map%
\begin{align*}
\left(  \mathbb{N}_{+}\right)  ^{V/\left(  \sim\right)  }  &  \rightarrow
\left(  \mathbb{N}_{+}\right)  ^{k},\\
f  &  \mapsto\left(  f\left(  C_{1}\right)  ,f\left(  C_{2}\right)
,\ldots,f\left(  C_{k}\right)  \right)
\end{align*}
is a bijection (by Lemma \ref{lem.function-count}, applied to $W=V/\left(
\sim\right)  $ and $Y=\mathbb{N}_{+}$).

For every $w\in V/\left(  \sim\right)  $, we have%
\begin{equation}
\pi_{V}^{-1}\left(  w\right)  =w \label{pf.lem.Eqs.sum.pi-1}%
\end{equation}
\footnote{\textit{Proof of (\ref{pf.lem.Eqs.sum.pi-1}):} Let $w\in V/\left(
\sim\right)  $. Thus, $w$ is an element of $V/\left(  \sim\right)  $. In other
words, $w$ is an $\sim$-equivalence class (since the elements of $V/\left(
\sim\right)  $ are the $\sim$-equivalence classes). In particular, $w\subseteq
V$.
\par
Let $x\in w$. Then, $x\in w\subseteq V$. The $\sim$-equivalence class of $x$
must be $w$ (since $x$ lies in the $\sim$-equivalence class $w$ (since $x\in
w$)). In other words, $\left[  x\right]  _{\sim}$ must be $w$ (since $\left[
x\right]  _{\sim}$ is the $\sim$-equivalence class of $x$). In other words,
$\left[  x\right]  _{\sim}=w$. Now, the definition of $\pi_{V}$ yields
$\pi_{V}\left(  x\right)  =\left[  x\right]  _{\sim}=w$. Hence, $x\in\pi
_{V}^{-1}\left(  w\right)  $.
\par
Let us now forget that we fixed $x$. We thus have shown that $x\in\pi_{V}%
^{-1}\left(  w\right)  $ for every $x\in w$. In other words, $w\subseteq
\pi_{V}^{-1}\left(  w\right)  $.
\par
Let now $y\in\pi_{V}^{-1}\left(  w\right)  $. Thus, $y\in V$ and $\pi
_{V}\left(  y\right)  =w$. The definition of $\pi_{V}$ yields $\pi_{V}\left(
y\right)  =\left[  y\right]  _{\sim}$. Thus, $w=\pi_{V}\left(  y\right)
=\left[  y\right]  _{\sim}$. Hence, $w$ is the $\sim$-equivalence class of $y$
(since $\left[  y\right]  _{\sim}$ is the $\sim$-equivalence class of $y$).
Consequently, $y$ must belong to $w$. In other words, $y\in w$.
\par
Let us now forget that we fixed $y$. We thus have shown that $y\in w$ for
every $y\in\pi_{V}^{-1}\left(  w\right)  $. In other words, $\pi_{V}%
^{-1}\left(  w\right)  \subseteq w$. Combining this with $w\subseteq\pi
_{V}^{-1}\left(  w\right)  $, we obtain $\pi_{V}^{-1}\left(  w\right)  =w$.
This proves (\ref{pf.lem.Eqs.sum.pi-1}).}.

Also, the map $\left\{  1,2,\ldots,k\right\}  \rightarrow V/\left(
\sim\right)  ,\ i\mapsto C_{i}$ is a bijection (since $\left(  C_{1}%
,C_{2},\ldots,C_{k}\right)  $ is a list of all elements of $V/\left(
\sim\right)  $, and contains each of these elements exactly once).

We have%
\begin{equation}
\mathbf{x}_{f\circ\pi_{V}}=\prod_{i=1}^{k}x_{f\left(  C_{i}\right)
}^{\left\vert C_{i}\right\vert }\ \ \ \ \ \ \ \ \ \ \text{for every }%
f\in\left(  \mathbb{N}_{+}\right)  ^{V/\left(  \sim\right)  }
\label{pf.lem.Eqs.sum.3}%
\end{equation}
\footnote{\textit{Proof of (\ref{pf.lem.Eqs.sum.3}):} Let $f\in\left(
\mathbb{N}_{+}\right)  ^{V/\left(  \sim\right)  }$. Then, the definition of
$\mathbf{x}_{f\circ\pi_{V}}$ yields%
\begin{align*}
\mathbf{x}_{f\circ\pi_{V}}  &  =\prod_{v\in V}x_{\left(  f\circ\pi_{V}\right)
\left(  v\right)  }=\prod_{w\in V/\left(  \sim\right)  }\underbrace{\prod
_{\substack{v\in V;\\\pi_{V}\left(  v\right)  =w}}}_{\substack{=\prod_{v\in
\pi_{V}^{-1}\left(  w\right)  }=\prod_{v\in w}\\\text{(since }\pi_{V}%
^{-1}\left(  w\right)  =w\\\text{(by (\ref{pf.lem.Eqs.sum.pi-1}))}%
}}\underbrace{x_{\left(  f\circ\pi_{V}\right)  \left(  v\right)  }%
}_{\substack{=x_{f\left(  w\right)  }\\\text{(since }\left(  f\circ\pi
_{V}\right)  \left(  v\right)  =f\left(  \pi_{V}\left(  v\right)  \right)
=f\left(  w\right)  \\\text{(since }\pi_{V}\left(  v\right)  =w\text{))}}}\\
&  \ \ \ \ \ \ \ \ \ \ \left(
\begin{array}
[c]{c}%
\text{because for every }v\in V\text{, there exists a unique }w\in V/\left(
\sim\right) \\
\text{such that }\pi_{V}\left(  v\right)  =w\text{ (since }\pi_{V}\text{ is a
map }V\rightarrow V/\left(  \sim\right)  \text{)}%
\end{array}
\right) \\
&  =\prod_{w\in V/\left(  \sim\right)  }\underbrace{\prod_{v\in w}x_{f\left(
w\right)  }}_{=x_{f\left(  w\right)  }^{\left\vert w\right\vert }}=\prod_{w\in
V/\left(  \sim\right)  }x_{f\left(  w\right)  }^{\left\vert w\right\vert
}=\prod_{i\in\left\{  1,2,\ldots,k\right\}  }x_{f\left(  C_{i}\right)
}^{\left\vert C_{i}\right\vert }%
\end{align*}
(here, we have substituted $C_{i}$ for $w$ in the product, since the map
$\left\{  1,2,\ldots,k\right\}  \rightarrow V/\left(  \sim\right)  ,\ i\mapsto
C_{i}$ is a bijection). Thus,%
\[
\mathbf{x}_{f\circ\pi_{V}}=\underbrace{\prod_{i\in\left\{  1,2,\ldots
,k\right\}  }}_{=\prod_{i=1}^{k}}x_{f\left(  C_{i}\right)  }^{\left\vert
C_{i}\right\vert }=\prod_{i=1}^{k}x_{f\left(  C_{i}\right)  }^{\left\vert
C_{i}\right\vert }.
\]
This proves (\ref{pf.lem.Eqs.sum.3}).}.

Now, (\ref{pf.lem.Eqs.sum.1}) becomes%
\begin{align*}
\sum_{\substack{f:V\rightarrow\mathbb{N}_{+};\\B\subseteq\operatorname*{Eqs}%
f}}\mathbf{x}_{f}  &  =\sum_{f\in\left(  \mathbb{N}_{+}\right)  ^{V/\left(
\sim\right)  }}\underbrace{\mathbf{x}_{f\circ\pi_{V}}}_{\substack{=\prod
_{i=1}^{k}x_{f\left(  C_{i}\right)  }^{\left\vert C_{i}\right\vert
}\\\text{(by (\ref{pf.lem.Eqs.sum.3}))}}}=\sum_{f\in\left(  \mathbb{N}%
_{+}\right)  ^{V/\left(  \sim\right)  }}\prod_{i=1}^{k}x_{f\left(
C_{i}\right)  }^{\left\vert C_{i}\right\vert }\\
&  =\sum_{\left(  s_{1},s_{2},\ldots,s_{k}\right)  \in\left(  \mathbb{N}%
_{+}\right)  ^{k}}\prod_{i=1}^{k}x_{s_{i}}^{\left\vert C_{i}\right\vert }%
\end{align*}
(here, we have substituted $\left(  s_{1},s_{2},\ldots,s_{k}\right)  $ for
$\left(  f\left(  C_{1}\right)  ,f\left(  C_{2}\right)  ,\ldots,f\left(
C_{k}\right)  \right)  $ in the sum, since the map $\left(  \mathbb{N}%
_{+}\right)  ^{V/\left(  \sim\right)  }\rightarrow\left(  \mathbb{N}%
_{+}\right)  ^{k},\ f\mapsto\left(  f\left(  C_{1}\right)  ,f\left(
C_{2}\right)  ,\ldots,f\left(  C_{k}\right)  \right)  $ is a bijection).
Comparing this with (\ref{pf.lem.Eqs.sum.p2}), we obtain $\sum
_{\substack{f:V\rightarrow\mathbb{N}_{+};\\B\subseteq\operatorname*{Eqs}%
f}}\mathbf{x}_{f}=p_{\lambda\left(  V,B\right)  }$. This proves Lemma
\ref{lem.Eqs.sum}.
\end{proof}
\end{verlong}

\begin{lemma}
\label{lem.BC.nonempty}Let $G=\left(  V,E\right)  $ be a finite graph. Let $X$
be a totally ordered set. Let $\ell:E\rightarrow X$ be a function. Let $K$ be
a broken circuit of $G$. Then, $K\neq\varnothing$.
\end{lemma}

\begin{vershort}
\begin{proof}
[Proof of Lemma \ref{lem.BC.nonempty}.]The set $K$ is a broken circuit of $G$,
and thus is a circuit of $G$ with an edge removed (by the definition of a
broken circuit). Thus, the set $K$ contains at least $1$ edge (since every
circuit of $G$ contains at least $2$ edges). This proves Lemma
\ref{lem.BC.nonempty}.
\end{proof}
\end{vershort}

\begin{verlong}
\begin{proof}
[Proof of Lemma \ref{lem.BC.nonempty}.]The set $K$ is a broken circuit of $G$.
In other words, the set $K$ is a subset of $E$ having the form $C\setminus
\left\{  e\right\}  $, where $C$ is a circuit of $G$, and where $e$ is the
unique edge in $C$ having maximum label (among the edges in $C$%
)\ \ \ \ \footnote{because a broken circuit of $G$ is the same as a subset of
$E$ having the form $C\setminus\left\{  e\right\}  $, where $C$ is a circuit
of $G$, and where $e$ is the unique edge in $C$ having maximum label (among
the edges in $C$) (by the definition of a \textquotedblleft broken
circuit\textquotedblright)}. Consider this $C$ and this $e$. Thus, we have the
following facts:

\begin{itemize}
\item The set $C$ is a circuit of $G$.

\item The element $e$ is the unique edge in $C$ having maximum label
(among the edges in $C$).

\item We have $K=C\setminus\left\{  e\right\}  $.
\end{itemize}

Now, assume (for the sake of contradiction) that $K=\varnothing$. Consider the
map $\operatorname*{id}:V\rightarrow V$. Then, $\operatorname*{Eqs}%
\operatorname*{id}=\varnothing$\ \ \ \ \footnote{\textit{Proof.} Let
$f\in\operatorname*{Eqs}\operatorname*{id}$. Thus,%
\begin{equation}
f\in\operatorname*{Eqs}\operatorname*{id}=\left\{  \left\{  s,t\right\}
\ \mid\ \left(  s,t\right)  \in V^{2},\ s\neq t\text{ and }\operatorname*{id}%
\left(  s\right)  =\operatorname*{id}\left(  t\right)  \right\} \nonumber
\end{equation}
(by the definition of $\operatorname*{Eqs}\operatorname*{id}$). In other
words, $f$ has the form $\left\{  s,t\right\}  $ for some $\left(  s,t\right)
\in V^{2}$ satisfying $s\neq t$ and $\operatorname*{id}\left(  s\right)
=\operatorname*{id}\left(  t\right)  $. Consider this $\left(  s,t\right)  $.
We have $s=\operatorname*{id}\left(  s\right)  =\operatorname*{id}\left(
t\right)  =t$; but this contradicts $s\neq t$.
\par
Now, let us forget that we fixed $f$. We thus have obtained a contradiction
for every $f\in\operatorname*{Eqs}\operatorname*{id}$. Thus, there exists no
$f\in\operatorname*{Eqs}\operatorname*{id}$. In other words,
$\operatorname*{Eqs}\operatorname*{id}$ is the empty set. Thus,
$\operatorname*{Eqs}\operatorname*{id}=\varnothing$, qed.}. But $C\setminus
\left\{  e\right\}  =K=\varnothing\subseteq\varnothing=\operatorname*{Eqs}%
\operatorname*{id}$. Hence, Lemma \ref{lem.Eqs.circuit} (applied to $V$ and
$\operatorname*{id}$ instead of $X$ and $f$) yields $e\in E\cap
\underbrace{\operatorname*{Eqs}\operatorname*{id}}_{=\varnothing}%
=E\cap\varnothing=\varnothing$. Thus, the set $\varnothing$ has at least one
element (namely, $e$). This contradicts the fact that this set $\varnothing$
is empty. This contradiction shows that our assumption (that $K=\varnothing$)
was wrong. Hence, we cannot have $K=\varnothing$. We thus have $K\neq
\varnothing$. This proves Lemma \ref{lem.BC.nonempty}.
\end{proof}
\end{verlong}

\subsection{Alternating sums}

We shall now come to less simple lemmas.

\begin{definition}
\label{def.iverson}We shall use the so-called \textit{Iverson bracket
notation}: If $\mathcal{S}$ is any logical statement, then $\left[
\mathcal{S}\right]  $ shall mean the integer $%
\begin{cases}
1, & \text{if }\mathcal{S}\text{ is true;}\\
0, & \text{if }\mathcal{S}\text{ is false}%
\end{cases}
$.
\end{definition}

The following lemma is probably the most crucial one in this note:

\begin{lemma}
\label{lem.NBCm.moeb}Let $G=\left(  V,E\right)  $ be a finite graph. Let $X$
be a totally ordered set. Let $\ell:E\rightarrow X$ be a function. Let
$\mathfrak{K}$ be some set of broken circuits of $G$ (not necessarily
containing all of them). Let $a_{K}$ be an element of $\mathbf{k}$ for every
$K\in\mathfrak{K}$.

Let $Y$ be any set. Let $f:V\rightarrow Y$ be any map. Then,%
\[
\sum_{B\subseteq E\cap\operatorname*{Eqs}f}\left(  -1\right)  ^{\left\vert
B\right\vert }\prod_{\substack{K\in\mathfrak{K};\\K\subseteq B}}a_{K}=\left[
E\cap\operatorname*{Eqs}f=\varnothing\right]  .
\]

\end{lemma}

\begin{vershort}
\begin{proof}
[Proof of Lemma \ref{lem.NBCm.moeb}.]We WLOG assume that $E\cap
\operatorname*{Eqs}f\neq\varnothing$ (since otherwise, the claim is
obvious\footnote{In (slightly) more detail: If $E\cap\operatorname*{Eqs}%
f=\varnothing$, then the sum $\sum_{B\subseteq E\cap\operatorname*{Eqs}%
f}\left(  -1\right)  ^{\left\vert B\right\vert }\prod_{\substack{K\in
\mathfrak{K};\\K\subseteq B}}a_{K}$ has only one addend (namely, the addend
for $B=\varnothing$), and thus simplifies to
\begin{align*}
\underbrace{\left(  -1\right)  ^{\left\vert \varnothing\right\vert }%
}_{=\left(  -1\right)  ^{0}=1}\underbrace{\prod_{\substack{K\in\mathfrak{K}%
;\\K\subseteq\varnothing}}}_{=\prod_{\substack{K\in\mathfrak{K}%
;\\K=\varnothing}}}a_{K}  &  =\prod_{\substack{K\in\mathfrak{K}%
;\\K=\varnothing}}a_{K}=\left(  \text{empty product}\right)
\ \ \ \ \ \ \ \ \ \ \left(
\begin{array}
[c]{c}%
\text{since no }K\in\mathfrak{K}\text{ satisfies }K=\varnothing\\
\text{(by Lemma \ref{lem.BC.nonempty})}%
\end{array}
\right) \\
&  =1=\left[  E\cap\operatorname*{Eqs}f=\varnothing\right]  .
\end{align*}
}). Thus, $\left[  E\cap\operatorname*{Eqs}f=\varnothing\right]  =0$.

Pick any $d\in E\cap\operatorname*{Eqs}f$ with maximum $\ell\left(  d\right)
$ (among all $d\in E\cap\operatorname*{Eqs}f$). (This is clearly possible,
since $E\cap\operatorname*{Eqs}f\neq\varnothing$.) Define two subsets
$\mathcal{U}$ and $\mathcal{V}$ of $\mathcal{P}\left(  E\cap
\operatorname*{Eqs}f\right)  $ as follows:%
\begin{align*}
\mathcal{U}  &  =\left\{  F\in\mathcal{P}\left(  E\cap\operatorname*{Eqs}%
f\right)  \ \mid\ d\notin F\right\}  ;\\
\mathcal{V}  &  =\left\{  F\in\mathcal{P}\left(  E\cap\operatorname*{Eqs}%
f\right)  \ \mid\ d\in F\right\}  .
\end{align*}
Thus, we have $\mathcal{P}\left(  E\cap\operatorname*{Eqs}f\right)
=\mathcal{U}\cup\mathcal{V}$, and the sets $\mathcal{U}$ and $\mathcal{V}$ are
disjoint. Now, we define a map $\Phi:\mathcal{U}\rightarrow\mathcal{V}$ by%
\[
\left(  \Phi\left(  B\right)  =B\cup\left\{  d\right\}
\ \ \ \ \ \ \ \ \ \ \text{for every }B\in\mathcal{U}\right)  .
\]
This map $\Phi$ is well-defined (because for every $B\in\mathcal{U}$, we have
$B\cup\left\{  d\right\}  \in\mathcal{V}$\ \ \ \ \footnote{This follows from
the fact that $d\in E\cap\operatorname*{Eqs}f$.}) and a bijection\footnote{Its
inverse is the map $\Psi:\mathcal{V}\rightarrow\mathcal{U}$ defined by
$\left(  \Psi\left(  B\right)  =B\setminus\left\{  d\right\}
\ \ \ \ \ \ \ \ \ \ \text{for every }B\in\mathcal{V}\right)  $.}. Moreover,
every $B\in\mathcal{U}$ satisfies%
\begin{equation}
\left(  -1\right)  ^{\left\vert \Phi\left(  B\right)  \right\vert }=-\left(
-1\right)  ^{\left\vert B\right\vert } \label{pf.lem.NBCm.moeb.short.Phi.-1}%
\end{equation}
\footnote{\textit{Proof.} Let $B\in\mathcal{U}$. Thus, $d\notin B$ (by the
definition of $\mathcal{U}$). Now, $\left\vert \underbrace{\Phi\left(
B\right)  }_{=B\cup\left\{  d\right\}  }\right\vert =\left\vert B\cup\left\{
d\right\}  \right\vert =\left\vert B\right\vert +1$ (since $d\notin B$), so
that $\left(  -1\right)  ^{\left\vert \Phi\left(  B\right)  \right\vert
}=-\left(  -1\right)  ^{\left\vert B\right\vert }$, qed.}.

Now, we claim that, for every $B\in\mathcal{U}$ and every $K\in\mathfrak{K}$,
we have the following logical equivalence:%
\begin{equation}
\left(  K\subseteq B\right)  \ \Longleftrightarrow\ \left(  K\subseteq
\Phi\left(  B\right)  \right)  . \label{pf.lem.NBCm.moeb.short.Phi.equiv}%
\end{equation}


\textit{Proof of (\ref{pf.lem.NBCm.moeb.short.Phi.equiv}):} Let $B\in
\mathcal{U}$ and $K\in\mathfrak{K}$. We must prove the equivalence
(\ref{pf.lem.NBCm.moeb.short.Phi.equiv}). The definition of $\Phi$ yields
$\Phi\left(  B\right)  =B\cup\left\{  d\right\}  \supseteq B$, so that
$B\subseteq\Phi\left(  B\right)  $. Hence, if $K\subseteq B$, then $K\subseteq
B\subseteq\Phi\left(  B\right)  $. Therefore, the forward implication of the
equivalence (\ref{pf.lem.NBCm.moeb.short.Phi.equiv}) is proven. It thus
remains to prove the backward implication of this equivalence. In other words,
it remains to prove that if $K\subseteq\Phi\left(  B\right)  $, then
$K\subseteq B$. So let us assume that $K\subseteq\Phi\left(  B\right)  $.

We want to prove that $K\subseteq B$. Assume the contrary. Thus,
$K\not \subseteq B$. We have $K\in\mathfrak{K}$. Thus, $K$ is a broken circuit
of $G$ (since $\mathfrak{K}$ is a set of broken circuits of $G$). In other
words, $K$ is a subset of $E$ having the form $C\setminus\left\{  e\right\}
$, where $C$ is a circuit of $G$, and where $e$ is the unique edge in $C$
having maximum label (among the edges in $C$) (because this is how a broken
circuit is defined). Consider these $C$ and $e$. Thus, $K=C\setminus\left\{
e\right\}  $.

The element $e$ is the unique edge in $C$ having maximum label (among
the edges in $C$). Thus, if $e^{\prime}$ is any edge in $C$ satisfying
$\ell\left(  e^{\prime}\right)  \geq\ell\left(  e\right)  $, then%
\begin{equation}
e^{\prime}=e. \label{pf.lem.NBCm.moeb.short.Phi.equiv.pf.e'}%
\end{equation}


But $\underbrace{K}_{\subseteq\Phi\left(  B\right)  =B\cup\left\{  d\right\}
}\setminus\left\{  d\right\}  \subseteq\left(  B\cup\left\{  d\right\}
\right)  \setminus\left\{  d\right\}  \subseteq B$.

If we had $d\notin K$, then we would have $K\setminus\left\{  d\right\}  =K$
and therefore $K=K\setminus\left\{  d\right\}  \subseteq B$; this would
contradict $K\not \subseteq B$. Hence, we cannot have $d\notin K$. We thus
must have $d\in K$. Hence, $d\in K=C\setminus\left\{  e\right\}  $. Hence,
$d\in C$ and $d\neq e$.

But $C\setminus\left\{  e\right\}  =K\subseteq\Phi\left(  B\right)  \subseteq
E\cap\operatorname*{Eqs}f$ (since $\Phi\left(  B\right)  \in\mathcal{P}\left(
E\cap\operatorname*{Eqs}f\right)  $), so that $C\setminus\left\{  e\right\}
\subseteq E\cap\operatorname*{Eqs}f\subseteq\operatorname*{Eqs}f$. Hence,
Lemma \ref{lem.Eqs.circuit} (applied to $Y$ instead of $X$) shows that $e\in
E\cap\operatorname*{Eqs}f$. Thus, $\ell\left(  d\right)  \geq\ell\left(
e\right)  $ (since $d$ was defined to be an element of $E\cap
\operatorname*{Eqs}f$ with maximum $\ell\left(  d\right)  $ among all $d\in
E\cap\operatorname*{Eqs}f$).

Also, $d\in C$. Since $\ell\left(  d\right)  \geq\ell\left(  e\right)  $, we
can therefore apply (\ref{pf.lem.NBCm.moeb.short.Phi.equiv.pf.e'}) to
$e^{\prime}=d$. We thus obtain $d=e$. This contradicts $d\neq e$. This
contradiction proves that our assumption was wrong. Hence, $K\subseteq B$ is
proven. Thus, we have proven the backward implication of the equivalence
(\ref{pf.lem.NBCm.moeb.short.Phi.equiv}); this completes the proof of
(\ref{pf.lem.NBCm.moeb.short.Phi.equiv}).

Now, recall that we have $\mathcal{P}\left(  E\cap\operatorname*{Eqs}f\right)
=\mathcal{U}\cup\mathcal{V}$, and the sets $\mathcal{U}$ and $\mathcal{V}$ are
disjoint. Hence, the sum $\sum_{B\subseteq E\cap\operatorname*{Eqs}f}\left(
-1\right)  ^{\left\vert B\right\vert }\prod_{\substack{K\in\mathfrak{K}%
;\\K\subseteq B}}a_{K}$ can be split into two sums as follows:
\begin{align}
&  \sum_{B\subseteq E\cap\operatorname*{Eqs}f}\left(  -1\right)  ^{\left\vert
B\right\vert }\prod_{\substack{K\in\mathfrak{K};\\K\subseteq B}}a_{K}%
\nonumber\\
&  =\sum_{B\in\mathcal{U}}\underbrace{\left(  -1\right)  ^{\left\vert
B\right\vert }}_{\substack{=-\left(  -1\right)  ^{\left\vert \Phi\left(
B\right)  \right\vert }\\\text{(by (\ref{pf.lem.NBCm.moeb.short.Phi.-1}))}%
}}\underbrace{\prod_{\substack{K\in\mathfrak{K};\\K\subseteq B}}}%
_{\substack{=\prod_{\substack{K\in\mathfrak{K};\\K\subseteq\Phi\left(
B\right)  }}\\\text{(because of the equivalence
(\ref{pf.lem.NBCm.moeb.short.Phi.equiv}))}}}a_{K}+\underbrace{\sum
_{B\in\mathcal{V}}\left(  -1\right)  ^{\left\vert B\right\vert }%
\prod_{\substack{K\in\mathfrak{K};\\K\subseteq B}}a_{K}}_{\substack{=\sum
_{B\in\mathcal{U}}\left(  -1\right)  ^{\left\vert \Phi\left(  B\right)
\right\vert }\prod_{\substack{K\in\mathfrak{K};\\K\subseteq\Phi\left(
B\right)  }}a_{K}\\\text{(here, we have substituted }\Phi\left(  B\right)
\text{ for }B\text{ in the sum,}\\\text{since the map }\Phi:\mathcal{U}%
\rightarrow\mathcal{V}\text{ is a bijection)}}}\nonumber\\
&  =\sum_{B\in\mathcal{U}}\left(  -\left(  -1\right)  ^{\left\vert \Phi\left(
B\right)  \right\vert }\right)  \prod_{\substack{K\in\mathfrak{K}%
;\\K\subseteq\Phi\left(  B\right)  }}a_{K}+\sum_{B\in\mathcal{U}}\left(
-1\right)  ^{\left\vert \Phi\left(  B\right)  \right\vert }\prod
_{\substack{K\in\mathfrak{K};\\K\subseteq\Phi\left(  B\right)  }%
}a_{K}\nonumber\\
&  =-\sum_{B\in\mathcal{U}}\left(  -1\right)  ^{\left\vert \Phi\left(
B\right)  \right\vert }\prod_{\substack{K\in\mathfrak{K};\\K\subseteq
\Phi\left(  B\right)  }}a_{K}+\sum_{B\in\mathcal{U}}\left(  -1\right)
^{\left\vert \Phi\left(  B\right)  \right\vert }\prod_{\substack{K\in
\mathfrak{K};\\K\subseteq\Phi\left(  B\right)  }}a_{K}\nonumber\\
&  =0=\left[  E\cap\operatorname*{Eqs}f=\varnothing\right]
\ \ \ \ \ \ \ \ \ \ \left(  \text{since }\left[  E\cap\operatorname*{Eqs}%
f=\varnothing\right]  =0\right)  . \label{pf.lem.NBCm.moeb.there}%
\end{align}
This proves Lemma \ref{lem.NBCm.moeb}.
\end{proof}
\end{vershort}

\begin{verlong}
\begin{proof}
[Proof of Lemma \ref{lem.NBCm.moeb}.]If $E\cap\operatorname*{Eqs}%
f=\varnothing$, then Lemma \ref{lem.NBCm.moeb} holds\footnote{\textit{Proof.}
Assume that $E\cap\operatorname*{Eqs}f=\varnothing$. We need to check that
Lemma \ref{lem.NBCm.moeb} holds.
\par
We have%
\[
\sum_{B\subseteq E\cap\operatorname*{Eqs}f}\left(  -1\right)  ^{\left\vert
B\right\vert }\prod_{\substack{K\in\mathfrak{K};\\K\subseteq B}}a_{K}%
=\sum_{B\subseteq\varnothing}\left(  -1\right)  ^{\left\vert B\right\vert
}\prod_{\substack{K\in\mathfrak{K};\\K\subseteq B}}a_{K}%
\ \ \ \ \ \ \ \ \ \ \left(  \text{since }E\cap\operatorname*{Eqs}%
f=\varnothing\right)  .
\]
But the only subset $B$ of $\varnothing$ is the set $\varnothing$. Thus, the
only addend of the sum $\sum_{B\subseteq\varnothing}\left(  -1\right)
^{\left\vert B\right\vert }\prod_{\substack{K\in\mathfrak{K};\\K\subseteq
B}}a_{K}$ is the addend for $B=\varnothing$. Hence,%
\[
\sum_{B\subseteq\varnothing}\left(  -1\right)  ^{\left\vert B\right\vert
}\prod_{\substack{K\in\mathfrak{K};\\K\subseteq B}}a_{K}=\underbrace{\left(
-1\right)  ^{\left\vert \varnothing\right\vert }}_{\substack{=1\\\text{(since
}\left\vert \varnothing\right\vert =0\text{ is even)}}}\prod_{\substack{K\in
\mathfrak{K};\\K\subseteq\varnothing}}a_{K}=\prod_{\substack{K\in
\mathfrak{K};\\K\subseteq\varnothing}}a_{K}.
\]
\par
But let $K\in\mathfrak{K}$ be such that $K\subseteq\varnothing$. Then, $K$ is
an element of $\mathfrak{K}$, and thus a broken circuit of $G$ (since
$\mathfrak{K}$ is a set of broken circuits of $G$). Hence, Lemma
\ref{lem.BC.nonempty} shows that $K\neq\varnothing$. But from $K\subseteq
\varnothing$, we obtain $K=\varnothing$; this contradicts $K\neq\varnothing$.
\par
Now, let us forget that we fixed $K$. Thus, we have obtained a contradiction
for every $K\in\mathfrak{K}$ satisfying $K\subseteq\varnothing$. Hence, there
exists no $K\in\mathfrak{K}$ satisfying $K\subseteq\varnothing$. Therefore,
the product $\prod_{\substack{K\in\mathfrak{K};\\K\subseteq\varnothing}}a_{K}$
is empty, and thus equals $1$. In other words, $\prod_{\substack{K\in
\mathfrak{K};\\K\subseteq\varnothing}}a_{K}=1$.
\par
Now,%
\[
\sum_{B\subseteq E\cap\operatorname*{Eqs}f}\left(  -1\right)  ^{\left\vert
B\right\vert }\prod_{\substack{K\in\mathfrak{K};\\K\subseteq B}}a_{K}%
=\sum_{B\subseteq\varnothing}\left(  -1\right)  ^{\left\vert B\right\vert
}\prod_{\substack{K\in\mathfrak{K};\\K\subseteq B}}a_{K}=\prod_{\substack{K\in
\mathfrak{K};\\K\subseteq\varnothing}}a_{K}=1=\left[  E\cap\operatorname*{Eqs}%
f=\varnothing\right]
\]
(since $\left[  E\cap\operatorname*{Eqs}f=\varnothing\right]  =1$ (since
$E\cap\operatorname*{Eqs}f=\varnothing$ holds)). Thus, Lemma
\ref{lem.NBCm.moeb} holds, qed.}. Thus, we can WLOG assume that we don't have
$E\cap\operatorname*{Eqs}f=\varnothing$. Assume this.

We don't have $E\cap\operatorname*{Eqs}f=\varnothing$. Thus, we have $\left[
E\cap\operatorname*{Eqs}f=\varnothing\right]  =0$.

We know that $V$ is finite (since the graph $\left(  V,E\right)  $ is finite).
Thus, $E$ is finite (since $E\subseteq\dbinom{V}{2}$), and therefore the set
$E\cap\operatorname*{Eqs}f$ is also finite.

We have $E\cap\operatorname*{Eqs}f\neq\varnothing$ (since we don't have
$E\cap\operatorname*{Eqs}f=\varnothing$). Thus, $E\cap\operatorname*{Eqs}f$ is
a nonempty finite set. Hence, there exists some $d\in E\cap\operatorname*{Eqs}%
f$ with maximum $\ell\left(  d\right)  $ (among all $d\in E\cap
\operatorname*{Eqs}f$). Pick such a $d$. (If there are several such $d$, then
it does not matter which one we pick.)

We have chosen $d$ to be the element of $E\cap\operatorname*{Eqs}f$ with
maximum $\ell\left(  d\right)  $ (among all $d\in E\cap\operatorname*{Eqs}f$).
Thus,%
\begin{equation}
\ell\left(  d\right)  \geq\ell\left(  e\right)  \ \ \ \ \ \ \ \ \ \ \text{for
every }e\in E\cap\operatorname*{Eqs}f. \label{pf.lem.NBCm.moeb.maxl}%
\end{equation}


As usual, we let $\mathcal{P}\left(  S\right)  $ denote the powerset of any
set $S$. We now define two subsets $\mathcal{U}$ and $\mathcal{V}$ of
$\mathcal{P}\left(  E\cap\operatorname*{Eqs}f\right)  $ as follows:%
\begin{align*}
\mathcal{U}  &  =\left\{  F\in\mathcal{P}\left(  E\cap\operatorname*{Eqs}%
f\right)  \ \mid\ d\notin F\right\}  ;\\
\mathcal{V}  &  =\left\{  F\in\mathcal{P}\left(  E\cap\operatorname*{Eqs}%
f\right)  \ \mid\ d\in F\right\}  .
\end{align*}
Every $B\in\mathcal{U}$ satisfies $B\cup\left\{  d\right\}  \in\mathcal{V}%
$\ \ \ \ \footnote{\textit{Proof.} Let $B\in\mathcal{U}$. Thus, $B\in
\mathcal{U}=\left\{  F\in\mathcal{P}\left(  E\cap\operatorname*{Eqs}f\right)
\ \mid\ d\notin F\right\}  $. In other words, $B$ is an element $F$ of
$\mathcal{P}\left(  E\cap\operatorname*{Eqs}f\right)  $ satisfying $d\notin
F$. In other words, $B$ is an element of $\mathcal{P}\left(  E\cap
\operatorname*{Eqs}f\right)  $ and satisfies $d\notin B$. We have
$B\in\mathcal{P}\left(  E\cap\operatorname*{Eqs}f\right)  $; in other words,
$B$ is a subset of $E\cap\operatorname*{Eqs}f$. Also, $\left\{  d\right\}
\subseteq E\cap\operatorname*{Eqs}f$ (since $d\in E\cap\operatorname*{Eqs}f$).
Thus, both $B$ and $\left\{  d\right\}  $ are subsets of $E\cap
\operatorname*{Eqs}f$. Hence, $B\cup\left\{  d\right\}  $ is a subset of
$E\cap\operatorname*{Eqs}f$. In other words, $B\cup\left\{  d\right\}
\in\mathcal{P}\left(  E\cap\operatorname*{Eqs}f\right)  $. Also, $d\in\left\{
d\right\}  \subseteq B\cup\left\{  d\right\}  $. Hence, $B\cup\left\{
d\right\}  $ is an element of $\mathcal{P}\left(  E\cap\operatorname*{Eqs}%
f\right)  $ and satisfies $d\in B\cup\left\{  d\right\}  $. In other words,
$B\cup\left\{  d\right\}  $ is an element $F$ of $\mathcal{P}\left(
E\cap\operatorname*{Eqs}f\right)  $ satisfying $d\in F$. In other words,
$B\cup\left\{  d\right\}  \in\left\{  F\in\mathcal{P}\left(  E\cap
\operatorname*{Eqs}f\right)  \ \mid\ d\in F\right\}  =\mathcal{V}$, qed.}.
Thus, we can define a map $\Phi:\mathcal{U}\rightarrow\mathcal{V}$ by%
\[
\left(  \Phi\left(  B\right)  =B\cup\left\{  d\right\}
\ \ \ \ \ \ \ \ \ \ \text{for every }B\in\mathcal{U}\right)  .
\]
Consider this map $\Phi$.

Every $B\in\mathcal{V}$ satisfies $B\setminus\left\{  d\right\}
\in\mathcal{U}$\ \ \ \ \footnote{\textit{Proof.} Let $B\in\mathcal{V}$. Thus,
$B\in\mathcal{V}=\left\{  F\in\mathcal{P}\left(  E\cap\operatorname*{Eqs}%
f\right)  \ \mid\ d\in F\right\}  $. In other words, $B$ is an element $F$ of
$\mathcal{P}\left(  E\cap\operatorname*{Eqs}f\right)  $ satisfying $d\in F$.
In other words, $B$ is an element of $\mathcal{P}\left(  E\cap
\operatorname*{Eqs}f\right)  $ and satisfies $d\in B$. We have $B\in
\mathcal{P}\left(  E\cap\operatorname*{Eqs}f\right)  $; in other words, $B$ is
a subset of $E\cap\operatorname*{Eqs}f$. Hence, $B\setminus\left\{  d\right\}
$ is a subset of $E\cap\operatorname*{Eqs}f$ (since $B\setminus\left\{
d\right\}  \subseteq B$). In other words, $B\setminus\left\{  d\right\}
\in\mathcal{P}\left(  E\cap\operatorname*{Eqs}f\right)  $. Also, $d\notin
B\cup\left\{  d\right\}  $ (since $d\in\left\{  d\right\}  $). Hence,
$B\setminus\left\{  d\right\}  $ is an element of $\mathcal{P}\left(
E\cap\operatorname*{Eqs}f\right)  $ and satisfies $d\notin B\setminus\left\{
d\right\}  $. In other words, $B\setminus\left\{  d\right\}  $ is an element
$F$ of $\mathcal{P}\left(  E\cap\operatorname*{Eqs}f\right)  $ satisfying
$d\notin F$. In other words, $B\setminus\left\{  d\right\}  \in\left\{
F\in\mathcal{P}\left(  E\cap\operatorname*{Eqs}f\right)  \ \mid\ d\notin
F\right\}  =\mathcal{U}$, qed.}. Thus, we can define a map $\Psi
:\mathcal{V}\rightarrow\mathcal{U}$ by%
\[
\left(  \Psi\left(  B\right)  =B\setminus\left\{  d\right\}
\ \ \ \ \ \ \ \ \ \ \text{for every }B\in\mathcal{V}\right)  .
\]
Consider this map $\Psi$.

We have $\Phi\circ\Psi=\operatorname*{id}$\ \ \ \ \footnote{\textit{Proof.}
Let $B\in\mathcal{V}$. We have%
\[
\left(  \Phi\circ\Psi\right)  \left(  B\right)  =\Phi\left(  \underbrace{\Psi
\left(  B\right)  }_{\substack{=B\setminus\left\{  d\right\}  \\\text{(by the
definition of }\Psi\text{)}}}\right)  =\Phi\left(  B\setminus\left\{
d\right\}  \right)  =\left(  B\setminus\left\{  d\right\}  \right)
\cup\left\{  d\right\}
\]
(by the definition of $\Phi$).
\par
We have $B\in\mathcal{V}=\left\{  F\in\mathcal{P}\left(  E\cap
\operatorname*{Eqs}f\right)  \ \mid\ d\in F\right\}  $. In other words, $B$ is
an element $F$ of $\mathcal{P}\left(  E\cap\operatorname*{Eqs}f\right)  $
satisfying $d\in F$. In other words, $B$ is an element of $\mathcal{P}\left(
E\cap\operatorname*{Eqs}f\right)  $ and satisfies $d\in B$. From $d\in B$, we
obtain $\left\{  d\right\}  \subseteq B$. Now, $\left(  \Phi\circ\Psi\right)
\left(  B\right)  =\left(  B\setminus\left\{  d\right\}  \right)  \cup\left\{
d\right\}  =B$ (since $\left\{  d\right\}  \subseteq B$). Thus, $\left(
\Phi\circ\Psi\right)  \left(  B\right)  =B=\operatorname*{id}\left(  B\right)
$.
\par
Now, let us forget that we fixed $B$. We thus have proven that $\left(
\Phi\circ\Psi\right)  \left(  B\right)  =\operatorname*{id}\left(  B\right)  $
for every $B\in\mathcal{V}$. In other words, $\Phi\circ\Psi=\operatorname*{id}%
$, qed.} and $\Psi\circ\Phi=\operatorname*{id}$%
\ \ \ \ \footnote{\textit{Proof.} Let $B\in\mathcal{U}$. We have%
\[
\left(  \Psi\circ\Phi\right)  \left(  B\right)  =\Psi\left(  \underbrace{\Phi
\left(  B\right)  }_{\substack{=B\cup\left\{  d\right\}  \\\text{(by the
definition of }\Phi\text{)}}}\right)  =\Psi\left(  B\cup\left\{  d\right\}
\right)  =\left(  B\cup\left\{  d\right\}  \right)  \setminus\left\{
d\right\}
\]
(by the definition of $\Psi$).
\par
We have $B\in\mathcal{U}=\left\{  F\in\mathcal{P}\left(  E\cap
\operatorname*{Eqs}f\right)  \ \mid\ d\notin F\right\}  $. In other words, $B$
is an element $F$ of $\mathcal{P}\left(  E\cap\operatorname*{Eqs}f\right)  $
satisfying $d\notin F$. In other words, $B$ is an element of $\mathcal{P}%
\left(  E\cap\operatorname*{Eqs}f\right)  $ and satisfies $d\notin B$. From
$d\notin B$, we see that the sets $\left\{  d\right\}  $ and $B$ are disjoint.
Now, $\left(  \Psi\circ\Phi\right)  \left(  B\right)  =\left(  B\cup\left\{
d\right\}  \right)  \setminus\left\{  d\right\}  =B$ (since the sets $\left\{
d\right\}  $ and $B$ are disjoint). Thus, $\left(  \Psi\circ\Phi\right)
\left(  B\right)  =B=\operatorname*{id}\left(  B\right)  $.
\par
Now, let us forget that we fixed $B$. We thus have proven that $\left(
\Psi\circ\Phi\right)  \left(  B\right)  =\operatorname*{id}\left(  B\right)  $
for every $B\in\mathcal{U}$. In other words, $\Psi\circ\Phi=\operatorname*{id}%
$, qed.}. Thus, the maps $\Phi$ and $\Psi$ are mutually inverse. Hence, the
map $\Phi$ is a bijection.

Moreover, for every $B\in\mathcal{U}$ and every $K\in\mathfrak{K}$, we have
the following logical equivalence:%
\begin{equation}
\left(  K\subseteq B\right)  \ \Longleftrightarrow\ \left(  K\subseteq
\Phi\left(  B\right)  \right)  . \label{pf.lem.NBCm.moeb.equiv}%
\end{equation}


\textit{Proof of (\ref{pf.lem.NBCm.moeb.equiv}):} Let $B\in\mathcal{U}$ and
$K\in\mathfrak{K}$. We need to prove the logical equivalence
(\ref{pf.lem.NBCm.moeb.equiv}).

The definition of $\Phi$ yields $\Phi\left(  B\right)  =B\cup\left\{
d\right\}  \supseteq B$, so that $B\subseteq\Phi\left(  B\right)  $.

We have $B\in\mathcal{U}=\left\{  F\in\mathcal{P}\left(  E\cap
\operatorname*{Eqs}f\right)  \ \mid\ d\notin F\right\}  $. In other words, $B$
is an element $F$ of $\mathcal{P}\left(  E\cap\operatorname*{Eqs}f\right)  $
satisfying $d\notin F$. In other words, $B$ is an element of $\mathcal{P}%
\left(  E\cap\operatorname*{Eqs}f\right)  $ and satisfies $d\notin B$. We have
$B\in\mathcal{P}\left(  E\cap\operatorname*{Eqs}f\right)  $; in other words,
$B$ is a subset of $E\cap\operatorname*{Eqs}f$.

Also, $\Phi\left(  B\right)  \in\mathcal{V}=\left\{  F\in\mathcal{P}\left(
E\cap\operatorname*{Eqs}f\right)  \ \mid\ d\notin F\right\}  \subseteq
\mathcal{P}\left(  E\cap\operatorname*{Eqs}f\right)  $. In other words,
$\Phi\left(  B\right)  \subseteq E\cap\operatorname*{Eqs}f$.

We have $K\in\mathfrak{K}$. Thus, $K$ is a broken circuit of $G$ (since
$\mathfrak{K}$ is a set of broken circuits of $G$). In other words, $K$ is a
subset of $E$ having the form $C\setminus\left\{  e\right\}  $, where $C$ is a
circuit of $G$, and where $e$ is the unique edge in $C$ having maximum label
(among the edges in $C$)\ \ \ \ \footnote{because a broken circuit of $G$ is
the same as a subset of $E$ having the form $C\setminus\left\{  e\right\}  $,
where $C$ is a circuit of $G$, and where $e$ is the unique edge in $C$ having
maximum label (among the edges in $C$) (by the definition of a
\textquotedblleft broken circuit\textquotedblright)}. Consider this $C$ and
this $e$. Thus, we have the following facts:

\begin{itemize}
\item The set $C$ is a circuit of $G$.

\item The element $e$ is the unique edge in $C$ having maximum label (among
the edges in $C$).

\item We have $K=C\setminus\left\{  e\right\}  $.
\end{itemize}

The element $e$ is the unique edge in $C$ having maximum label (among
the edges in $C$). Thus, the only edge in $C$ whose label is greater or equal
to the label of $e$ is $e$ itself. In other words, if $e^{\prime}$ is any edge
in $C$ satisfying $\ell\left(  e^{\prime}\right)  \geq\ell\left(  e\right)  $,
then%
\begin{equation}
e^{\prime}=e. \label{pf.lem.NBCm.moeb.equiv.pf.maxlab}%
\end{equation}


Let us now assume that $K\subseteq\Phi\left(  B\right)  $. Thus,
$K\subseteq\Phi\left(  B\right)  =B\cup\left\{  d\right\}  $ (by the
definition of $\Phi$). Hence, $\underbrace{K}_{\subseteq B\cup\left\{
d\right\}  }\setminus\left\{  d\right\}  \subseteq\left(  B\cup\left\{
d\right\}  \right)  \setminus\left\{  d\right\}  \subseteq B$.

We shall now prove that $K\subseteq B$.

Indeed, assume the contrary. Thus, $K\not \subseteq B$. If we had $d\notin K$,
then we would have $K\setminus\left\{  d\right\}  =K$ and therefore
$K=K\setminus\left\{  d\right\}  \subseteq B$; this would contradict
$K\not \subseteq B$. Hence, we cannot have $d\notin K$. We thus must have
$d\in K$. Hence, $d\in K=C\setminus\left\{  e\right\}  $. Hence, $d\in C$ and
$d\notin\left\{  e\right\}  $. From $d\notin\left\{  e\right\}  $, we obtain
$d\neq e$.

But $C\setminus\left\{  e\right\}  =K\subseteq\Phi\left(  B\right)  \subseteq
E\cap\operatorname*{Eqs}f\subseteq\operatorname*{Eqs}f$. Hence, Lemma
\ref{lem.Eqs.circuit} (applied to $Y$ instead of $X$) shows that $e\in
E\cap\operatorname*{Eqs}f$. Thus, (\ref{pf.lem.NBCm.moeb.maxl}) shows that
$\ell\left(  d\right)  \geq\ell\left(  e\right)  $.

Also, $d\in C$. In other words, $d$ is an edge in $C$. Since $\ell\left(
d\right)  \geq\ell\left(  e\right)  $, we can therefore apply
(\ref{pf.lem.NBCm.moeb.equiv.pf.maxlab}) to $e^{\prime}=d$. We thus obtain
$d=e$. This contradicts $d\neq e$. This contradiction proves that our
assumption was wrong. Hence, $K\subseteq B$ is proven.

Now, let us forget that we assumed that $K\subseteq\Phi\left(  B\right)  $. We
thus have proven that $K\subseteq B$ under the assumption that $K\subseteq
\Phi\left(  B\right)  $. In other words, we have proven the implication%
\begin{equation}
\left(  K\subseteq\Phi\left(  B\right)  \right)  \ \Longrightarrow\ \left(
K\subseteq B\right)  . \label{pf.lem.NBCm.moeb.equiv.pf.dir1}%
\end{equation}


On the other hand, if $K\subseteq B$, then $K\subseteq B\subseteq\Phi\left(
B\right)  $. Hence, the implication%
\[
\left(  K\subseteq B\right)  \ \Longrightarrow\ \left(  K\subseteq\Phi\left(
B\right)  \right)
\]
holds. Combining this implication with (\ref{pf.lem.NBCm.moeb.equiv.pf.dir1}),
we obtain the logical equivalence $\left(  K\subseteq B\right)
\ \Longleftrightarrow\ \left(  K\subseteq\Phi\left(  B\right)  \right)  $.
Thus, (\ref{pf.lem.NBCm.moeb.equiv}) is proven.

Also, every $B\in\mathcal{U}$ satisfies%
\begin{equation}
\left(  -1\right)  ^{\left\vert B\right\vert }=-\left(  -1\right)
^{\left\vert \Phi\left(  B\right)  \right\vert }
\label{pf.lem.NBCm.moeb.phisize}%
\end{equation}
\footnote{\textit{Proof of (\ref{pf.lem.NBCm.moeb.phisize}):} Let
$B\in\mathcal{U}$. We have $B\in\mathcal{U}=\left\{  F\in\mathcal{P}\left(
E\cap\operatorname*{Eqs}f\right)  \ \mid\ d\notin F\right\}  $. In other
words, $B$ is an element $F$ of $\mathcal{P}\left(  E\cap\operatorname*{Eqs}%
f\right)  $ satisfying $d\notin F$. In other words, $B$ is an element of
$\mathcal{P}\left(  E\cap\operatorname*{Eqs}f\right)  $ and satisfies $d\notin
B$. From $d\notin B$, we see that $\left\vert B\cup\left\{  d\right\}
\right\vert =\left\vert B\right\vert +1$. Now, the definition of $\Phi$ yields
$\Phi\left(  B\right)  =B\cup\left\{  d\right\}  $. Hence, $\left\vert
\Phi\left(  B\right)  \right\vert =\left\vert B\cup\left\{  d\right\}
\right\vert =\left\vert B\right\vert +1$. Thus, $\left(  -1\right)
^{\left\vert \Phi\left(  B\right)  \right\vert }=\left(  -1\right)
^{\left\vert B\right\vert +1}=-\left(  -1\right)  ^{\left\vert B\right\vert }%
$. Therefore, $\left(  -1\right)  ^{\left\vert B\right\vert }=-\left(
-1\right)  ^{\left\vert \Phi\left(  B\right)  \right\vert }$. This proves
(\ref{pf.lem.NBCm.moeb.phisize}).}.

Now,%
\begin{align*}
&  \sum_{B\subseteq E\cap\operatorname*{Eqs}f}\left(  -1\right)  ^{\left\vert
B\right\vert }\prod_{\substack{K\in\mathfrak{K};\\K\subseteq B}}a_{K}\\
&  =\underbrace{\sum_{\substack{B\subseteq E\cap\operatorname*{Eqs}f;\\d\in
B}}}_{\substack{=\sum_{B\in\left\{  F\in\mathcal{P}\left(  E\cap
\operatorname*{Eqs}f\right)  \ \mid\ d\in F\right\}  }=\sum_{B\in\mathcal{V}%
}\\\text{(since }\left\{  F\in\mathcal{P}\left(  E\cap\operatorname*{Eqs}%
f\right)  \ \mid\ d\in F\right\}  =\mathcal{V}\text{)}}}\left(  -1\right)
^{\left\vert B\right\vert }\prod_{\substack{K\in\mathfrak{K};\\K\subseteq
B}}a_{K}+\underbrace{\sum_{\substack{B\subseteq E\cap\operatorname*{Eqs}%
f;\\d\notin B}}}_{\substack{=\sum_{B\in\left\{  F\in\mathcal{P}\left(
E\cap\operatorname*{Eqs}f\right)  \ \mid\ d\notin F\right\}  }=\sum
_{B\in\mathcal{U}}\\\text{(since }\left\{  F\in\mathcal{P}\left(
E\cap\operatorname*{Eqs}f\right)  \ \mid\ d\notin F\right\}  =\mathcal{U}%
\text{)}}}\left(  -1\right)  ^{\left\vert B\right\vert }\prod_{\substack{K\in
\mathfrak{K};\\K\subseteq B}}a_{K}\\
&  =\underbrace{\sum_{B\in\mathcal{V}}\left(  -1\right)  ^{\left\vert
B\right\vert }\prod_{\substack{K\in\mathfrak{K};\\K\subseteq B}}a_{K}%
}_{\substack{=\sum_{B\in\mathcal{U}}\left(  -1\right)  ^{\left\vert
\Phi\left(  B\right)  \right\vert }\prod_{\substack{K\in\mathfrak{K}%
;\\K\subseteq\Phi\left(  B\right)  }}a_{K}\\\text{(here, we have substituted
}\Phi\left(  B\right)  \text{ for }B\text{ in the sum,}\\\text{since the map
}\Phi:\mathcal{U}\rightarrow\mathcal{V}\text{ is a bijection)}}}+\sum
_{B\in\mathcal{U}}\underbrace{\left(  -1\right)  ^{\left\vert B\right\vert }%
}_{\substack{=-\left(  -1\right)  ^{\left\vert \Phi\left(  B\right)
\right\vert }\\\text{(by (\ref{pf.lem.NBCm.moeb.phisize}))}}}\underbrace{\prod
_{\substack{K\in\mathfrak{K};\\K\subseteq B}}}_{\substack{=\prod
_{\substack{K\in\mathfrak{K};\\K\subseteq\Phi\left(  B\right)  }%
}\\\text{(because for every }K\in\mathfrak{K}\text{, the condition}\\\left(
K\subseteq B\right)  \text{ is equivalent to }\left(  K\subseteq\Phi\left(
B\right)  \right)  \\\text{(by (\ref{pf.lem.NBCm.moeb.equiv})))}}}a_{K}\\
&  =\sum_{B\in\mathcal{U}}\left(  -1\right)  ^{\left\vert \Phi\left(
B\right)  \right\vert }\prod_{\substack{K\in\mathfrak{K};\\K\subseteq
\Phi\left(  B\right)  }}a_{K}+\sum_{B\in\mathcal{U}}\left(  -\left(
-1\right)  ^{\left\vert \Phi\left(  B\right)  \right\vert }\right)
\prod_{\substack{K\in\mathfrak{K};\\K\subseteq\Phi\left(  B\right)  }}a_{K}\\
&  =\sum_{B\in\mathcal{U}}\left(  -1\right)  ^{\left\vert \Phi\left(
B\right)  \right\vert }\prod_{\substack{K\in\mathfrak{K};\\K\subseteq
\Phi\left(  B\right)  }}a_{K}-\sum_{B\in\mathcal{U}}\left(  -1\right)
^{\left\vert \Phi\left(  B\right)  \right\vert }\prod_{\substack{K\in
\mathfrak{K};\\K\subseteq\Phi\left(  B\right)  }}a_{K}\\
&  =0=\left[  E\cap\operatorname*{Eqs}f=\varnothing\right]
\ \ \ \ \ \ \ \ \ \ \left(  \text{since }\left[  E\cap\operatorname*{Eqs}%
f=\varnothing\right]  =0\right)  .
\end{align*}
This proves Lemma \ref{lem.NBCm.moeb}.
\end{proof}
\end{verlong}

We now finally proceed to the proof of Theorem \ref{thm.chromsym.varis}:

\begin{vershort}
\begin{proof}
[Proof of Theorem \ref{thm.chromsym.varis}.]The definition of $X_{G}$ shows
that
\begin{align*}
X_{G}  &  =\sum_{\substack{f:V\rightarrow\mathbb{N}_{+}\text{ is
a}\\\text{proper }\mathbb{N}_{+}\text{-coloring of }G}}\mathbf{x}_{f}\\
&  =\sum_{f:V\rightarrow\mathbb{N}_{+}}\left[  \underbrace{f\text{ is a proper
}\mathbb{N}_{+}\text{-coloring of }G}_{\substack{\Longleftrightarrow\ \left(
\text{the }\mathbb{N}_{+}\text{-coloring }f\text{ of }G\text{ is
proper}\right)  \\\Longleftrightarrow\ \left(  E\cap\operatorname*{Eqs}%
f=\varnothing\right)  \\\text{(by Lemma \ref{lem.Eqs.proper}, applied to
}\mathbb{N}_{+}\text{ instead of }X\text{)}}}\right]  \mathbf{x}_{f}\\
&  =\sum_{f:V\rightarrow\mathbb{N}_{+}}\underbrace{\left[  E\cap
\operatorname*{Eqs}f=\varnothing\right]  }_{\substack{=\sum_{B\subseteq
E\cap\operatorname*{Eqs}f}\left(  -1\right)  ^{\left\vert B\right\vert }%
\prod_{\substack{K\in\mathfrak{K};\\K\subseteq B}}a_{K}\\\text{(by Lemma
\ref{lem.NBCm.moeb}, applied to }Y=\mathbb{N}_{+}\text{)}}}\mathbf{x}_{f}\\
&  =\sum_{f:V\rightarrow\mathbb{N}_{+}}\underbrace{\sum_{B\subseteq
E\cap\operatorname*{Eqs}f}}_{=\sum_{\substack{B\subseteq E;\\B\subseteq
\operatorname*{Eqs}f}}}\left(  -1\right)  ^{\left\vert B\right\vert }\left(
\prod_{\substack{K\in\mathfrak{K};\\K\subseteq B}}a_{K}\right)  \mathbf{x}%
_{f}=\underbrace{\sum_{f:V\rightarrow\mathbb{N}_{+}}\sum_{\substack{B\subseteq
E;\\B\subseteq\operatorname*{Eqs}f}}}_{=\sum_{B\subseteq E}\sum
_{\substack{f:V\rightarrow\mathbb{N}_{+};\\B\subseteq\operatorname*{Eqs}f}%
}}\left(  -1\right)  ^{\left\vert B\right\vert }\left(  \prod_{\substack{K\in
\mathfrak{K};\\K\subseteq B}}a_{K}\right)  \mathbf{x}_{f}\\
&  =\sum_{B\subseteq E}\sum_{\substack{f:V\rightarrow\mathbb{N}_{+}%
;\\B\subseteq\operatorname*{Eqs}f}}\left(  -1\right)  ^{\left\vert
B\right\vert }\left(  \prod_{\substack{K\in\mathfrak{K};\\K\subseteq B}%
}a_{K}\right)  \mathbf{x}_{f}=\sum_{B\subseteq E}\left(  -1\right)
^{\left\vert B\right\vert }\left(  \prod_{\substack{K\in\mathfrak{K}%
;\\K\subseteq B}}a_{K}\right)  \underbrace{\sum_{\substack{f:V\rightarrow
\mathbb{N}_{+};\\B\subseteq\operatorname*{Eqs}f}}\mathbf{x}_{f}}%
_{\substack{=p_{\lambda\left(  V,B\right)  }\\\text{(by Lemma
\ref{lem.Eqs.sum})}}}\\
&  =\sum_{B\subseteq E}\left(  -1\right)  ^{\left\vert B\right\vert }\left(
\prod_{\substack{K\in\mathfrak{K};\\K\subseteq B}}a_{K}\right)  p_{\lambda
\left(  V,B\right)  }=\sum_{F\subseteq E}\left(  -1\right)  ^{\left\vert
F\right\vert }\left(  \prod_{\substack{K\in\mathfrak{K};\\K\subseteq F}%
}a_{K}\right)  p_{\lambda\left(  V,F\right)  }%
\end{align*}
(here, we have renamed the summation index $B$ as $F$). This proves Theorem
\ref{thm.chromsym.varis}.
\end{proof}
\end{vershort}

\begin{verlong}
\begin{proof}
[Proof of Theorem \ref{thm.chromsym.varis}.]We have
\begin{equation}
X_{G}=\sum_{\substack{f:V\rightarrow\mathbb{N}_{+}\text{ is a}\\\text{proper
}\mathbb{N}_{+}\text{-coloring of }G}}\mathbf{x}_{f}
\label{pf.thm.chromsym.varis.XG-def}%
\end{equation}
(by the definition of $X_{G}$). Now, if $f:V\rightarrow\mathbb{N}_{+}$ is a
map, then we have the following logical equivalence:%
\begin{equation}
\left(  \text{the }\mathbb{N}_{+}\text{-coloring }f\text{ of }G\text{ is
proper}\right)  \ \Longleftrightarrow\ \left(  E\cap\operatorname*{Eqs}%
f=\varnothing\right)  \label{pf.thm.chromsym.varis.equiv}%
\end{equation}
(because the $\mathbb{N}_{+}$-coloring $f$ of $G$ is proper if and only if
$E\cap\operatorname*{Eqs}f=\varnothing$\ \ \ \ \footnote{by Lemma
\ref{lem.Eqs.proper} (applied to $\mathbb{N}_{+}$ instead of $X$)}). Now,%
\begin{align}
&  \sum_{f:V\rightarrow\mathbb{N}_{+}}\left[  \underbrace{E\cap
\operatorname*{Eqs}f=\varnothing}_{\substack{\Longleftrightarrow\ \left(
\text{the }\mathbb{N}_{+}\text{-coloring }f\text{ of }G\text{ is
proper}\right)  \\\text{(by (\ref{pf.thm.chromsym.varis.equiv}))}}}\right]
\mathbf{x}_{f}\nonumber\\
&  =\sum_{f:V\rightarrow\mathbb{N}_{+}}\left[  \underbrace{\text{the
}\mathbb{N}_{+}\text{-coloring }f\text{ of }G\text{ is proper}}%
_{\Longleftrightarrow\ \left(  f\text{ is a proper }\mathbb{N}_{+}%
\text{-coloring of }G\right)  }\right]  \mathbf{x}_{f}\nonumber\\
&  =\sum_{f:V\rightarrow\mathbb{N}_{+}}\left[  f\text{ is a proper }%
\mathbb{N}_{+}\text{-coloring of }G\right]  \mathbf{x}_{f}\nonumber\\
&  =\sum_{\substack{f:V\rightarrow\mathbb{N}_{+}\text{ is a}\\\text{proper
}\mathbb{N}_{+}\text{-coloring of }G}}\underbrace{\left[  f\text{ is a proper
}\mathbb{N}_{+}\text{-coloring of }G\right]  }_{\substack{=1\\\text{(since
}f\text{ is a proper }\mathbb{N}_{+}\text{-coloring of }G\text{)}}%
}\mathbf{x}_{f}\nonumber\\
&  \ \ \ \ \ \ \ \ \ \ +\sum_{\substack{f:V\rightarrow\mathbb{N}_{+}\text{ is
not a}\\\text{proper }\mathbb{N}_{+}\text{-coloring of }G}}\underbrace{\left[
f\text{ is a proper }\mathbb{N}_{+}\text{-coloring of }G\right]
}_{\substack{=0\\\text{(since }f\text{ is not a proper }\mathbb{N}%
_{+}\text{-coloring of }G\text{)}}}\mathbf{x}_{f}\nonumber\\
&  =\sum_{\substack{f:V\rightarrow\mathbb{N}_{+}\text{ is a}\\\text{proper
}\mathbb{N}_{+}\text{-coloring of }G}}\mathbf{x}_{f}+\underbrace{\sum
_{\substack{f:V\rightarrow\mathbb{N}_{+}\text{ is not a}\\\text{proper
}\mathbb{N}_{+}\text{-coloring of }G}}0\mathbf{x}_{f}}_{=0}=\sum
_{\substack{f:V\rightarrow\mathbb{N}_{+}\text{ is a}\\\text{proper }%
\mathbb{N}_{+}\text{-coloring of }G}}\mathbf{x}_{f}=X_{G}
\label{pf.thm.chromsym.varis.step1}%
\end{align}
(by (\ref{pf.thm.chromsym.varis.XG-def})).

But for every $f:V\rightarrow\mathbb{N}_{+}$, we have%
\begin{equation}
\sum_{B\subseteq E\cap\operatorname*{Eqs}f}\left(  -1\right)  ^{\left\vert
B\right\vert }\prod_{\substack{K\in\mathfrak{K};\\K\subseteq B}}a_{K}=\left[
E\cap\operatorname*{Eqs}f=\varnothing\right]
\label{pf.thm.chromsym.varis.moeb}%
\end{equation}
(by Lemma \ref{lem.NBCm.moeb} (applied to $\mathbb{N}_{+}$ instead of $Y$)).

For every $f:V\rightarrow\mathbb{N}_{+}$, we have%
\begin{equation}
\left\{  F\subseteq E\ \mid\ F\subseteq\operatorname*{Eqs}f\right\}
=\mathcal{P}\left(  E\cap\operatorname*{Eqs}f\right)
\label{pf.thm.chromsym.varis.cut}%
\end{equation}
\footnote{\textit{Proof of (\ref{pf.thm.chromsym.varis.cut}):} Let
$f:V\rightarrow\mathbb{N}_{+}$.
\par
Let $B\in\left\{  F\subseteq E\ \mid\ F\subseteq\operatorname*{Eqs}f\right\}
$. Thus, $B$ is a subset $F$ of $E$ satisfying $F\subseteq\operatorname*{Eqs}%
f$. In other words, $B$ is a subset of $E$ and satisfies $B\subseteq
\operatorname*{Eqs}f$. Since $B$ is a subset of $E$, we have $B\subseteq E$.
Combining this with $B\subseteq\operatorname*{Eqs}f$, we obtain $B\subseteq
E\cap\operatorname*{Eqs}f$. In other words, $B\in\mathcal{P}\left(
E\cap\operatorname*{Eqs}f\right)  $.
\par
Let us now forget that we fixed $B$. We thus have proven that every
$B\in\left\{  F\subseteq E\ \mid\ F\subseteq\operatorname*{Eqs}f\right\}  $
satisfies $B\in\mathcal{P}\left(  E\cap\operatorname*{Eqs}f\right)  $. In
other words,%
\begin{equation}
\left\{  F\subseteq E\ \mid\ F\subseteq\operatorname*{Eqs}f\right\}
\subseteq\mathcal{P}\left(  E\cap\operatorname*{Eqs}f\right)  .
\label{pf.thm.chromsym.varis.cut.pf.1}%
\end{equation}
\par
On the other hand, let $C\in\mathcal{P}\left(  E\cap\operatorname*{Eqs}%
f\right)  $. Thus, $C$ is a subset of $E\cap\operatorname*{Eqs}f$. Hence,
$C\subseteq E\cap\operatorname*{Eqs}f\subseteq E$, so that $C$ is a subset of
$E$. Also, $C\subseteq E\cap\operatorname*{Eqs}f\subseteq\operatorname*{Eqs}%
f$. Thus, $C$ is a subset of $E$ and satisfies $C\subseteq\operatorname*{Eqs}%
f$. In other words, $C$ is a subset $F$ of $E$ satisfying $F\subseteq
\operatorname*{Eqs}f$. In other words, $C\in\left\{  F\subseteq E\ \mid
\ F\subseteq\operatorname*{Eqs}f\right\}  $.
\par
Let us now forget that we fixed $C$. We thus have proven that every
$C\in\mathcal{P}\left(  E\cap\operatorname*{Eqs}f\right)  $ satisfies
$C\in\left\{  F\subseteq E\ \mid\ F\subseteq\operatorname*{Eqs}f\right\}  $.
In other words,%
\[
\mathcal{P}\left(  E\cap\operatorname*{Eqs}f\right)  \subseteq\left\{
F\subseteq E\ \mid\ F\subseteq\operatorname*{Eqs}f\right\}  .
\]
Combining this inclusion with (\ref{pf.thm.chromsym.varis.cut.pf.1}), we
obtain $\left\{  F\subseteq E\ \mid\ F\subseteq\operatorname*{Eqs}f\right\}
=\mathcal{P}\left(  E\cap\operatorname*{Eqs}f\right)  $. This proves
(\ref{pf.thm.chromsym.varis.cut}).}.

For every $f:V\rightarrow\mathbb{N}_{+}$, we have%
\begin{align}
&  \underbrace{\sum_{\substack{B\subseteq E;\\B\subseteq\operatorname*{Eqs}%
f}}}_{\substack{=\sum_{B\in\left\{  F\subseteq E\ \mid\ F\subseteq
\operatorname*{Eqs}f\right\}  }=\sum_{B\in\mathcal{P}\left(  E\cap
\operatorname*{Eqs}f\right)  }\\\text{(because }\left\{  F\subseteq
E\ \mid\ F\subseteq\operatorname*{Eqs}f\right\}  =\mathcal{P}\left(
E\cap\operatorname*{Eqs}f\right)  \\\text{(by (\ref{pf.thm.chromsym.varis.cut}%
)))}}}\left(  -1\right)  ^{\left\vert B\right\vert }\prod_{\substack{K\in
\mathfrak{K};\\K\subseteq B}}a_{K}\nonumber\\
&  =\underbrace{\sum_{B\in\mathcal{P}\left(  E\cap\operatorname*{Eqs}f\right)
}}_{=\sum_{B\subseteq E\cap\operatorname*{Eqs}f}}\left(  -1\right)
^{\left\vert B\right\vert }\prod_{\substack{K\in\mathfrak{K};\\K\subseteq
B}}a_{K}=\sum_{B\subseteq E\cap\operatorname*{Eqs}f}\left(  -1\right)
^{\left\vert B\right\vert }\prod_{\substack{K\in\mathfrak{K};\\K\subseteq
B}}a_{K}\nonumber\\
&  =\left[  E\cap\operatorname*{Eqs}f=\varnothing\right]
\label{pf.thm.chromsym.varis.moeb2}%
\end{align}
(by (\ref{pf.thm.chromsym.varis.moeb})).

Now, (\ref{pf.thm.chromsym.varis.step1}) yields%
\begin{align*}
X_{G}  &  =\sum_{f:V\rightarrow\mathbb{N}_{+}}\underbrace{\left[
E\cap\operatorname*{Eqs}f=\varnothing\right]  }_{\substack{=\sum
_{\substack{B\subseteq E;\\B\subseteq\operatorname*{Eqs}f}}\left(  -1\right)
^{\left\vert B\right\vert }\prod_{\substack{K\in\mathfrak{K};\\K\subseteq
B}}a_{K}\\\text{(by (\ref{pf.thm.chromsym.varis.moeb2}))}}}\mathbf{x}_{f}\\
&  =\sum_{f:V\rightarrow\mathbb{N}_{+}}\left(  \sum_{\substack{B\subseteq
E;\\B\subseteq\operatorname*{Eqs}f}}\left(  -1\right)  ^{\left\vert
B\right\vert }\prod_{\substack{K\in\mathfrak{K};\\K\subseteq B}}a_{K}\right)
\mathbf{x}_{f}=\underbrace{\sum_{f:V\rightarrow\mathbb{N}_{+}}\sum
_{\substack{B\subseteq E;\\B\subseteq\operatorname*{Eqs}f}}}_{=\sum
_{B\subseteq E}\sum_{\substack{f:V\rightarrow\mathbb{N}_{+};\\B\subseteq
\operatorname*{Eqs}f}}}\left(  -1\right)  ^{\left\vert B\right\vert }\left(
\prod_{\substack{K\in\mathfrak{K};\\K\subseteq B}}a_{K}\right)  \mathbf{x}%
_{f}\\
&  =\sum_{B\subseteq E}\sum_{\substack{f:V\rightarrow\mathbb{N}_{+}%
;\\B\subseteq\operatorname*{Eqs}f}}\left(  -1\right)  ^{\left\vert
B\right\vert }\left(  \prod_{\substack{K\in\mathfrak{K};\\K\subseteq B}%
}a_{K}\right)  \mathbf{x}_{f}\\
&  =\sum_{B\subseteq E}\left(  -1\right)  ^{\left\vert B\right\vert }\left(
\prod_{\substack{K\in\mathfrak{K};\\K\subseteq B}}a_{K}\right)
\underbrace{\sum_{\substack{f:V\rightarrow\mathbb{N}_{+};\\B\subseteq
\operatorname*{Eqs}f}}\mathbf{x}_{f}}_{\substack{=p_{\lambda\left(
V,B\right)  }\\\text{(by Lemma \ref{lem.Eqs.sum}}\\\text{(since }\left(
V,B\right)  \text{ is a finite graph}\\\text{(since }V\text{ is a finite set
and }B\subseteq E\subseteq\dbinom{V}{2}\text{)))}}}\\
&  =\sum_{B\subseteq E}\left(  -1\right)  ^{\left\vert B\right\vert }\left(
\prod_{\substack{K\in\mathfrak{K};\\K\subseteq B}}a_{K}\right)  p_{\lambda
\left(  V,B\right)  }=\sum_{F\subseteq E}\left(  -1\right)  ^{\left\vert
F\right\vert }\left(  \prod_{\substack{K\in\mathfrak{K};\\K\subseteq F}%
}a_{K}\right)  p_{\lambda\left(  V,F\right)  }%
\end{align*}
(here, we have renamed the summation index $B$ as $F$). This proves Theorem
\ref{thm.chromsym.varis}.
\end{proof}
\end{verlong}

Thus, Theorem \ref{thm.chromsym.varis} is proven; as we know, this entails the
correctness of Theorem \ref{thm.chromsym.empty}, Corollary
\ref{cor.chromsym.K-free} and Corollary \ref{cor.chromsym.NBC}.

\section{The chromatic polynomial}

\subsection{Definition}

We have so far studied the chromatic symmetric function. We shall now apply
the above results to the chromatic polynomial. The definition of the chromatic
polynomial rests upon the following fact:

\begin{theorem}
\label{thm.chrompol.exist}Let $G=\left(  V,E\right)  $ be a finite graph.
Then, there exists a unique polynomial $P\in\mathbb{Z}\left[  x\right]  $ such
that every $q\in\mathbb{N}$ satisfies%
\[
P\left(  q\right)  =\left(  \text{the number of all proper }\left\{
1,2,\ldots,q\right\}  \text{-colorings of }G\right)  .
\]

\end{theorem}

\begin{definition}
\label{def.chrompol}Let $G=\left(  V,E\right)  $ be a finite graph. Theorem
\ref{thm.chrompol.exist} shows that there exists a polynomial $P\in
\mathbb{Z}\left[  x\right]  $ such that every $q\in\mathbb{N}$ satisfies
$P\left(  q\right)  =\left(  \text{the number of all proper }\left\{
1,2,\ldots,q\right\}  \text{-colorings of }G\right)  $. This polynomial $P$ is
called the \textit{chromatic polynomial} of $G$, and will be denoted by
$\chi_{G}$.
\end{definition}

We shall later prove Theorem \ref{thm.chrompol.exist} (as a consequence of
something stronger that we show). First, we shall state some formulas for the
chromatic polynomial which are analogues of results proven before for the
chromatic symmetric function.

\subsection{Formulas for $\chi_{G}$}

Before we state several formulas for $\chi_{G}$, we need to introduce one more notation:

\begin{definition}
\label{def.conn}Let $G$ be a finite graph. We let $\operatorname*{conn}G$
denote the number of connected components of $G$.
\end{definition}

The following results are analogues of Theorem \ref{thm.chromsym.empty},
Theorem \ref{thm.chromsym.varis}, Corollary \ref{cor.chromsym.K-free} and
Corollary \ref{cor.chromsym.NBC}, respectively:

\begin{theorem}
\label{thm.chrompol.empty}Let $G=\left(  V,E\right)  $ be a finite graph.
Then,%
\[
\chi_{G}=\sum_{F\subseteq E}\left(  -1\right)  ^{\left\vert F\right\vert
}x^{\operatorname*{conn}\left(  V,F\right)  }.
\]
(Here, of course, the pair $\left(  V,F\right)  $ is regarded as a graph, and
the expression $\operatorname*{conn}\left(  V,F\right)  $ is understood
according to Definition \ref{def.conn}.)
\end{theorem}

\begin{theorem}
\label{thm.chrompol.varis}Let $G=\left(  V,E\right)  $ be a finite graph. Let
$X$ be a totally ordered set. Let $\ell:E\rightarrow X$ be a function. Let
$\mathfrak{K}$ be some set of broken circuits of $G$ (not necessarily
containing all of them). Let $a_{K}$ be an element of $\mathbf{k}$ for every
$K\in\mathfrak{K}$. Then,%
\[
\chi_{G}=\sum_{F\subseteq E}\left(  -1\right)  ^{\left\vert F\right\vert
}\left(  \prod_{\substack{K\in\mathfrak{K};\\K\subseteq F}}a_{K}\right)
x^{\operatorname*{conn}\left(  V,F\right)  }.
\]
(Here, of course, the pair $\left(  V,F\right)  $ is regarded as a graph, and
the expression $\operatorname*{conn}\left(  V,F\right)  $ is understood
according to Definition \ref{def.conn}.)
\end{theorem}

\begin{corollary}
\label{cor.chrompol.K-free}Let $G=\left(  V,E\right)  $ be a finite graph. Let
$X$ be a totally ordered set. Let $\ell:E\rightarrow X$ be a function. Let
$\mathfrak{K}$ be some set of broken circuits of $G$ (not necessarily
containing all of them). Then,%
\[
\chi_{G}=\sum_{\substack{F\subseteq E;\\F\text{ is }\mathfrak{K}\text{-free}%
}}\left(  -1\right)  ^{\left\vert F\right\vert }x^{\operatorname*{conn}\left(
V,F\right)  }.
\]

\end{corollary}

\begin{corollary}
\label{cor.chrompol.NBC}Let $G=\left(  V,E\right)  $ be a finite graph. Let
$X$ be a totally ordered set. Let $\ell:E\rightarrow X$ be a function. Then,%
\[
\chi_{G}=\sum_{\substack{F\subseteq E;\\F\text{ contains no broken}%
\\\text{circuit of }G\text{ as a subset}}}\left(  -1\right)  ^{\left\vert
F\right\vert }x^{\operatorname*{conn}\left(  V,F\right)  }.
\]

\end{corollary}

\subsection{Proofs}

There are two approaches to these results: One is to derive them similarly to
how we derived the analogous results about $X_{G}$; the other is to derive
them from the latter. We shall take the first approach, since it yields a
proof of the classical Theorem \ref{thm.chrompol.exist} \textquotedblleft for
free\textquotedblright. We begin with an analogue of Lemma \ref{lem.Eqs.sum}:

\begin{lemma}
\label{lem.Eqs.sum1}Let $\left(  V,B\right)  $ be a finite graph. Let
$q\in\mathbb{N}$. Then,%
\[
\sum_{\substack{f:V\rightarrow\left\{  1,2,\ldots,q\right\}  ;\\B\subseteq
\operatorname*{Eqs}f}}1=q^{\operatorname*{conn}\left(  V,B\right)  }.
\]
(Here, the expression $\operatorname*{conn}\left(  V,B\right)  $ is understood
according to Definition \ref{def.connectedness} \textbf{(b)}.)
\end{lemma}

One way to prove Lemma \ref{lem.Eqs.sum1} is to evaluate the equality given by
Lemma \ref{lem.Eqs.sum} at $x_{k}=%
\begin{cases}
1, & \text{if }k\leq q\\
0, & \text{if }k>q
\end{cases}
$. Another proof can be obtained by mimicking our proof of Lemma
\ref{lem.Eqs.sum}:

\begin{vershort}
\begin{proof}
[Proof of Lemma \ref{lem.Eqs.sum1}.]Define $\left(  C_{1},C_{2},\ldots
,C_{k}\right)  $ as in the proof of Lemma \ref{lem.Eqs.sum}. Thus,
$\operatorname*{conn}\left(  V,B\right)  =k$. Define a map $\Phi$ as in the
proof of Lemma \ref{lem.Eqs.sum}, but with $\mathbb{N}_{+}$ replaced by
$\left\{  1,2,\ldots,q\right\}  $. Then,
\[
\Phi:\left\{  1,2,\ldots,q\right\}  ^{k}\rightarrow\left\{  f:V\rightarrow
\left\{  1,2,\ldots,q\right\}  \ \mid\ B\subseteq\operatorname*{Eqs}f\right\}
\]
is a bijection\footnote{This can be shown in the same way as for the map
$\Phi$ in the proof of Lemma \ref{lem.Eqs.sum}; we just have to replace every
$\mathbb{N}_{+}$ by $\left\{  1,2,\ldots,q\right\}  $.}. Now,%
\begin{align*}
&  \sum_{\substack{f:V\rightarrow\left\{  1,2,\ldots,q\right\}  ;\\B\subseteq
\operatorname*{Eqs}f}}1\\
&  =\sum_{\left(  s_{1},s_{2},\ldots,s_{k}\right)  \in\left\{  1,2,\ldots
,q\right\}  ^{k}}1\\
&  \ \ \ \ \ \ \ \ \ \ \left(
\begin{array}
[c]{c}%
\text{here, we have substituted }\Phi\left(  s_{1},s_{2},\ldots,s_{k}\right)
\text{ for }f\text{ in the sum,}\\
\text{since the map }\Phi:\left\{  1,2,\ldots,q\right\}  ^{k}\rightarrow
\left\{  f:V\rightarrow\left\{  1,2,\ldots,q\right\}  \ \mid\ B\subseteq
\operatorname*{Eqs}f\right\} \\
\text{is a bijection}%
\end{array}
\right) \\
&  =\left(  \text{the number of all }\left(  s_{1},s_{2},\ldots,s_{k}\right)
\in\left\{  1,2,\ldots,q\right\}  ^{k}\right) \\
&  =q^{k}=q^{\operatorname*{conn}\left(  V,B\right)  }%
\ \ \ \ \ \ \ \ \ \ \left(  \text{since }k=\operatorname*{conn}\left(
V,B\right)  \right)  .
\end{align*}
This proves Lemma \ref{lem.Eqs.sum1}.
\end{proof}
\end{vershort}

\begin{verlong}
\begin{proof}
[Proof of Lemma \ref{lem.Eqs.sum1}.]Let $\sim$ denote the equivalence relation
$\sim_{\left(  V,B\right)  }$ (defined as in Definition
\ref{def.connectedness} \textbf{(a)}). The connected components of $\left(
V,B\right)  $ are the $\sim_{\left(  V,B\right)  }$-equivalence classes
(because this is how the connected components of $\left(  V,B\right)  $ are
defined). In other words, the connected components of $\left(  V,B\right)  $
are the $\sim$-equivalence classes (since $\sim$ is the relation
$\sim_{\left(  V,B\right)  }$). In other words, the connected components of
$\left(  V,B\right)  $ are the elements of $V/\left(  \sim\right)  $ (since
the elements of $V/\left(  \sim\right)  $ are the $\sim$-equivalence classes
(by the definition of $V/\left(  \sim\right)  $)).

Let $Y$ be the set $\left\{  1,2,\ldots,q\right\}  $. Thus, $\left\vert
Y\right\vert =q$. A set $Y_{\sim}^{V}$ is defined (according to Definition
\ref{def.relquot.maps} \textbf{(b)}).

Proposition \ref{prop.relquot.uniprop} \textbf{(b)} (applied to $X=V$) shows
that the map%
\[
Y^{V/\left(  \sim\right)  }\rightarrow Y_{\sim}^{V}%
,\ \ \ \ \ \ \ \ \ \ f\mapsto f\circ\pi_{V}%
\]
is a bijection. Thus, there exists a bijection $Y^{V/\left(  \sim\right)
}\rightarrow Y_{\sim}^{V}$ (namely, this map). Hence, $\left\vert Y_{\sim}%
^{V}\right\vert =\left\vert Y^{V/\left(  \sim\right)  }\right\vert $.

For every map $f:V\rightarrow Y$, we have the following equivalence:%
\begin{equation}
\left(  B\subseteq\operatorname*{Eqs}f\right)  \ \Longleftrightarrow\ \left(
f\in Y_{\sim}^{V}\right)  \label{pf.lem.Eqs.sum1.equivalence}%
\end{equation}
(according to Lemma \ref{lem.Eqs.sum-aux}). Thus, we have the following
equality of summation signs:
\[
\sum_{\substack{f:V\rightarrow Y;\\B\subseteq\operatorname*{Eqs}f}%
}=\sum_{\substack{f:V\rightarrow Y;\\f\in Y_{\sim}^{V}}}=\sum_{\substack{f\in
Y^{V};\\f\in Y_{\sim}^{V}}}=\sum_{f\in Y_{\sim}^{V}}%
\]
(since $Y_{\sim}^{V}$ is a subset of $Y^{V}$). Hence,%
\begin{equation}
\underbrace{\sum_{\substack{f:V\rightarrow Y;\\B\subseteq\operatorname*{Eqs}%
f}}}_{=\sum_{f\in Y_{\sim}^{V}}}1=\sum_{f\in Y_{\sim}^{V}}1=\left\vert
Y_{\sim}^{V}\right\vert \cdot1=\left\vert Y_{\sim}^{V}\right\vert =\left\vert
Y^{V/\left(  \sim\right)  }\right\vert . \label{pf.lem.Eqs.sum1.1}%
\end{equation}


Now, let $\left(  C_{1},C_{2},\ldots,C_{k}\right)  $ be a list of all
connected components of $\left(  V,B\right)  $.\ \ \ \ \footnote{Every
connected component of $\left(  V,B\right)  $ should appear exactly once in
this list.} Thus, $k$ is the number of connected components of $\left(
V,B\right)  $. In other words, $k$ is $\operatorname*{conn}\left(  V,B\right)
$ (since $\operatorname*{conn}\left(  V,B\right)  $ is the number of connected
components of $\left(  V,B\right)  $ (by the definition of
$\operatorname*{conn}\left(  V,B\right)  $)). In other words,
$k=\operatorname*{conn}\left(  V,B\right)  $.

Recall that $\left(  C_{1},C_{2},\ldots,C_{k}\right)  $ is a list of all
connected components of $\left(  V,B\right)  $. In other words, $\left(
C_{1},C_{2},\ldots,C_{k}\right)  $ is a list of all elements of $V/\left(
\sim\right)  $ (since the elements of $V/\left(  \sim\right)  $ are the
connected components of $\left(  V,B\right)  $). Moreover, every element of
$V/\left(  \sim\right)  $ appears exactly once in this list $\left(
C_{1},C_{2},\ldots,C_{k}\right)  $ (since the entries of the list $\left(
C_{1},C_{2},\ldots,C_{k}\right)  $ are pairwise distinct\footnote{since every
connected component of $\left(  V,B\right)  $ appears exactly once in this
list}). Thus, $\left(  C_{1},C_{2},\ldots,C_{k}\right)  $ is a list of all
elements of $V/\left(  \sim\right)  $, and contains each of these elements
exactly once. Hence, the map%
\begin{align*}
Y^{V/\left(  \sim\right)  }  &  \rightarrow Y^{k},\\
f  &  \mapsto\left(  f\left(  C_{1}\right)  ,f\left(  C_{2}\right)
,\ldots,f\left(  C_{k}\right)  \right)
\end{align*}
is a bijection (by Lemma \ref{lem.function-count}, applied to $W=V/\left(
\sim\right)  $). Hence, there exists a bijection $Y^{V/\left(  \sim\right)
}\rightarrow Y^{k}$ (namely, this map). Thus, $\left\vert Y^{k}\right\vert
=\left\vert Y^{V/\left(  \sim\right)  }\right\vert $.

Now, (\ref{pf.lem.Eqs.sum.1}) becomes%
\begin{align*}
\sum_{\substack{f:V\rightarrow Y;\\B\subseteq\operatorname*{Eqs}f}}1  &
=\left\vert Y^{V/\left(  \sim\right)  }\right\vert =\left\vert Y^{k}%
\right\vert =\left\vert Y\right\vert ^{k}=q^{k}\ \ \ \ \ \ \ \ \ \ \left(
\text{since }\left\vert Y\right\vert =q\right) \\
&  =q^{\operatorname*{conn}\left(  V,B\right)  }\ \ \ \ \ \ \ \ \ \ \left(
\text{since }k=\operatorname*{conn}\left(  V,B\right)  \right)  .
\end{align*}
This proves Lemma \ref{lem.Eqs.sum1}.
\end{proof}
\end{verlong}

We shall now show a weaker version of Theorem \ref{thm.chrompol.varis} (as a
stepping stone to the actual theorem):

\begin{lemma}
\label{lem.chrompol.weak.varis}Let $G=\left(  V,E\right)  $ be a finite graph.
Let $X$ be a totally ordered set. Let $\ell:E\rightarrow X$ be a function. Let
$\mathfrak{K}$ be some set of broken circuits of $G$ (not necessarily
containing all of them). Let $a_{K}$ be an element of $\mathbf{k}$ for every
$K\in\mathfrak{K}$. Let $q\in\mathbb{N}$. Then,%
\begin{align*}
&  \left(  \text{the number of all proper }\left\{  1,2,\ldots,q\right\}
\text{-colorings of }G\right) \\
&  =\sum_{F\subseteq E}\left(  -1\right)  ^{\left\vert F\right\vert }\left(
\prod_{\substack{K\in\mathfrak{K};\\K\subseteq F}}a_{K}\right)
q^{\operatorname*{conn}\left(  V,F\right)  }.
\end{align*}
(Here, of course, the pair $\left(  V,F\right)  $ is regarded as a graph, and
the expression $\operatorname*{conn}\left(  V,F\right)  $ is understood
according to Definition \ref{def.conn}.)
\end{lemma}

\begin{vershort}
\begin{proof}
[Proof of Lemma \ref{lem.chrompol.weak.varis}.]We have\footnote{We are again
using the Iverson bracket notation, as defined in Definition \ref{def.iverson}%
.}%
\begin{align*}
&  \left(  \text{the number of all proper }\left\{  1,2,\ldots,q\right\}
\text{-colorings of }G\right) \\
&  =\sum_{f:V\rightarrow\left\{  1,2,\ldots,q\right\}  }\left[
\underbrace{f\text{ is a proper }\left\{  1,2,\ldots,q\right\}
\text{-coloring of }G}_{\substack{\Longleftrightarrow\ \left(  \text{the
}\left\{  1,2,\ldots,q\right\}  \text{-coloring }f\text{ of }G\text{ is
proper}\right)  \\\Longleftrightarrow\ \left(  E\cap\operatorname*{Eqs}%
f=\varnothing\right)  \\\text{(by Lemma \ref{lem.Eqs.proper}, applied to
}\left\{  1,2,\ldots,q\right\}  \text{ instead of }X\text{)}}}\right] \\
&  =\sum_{f:V\rightarrow\left\{  1,2,\ldots,q\right\}  }\underbrace{\left[
E\cap\operatorname*{Eqs}f=\varnothing\right]  }_{\substack{=\sum_{B\subseteq
E\cap\operatorname*{Eqs}f}\left(  -1\right)  ^{\left\vert B\right\vert }%
\prod_{\substack{K\in\mathfrak{K};\\K\subseteq B}}a_{K}\\\text{(by Lemma
\ref{lem.NBCm.moeb}, applied to }Y=\mathbb{N}_{+}\text{)}}}\\
&  =\sum_{f:V\rightarrow\left\{  1,2,\ldots,q\right\}  }\underbrace{\sum
_{B\subseteq E\cap\operatorname*{Eqs}f}}_{=\sum_{\substack{B\subseteq
E;\\B\subseteq\operatorname*{Eqs}f}}}\left(  -1\right)  ^{\left\vert
B\right\vert }\left(  \prod_{\substack{K\in\mathfrak{K};\\K\subseteq B}%
}a_{K}\right)  =\underbrace{\sum_{f:V\rightarrow\left\{  1,2,\ldots,q\right\}
}\sum_{\substack{B\subseteq E;\\B\subseteq\operatorname*{Eqs}f}}}%
_{=\sum_{B\subseteq E}\sum_{\substack{f:V\rightarrow\mathbb{N}_{+}%
;\\B\subseteq\operatorname*{Eqs}f}}}\left(  -1\right)  ^{\left\vert
B\right\vert }\left(  \prod_{\substack{K\in\mathfrak{K};\\K\subseteq B}%
}a_{K}\right) \\
&  =\sum_{B\subseteq E}\sum_{\substack{f:V\rightarrow\left\{  1,2,\ldots
,q\right\}  ;\\B\subseteq\operatorname*{Eqs}f}}\left(  -1\right)  ^{\left\vert
B\right\vert }\left(  \prod_{\substack{K\in\mathfrak{K};\\K\subseteq B}%
}a_{K}\right)  =\sum_{B\subseteq E}\left(  -1\right)  ^{\left\vert
B\right\vert }\left(  \prod_{\substack{K\in\mathfrak{K};\\K\subseteq B}%
}a_{K}\right)  \underbrace{\sum_{\substack{f:V\rightarrow\left\{
1,2,\ldots,q\right\}  ;\\B\subseteq\operatorname*{Eqs}f}}1}%
_{\substack{=q^{\operatorname*{conn}\left(  V,B\right)  }\\\text{(by Lemma
\ref{lem.Eqs.sum1})}}}\\
&  =\sum_{B\subseteq E}\left(  -1\right)  ^{\left\vert B\right\vert }\left(
\prod_{\substack{K\in\mathfrak{K};\\K\subseteq B}}a_{K}\right)
q^{\operatorname*{conn}\left(  V,B\right)  }=\sum_{F\subseteq E}\left(
-1\right)  ^{\left\vert F\right\vert }\left(  \prod_{\substack{K\in
\mathfrak{K};\\K\subseteq F}}a_{K}\right)  q^{\operatorname*{conn}\left(
V,F\right)  }%
\end{align*}
(here, we have renamed the summation index $B$ as $F$). This proves Theorem
\ref{thm.chromsym.varis}.
\end{proof}
\end{vershort}

\begin{verlong}
\begin{proof}
[Proof of Lemma \ref{lem.chrompol.weak.varis}.]Let $Q=\left\{  1,2,\ldots
,q\right\}  $. If $f:V\rightarrow Q$ is a map, then we have the following
logical equivalence:%
\begin{equation}
\left(  \text{the }Q\text{-coloring }f\text{ of }G\text{ is proper}\right)
\ \Longleftrightarrow\ \left(  E\cap\operatorname*{Eqs}f=\varnothing\right)
\label{pf.lem.chrompol.weak.varis.equiv}%
\end{equation}
(because the $Q$-coloring $f$ of $G$ is proper if and only if $E\cap
\operatorname*{Eqs}f=\varnothing$\ \ \ \ \footnote{by Lemma
\ref{lem.Eqs.proper} (applied to $Q$ instead of $X$)}). Now,\footnote{We are
again using the Iverson bracket notation, as defined in Definition
\ref{def.iverson}.}%
\begin{align}
&  \sum_{f:V\rightarrow Q}\left[  \underbrace{E\cap\operatorname*{Eqs}%
f=\varnothing}_{\substack{\Longleftrightarrow\ \left(  \text{the
}Q\text{-coloring }f\text{ of }G\text{ is proper}\right)  \\\text{(by
(\ref{pf.lem.chrompol.weak.varis.equiv}))}}}\right] \nonumber\\
&  =\sum_{f:V\rightarrow Q}\left[  \underbrace{\text{the }Q\text{-coloring
}f\text{ of }G\text{ is proper}}_{\Longleftrightarrow\ \left(  f\text{ is a
proper }Q\text{-coloring of }G\right)  }\right] \nonumber\\
&  =\sum_{f:V\rightarrow Q}\left[  f\text{ is a proper }Q\text{-coloring of
}G\right] \nonumber\\
&  =\sum_{\substack{f:V\rightarrow Q\text{ is a}\\\text{proper }%
Q\text{-coloring of }G}}\underbrace{\left[  f\text{ is a proper }%
Q\text{-coloring of }G\right]  }_{\substack{=1\\\text{(since }f\text{ is a
proper }Q\text{-coloring of }G\text{)}}}\nonumber\\
&  \ \ \ \ \ \ \ \ \ \ +\sum_{\substack{f:V\rightarrow Q\text{ is not
a}\\\text{proper }Q\text{-coloring of }G}}\underbrace{\left[  f\text{ is a
proper }Q\text{-coloring of }G\right]  }_{\substack{=0\\\text{(since }f\text{
is not a proper }Q\text{-coloring of }G\text{)}}}\nonumber\\
&  =\sum_{\substack{f:V\rightarrow Q\text{ is a}\\\text{proper }%
Q\text{-coloring of }G}}1+\underbrace{\sum_{\substack{f:V\rightarrow Q\text{
is not a}\\\text{proper }Q\text{-coloring of }G}}0}_{=0}=\sum
_{\substack{f:V\rightarrow Q\text{ is a}\\\text{proper }Q\text{-coloring of
}G}}1\nonumber\\
&  =\left(  \text{the number of all proper }Q\text{-colorings of }G\right)
\cdot1\nonumber\\
&  =\left(  \text{the number of all proper }Q\text{-colorings of }G\right)
\nonumber\\
&  =\left(  \text{the number of all proper }\left\{  1,2,\ldots,q\right\}
\text{-colorings of }G\right)  \label{pf.lem.chrompol.weak.varis.1}%
\end{align}
(since $Q=\left\{  1,2,\ldots,q\right\}  $).

But for every $f:V\rightarrow Q$, we have%
\begin{equation}
\sum_{B\subseteq E\cap\operatorname*{Eqs}f}\left(  -1\right)  ^{\left\vert
B\right\vert }\prod_{\substack{K\in\mathfrak{K};\\K\subseteq B}}a_{K}=\left[
E\cap\operatorname*{Eqs}f=\varnothing\right]
\label{pf.lem.chrompol.weak.varis.moeb}%
\end{equation}
(by Lemma \ref{lem.NBCm.moeb} (applied to $Q$ instead of $Y$)).

For every $f:V\rightarrow Q$, we have%
\begin{equation}
\left\{  F\subseteq E\ \mid\ F\subseteq\operatorname*{Eqs}f\right\}
=\mathcal{P}\left(  E\cap\operatorname*{Eqs}f\right)
\label{pf.lem.chrompol.weak.varis.cut}%
\end{equation}
\footnote{\textit{Proof of (\ref{pf.lem.chrompol.weak.varis.cut}):} Let
$f:V\rightarrow Q$.
\par
Let $B\in\left\{  F\subseteq E\ \mid\ F\subseteq\operatorname*{Eqs}f\right\}
$. Thus, $B$ is a subset $F$ of $E$ satisfying $F\subseteq\operatorname*{Eqs}%
f$. In other words, $B$ is a subset of $E$ and satisfies $B\subseteq
\operatorname*{Eqs}f$. Since $B$ is a subset of $E$, we have $B\subseteq E$.
Combining this with $B\subseteq\operatorname*{Eqs}f$, we obtain $B\subseteq
E\cap\operatorname*{Eqs}f$. In other words, $B\in\mathcal{P}\left(
E\cap\operatorname*{Eqs}f\right)  $.
\par
Let us now forget that we fixed $B$. We thus have proven that every
$B\in\left\{  F\subseteq E\ \mid\ F\subseteq\operatorname*{Eqs}f\right\}  $
satisfies $B\in\mathcal{P}\left(  E\cap\operatorname*{Eqs}f\right)  $. In
other words,%
\begin{equation}
\left\{  F\subseteq E\ \mid\ F\subseteq\operatorname*{Eqs}f\right\}
\subseteq\mathcal{P}\left(  E\cap\operatorname*{Eqs}f\right)  .
\label{pf.lem.chrompol.weak.varis.cut.pf.1}%
\end{equation}
\par
On the other hand, let $C\in\mathcal{P}\left(  E\cap\operatorname*{Eqs}%
f\right)  $. Thus, $C$ is a subset of $E\cap\operatorname*{Eqs}f$. Hence,
$C\subseteq E\cap\operatorname*{Eqs}f\subseteq E$, so that $C$ is a subset of
$E$. Also, $C\subseteq E\cap\operatorname*{Eqs}f\subseteq\operatorname*{Eqs}%
f$. Thus, $C$ is a subset of $E$ and satisfies $C\subseteq\operatorname*{Eqs}%
f$. In other words, $C$ is a subset $F$ of $E$ satisfying $F\subseteq
\operatorname*{Eqs}f$. In other words, $C\in\left\{  F\subseteq E\ \mid
\ F\subseteq\operatorname*{Eqs}f\right\}  $.
\par
Let us now forget that we fixed $C$. We thus have proven that every
$C\in\mathcal{P}\left(  E\cap\operatorname*{Eqs}f\right)  $ satisfies
$C\in\left\{  F\subseteq E\ \mid\ F\subseteq\operatorname*{Eqs}f\right\}  $.
In other words,%
\[
\mathcal{P}\left(  E\cap\operatorname*{Eqs}f\right)  \subseteq\left\{
F\subseteq E\ \mid\ F\subseteq\operatorname*{Eqs}f\right\}  .
\]
Combining this inclusion with (\ref{pf.lem.chrompol.weak.varis.cut.pf.1}), we
obtain $\left\{  F\subseteq E\ \mid\ F\subseteq\operatorname*{Eqs}f\right\}
=\mathcal{P}\left(  E\cap\operatorname*{Eqs}f\right)  $. This proves
(\ref{pf.lem.chrompol.weak.varis.cut}).}.

For every $f:V\rightarrow Q$, we have%
\begin{align}
&  \underbrace{\sum_{\substack{B\subseteq E;\\B\subseteq\operatorname*{Eqs}%
f}}}_{\substack{=\sum_{B\in\left\{  F\subseteq E\ \mid\ F\subseteq
\operatorname*{Eqs}f\right\}  }=\sum_{B\in\mathcal{P}\left(  E\cap
\operatorname*{Eqs}f\right)  }\\\text{(because }\left\{  F\subseteq
E\ \mid\ F\subseteq\operatorname*{Eqs}f\right\}  =\mathcal{P}\left(
E\cap\operatorname*{Eqs}f\right)  \\\text{(by
(\ref{pf.lem.chrompol.weak.varis.cut})))}}}\left(  -1\right)  ^{\left\vert
B\right\vert }\prod_{\substack{K\in\mathfrak{K};\\K\subseteq B}}a_{K}%
\nonumber\\
&  =\underbrace{\sum_{B\in\mathcal{P}\left(  E\cap\operatorname*{Eqs}f\right)
}}_{=\sum_{B\subseteq E\cap\operatorname*{Eqs}f}}\left(  -1\right)
^{\left\vert B\right\vert }\prod_{\substack{K\in\mathfrak{K};\\K\subseteq
B}}a_{K}=\sum_{B\subseteq E\cap\operatorname*{Eqs}f}\left(  -1\right)
^{\left\vert B\right\vert }\prod_{\substack{K\in\mathfrak{K};\\K\subseteq
B}}a_{K}\nonumber\\
&  =\left[  E\cap\operatorname*{Eqs}f=\varnothing\right]
\label{pf.lem.chrompol.weak.varis.moeb2}%
\end{align}
(by (\ref{pf.lem.chrompol.weak.varis.moeb})).

Now, (\ref{pf.lem.chrompol.weak.varis.1}) yields%
\begin{align*}
&  \left(  \text{the number of all proper }\left\{  1,2,\ldots,q\right\}
\text{-colorings of }G\right) \\
&  =\sum_{f:V\rightarrow Q}\underbrace{\left[  E\cap\operatorname*{Eqs}%
f=\varnothing\right]  }_{\substack{=\sum_{\substack{B\subseteq E;\\B\subseteq
\operatorname*{Eqs}f}}\left(  -1\right)  ^{\left\vert B\right\vert }%
\prod_{\substack{K\in\mathfrak{K};\\K\subseteq B}}a_{K}\\\text{(by
(\ref{pf.lem.chrompol.weak.varis.moeb2}))}}}\\
&  =\sum_{f:V\rightarrow Q}\left(  \sum_{\substack{B\subseteq E;\\B\subseteq
\operatorname*{Eqs}f}}\left(  -1\right)  ^{\left\vert B\right\vert }%
\prod_{\substack{K\in\mathfrak{K};\\K\subseteq B}}a_{K}\right)
=\underbrace{\sum_{f:V\rightarrow Q}\sum_{\substack{B\subseteq E;\\B\subseteq
\operatorname*{Eqs}f}}}_{=\sum_{B\subseteq E}\sum_{\substack{f:V\rightarrow
Q;\\B\subseteq\operatorname*{Eqs}f}}}\left(  -1\right)  ^{\left\vert
B\right\vert }\underbrace{\left(  \prod_{\substack{K\in\mathfrak{K}%
;\\K\subseteq B}}a_{K}\right)  }_{=\left(  \prod_{\substack{K\in
\mathfrak{K};\\K\subseteq B}}a_{K}\right)  1}\\
&  =\sum_{B\subseteq E}\sum_{\substack{f:V\rightarrow Q;\\B\subseteq
\operatorname*{Eqs}f}}\left(  -1\right)  ^{\left\vert B\right\vert }\left(
\prod_{\substack{K\in\mathfrak{K};\\K\subseteq B}}a_{K}\right)  1=\sum
_{B\subseteq E}\sum_{\substack{f:V\rightarrow\left\{  1,2,\ldots,q\right\}
;\\B\subseteq\operatorname*{Eqs}f}}\left(  -1\right)  ^{\left\vert
B\right\vert }\left(  \prod_{\substack{K\in\mathfrak{K};\\K\subseteq B}%
}a_{K}\right)  1\\
&  \ \ \ \ \ \ \ \ \ \ \left(  \text{since }Q=\left\{  1,2,\ldots,q\right\}
\right) \\
&  =\sum_{B\subseteq E}\left(  -1\right)  ^{\left\vert B\right\vert }\left(
\prod_{\substack{K\in\mathfrak{K};\\K\subseteq B}}a_{K}\right)
\underbrace{\sum_{\substack{f:V\rightarrow\left\{  1,2,\ldots,q\right\}
;\\B\subseteq\operatorname*{Eqs}f}}1}_{\substack{=q^{\operatorname*{conn}%
\left(  V,B\right)  }\\\text{(by Lemma \ref{lem.Eqs.sum1}}\\\text{(since
}\left(  V,B\right)  \text{ is a finite graph}\\\text{(since }V\text{ is a
finite set and }B\subseteq E\subseteq\dbinom{V}{2}\text{)))}}}\\
&  =\sum_{B\subseteq E}\left(  -1\right)  ^{\left\vert B\right\vert }\left(
\prod_{\substack{K\in\mathfrak{K};\\K\subseteq B}}a_{K}\right)
q^{\operatorname*{conn}\left(  V,B\right)  }=\sum_{F\subseteq E}\left(
-1\right)  ^{\left\vert F\right\vert }\left(  \prod_{\substack{K\in
\mathfrak{K};\\K\subseteq F}}a_{K}\right)  q^{\operatorname*{conn}\left(
V,F\right)  }%
\end{align*}
(here, we have renamed the summation index $B$ as $F$). This proves Lemma
\ref{lem.chrompol.weak.varis}.
\end{proof}
\end{verlong}

From Lemma \ref{lem.chrompol.weak.varis}, we obtain the following consequence:

\begin{lemma}
\label{lem.chrompol.weak.empty}Let $G=\left(  V,E\right)  $ be a finite graph.
Let $q\in\mathbb{N}$. Then,%
\begin{align*}
&  \left(  \text{the number of all proper }\left\{  1,2,\ldots,q\right\}
\text{-colorings of }G\right) \\
&  =\sum_{F\subseteq E}\left(  -1\right)  ^{\left\vert F\right\vert
}q^{\operatorname*{conn}\left(  V,F\right)  }.
\end{align*}
(Here, of course, the pair $\left(  V,F\right)  $ is regarded as a graph, and
the expression $\operatorname*{conn}\left(  V,F\right)  $ is understood
according to Definition \ref{def.conn}.)
\end{lemma}

\begin{vershort}
\begin{proof}
[Proof of Lemma \ref{lem.chrompol.weak.empty}.]This is derived from Lemma
\ref{lem.chrompol.weak.varis} in the same way as Theorem
\ref{thm.chromsym.empty} was derived from Theorem \ref{thm.chrompol.varis}.
\end{proof}
\end{vershort}

\begin{verlong}
\begin{proof}
[Proof of Lemma \ref{lem.chrompol.weak.empty}.]Let $X$ be the totally ordered
set $\left\{  1\right\}  $ (equipped with the only possible order on this
set). Let $\ell:E\rightarrow X$ be the function sending each $e\in E$ to $1\in
X$. Let $\mathfrak{K}$ be the empty set. Clearly, $\mathfrak{K}$ is a set of
broken circuits of $G$. Lemma \ref{lem.chrompol.weak.varis} (applied to $0$
instead of $a_{K}$) yields%
\begin{align*}
&  \left(  \text{the number of all proper }\left\{  1,2,\ldots,q\right\}
\text{-colorings of }G\right) \\
&  =\sum_{F\subseteq E}\left(  -1\right)  ^{\left\vert F\right\vert
}\underbrace{\left(  \prod_{\substack{K\in\mathfrak{K};\\K\subseteq
F}}0\right)  }_{\substack{=\left(  \text{empty product}\right)  \\\text{(since
}\mathfrak{K}\text{ is the empty set)}}}q^{\operatorname*{conn}\left(
V,F\right)  }\\
&  =\sum_{F\subseteq E}\left(  -1\right)  ^{\left\vert F\right\vert
}\underbrace{\left(  \text{empty product}\right)  }_{=1}%
q^{\operatorname*{conn}\left(  V,F\right)  }=\sum_{F\subseteq E}\left(
-1\right)  ^{\left\vert F\right\vert }q^{\operatorname*{conn}\left(
V,F\right)  }.
\end{align*}
This proves Lemma \ref{lem.chrompol.weak.empty}.
\end{proof}
\end{verlong}

Next, we recall a classical fact about a polynomials over fields: Namely, if a
polynomial (in one variable) over a field has infinitely many roots, then this
polynomial is $0$. Let us state this more formally:

\begin{proposition}
\label{prop.poly0.field}Let $K$ be a field. Let $P\in K\left[  x\right]  $ be
a polynomial over $K$. Assume that there are infinitely many $\lambda\in K$
satisfying $P\left(  \lambda\right)  =0$. Then, $P=0$.
\end{proposition}

We shall use the following consequence of this proposition:

\begin{corollary}
\label{cor.poly0.intdom}Let $R$ be an integral domain. Assume that the
canonical ring homomorphism from the ring $\mathbb{Z}$ to the ring $R$ is
injective. Let $P\in R\left[  x\right]  $ be a polynomial over $R$. Assume
that $P\left(  q\cdot1_{R}\right)  =0$ for every $q\in\mathbb{N}$ (where
$1_{R}$ denotes the unity of $R$). Then, $P=0$.
\end{corollary}

\begin{vershort}
\begin{proof}
[Proof of Corollary \ref{cor.poly0.intdom}.]Let $K$ denote the fraction field
of the integral domain $R$. We regard $R$ and $R\left[  x\right]  $ as
subrings of $K$ and $K\left[  x\right]  $, respectively. By assumption, we
have $P\left(  q\cdot1_{R}\right)  =0$ for every $q\in\mathbb{N}$. But the
elements $q\cdot1_{R}$ of $R$ for $q\in\mathbb{N}$ are pairwise distinct
(since the canonical ring homomorphism from the ring $\mathbb{Z}$ to the ring
$R$ is injective). Hence, there are infinitely many $\lambda\in K$ satisfying
$P\left(  \lambda\right)  =0$ (namely, $\lambda=q\cdot1_{R}$ for all
$q\in\mathbb{N}$). Thus, Proposition \ref{prop.poly0.field} shows that $P=0$.
This proves Corollary \ref{cor.poly0.intdom}.
\end{proof}
\end{vershort}

\begin{verlong}
\begin{proof}
[Proof of Corollary \ref{cor.poly0.intdom}.]Let $K$ denote the fraction field
of the integral domain $R$. Then, there is a canonical injective ring
homorphism $R\rightarrow K$. We use this homomorphism to regard $R$ as a
subring of $K$. Consequently, $R\left[  x\right]  $ will be regarded as a
subring of $K\left[  x\right]  $. In particular, the polynomial $P\in R\left[
x\right]  $ will thus be regarded as a polynomial in $K\left[  x\right]  $.

Let $\iota:\mathbb{Z}\rightarrow R$ be the canonical ring homomorphism from
the ring $\mathbb{Z}$ to the ring $R$. (Thus, $\iota$ sends every
$q\in\mathbb{Z}$ to $q\cdot1_{R}\in R$.)

We have assumed that the canonical ring homomorphism from the ring
$\mathbb{Z}$ to the ring $R$ is injective. In other words, the map $\iota$ is
injective (since the map $\iota$ is the canonical ring homomorphism from the
ring $\mathbb{Z}$ to the ring $R$). Hence, $\left\vert \iota\left(
\mathbb{N}\right)  \right\vert =\left\vert \mathbb{N}\right\vert =\infty$.
Moreover, every $\lambda\in\iota\left(  \mathbb{N}\right)  $ satisfies
$P\left(  \lambda\right)  =0$\ \ \ \ \footnote{\textit{Proof.} Let $\lambda
\in\iota\left(  \mathbb{N}\right)  $. Thus, there exists some $h\in\mathbb{N}$
such that $\lambda=\iota\left(  h\right)  $. Consider this $h$.
\par
We have assumed that $P\left(  q\cdot1_{R}\right)  =0$ for every
$q\in\mathbb{N}$. Applying this to $q=h$, we obtain $P\left(  h\cdot
1_{R}\right)  =0$. But $\lambda=\iota\left(  h\right)  =h\cdot1_{R}$ (by the
definition of $\iota$). Hence, $P\left(  \underbrace{\lambda}_{=h\cdot1_{R}%
}\right)  =P\left(  h\cdot1_{R}\right)  =0$, qed.}. Hence, there are
infinitely many $\lambda\in K$ satisfying $P\left(  \lambda\right)  =0$
(because there are infinitely many $\lambda\in\iota\left(  \mathbb{N}\right)
$ (since $\left\vert \iota\left(  \mathbb{N}\right)  \right\vert =\infty$),
and because they all are elements of $K$ (since $\iota\left(  \mathbb{N}%
\right)  \subseteq R\subseteq K$)). Hence, Proposition \ref{prop.poly0.field}
shows that $P=0$. This proves Corollary \ref{cor.poly0.intdom}.
\end{proof}
\end{verlong}

We can now prove the classical Theorem \ref{thm.chrompol.exist}:

\begin{vershort}
\begin{proof}
[Proof of Theorem \ref{thm.chrompol.exist}.]We need to show that there exists
a unique polynomial $P\in\mathbb{Z}\left[  x\right]  $ such that every
$q\in\mathbb{N}$ satisfies%
\[
P\left(  q\right)  =\left(  \text{the number of all proper }\left\{
1,2,\ldots,q\right\}  \text{-colorings of }G\right)  .
\]
To see that such a polynomial exists, we notice that $P=\sum_{F\subseteq
E}\left(  -1\right)  ^{\left\vert F\right\vert }x^{\operatorname*{conn}\left(
V,F\right)  }$ is such a polynomial (by Lemma \ref{lem.chrompol.weak.empty}).
It remains to prove that such a polynomial is unique. This follows from the
fact that if two polynomials $P_{1}\in\mathbb{Z}\left[  x\right]  $ and
$P_{2}\in\mathbb{Z}\left[  x\right]  $ satisfy
\[
P_{1}\left(  q\right)  =P_{2}\left(  q\right)  \ \ \ \ \ \ \ \ \ \ \text{for
all }q\in\mathbb{N},
\]
then $P_{1}=P_{2}$\ \ \ \ \footnote{This fact follows from Corollary
\ref{cor.poly0.intdom} (applied to $R=\mathbb{Z}$ and $P=P_{1}-P_{2}$).}.
Theorem \ref{thm.chrompol.exist} is therefore proven.
\end{proof}
\end{vershort}

\begin{verlong}
\begin{proof}
[Proof of Theorem \ref{thm.chrompol.exist}.]We need to show that there exists
a unique polynomial $P\in\mathbb{Z}\left[  x\right]  $ such that every
$q\in\mathbb{N}$ satisfies%
\begin{equation}
P\left(  q\right)  =\left(  \text{the number of all proper }\left\{
1,2,\ldots,q\right\}  \text{-colorings of }G\right)  .
\label{pf.thm.chrompol.exist.P(q)=}%
\end{equation}


Let us first show that there exists at most one such polynomial. Indeed, let
$P_{1}$ and $P_{2}$ be two polynomials $P\in\mathbb{Z}\left[  x\right]  $ such
that every $q\in\mathbb{N}$ satisfies (\ref{pf.thm.chrompol.exist.P(q)=}). We
shall show that $P_{1}=P_{2}$.

We know that $P_{1}$ is a polynomial $P\in\mathbb{Z}\left[  x\right]  $ such
that every $q\in\mathbb{N}$ satisfies (\ref{pf.thm.chrompol.exist.P(q)=}). In
other words, $P_{1}$ is a polynomial in $\mathbb{Z}\left[  x\right]  $ and
every $q\in\mathbb{N}$ satisfies
\begin{equation}
P_{1}\left(  q\right)  =\left(  \text{the number of all proper }\left\{
1,2,\ldots,q\right\}  \text{-colorings of }G\right)  .
\label{pf.thm.chrompol.exist.P1(q)=}%
\end{equation}
The same argument (applied to $P_{2}$ instead of $P_{1}$) shows that $P_{2}$
is a polynomial in $\mathbb{Z}\left[  x\right]  $ and every $q\in\mathbb{N}$
satisfies
\begin{equation}
P_{2}\left(  q\right)  =\left(  \text{the number of all proper }\left\{
1,2,\ldots,q\right\}  \text{-colorings of }G\right)  .
\label{pf.thm.chrompol.exist.P2(q)=}%
\end{equation}


The ring $\mathbb{Z}$ is an integral domain. The canonical ring homomorphism
from the ring $\mathbb{Z}$ to the ring $\mathbb{Z}$ is
injective\footnote{Indeed, this homomorphism is the identity map.}. Now, every
$q\in\mathbb{N}$ satisfies%
\begin{align}
P_{1}\left(  q\right)   &  =\left(  \text{the number of all proper }\left\{
1,2,\ldots,q\right\}  \text{-colorings of }G\right) \nonumber\\
&  \ \ \ \ \ \ \ \ \ \ \left(  \text{by (\ref{pf.thm.chrompol.exist.P1(q)=}%
)}\right) \nonumber\\
&  =P_{2}\left(  q\right)  \ \ \ \ \ \ \ \ \ \ \left(  \text{by
(\ref{pf.thm.chrompol.exist.P2(q)=})}\right)  \label{pf.thm.chrompol.exist.4}%
\end{align}
Thus, every $q\in\mathbb{N}$ satisfies%
\begin{align*}
\left(  P_{1}-P_{2}\right)  \left(  \underbrace{q\cdot1_{\mathbb{Z}}}%
_{=q}\right)   &  =\left(  P_{1}-P_{2}\right)  \left(  q\right)
=\underbrace{P_{1}\left(  q\right)  }_{\substack{=P_{2}\left(  q\right)
\\\text{(by (\ref{pf.thm.chrompol.exist.4}))}}}-P_{2}\left(  q\right) \\
&  =P_{2}\left(  q\right)  -P_{2}\left(  q\right)  =0.
\end{align*}
Hence, Corollary \ref{cor.poly0.intdom} (applied to $\mathbb{Z}$ and
$P_{1}-P_{2}$ instead of $R$ and $P$) shows that $P_{1}-P_{2}=0$. In other
words, $P_{1}=P_{2}$.

Now, let us forget that we fixed $P_{1}$ and $P_{2}$. We thus have showed that
if $P_{1}$ and $P_{2}$ are two polynomials $P\in\mathbb{Z}\left[  x\right]  $
such that every $q\in\mathbb{N}$ satisfies (\ref{pf.thm.chrompol.exist.P(q)=}%
), then $P_{1}=P_{2}$. In other words, there exists \textbf{at most one}
polynomial $P\in\mathbb{Z}\left[  x\right]  $ such that every $q\in\mathbb{N}$
satisfies (\ref{pf.thm.chrompol.exist.P(q)=}).

Let us now prove that there exists at least one such polynomial. Indeed,
define a polynomial $Q\in\mathbb{Z}\left[  x\right]  $ by%
\begin{equation}
Q=\sum_{F\subseteq E}\left(  -1\right)  ^{\left\vert F\right\vert
}x^{\operatorname*{conn}\left(  V,F\right)  }.
\label{pf.thm.chrompol.exist.Q=}%
\end{equation}
Then, every $q\in\mathbb{N}$ satisfies%
\begin{align*}
Q\left(  q\right)   &  =\sum_{F\subseteq E}\left(  -1\right)  ^{\left\vert
F\right\vert }q^{\operatorname*{conn}\left(  V,F\right)  }%
\ \ \ \ \ \ \ \ \ \ \left(  \text{here, we have substituted }q\text{ for
}x\text{ in (\ref{pf.thm.chrompol.exist.Q=})}\right) \\
&  =\left(  \text{the number of all proper }\left\{  1,2,\ldots,q\right\}
\text{-colorings of }G\right)
\end{align*}
(by Lemma \ref{lem.chrompol.weak.empty}). In other words, every $q\in
\mathbb{N}$ satisfies (\ref{pf.thm.chrompol.exist.P(q)=}) for $P=Q$. Thus, $Q$
is a polynomial $P\in\mathbb{Z}\left[  x\right]  $ such that every
$q\in\mathbb{N}$ satisfies (\ref{pf.thm.chrompol.exist.P(q)=}). Hence, there
exists \textbf{at least} one polynomial $P\in\mathbb{Z}\left[  x\right]  $
such that every $q\in\mathbb{N}$ satisfies (\ref{pf.thm.chrompol.exist.P(q)=})
(namely, $P=Q$). Consequently, there exists \textbf{exactly} one polynomial
$P\in\mathbb{Z}\left[  x\right]  $ such that every $q\in\mathbb{N}$ satisfies
(\ref{pf.thm.chrompol.exist.P(q)=}) (because we have already shown that there
exists \textbf{at most} one such polynomial). This proves Theorem
\ref{thm.chrompol.exist}.
\end{proof}
\end{verlong}

Next, it is the turn of Theorem \ref{thm.chrompol.varis}:

\begin{vershort}
\begin{proof}
[Proof of Theorem \ref{thm.chrompol.varis}.]Let $R$ be the polynomial ring
$\mathbb{Z}\left[  y_{K}\ \mid\ K\in\mathfrak{K}\right]  $, where $y_{K}$ is a
new indeterminate for each $K\in\mathfrak{K}$.

The claim of Theorem \ref{thm.chrompol.varis} is a polynomial identity in the
elements $a_{K}$ of $\mathbf{k}$. Hence, we can WLOG assume that
$\mathbf{k}=R$ and $a_{K}=y_{K}$ for each $K\in\mathfrak{K}$. Assume this.
Thus, $\mathbf{k}$ is an integral domain, and the canonical ring homomorphism
from the ring $\mathbb{Z}$ to the ring $\mathbf{k}$ is injective.

For every $q\in\mathbb{N}$, we have%
\begin{align}
\chi_{G}\left(  q\right)   &  =\left(  \text{the number of all proper
}\left\{  1,2,\ldots,q\right\}  \text{-colorings of }G\right) \nonumber\\
&  \ \ \ \ \ \ \ \ \ \ \left(  \text{by the definition of the chromatic
polynomial }\chi_{G}\right) \nonumber\\
&  =\sum_{F\subseteq E}\left(  -1\right)  ^{\left\vert F\right\vert }\left(
\prod_{\substack{K\in\mathfrak{K};\\K\subseteq F}}a_{K}\right)
q^{\operatorname*{conn}\left(  V,F\right)  }
\label{pf.thm.chrompol.varis.short.eq1}%
\end{align}
(by Lemma \ref{lem.chrompol.weak.varis}). Define a polynomial $P\in
\mathbf{k}\left[  x\right]  $ by%
\begin{equation}
P=\chi_{G}-\sum_{F\subseteq E}\left(  -1\right)  ^{\left\vert F\right\vert
}\left(  \prod_{\substack{K\in\mathfrak{K};\\K\subseteq F}}a_{K}\right)
x^{\operatorname*{conn}\left(  V,F\right)  }.
\label{pf.thm.chrompol.varis.short.P=}%
\end{equation}
Then, for every $q\in\mathbb{N}$, we have%
\begin{align*}
P\left(  \underbrace{q\cdot1_{\mathbf{k}}}_{=q}\right)   &  =P\left(
q\right)  =\chi_{G}\left(  q\right)  -\sum_{F\subseteq E}\left(  -1\right)
^{\left\vert F\right\vert }\left(  \prod_{\substack{K\in\mathfrak{K}%
;\\K\subseteq F}}a_{K}\right)  q^{\operatorname*{conn}\left(  V,F\right)
}\ \ \ \ \ \ \ \ \ \ \left(  \text{by (\ref{pf.thm.chrompol.varis.short.P=}%
)}\right) \\
&  =0\ \ \ \ \ \ \ \ \ \ \left(  \text{by
(\ref{pf.thm.chrompol.varis.short.eq1})}\right)  .
\end{align*}
Thus, Corollary \ref{cor.poly0.intdom} (applied to $R=\mathbf{k}$) shows that
$P=0$. In light of (\ref{pf.thm.chrompol.varis.short.P=}), this rewrites as
follows:%
\[
\chi_{G}=\sum_{F\subseteq E}\left(  -1\right)  ^{\left\vert F\right\vert
}\left(  \prod_{\substack{K\in\mathfrak{K};\\K\subseteq F}}a_{K}\right)
x^{\operatorname*{conn}\left(  V,F\right)  }.
\]
This proves Theorem \ref{thm.chrompol.varis}.
\end{proof}
\end{vershort}

\begin{verlong}
\begin{proof}
[Proof of Theorem \ref{thm.chrompol.varis}.]Let $R$ be the polynomial ring
$\mathbb{Z}\left[  y_{K}\ \mid\ K\in\mathfrak{K}\right]  $, where $y_{K}$ is a
new indeterminate for each $K\in\mathfrak{K}$. Clearly, $R$ is an integral
domain (since $R$ is a polynomial ring over $\mathbb{Z}$). Moreover, the
canonical ring homomorphism from the ring $\mathbb{Z}$ to the ring $R$ is
injective (since $R$ is a polynomial ring over $\mathbb{Z}$). We regard
$\mathbb{Z}\left[  x\right]  $ as a subring of $R\left[  x\right]  $ (since
$\mathbb{Z}$ is a subring of $R$).

The chromatic polynomial $\chi_{G}$ is the unique polynomial $P\in
\mathbb{Z}\left[  x\right]  $ such that every $q\in\mathbb{N}$ satisfies
$P\left(  q\right)  =\left(  \text{the number of all proper }\left\{
1,2,\ldots,q\right\}  \text{-colorings of }G\right)  $ (because this is how
$\chi_{G}$ is defined). Thus, $\chi_{G}$ is a polynomial in $\mathbb{Z}\left[
x\right]  $ and has the property that every $q\in\mathbb{N}$ satisfies%
\begin{equation}
\chi_{G}\left(  q\right)  =\left(  \text{the number of all proper }\left\{
1,2,\ldots,q\right\}  \text{-colorings of }G\right)  .
\label{pf.thm.chrompol.varis.chiG(q)=}%
\end{equation}


Lemma \ref{lem.chrompol.weak.varis} (applied to $R$ and $y_{K}$ instead of
$\mathbf{k}$ and $a_{K}$) yields that%
\begin{align*}
&  \left(  \text{the number of all proper }\left\{  1,2,\ldots,q\right\}
\text{-colorings of }G\right) \\
&  =\sum_{F\subseteq E}\left(  -1\right)  ^{\left\vert F\right\vert }\left(
\prod_{\substack{K\in\mathfrak{K};\\K\subseteq F}}y_{K}\right)
q^{\operatorname*{conn}\left(  V,F\right)  }%
\end{align*}
for every $q\in\mathbb{N}$. Thus, for every $q\in\mathbb{N}$, we have%
\begin{align}
&  \sum_{F\subseteq E}\left(  -1\right)  ^{\left\vert F\right\vert }\left(
\prod_{\substack{K\in\mathfrak{K};\\K\subseteq F}}y_{K}\right)
q^{\operatorname*{conn}\left(  V,F\right)  }\nonumber\\
&  =\left(  \text{the number of all proper }\left\{  1,2,\ldots,q\right\}
\text{-colorings of }G\right) \nonumber\\
&  =\chi_{G}\left(  q\right)  \ \ \ \ \ \ \ \ \ \ \left(  \text{by
(\ref{pf.thm.chrompol.varis.chiG(q)=})}\right)  .
\label{pf.thm.chrompol.varis.2}%
\end{align}


We have $\chi_{G}\in\mathbb{Z}\left[  x\right]  \subseteq R\left[  x\right]
$. On the other hand, define a polynomial $\widetilde{P}\in R\left[  x\right]
$ by%
\begin{equation}
\widetilde{P}=\sum_{F\subseteq E}\left(  -1\right)  ^{\left\vert F\right\vert
}\left(  \prod_{\substack{K\in\mathfrak{K};\\K\subseteq F}}y_{K}\right)
x^{\operatorname*{conn}\left(  V,F\right)  }.
\label{pf.thm.chrompol.varis.Ptilde=}%
\end{equation}
Every $q\in\mathbb{N}$ satisfies%
\begin{align*}
\widetilde{P}\left(  q\right)   &  =\sum_{F\subseteq E}\left(  -1\right)
^{\left\vert F\right\vert }\left(  \prod_{\substack{K\in\mathfrak{K}%
;\\K\subseteq F}}y_{K}\right)  q^{\operatorname*{conn}\left(  V,F\right)  }\\
&  \ \ \ \ \ \ \ \ \ \ \left(  \text{here, we have substituted }q\text{ for
}x\text{ in (\ref{pf.thm.chrompol.varis.Ptilde=})}\right) \\
&  =\chi_{G}\left(  q\right)  \ \ \ \ \ \ \ \ \ \ \left(  \text{by
(\ref{pf.thm.chrompol.varis.2})}\right)  .
\end{align*}
Hence, every $q\in\mathbb{N}$ satisfies%
\[
\left(  \chi_{G}-\widetilde{P}\right)  \left(  \underbrace{q\cdot1_{R}}%
_{=q}\right)  =\left(  \chi_{G}-\widetilde{P}\right)  \left(  q\right)
=\chi_{G}\left(  q\right)  -\underbrace{\widetilde{P}\left(  q\right)
}_{=\chi_{G}\left(  q\right)  }=\chi_{G}\left(  q\right)  -\chi_{G}\left(
q\right)  =0.
\]
Thus, Corollary \ref{cor.poly0.intdom} (applied to $\chi_{G}-\widetilde{P}$
instead of $P$) shows that $\chi_{G}-\widetilde{P}=0$. In other words,
$\chi_{G}=\widetilde{P}$.

Let $\mathbf{a}$ denote the family $\left(  a_{K}\right)  _{K\in\mathfrak{K}%
}\in\mathbf{k}^{\mathfrak{K}}$ of elements of $\mathbf{k}$.

Now, recall that $R$ is the polynomial ring $\mathbb{Z}\left[  y_{K}%
\ \mid\ K\in\mathfrak{K}\right]  $. Hence, $R$ satisfies the following
universal property (the well-known universal property of a polynomial ring):
For any commutative $\mathbb{Z}$-algebra $B$ and any family $\mathbf{b}%
=\left(  b_{K}\right)  _{K\in\mathfrak{K}}\in B^{\mathfrak{K}}$ of elements of
$B$, there exists a unique $\mathbb{Z}$-algebra homomorphism $\phi
:R\rightarrow B$ satisfying%
\[
\left(  \phi\left(  y_{K}\right)  =b_{K}\ \ \ \ \ \ \ \ \ \ \text{for every
}K\in\mathfrak{K}\right)  .
\]
This $\mathbb{Z}$-algebra homomorphism $\phi$ is denoted by
$\operatorname*{ev}\nolimits_{\mathbf{b}}$, and is called the
\textit{evaluation homomorphism} at the family $\mathbf{b}$.

Thus we have constructed a $\mathbb{Z}$-algebra homomorphism
$\operatorname*{ev}\nolimits_{\mathbf{b}}:R\rightarrow B$ for every
commutative $\mathbb{Z}$-algebra $B$ and every family $\mathbf{b}=\left(
b_{K}\right)  _{K\in\mathfrak{K}}\in B^{\mathfrak{K}}$ of elements of $B$.
Applying this construction to $B=\mathbf{k}$, $\mathbf{b}=\mathbf{a}$ and
$b_{K}=a_{K}$, we obtain a $\mathbb{Z}$-algebra homomorphism
$\operatorname*{ev}\nolimits_{\mathbf{a}}:R\rightarrow\mathbf{k}$. This
homomorphism, in turn, induces a $\mathbb{Z}\left[  x\right]  $-algebra
homomorphism $\operatorname*{ev}\nolimits_{\mathbf{a}}\left[  x\right]
:R\left[  x\right]  \rightarrow\mathbf{k}\left[  x\right]  $ satisfying
$\left(  \operatorname*{ev}\nolimits_{\mathbf{a}}\left[  x\right]  \right)
\left(  x\right)  =x$.\ \ \ \ \footnote{This homomorphism $\operatorname*{ev}%
\nolimits_{\mathbf{a}}\left[  x\right]  $ is explicitly given by the formula%
\[
\left(  \operatorname*{ev}\nolimits_{\mathbf{a}}\left[  x\right]  \right)
\left(  \sum_{i=0}^{\infty}r_{i}x^{i}\right)  =\sum_{i=0}^{\infty
}\operatorname*{ev}\nolimits_{\mathbf{a}}\left(  r_{i}\right)  \cdot x^{i}%
\]
for every sequence $\left(  r_{0},r_{1},r_{2},\ldots\right)  \in R^{\infty}$
of elements of $R$ which satisfies $r_{i}=0$ for all sufficiently high $i$.}.
Consider this homomorphism $\operatorname*{ev}\nolimits_{\mathbf{a}}\left[
x\right]  $.

Recall that $\mathbf{a}=\left(  a_{K}\right)  _{K\in\mathfrak{K}}$. Hence,
$\operatorname*{ev}\nolimits_{\mathbf{a}}$ is the unique $\mathbb{Z}$-algebra
homomorphism $\phi:R\rightarrow\mathbf{k}$ satisfying%
\[
\left(  \phi\left(  y_{K}\right)  =a_{K}\ \ \ \ \ \ \ \ \ \ \text{for every
}K\in\mathfrak{K}\right)
\]
(by the definition of $\operatorname*{ev}\nolimits_{\mathbf{a}}$). Thus,
$\operatorname*{ev}\nolimits_{\mathbf{a}}$ is a $\mathbb{Z}$-algebra
homomorphism and satisfies%
\begin{equation}
\operatorname*{ev}\nolimits_{\mathbf{a}}\left(  y_{K}\right)  =a_{K}%
\ \ \ \ \ \ \ \ \ \ \text{for every }K\in\mathfrak{K}.
\label{pf.thm.chrompol.varis.eva}%
\end{equation}


The construction of $\operatorname*{ev}\nolimits_{\mathbf{a}}\left[  x\right]
$ shows that
\begin{equation}
\left(  \operatorname*{ev}\nolimits_{\mathbf{a}}\left[  x\right]  \right)
\left(  u\right)  =\operatorname*{ev}\nolimits_{\mathbf{a}}\left(  u\right)
\ \ \ \ \ \ \ \ \ \ \text{for every }u\in R.
\label{pf.thm.chrompol.varis.eva(u)}%
\end{equation}
Now, for every $K\in\mathfrak{K}$, we have%
\begin{align}
\left(  \operatorname*{ev}\nolimits_{\mathbf{a}}\left[  x\right]  \right)
\left(  y_{K}\right)   &  =\operatorname*{ev}\nolimits_{\mathbf{a}}\left(
y_{K}\right)  \ \ \ \ \ \ \ \ \ \ \left(  \text{by
(\ref{pf.thm.chrompol.varis.eva(u)}), applied to }u=y_{K}\right) \nonumber\\
&  =a_{K}\ \ \ \ \ \ \ \ \ \ \left(  \text{by (\ref{pf.thm.chrompol.varis.eva}%
)}\right)  . \label{pf.thm.chrompol.varis.evax}%
\end{align}


By its definition, the homomorphism $\operatorname*{ev}\nolimits_{\mathbf{a}%
}\left[  x\right]  $ preserves $\mathbb{Z}\left[  x\right]  $ (or, rather,
sends every element of $\mathbb{Z}\left[  x\right]  \subseteq R\left[
x\right]  $ to the corresponding element of $\mathbb{Z}\left[  x\right]
\subseteq R\left[  x\right]  $). In other words, $\left(  \operatorname*{ev}%
\nolimits_{\mathbf{a}}\left[  x\right]  \right)  \left(  Q\right)  =Q$ for
every $Q\in\mathbb{Z}\left[  x\right]  $. Applying this to $Q=\chi_{G}$, we
obtain%
\[
\left(  \operatorname*{ev}\nolimits_{\mathbf{a}}\left[  x\right]  \right)
\left(  \chi_{G}\right)  =\chi_{G}.
\]
Thus,%
\begin{align*}
\chi_{G}  &  =\left(  \operatorname*{ev}\nolimits_{\mathbf{a}}\left[
x\right]  \right)  \left(  \underbrace{\chi_{G}}_{=\widetilde{P}%
=\sum_{F\subseteq E}\left(  -1\right)  ^{\left\vert F\right\vert }\left(
\prod_{\substack{K\in\mathfrak{K};\\K\subseteq F}}y_{K}\right)
x^{\operatorname*{conn}\left(  V,F\right)  }}\right) \\
&  =\left(  \operatorname*{ev}\nolimits_{\mathbf{a}}\left[  x\right]  \right)
\left(  \sum_{F\subseteq E}\left(  -1\right)  ^{\left\vert F\right\vert
}\left(  \prod_{\substack{K\in\mathfrak{K};\\K\subseteq F}}y_{K}\right)
x^{\operatorname*{conn}\left(  V,F\right)  }\right) \\
&  =\sum_{F\subseteq E}\left(  -1\right)  ^{\left\vert F\right\vert }\left(
\prod_{\substack{K\in\mathfrak{K};\\K\subseteq F}}\underbrace{\left(
\operatorname*{ev}\nolimits_{\mathbf{a}}\left[  x\right]  \right)  \left(
y_{K}\right)  }_{\substack{=a_{K}\\\text{(by (\ref{pf.thm.chrompol.varis.evax}%
))}}}\right)  \left(  \underbrace{\left(  \operatorname*{ev}%
\nolimits_{\mathbf{a}}\left[  x\right]  \right)  \left(  x\right)  }%
_{=x}\right)  ^{\operatorname*{conn}\left(  V,F\right)  }\\
&  \ \ \ \ \ \ \ \ \ \ \left(  \text{since }\operatorname*{ev}%
\nolimits_{\mathbf{a}}\left[  x\right]  \text{ is a }\mathbb{Z}\text{-algebra
homomorphism}\right) \\
&  =\sum_{F\subseteq E}\left(  -1\right)  ^{\left\vert F\right\vert }\left(
\prod_{\substack{K\in\mathfrak{K};\\K\subseteq F}}a_{K}\right)
x^{\operatorname*{conn}\left(  V,F\right)  }.
\end{align*}
This proves Theorem \ref{thm.chrompol.varis}.
\end{proof}
\end{verlong}

\begin{vershort}
Now that Theorem \ref{thm.chrompol.varis} is proven, we could derive Theorem
\ref{thm.chrompol.empty}, Corollary \ref{cor.chrompol.K-free} and Corollary
\ref{cor.chrompol.NBC} from it in the same way as we have derived Theorem
\ref{thm.chromsym.empty}, Corollary \ref{cor.chromsym.K-free} and Corollary
\ref{cor.chromsym.NBC} from Theorem \ref{thm.chromsym.varis}. We leave the
details to the reader.
\end{vershort}

\begin{verlong}
Now that Theorem \ref{thm.chrompol.varis} is proven, we can derive Theorem
\ref{thm.chrompol.empty}, Corollary \ref{cor.chrompol.K-free} and Corollary
\ref{cor.chrompol.NBC} from it in the same way as we have derived Theorem
\ref{thm.chromsym.empty}, Corollary \ref{cor.chromsym.K-free} and Corollary
\ref{cor.chromsym.NBC} from Theorem \ref{thm.chromsym.varis}. Here are the details:
\end{verlong}

\begin{verlong}
\begin{proof}
[Proof of Corollary \ref{cor.chrompol.K-free}.]We can apply Theorem
\ref{thm.chrompol.varis} to $0$ instead of $a_{K}$. As a result, we obtain%
\begin{equation}
\chi_{G}=\sum_{F\subseteq E}\left(  -1\right)  ^{\left\vert F\right\vert
}\left(  \prod_{\substack{K\in\mathfrak{K};\\K\subseteq F}}0\right)
x^{\operatorname*{conn}\left(  V,F\right)  }. \label{pf.cor.chrompol.K-free.0}%
\end{equation}
Now, if $F$ is any subset of $E$, then%
\begin{equation}
\prod_{\substack{K\in\mathfrak{K};\\K\subseteq F}}0=%
\begin{cases}
1, & \text{if }F\text{ is }\mathfrak{K}\text{-free;}\\
0, & \text{if }F\text{ is not }\mathfrak{K}\text{-free}%
\end{cases}
\label{pf.cor.chrompol.K-free.1}%
\end{equation}
\footnote{\textit{Proof of (\ref{pf.cor.chrompol.K-free.1}):} The equality
(\ref{pf.cor.chrompol.K-free.1}) has already been proven in our proof of
Corollary \ref{cor.chromsym.K-free}.}. Thus, (\ref{pf.cor.chrompol.K-free.0})
becomes%
\begin{align*}
\chi_{G}  &  =\sum_{F\subseteq E}\left(  -1\right)  ^{\left\vert F\right\vert
}\underbrace{\left(  \prod_{\substack{K\in\mathfrak{K};\\K\subseteq
F}}0\right)  }_{\substack{=%
\begin{cases}
1, & \text{if }F\text{ is }\mathfrak{K}\text{-free;}\\
0, & \text{if }F\text{ is not }\mathfrak{K}\text{-free}%
\end{cases}
\\\text{(by (\ref{pf.cor.chrompol.K-free.1}))}}}x^{\operatorname*{conn}\left(
V,F\right)  }\\
&  =\sum_{F\subseteq E}\left(  -1\right)  ^{\left\vert F\right\vert }%
\begin{cases}
1, & \text{if }F\text{ is }\mathfrak{K}\text{-free;}\\
0, & \text{if }F\text{ is not }\mathfrak{K}\text{-free}%
\end{cases}
x^{\operatorname*{conn}\left(  V,F\right)  }\\
&  =\sum_{\substack{F\subseteq E;\\F\text{ is }\mathfrak{K}\text{-free}%
}}\left(  -1\right)  ^{\left\vert F\right\vert }\underbrace{%
\begin{cases}
1, & \text{if }F\text{ is }\mathfrak{K}\text{-free;}\\
0, & \text{if }F\text{ is not }\mathfrak{K}\text{-free}%
\end{cases}
}_{\substack{=1\\\text{(since }F\text{ is }\mathfrak{K}\text{-free)}%
}}x^{\operatorname*{conn}\left(  V,F\right)  }\\
&  \ \ \ \ \ \ \ \ \ \ +\sum_{\substack{F\subseteq E;\\F\text{ is not
}\mathfrak{K}\text{-free}}}\left(  -1\right)  ^{\left\vert F\right\vert
}\underbrace{%
\begin{cases}
1, & \text{if }F\text{ is }\mathfrak{K}\text{-free;}\\
0, & \text{if }F\text{ is not }\mathfrak{K}\text{-free}%
\end{cases}
}_{\substack{=0\\\text{(since }F\text{ is not }\mathfrak{K}\text{-free)}%
}}x^{\operatorname*{conn}\left(  V,F\right)  }\\
&  =\sum_{\substack{F\subseteq E;\\F\text{ is }\mathfrak{K}\text{-free}%
}}\left(  -1\right)  ^{\left\vert F\right\vert }x^{\operatorname*{conn}\left(
V,F\right)  }+\underbrace{\sum_{\substack{F\subseteq E;\\F\text{ is not
}\mathfrak{K}\text{-free}}}\left(  -1\right)  ^{\left\vert F\right\vert
}0x^{\operatorname*{conn}\left(  V,F\right)  }}_{=0}\\
&  =\sum_{\substack{F\subseteq E;\\F\text{ is }\mathfrak{K}\text{-free}%
}}\left(  -1\right)  ^{\left\vert F\right\vert }x^{\operatorname*{conn}\left(
V,F\right)  }.
\end{align*}
This proves Corollary \ref{cor.chrompol.K-free}.
\end{proof}
\end{verlong}

\begin{verlong}
\begin{proof}
[Proof of Corollary \ref{cor.chrompol.NBC}.]Let $\mathfrak{K}$ be the set of
all broken circuits of $G$. Thus, the elements of $\mathfrak{K}$ are the
broken circuits of $G$.

Now, for every subset $F$ of $E$, we have the following equivalence of
statements:%
\begin{align*}
&  \ \left(  F\text{ is }\mathfrak{K}\text{-free}\right) \\
&  \Longleftrightarrow\ \left(  F\text{ contains no }K\in\mathfrak{K}\text{ as
a subset}\right) \\
&  \ \ \ \ \ \ \ \ \ \ \left(
\begin{array}
[c]{c}%
\text{because }F\text{ is }\mathfrak{K}\text{-free if and only if }F\text{
contains no }K\in\mathfrak{K}\text{ as a subset}\\
\text{(by the definition of \textquotedblleft}\mathfrak{K}%
\text{-free\textquotedblright)}%
\end{array}
\right) \\
&  \Longleftrightarrow\ \left(  F\text{ contains no element of }%
\mathfrak{K}\text{ as a subset}\right) \\
&  \Longleftrightarrow\ \left(  F\text{ contains no broken circuit of }G\text{
as a subset}\right)
\end{align*}
(since the elements of $\mathfrak{K}$ are the broken circuits of $G$). Hence,
$\sum_{\substack{F\subseteq E;\\F\text{ is }\mathfrak{K}\text{-free}}%
}=\sum_{\substack{F\subseteq E;\\F\text{ contains no broken}\\\text{circuit of
}G\text{ as a subset}}}$ (an equality between summation signs). Now, Corollary
\ref{cor.chrompol.K-free} yields%
\[
\chi_{G}=\underbrace{\sum_{\substack{F\subseteq E;\\F\text{ is }%
\mathfrak{K}\text{-free}}}}_{=\sum_{\substack{F\subseteq E;\\F\text{ contains
no broken}\\\text{circuit of }G\text{ as a subset}}}}\left(  -1\right)
^{\left\vert F\right\vert }x^{\operatorname*{conn}\left(  V,F\right)  }%
=\sum_{\substack{F\subseteq E;\\F\text{ contains no broken}\\\text{circuit of
}G\text{ as a subset}}}\left(  -1\right)  ^{\left\vert F\right\vert
}x^{\operatorname*{conn}\left(  V,F\right)  }.
\]
This proves Corollary \ref{cor.chrompol.NBC}.
\end{proof}
\end{verlong}

\begin{verlong}
\begin{proof}
[Proof of Theorem \ref{thm.chrompol.empty}.]Let $X$ be the totally ordered set
$\left\{  1\right\}  $ (equipped with the only possible order on this set).
Let $\ell:E\rightarrow X$ be the function sending each $e\in E$ to $1\in X$.
Let $\mathfrak{K}$ be the empty set. Clearly, $\mathfrak{K}$ is a set of
broken circuits of $G$. Theorem \ref{thm.chrompol.varis} (applied to $0$
instead of $a_{K}$) yields%
\begin{align*}
\chi_{G}  &  =\sum_{F\subseteq E}\left(  -1\right)  ^{\left\vert F\right\vert
}\underbrace{\left(  \prod_{\substack{K\in\mathfrak{K};\\K\subseteq
F}}0\right)  }_{\substack{=\left(  \text{empty product}\right)  \\\text{(since
}\mathfrak{K}\text{ is the empty set)}}}x^{\operatorname*{conn}\left(
V,F\right)  }\\
&  =\sum_{F\subseteq E}\left(  -1\right)  ^{\left\vert F\right\vert
}\underbrace{\left(  \text{empty product}\right)  }_{=1}%
x^{\operatorname*{conn}\left(  V,F\right)  }=\sum_{F\subseteq E}\left(
-1\right)  ^{\left\vert F\right\vert }x^{\operatorname*{conn}\left(
V,F\right)  }.
\end{align*}
This proves Theorem \ref{thm.chrompol.empty}.
\end{proof}
\end{verlong}

\subsection{Special case: Whitney's Broken-Circuit Theorem}

Corollary \ref{cor.chrompol.NBC} is commonly stated in the following
simplified (if less general) form:

\Needspace{4cm}

\begin{corollary}
\label{cor.chrompol.NBCfor}Let $G=\left(  V,E\right)  $ be a finite graph. Let
$X$ be a totally ordered set. Let $\ell:E\rightarrow X$ be an injective
function. Then,%
\[
\chi_{G}=\sum_{\substack{F\subseteq E;\\F\text{ contains no broken}%
\\\text{circuit of }G\text{ as a subset}}}\left(  -1\right)  ^{\left\vert
F\right\vert }x^{\left\vert V\right\vert -\left\vert F\right\vert }.
\]

\end{corollary}

Corollary \ref{cor.chrompol.NBCfor} is known as \textit{Whitney's
Broken-Circuit theorem} (see, e.g., \cite{BlaSag86}).

Notice that $\ell$ is required to be injective in Corollary
\ref{cor.chrompol.NBCfor}; the purpose of this requirement is to ensure that
every circuit of $G$ has a unique edge $e$ with maximum $\ell\left(  e\right)
$, and thus induces a broken circuit of $G$. The proof of Corollary
\ref{cor.chrompol.NBCfor} relies on the following standard result:

\begin{lemma}
\label{lem.conn.forest}Let $\left(  V,F\right)  $ be a finite graph. Assume
that $\left(  V,F\right)  $ has no circuits. Then, $\operatorname*{conn}%
\left(  V,F\right)  =\left\vert V\right\vert -\left\vert F\right\vert $.
\end{lemma}

(A graph which has no circuits is commonly known as a \textit{forest}.)

\begin{verlong}
[TODO: Write a proof of Lemma \ref{lem.conn.forest}.]
\end{verlong}

Lemma \ref{lem.conn.forest} is both extremely elementary and well-known; for
example, it appears in \cite[Proposition 10.6]{Bona11} and in \cite[\S I.2,
Corollary 6]{Bollobas}. Let us now see how it entails Corollary
\ref{cor.chrompol.NBCfor}:

\begin{vershort}
\begin{proof}
[Proof of Corollary \ref{cor.chrompol.NBCfor}.]Corollary
\ref{cor.chrompol.NBCfor} follows from Corollary \ref{cor.chrompol.NBC}.
Indeed, the injectivity of $\ell$ shows that every circuit of $G$ has a unique
edge $e$ with maximum $\ell\left(  e\right)  $, and thus contains a broken
circuit of $G$. Therefore, if a subset $F$ of $E$ contains no broken circuit
of $G$ as a subset, then $F$ contains no circuit of $G$ either, and therefore
the graph $\left(  V,F\right)  $ contains no circuits; but this entails that
$\operatorname*{conn}\left(  V,F\right)  =\left\vert V\right\vert -\left\vert
F\right\vert $ (by Lemma \ref{lem.conn.forest}). Hence, Corollary
\ref{cor.chrompol.NBC} immediately yields Corollary \ref{cor.chrompol.NBCfor}.
\end{proof}
\end{vershort}

\begin{verlong}
\begin{proof}
[Proof of Corollary \ref{cor.chrompol.NBCfor}.]We first claim that if $F$ is a
subset of $E$ such that $F$ contains no broken circuit of $G$ as a subset,
then%
\begin{equation}
\operatorname*{conn}\left(  V,F\right)  =\left\vert V\right\vert -\left\vert
F\right\vert . \label{pf.cor.chrompol.NBCfor.1}%
\end{equation}


\textit{Proof of (\ref{pf.cor.chrompol.NBCfor.1}):} Let $F$ be a subset of $E$
such that $F$ contains no broken circuit of $G$ as a subset. We shall now show
that the graph $\left(  V,F\right)  $ has no circuits.

Indeed, assume the contrary (for the sake of contradiction). Thus, the graph
$\left(  V,F\right)  $ has a circuit. In other words, there exists a circuit
$D$ of the graph $\left(  V,F\right)  $. Consider this $D$.

The set $D$ is a circuit of $\left(  V,F\right)  $. In other words, the set
$D$ has the form $\left\{  \left\{  v_{1},v_{2}\right\}  ,\left\{  v_{2}%
,v_{3}\right\}  ,\ldots,\left\{  v_{m},v_{m+1}\right\}  \right\}  $, where
$\left(  v_{1},v_{2},\ldots,v_{m+1}\right)  $ is a cycle of $\left(
V,F\right)  $ (by the definition of a \textquotedblleft
circuit\textquotedblright). Consider this cycle $\left(  v_{1},v_{2}%
,\ldots,v_{m+1}\right)  $. We thus have
\begin{equation}
D=\left\{  \left\{  v_{1},v_{2}\right\}  ,\left\{  v_{2},v_{3}\right\}
,\ldots,\left\{  v_{m},v_{m+1}\right\}  \right\}  .
\label{pf.cor.chrompol.NBCfor.5}%
\end{equation}


The list $\left(  v_{1},v_{2},\ldots,v_{m+1}\right)  $ is a cycle of $\left(
V,F\right)  $. According to the definition of a \textquotedblleft
cycle\textquotedblright, this means that this list is a list of elements of
$V$ satisfying the following four properties:

\begin{itemize}
\item We have $m>1$.

\item We have $v_{m+1}=v_{1}$.

\item The vertices $v_{1},v_{2},\ldots,v_{m}$ are pairwise distinct.

\item We have $\left\{  v_{i},v_{i+1}\right\}  \in F$ for every $i\in\left\{
1,2,\ldots,m\right\}  $.
\end{itemize}

Thus, $\left(  v_{1},v_{2},\ldots,v_{m+1}\right)  $ is a list of elements of
$V$ satisfying the four properties that we have just mentioned. Notice that%
\begin{align*}
D  &  =\left\{  \left\{  v_{1},v_{2}\right\}  ,\left\{  v_{2},v_{3}\right\}
,\ldots,\left\{  v_{m},v_{m+1}\right\}  \right\} \\
&  =\left\{  \left\{  v_{i},v_{i+1}\right\}  \ \mid\ i\in\left\{
1,2,\ldots,m\right\}  \right\}  \subseteq F
\end{align*}
(since $\left\{  v_{i},v_{i+1}\right\}  \in F$ for every $i\in\left\{
1,2,\ldots,m\right\}  $).

Now, we have $\left\{  v_{i},v_{i+1}\right\}  \in F\subseteq E$ for every
$i\in\left\{  1,2,\ldots,m\right\}  $. Hence, $\left(  v_{1},v_{2}%
,\ldots,v_{m+1}\right)  $ is a list of elements of $V$ satisfying the
following four properties:

\begin{itemize}
\item We have $m>1$.

\item We have $v_{m+1}=v_{1}$.

\item The vertices $v_{1},v_{2},\ldots,v_{m}$ are pairwise distinct.

\item We have $\left\{  v_{i},v_{i+1}\right\}  \in E$ for every $i\in\left\{
1,2,\ldots,m\right\}  $.
\end{itemize}

According to the definition of a \textquotedblleft cycle\textquotedblright,
this means that the list $\left(  v_{1},v_{2},\ldots,v_{m+1}\right)  $ is a
cycle of $\left(  V,E\right)  $.

Thus, we conclude that the list $\left(  v_{1},v_{2},\ldots,v_{m+1}\right)  $
is a cycle of $\left(  V,E\right)  $. In other words, the list $\left(
v_{1},v_{2},\ldots,v_{m+1}\right)  $ is a cycle of $G$ (since $G=\left(
V,E\right)  $). Thus, the set $\left\{  \left\{  v_{1},v_{2}\right\}
,\left\{  v_{2},v_{3}\right\}  ,\ldots,\left\{  v_{m},v_{m+1}\right\}
\right\}  $ is a circuit of $G$ (by the definition of a \textquotedblleft
circuit\textquotedblright). In view of (\ref{pf.cor.chrompol.NBCfor.5}), this
rewrites as follows: The set $D$ is a circuit of $G$.

No two distinct edges in $D$ have the same label\footnote{\textit{Proof.}
Assume the contrary. Thus, two distinct edges in $D$ have the same label. In
other words, there exist two distinct edges $e$ and $e^{\prime}$ in $D$ such
that $e$ and $e^{\prime}$ have the same label. Consider these $e$ and
$e^{\prime}$.
\par
The edges $e$ and $e^{\prime}$ have the same label. In other words, the label
of $e$ equals the label of $e^{\prime}$. In other words, $\ell\left(
e\right)  $ equals $\ell\left(  e^{\prime}\right)  $ (since the label of $e$
is $\ell\left(  e\right)  $ (by the definition of \textquotedblleft
label\textquotedblright), whereas the label of $e^{\prime}$ is $\ell\left(
e^{\prime}\right)  $ (by the definition of \textquotedblleft
label\textquotedblright)). In other words, $\ell\left(  e\right)  =\ell\left(
e^{\prime}\right)  $. Since the map $\ell$ is injective, this shows that
$e=e^{\prime}$. This contradicts the assumption that $e$ and $e^{\prime}$ are
distinct. This contradiction proves that our assumption was wrong. Qed.}.

From $m>1$, we obtain%
\[
\left\{  v_{1},v_{2}\right\}  \in\left\{  \left\{  v_{1},v_{2}\right\}
,\left\{  v_{2},v_{3}\right\}  ,\ldots,\left\{  v_{m},v_{m+1}\right\}
\right\}  =D
\]
(by (\ref{pf.cor.chrompol.NBCfor.5})). Hence, the set $D$ is nonempty (since
it contains $\left\{  v_{1},v_{2}\right\}  $). Thus, $D$ is a nonempty finite
set. Hence, there exists an edge in $D$ having maximum label. Let $f$ be this
edge. Clearly, $D\setminus\left\{  f\right\}  \subseteq D\subseteq F\subseteq
E$. Thus, $D\setminus\left\{  f\right\}  $ is a subset of $E$.

The edge $f$ is an edge in $D$ having maximum label. Since no other edge in
$D$ has the same label as $f$ (because no two distinct edges in $D$ have the
same label), this shows that the edge $f$ is the \textbf{unique} edge in $D$
having maximum label. Therefore, $D\setminus\left\{  f\right\}  $ is a subset
of $E$ having the form $C\setminus\left\{  e\right\}  $, where $C$ is a
circuit of $G$, and where $e$ is the unique edge in $C$ having maximum label
(among the edges in $C$)\ \ \ \ \footnote{Namely, $D\setminus\left\{
f\right\}  $ has this form for $C=D$ and $e=f$.}. In other words,
$D\setminus\left\{  f\right\}  $ is a broken circuit of $G$ (since
$D\setminus\left\{  f\right\}  $ is a broken circuit of $G$ if and only if
$D\setminus\left\{  f\right\}  $ is a subset of $E$ having the form
$C\setminus\left\{  e\right\}  $, where $C$ is a circuit of $G$, and where $e$
is the unique edge in $C$ having maximum label (among the edges in
$C$)\ \ \ \ \footnote{by the definition of a \textquotedblleft broken
circuit\textquotedblright}). This broken circuit $D\setminus\left\{
f\right\}  $ satisfies $D\setminus\left\{  f\right\}  \subseteq F$. Thus,
there exists a broken circuit $B$ of $G$ such that $B\subseteq F$ (namely,
$B=D\setminus\left\{  f\right\}  $).

But the set $F$ contains no broken circuit of $G$ as a subset. In other words,
there exists no broken circuit $B$ of $G$ such that $B\subseteq F$. This
contradicts the fact that there exists a broken circuit $B$ of $G$ such that
$B\subseteq F$. This contradiction proves that our assumption was wrong.
Hence, the graph $\left(  V,F\right)  $ has no circuits. Lemma
\ref{lem.conn.forest} thus shows that $\operatorname*{conn}\left(  V,F\right)
=\left\vert V\right\vert -\left\vert F\right\vert $. This proves
(\ref{pf.cor.chrompol.NBCfor.1}).

Now, Corollary \ref{cor.chrompol.NBC} yields%
\[
\chi_{G}=\sum_{\substack{F\subseteq E;\\F\text{ contains no broken}%
\\\text{circuit of }G\text{ as a subset}}}\left(  -1\right)  ^{\left\vert
F\right\vert }\underbrace{x^{\operatorname*{conn}\left(  V,F\right)  }%
}_{\substack{=x^{\left\vert V\right\vert -\left\vert F\right\vert
}\\\text{(since }\operatorname*{conn}\left(  V,F\right)  =\left\vert
V\right\vert -\left\vert F\right\vert \\\text{(by
(\ref{pf.cor.chrompol.NBCfor.1})))}}}=\sum_{\substack{F\subseteq E;\\F\text{
contains no broken}\\\text{circuit of }G\text{ as a subset}}}\left(
-1\right)  ^{\left\vert F\right\vert }x^{\left\vert V\right\vert -\left\vert
F\right\vert }.
\]
This proves Corollary \ref{cor.chrompol.NBCfor}.
\end{proof}
\end{verlong}

\section{Application to transitive directed graphs}

We shall now see an application of Corollary \ref{cor.chrompol.K-free} to
graphs which are obtained from certain directed graphs by \textquotedblleft
forgetting the directions of the edges\textquotedblright. Let us first
introduce the notations involved:

\begin{definition}
\label{def.digraph}\textbf{(a)} A \textit{digraph} means a pair $\left(
V,A\right)  $, where $V$ is a set, and where $A$ is a subset of $V^{2}$.
Digraphs are also called \textit{directed graphs}. A digraph $\left(
V,A\right)  $ is said to be \textit{finite} if the set $V$ is finite. If
$D=\left(  V,A\right)  $ is a digraph, then the elements of $V$ are called the
\textit{vertices} of the digraph $D$, while the elements of $A$ are called the
\textit{arcs} (or the \textit{directed edges}) of the digraph $D$. If
$a=\left(  v,w\right)  $ is an arc of a digraph $D$, then $v$ is called the
\textit{source} of $a$, whereas $w$ is called the \textit{target} of $a$.

\textbf{(b)} A digraph $\left(  V,A\right)  $ is said to be \textit{loopless}
if every $v\in V$ satisfies $\left(  v,v\right)  \notin A$. (In other words, a
digraph is loopless if and only if it has no arc whose source and target are identical.)

\textbf{(c)} A digraph $\left(  V,A\right)  $ is said to be
\textit{transitive} if it has the following property: For any $u\in V$, $v\in
V$ and $w\in V$ satisfying $\left(  u,v\right)  \in A$ and $\left(
v,w\right)  \in A$, we have $\left(  u,w\right)  \in A$.

\textbf{(d)} A digraph $\left(  V,A\right)  $ is said to be $2$%
\textit{-path-free} if there exist no three elements $u$, $v$ and $w$ of $V$
satisfying $\left(  u,v\right)  \in A$ and $\left(  v,w\right)  \in A$.

\textbf{(e)} Let $D=\left(  V,A\right)  $ be a loopless digraph. Define a map
$\operatorname*{set}:A\rightarrow\dbinom{V}{2}$ by setting%
\[
\left(  \operatorname*{set}\left(  v,w\right)  =\left\{  v,w\right\}
\ \ \ \ \ \ \ \ \ \ \text{for every }\left(  v,w\right)  \in A\right)  .
\]
(It is easy to see that $\operatorname*{set}$ is well-defined, because
$\left(  V,A\right)  $ is loopless.) The graph $\left(  V,\operatorname*{set}%
A\right)  $ will be denoted by $\underline{D}$.
\end{definition}

We can now state our application of Corollary \ref{cor.chrompol.K-free},
answering a question suggested by Alexander Postnikov:

\begin{proposition}
\label{prop.digraph.2pf-chrom}Let $D=\left(  V,A\right)  $ be a finite
transitive loopless digraph. Then,%
\[
\chi_{\underline{D}}=\sum_{\substack{F\subseteq A;\\\text{the digraph }\left(
V,F\right)  \text{ is }2\text{-path-free}}}\left(  -1\right)  ^{\left\vert
F\right\vert }x^{\operatorname*{conn}\left(  V,\operatorname*{set}F\right)
}.
\]

\end{proposition}

\begin{vershort}
\begin{proof}
[Proof of Proposition \ref{prop.digraph.2pf-chrom}.]Let $E=\operatorname*{set}%
A$. Then, the definition of $\underline{D}$ yields $\underline{D}=\left(
V,\underbrace{\operatorname*{set}A}_{=E}\right)  =\left(  V,E\right)  $.

The map $\operatorname*{set}:A\rightarrow\dbinom{V}{2}$ (which sends every arc
$\left(  v,w\right)  \in A$ to $\left\{  v,w\right\}  \in\dbinom{V}{2}$)
restricts to a surjection $A\rightarrow E$ (since $E=\operatorname*{set}A$).
Let us denote this surjection by $\pi$. Thus, $\pi$ is a map from $A$ to $E$
sending each arc $\left(  v,w\right)  \in A$ to $\left\{  v,w\right\}  \in E$.
We shall soon see that $\pi$ is a bijection.

We define a partial order on the set $V$ as follows: For $i\in V$ and $j\in
V$, we set $i<j$ if and only if $\left(  i,j\right)  \in A$ (that is, if and
only if there is an arc from $i$ to $j$ in $D$). This is a well-defined
partial order\footnote{Indeed, the relation $<$ that we have just defined is
transitive (since the digraph $\left(  V,A\right)  $ is transitive) and
antisymmetric (since the digraph $\left(  V,A\right)  $ is loopless).}. Thus,
$V$ becomes a poset. For every $i\in V$ and $j\in V$ satisfying $i\leq j$, we
let $\left[  i,j\right]  $ denote the interval $\left\{  k\in V\ \mid\ i\leq
k\leq j\right\}  $ of the poset $V$.

There exist no $i,j\in V$ such that both $\left(  i,j\right)  $ and $\left(
j,i\right)  $ belong to $A$ (because if such $i$ and $j$ would exist, then
they would satisfy $i<j$ and $j<i$, but this would contradict the fact that
$V$ is a poset). Hence, the projection $\pi:A\rightarrow E$ is injective, and
thus bijective (since we already know that $\pi$ is surjective). Hence, its
inverse map $\pi^{-1}:E\rightarrow A$ is well-defined. For every subset $F$ of
$E$, we have%
\begin{align}
F  &  =\pi\left(  \pi^{-1}\left(  F\right)  \right)
\ \ \ \ \ \ \ \ \ \ \left(  \text{since }\pi\text{ is bijective}\right)
\nonumber\\
&  =\operatorname*{set}\left(  \pi^{-1}\left(  F\right)  \right)
\label{pf.prop.digraph.2pf-chrom.short.F=set}%
\end{align}
(since $\pi$ is a restriction of the map $\operatorname*{set}$).

For any $\left(  u,v\right)  \in A$ and any subset $F$ of $E$, we have the
following logical equivalence:%
\begin{equation}
\left(  \left\{  u,v\right\}  \in F\right)  \ \Longleftrightarrow\ \left(
\left(  u,v\right)  \in\pi^{-1}\left(  F\right)  \right)
\label{pf.prop.digraph.2pf-chrom.short.equiv}%
\end{equation}
\footnote{\textit{Proof of (\ref{pf.prop.digraph.2pf-chrom.short.equiv}):} Let
$\left(  u,v\right)  \in A$, and let $F$ be a subset of $E$. We need to prove
the equivalence (\ref{pf.prop.digraph.2pf-chrom.short.equiv}).
\par
From $\left(  u,v\right)  \in A$, we see that $\pi\left(  u,v\right)  $ is
well-defined. The definition of $\pi$ shows that $\pi\left(  u,v\right)
=\left\{  u,v\right\}  $. Hence, we have the following chain of equivalences:%
\[
\left(  \underbrace{\left\{  u,v\right\}  }_{=\pi\left(  u,v\right)  }\in
F\right)  \ \Longleftrightarrow\ \left(  \pi\left(  u,v\right)  \in F\right)
\ \Longleftrightarrow\ \left(  \left(  u,v\right)  \in\pi^{-1}\left(
F\right)  \right)  .
\]
This proves (\ref{pf.prop.digraph.2pf-chrom.short.equiv}).}.

Define a function $\ell^{\prime}:A\rightarrow\mathbb{N}$ by%
\[
\ell^{\prime}\left(  i,j\right)  =\left\vert \left[  i,j\right]  \right\vert
\ \ \ \ \ \ \ \ \ \ \text{for all }\left(  i,j\right)  \in A.
\]
Define a function $\ell:E\rightarrow\mathbb{N}$ by $\ell=\ell^{\prime}\circ
\pi^{-1}$. Thus, $\ell\circ\pi=\ell^{\prime}$. Therefore,%
\begin{equation}
\ell\left(  \underbrace{\left\{  i,j\right\}  }_{=\pi\left(  i,j\right)
}\right)  =\underbrace{\left(  \ell\circ\pi\right)  }_{=\ell^{\prime}}\left(
i,j\right)  =\ell^{\prime}\left(  i,j\right)  =\left\vert \left[  i,j\right]
\right\vert \label{pf.prop.digraph.2pf-chrom.short.l}%
\end{equation}
for all $\left(  i,j\right)  \in A$.

Let $\mathfrak{K}$ be the set
\[
\left\{  \left\{  \left\{  i,k\right\}  ,\left\{  k,j\right\}  \right\}
\ \mid\ \left(  i,k\right)  \in A\text{ and }\left(  k,j\right)  \in
A\right\}  .
\]
Each $K\in\mathfrak{K}$ is a broken circuit of $\underline{D}$%
\ \ \ \ \footnote{\textit{Proof.} Let $K\in\mathfrak{K}$. Then, $K=\left\{
\left\{  i,k\right\}  ,\left\{  k,j\right\}  \right\}  $ for some $\left(
i,k\right)  \in A$ and $\left(  k,j\right)  \in A$ (by the definition of
$\mathfrak{K}$). Consider these $\left(  i,k\right)  $ and $\left(
k,j\right)  $. Since $\left(  V,A\right)  $ is transitive, we have $\left(
i,j\right)  \in A$. Thus, $\left\{  i,k\right\}  $, $\left\{  k,j\right\}  $
and $\left\{  i,j\right\}  $ are edges of $\underline{D}$. These edges form a
circuit of $\underline{D}$. In particular, $i$, $j$ and $k$ are pairwise
distinct.
\par
Applications of (\ref{pf.prop.digraph.2pf-chrom.short.l}) yield $\ell\left(
\left\{  i,j\right\}  \right)  =\left\vert \left[  i,j\right]  \right\vert $,
$\ell\left(  \left\{  i,k\right\}  \right)  =\left\vert \left[  i,k\right]
\right\vert $ and $\ell\left(  \left\{  k,j\right\}  \right)  =\left\vert
\left[  k,j\right]  \right\vert $.
\par
But we have $i<k$ (since $\left(  i,k\right)  \in A$) and $k<j$ (since
$\left(  k,j\right)  \in A$). Hence, $\left[  i,k\right]  $ is a proper subset
of $\left[  i,j\right]  $. (It is proper because it does not contain $j$,
whereas $\left[  i,j\right]  $ does.) Hence, $\left\vert \left[  i,k\right]
\right\vert <\left\vert \left[  i,j\right]  \right\vert $. Thus, $\ell\left(
\left\{  i,j\right\}  \right)  =\left\vert \left[  i,j\right]  \right\vert
>\left\vert \left[  i,k\right]  \right\vert =\ell\left(  \left\{  i,k\right\}
\right)  $. Similarly, $\ell\left(  \left\{  i,j\right\}  \right)
>\ell\left(  \left\{  k,j\right\}  \right)  $. The last two inequalities show
that $\left\{  i,j\right\}  $ is the unique edge of the circuit $\left\{
\left\{  i,k\right\}  ,\left\{  k,j\right\}  ,\left\{  i,j\right\}  \right\}
$ having maximum label. Hence, $\left\{  \left\{  i,k\right\}  ,\left\{
k,j\right\}  ,\left\{  i,j\right\}  \right\}  \setminus\left\{  \left\{
i,j\right\}  \right\}  $ is a broken circuit of $\underline{D}$. Since%
\begin{align*}
\left\{  \left\{  i,k\right\}  ,\left\{  k,j\right\}  ,\left\{  i,j\right\}
\right\}  \setminus\left\{  \left\{  i,j\right\}  \right\}   &  =\left\{
\left\{  i,k\right\}  ,\left\{  k,j\right\}  \right\}
\ \ \ \ \ \ \ \ \ \ \left(  \text{since }i\text{, }j\text{ and }k\text{ are
pairwise distinct}\right) \\
&  =K,
\end{align*}
this shows that $K$ is a broken circuit of $\underline{D}$, qed.}. Thus,
$\mathfrak{K}$ is a set of broken circuits of $\underline{D}$.

A subset $F$ of $E$ is $\mathfrak{K}$-free if and only if the digraph $\left(
V,\pi^{-1}\left(  F\right)  \right)  $ is $2$%
-path-free\footnote{\textit{Proof.} Let $F$ be a subset of $E$. Then, we have
the following equivalence of statements:%
\begin{align*}
&  \ \left(  F\text{ is }\mathfrak{K}\text{-free}\right) \\
&  \Longleftrightarrow\ \left(  \left\{  \left\{  i,k\right\}  ,\left\{
k,j\right\}  \right\}  \not \subseteq F\text{ whenever }\left(  i,k\right)
\in A\text{ and }\left(  k,j\right)  \in A\right) \\
&  \ \ \ \ \ \ \ \ \ \ \left(  \text{by the definition of }\mathfrak{K}\right)
\\
&  \Longleftrightarrow\ \left(  \text{no }\left(  i,k\right)  \in A\text{ and
}\left(  k,j\right)  \in A\text{ satisfy }\left\{  \left\{  i,k\right\}
,\left\{  k,j\right\}  \right\}  \subseteq F\right) \\
&  \Longleftrightarrow\ \left(  \text{no }\left(  i,k\right)  \in A\text{ and
}\left(  k,j\right)  \in A\text{ satisfy }\left\{  i,k\right\}  \in F\text{
and }\left\{  k,j\right\}  \in F\right) \\
&  \Longleftrightarrow\ \left(  \text{no }\left(  i,k\right)  \in A\text{ and
}\left(  k,j\right)  \in A\text{ satisfy }\left(  i,k\right)  \in\pi
^{-1}\left(  F\right)  \text{ and }\left\{  k,j\right\}  \in F\right) \\
&  \ \ \ \ \ \ \ \ \ \ \left(
\begin{array}
[c]{c}%
\text{because for }\left(  i,k\right)  \in A\text{, we have }\left\{
i,k\right\}  \in F\text{ if and only if }\left(  i,k\right)  \in\pi
^{-1}\left(  F\right) \\
\text{(by (\ref{pf.prop.digraph.2pf-chrom.short.equiv}), applied to }u=i\text{
and }v=k\text{)}%
\end{array}
\right) \\
&  \Longleftrightarrow\ \left(  \text{no }\left(  i,k\right)  \in A\text{ and
}\left(  k,j\right)  \in A\text{ satisfy }\left(  i,k\right)  \in\pi
^{-1}\left(  F\right)  \text{ and }\left(  k,j\right)  \in\pi^{-1}\left(
F\right)  \right) \\
&  \ \ \ \ \ \ \ \ \ \ \left(
\begin{array}
[c]{c}%
\text{because for }\left(  k,j\right)  \in A\text{, we have }\left\{
k,j\right\}  \in F\text{ if and only if }\left(  k,j\right)  \in\pi
^{-1}\left(  F\right) \\
\text{(by (\ref{pf.prop.digraph.2pf-chrom.short.equiv}), applied to }u=k\text{
and }v=j\text{)}%
\end{array}
\right) \\
&  \Longleftrightarrow\ \left(  \text{the digraph }\left(  V,\pi^{-1}\left(
F\right)  \right)  \text{ is }2\text{-path-free}\right)
\ \ \ \ \ \ \ \ \ \ \left(  \text{by the definition of \textquotedblleft%
}2\text{-path-free\textquotedblright}\right)  ,
\end{align*}
qed.}. Now, Corollary \ref{cor.chrompol.K-free} (applied to $X=\mathbb{N}$ and
$G=\underline{D}$) shows that%
\begin{align*}
\chi_{\underline{D}}  &  =\underbrace{\sum_{\substack{F\subseteq E;\\F\text{
is }\mathfrak{K}\text{-free}}}}_{\substack{=\sum_{\substack{F\subseteq
E;\\\text{the digraph }\left(  V,\pi^{-1}\left(  F\right)  \right)  \text{ is
}2\text{-path-free}}}\\\text{(since we have just shown that}\\\text{a subset
}F\text{ of }E\text{ is }\mathfrak{K}\text{-free if and only if}\\\text{the
digraph }\left(  V,\pi^{-1}\left(  F\right)  \right)  \text{ is }%
2\text{-path-free)}}}\underbrace{\left(  -1\right)  ^{\left\vert F\right\vert
}}_{\substack{=\left(  -1\right)  ^{\left\vert \pi^{-1}\left(  F\right)
\right\vert }\\\text{(since }\pi\text{ is bijective)}}%
}\underbrace{x^{\operatorname*{conn}\left(  V,F\right)  }}%
_{\substack{=x^{\operatorname*{conn}\left(  V,\operatorname*{set}\left(
\pi^{-1}\left(  F\right)  \right)  \right)  }\\\text{(by
(\ref{pf.prop.digraph.2pf-chrom.short.F=set}))}}}\\
&  =\sum_{\substack{F\subseteq E;\\\text{the digraph }\left(  V,\pi
^{-1}\left(  F\right)  \right)  \text{ is }2\text{-path-free}}}\left(
-1\right)  ^{\left\vert \pi^{-1}\left(  F\right)  \right\vert }%
x^{\operatorname*{conn}\left(  V,\operatorname*{set}\left(  \pi^{-1}\left(
F\right)  \right)  \right)  }\\
&  =\sum_{\substack{B\subseteq A;\\\text{the digraph }\left(  V,B\right)
\text{ is }2\text{-path-free}}}\left(  -1\right)  ^{\left\vert B\right\vert
}x^{\operatorname*{conn}\left(  V,\operatorname*{set}B\right)  }\\
&  \ \ \ \ \ \ \ \ \ \ \left(
\begin{array}
[c]{c}%
\text{here, we have substituted }B\text{ for }\pi^{-1}\left(  F\right)  \text{
in the sum,}\\
\text{since the map }\pi:A\rightarrow E\text{ is bijective and thus induces}\\
\text{a bijection from the subsets of }E\text{ to the subsets of }A\\
\text{sending each }F\subseteq E\text{ to }\pi^{-1}\left(  F\right)
\end{array}
\right) \\
&  =\sum_{\substack{F\subseteq A;\\\text{the digraph }\left(  V,F\right)
\text{ is }2\text{-path-free}}}\left(  -1\right)  ^{\left\vert F\right\vert
}x^{\operatorname*{conn}\left(  V,\operatorname*{set}F\right)  }%
\end{align*}
(here, we have renamed the summation index $B$ as $F$). This proves
Proposition \ref{prop.digraph.2pf-chrom}.
\end{proof}
\end{vershort}

\begin{verlong}
We prepare for the proof of this proposition by stating two simple lemmas:

\begin{lemma}
\label{lem.digraph.2pf-chrom.3cyc}Let $G=\left(  V,E\right)  $ be a graph. Let
$u$, $v$ and $w$ be three elements of $V$ such that $\left\{  u,v\right\}  \in
E$, $\left\{  v,w\right\}  \in E$ and $\left\{  u,w\right\}  \in E$. Let $C$
be the set $\left\{  \left\{  u,v\right\}  ,\left\{  v,w\right\}  ,\left\{
u,w\right\}  \right\}  $. Then, the set $C$ is a circuit of $G$ and satisfies
$C\setminus\left\{  \left\{  u,w\right\}  \right\}  =\left\{  \left\{
u,v\right\}  ,\left\{  v,w\right\}  \right\}  $.
\end{lemma}

\begin{proof}
[Proof of Lemma \ref{lem.digraph.2pf-chrom.3cyc}.]We have $E\subseteq
\dbinom{V}{2}$ (since $\left(  V,E\right)  $ is a graph). For any $a\in V$ and
$b\in V$ satisfying $\left\{  a,b\right\}  \in E$, we have%
\begin{equation}
a\neq b \label{pf.lem.digraph.2pf-chrom.3cyc.1}%
\end{equation}
\footnote{\textit{Proof of (\ref{pf.lem.digraph.2pf-chrom.3cyc.1}):} Let $a\in
V$ and $b\in V$ be such that $\left\{  a,b\right\}  \in E$. We need to prove
that $a\neq b$.
\par
Assume the contrary (for the sake of contradiction). Thus, $a=b$. But
$\left\{  a,b\right\}  \in E\subseteq\dbinom{V}{2}$. In other words, $\left\{
a,b\right\}  $ is a $2$-element subset of $V$ (since $\dbinom{V}{2}$ is the
set of all $2$-element subsets of $V$). Thus, $\left\vert \left\{
a,b\right\}  \right\vert =2$. This contradicts $\left\vert \left\{
\underbrace{a}_{=b},b\right\}  \right\vert =\left\vert \underbrace{\left\{
b,b\right\}  }_{=\left\{  b\right\}  }\right\vert =\left\vert \left\{
b\right\}  \right\vert =1<2$. This contradiction proves that our assumption
was wrong. Hence, $a\neq b$ holds. This proves
(\ref{pf.lem.digraph.2pf-chrom.3cyc.1}).}.

Now, (\ref{pf.lem.digraph.2pf-chrom.3cyc.1}) (applied to $a=u$ and $b=v$)
yields $u\neq v$ (since $\left\{  u,v\right\}  \in E$). Also,
(\ref{pf.lem.digraph.2pf-chrom.3cyc.1}) (applied to $a=v$ and $b=w$) yields
$v\neq w$ (since $\left\{  v,w\right\}  \in E$). Also,
(\ref{pf.lem.digraph.2pf-chrom.3cyc.1}) (applied to $a=u$ and $b=w$) yields
$u\neq w$ (since $\left\{  u,w\right\}  \in E$).

Now, set $m=3$. Thus, $m+1=4$. Now, $u$, $v$, $w$ and $u$ are elements of $V$.
Hence, $\left(  u,v,w,u\right)  \in V^{4}=V^{m+1}$ (since $4=m+1$). In other
words, $\left(  u,v,w,u\right)  $ is an $\left(  m+1\right)  $-tuple of
elements of $V$. Denote this $\left(  m+1\right)  $-tuple by $\left(
v_{1},v_{2},\ldots,v_{m+1}\right)  $. Hence, $\left(  v_{1},v_{2}%
,\ldots,v_{m+1}\right)  =\left(  u,v,w,u\right)  $. Therefore, $v_{1}=u$,
$v_{2}=v$, $v_{3}=w$ and $v_{4}=u$. Now, we have $m=3>1$ and%
\begin{align*}
v_{m+1}  &  =v_{4}\ \ \ \ \ \ \ \ \ \ \left(  \text{since }m+1=4\right) \\
&  =u=v_{1}.
\end{align*}
Also, the vertices $v_{1},v_{2},\ldots,v_{m}$ are pairwise
distinct\footnote{\textit{Proof.} We have $m=3$, and thus $\left(  v_{1}%
,v_{2},\ldots,v_{m}\right)  =\left(  \underbrace{v_{1}}_{=u},\underbrace{v_{2}%
}_{=v},\underbrace{v_{3}}_{=w}\right)  =\left(  u,v,w\right)  $.
\par
The vertices $u,v,w$ are pairwise distinct (since $u\neq v$, $u\neq w$ and
$v\neq w$). Since $\left(  u,v,w\right)  =\left(  v_{1},v_{2},\ldots
,v_{m}\right)  $, this rewrites as follows: The vertices $v_{1},v_{2}%
,\ldots,v_{m}$ are pairwise distinct. Qed.}. Finally, we have $\left\{
v_{i},v_{i+1}\right\}  \in E$ for every $i\in\left\{  1,2,\ldots,m\right\}
$\ \ \ \ \footnote{\textit{Proof.} Let $i\in\left\{  1,2,\ldots,m\right\}  $.
We want to prove that $\left\{  v_{i},v_{i+1}\right\}  \in E$. We have%
\begin{align*}
i  &  \in\left\{  1,2,\ldots,m\right\}  =\left\{  1,2,\ldots,3\right\}
\ \ \ \ \ \ \ \ \ \ \left(  \text{since }m=3\right) \\
&  =\left\{  1,2,3\right\}  .
\end{align*}
In other words, $i=1$ or $i=2$ or $i=3$. We are thus in one of the following
three cases:
\par
\textit{Case 1:} We have $i=1$.
\par
\textit{Case 2:} We have $i=2$.
\par
\textit{Case 3:} We have $i=3$.
\par
Let us first consider Case 1. In this case, we have $i=1$. Hence, $v_{i}%
=v_{1}=u$ and $v_{i+1}=v_{1+1}=v_{2}=v$. Now, $\left\{  \underbrace{v_{i}%
}_{=u},\underbrace{v_{i+1}}_{=v}\right\}  =\left\{  u,v\right\}  \in E$. Thus,
$\left\{  v_{i},v_{i+1}\right\}  \in E$ holds in Case 1.
\par
Let us first consider Case 2. In this case, we have $i=2$. Hence, $v_{i}%
=v_{2}=v$ and $v_{i+1}=v_{2+1}=v_{3}=w$. Now, $\left\{  \underbrace{v_{i}%
}_{=v},\underbrace{v_{i+1}}_{=w}\right\}  =\left\{  v,w\right\}  \in E$. Thus,
$\left\{  v_{i},v_{i+1}\right\}  \in E$ holds in Case 2.
\par
Let us first consider Case 3. In this case, we have $i=3$. Hence, $v_{i}%
=v_{3}=w$ and $v_{i+1}=v_{3+1}=v_{4}=u$. Now, $\left\{  \underbrace{v_{i}%
}_{=w},\underbrace{v_{i+1}}_{=u}\right\}  =\left\{  w,u\right\}  =\left\{
u,w\right\}  \in E$. Thus, $\left\{  v_{i},v_{i+1}\right\}  \in E$ holds in
Case 3.
\par
We have now proven $\left\{  v_{i},v_{i+1}\right\}  \in E$ in each of the
three Cases 1, 2 and 3. Thus, $\left\{  v_{i},v_{i+1}\right\}  \in E$ always
holds (since the Cases 1, 2 and 3 cover all possibilities). This completes our
proof.}.

Altogether, we now know that $\left(  v_{1},v_{2},\ldots,v_{m+1}\right)  $ is
a list of elements of $V$ satisfying the following four properties:

\begin{itemize}
\item We have $m>1$.

\item We have $v_{m+1}=v_{1}$.

\item The vertices $v_{1},v_{2},\ldots,v_{m}$ are pairwise distinct.

\item We have $\left\{  v_{i},v_{i+1}\right\}  \in E$ for every $i\in\left\{
1,2,\ldots,m\right\}  $.
\end{itemize}

According to the definition of a \textquotedblleft cycle\textquotedblright,
this means that the list $\left(  v_{1},v_{2},\ldots,v_{m+1}\right)  $ is a
cycle of $\left(  V,E\right)  $.

Thus, we conclude that the list $\left(  v_{1},v_{2},\ldots,v_{m+1}\right)  $
is a cycle of $\left(  V,E\right)  $. In other words, the list $\left(
v_{1},v_{2},\ldots,v_{m+1}\right)  $ is a cycle of $G$ (since $G=\left(
V,E\right)  $). Thus, the set $\left\{  \left\{  v_{1},v_{2}\right\}
,\left\{  v_{2},v_{3}\right\}  ,\ldots,\left\{  v_{m},v_{m+1}\right\}
\right\}  $ is a circuit of $G$ (by the definition of a \textquotedblleft
circuit\textquotedblright). Since%
\begin{align*}
&  \left\{  \left\{  v_{1},v_{2}\right\}  ,\left\{  v_{2},v_{3}\right\}
,\ldots,\left\{  v_{m},v_{m+1}\right\}  \right\} \\
&  =\left\{  \left\{  v_{1},v_{2}\right\}  ,\left\{  v_{2},v_{3}\right\}
,\ldots,\left\{  v_{3},v_{3+1}\right\}  \right\}  \ \ \ \ \ \ \ \ \ \ \left(
\text{since }m=3\right) \\
&  =\left\{  \left\{  \underbrace{v_{1}}_{=u},\underbrace{v_{2}}_{=v}\right\}
,\left\{  \underbrace{v_{2}}_{=v},\underbrace{v_{3}}_{=w}\right\}  ,\left\{
\underbrace{v_{3}}_{=w},\underbrace{v_{4}}_{=u}\right\}  \right\} \\
&  =\left\{  \left\{  u,v\right\}  ,\left\{  v,w\right\}
,\underbrace{\left\{  w,u\right\}  }_{=\left\{  u,w\right\}  }\right\}
=\left\{  \left\{  u,v\right\}  ,\left\{  v,w\right\}  ,\left\{  u,w\right\}
\right\}  =C
\end{align*}
(because $C=\left\{  \left\{  u,v\right\}  ,\left\{  v,w\right\}  ,\left\{
u,w\right\}  \right\}  $), this rewrites as follows: The set $C$ is a circuit
of $G$.

It remains to prove that $C\setminus\left\{  \left\{  u,w\right\}  \right\}
=\left\{  \left\{  u,v\right\}  ,\left\{  v,w\right\}  \right\}  $. Indeed, we
have $\left\{  u,w\right\}  \notin\left\{  \left\{  u,v\right\}  ,\left\{
v,w\right\}  \right\}  $\ \ \ \ \footnote{\textit{Proof.} Assume the contrary.
Thus, $\left\{  u,w\right\}  \in\left\{  \left\{  u,v\right\}  ,\left\{
v,w\right\}  \right\}  $. In other words, $\left\{  u,w\right\}  $ equals
either $\left\{  u,v\right\}  $ or $\left\{  v,w\right\}  $. In other words,
we must be in one of the following two cases:
\par
\textit{Case 1:} We have $\left\{  u,w\right\}  =\left\{  u,v\right\}  $.
\par
\textit{Case 2:} We have $\left\{  u,w\right\}  =\left\{  v,w\right\}  $.
\par
Let us first consider Case 1. In this case, we have $\left\{  u,w\right\}
=\left\{  u,v\right\}  $. Hence, $w\in\left\{  u,w\right\}  =\left\{
u,v\right\}  $. Combining this with $w\neq u$ (since $u\neq w$), we obtain
$w\in\left\{  u,v\right\}  \setminus\left\{  u\right\}  \subseteq\left\{
v\right\}  $. In other words, $w=v$. Hence, $v=w$. But this contradicts $v\neq
w$. Thus, we have found a contradiction in Case 1.
\par
Let us now consider Case 2. In this case, we have $\left\{  u,w\right\}
=\left\{  v,w\right\}  $. Hence, $u\in\left\{  u,w\right\}  =\left\{
v,w\right\}  $. Combining this with $u\neq v$, we obtain $u\in\left\{
v,w\right\}  \setminus\left\{  v\right\}  \subseteq\left\{  w\right\}  $. In
other words, $u=w$. This contradicts $u\neq w$. Thus, we have found a
contradiction in Case 2.
\par
We have now found a contradiction in each of the two Cases 1 and 2. Since
these two Cases cover all possibilities, we thus always have a contradiction.
This contradiction shows that our assumption was wrong, qed.}. Now,%
\[
C=\left\{  \left\{  u,v\right\}  ,\left\{  v,w\right\}  ,\left\{  u,w\right\}
\right\}  =\left\{  \left\{  u,v\right\}  ,\left\{  v,w\right\}  \right\}
\cup\left\{  \left\{  u,w\right\}  \right\}  .
\]
Hence,%
\begin{align*}
&  \underbrace{C}_{=\left\{  \left\{  u,v\right\}  ,\left\{  v,w\right\}
\right\}  \cup\left\{  \left\{  u,w\right\}  \right\}  }\setminus\left\{
\left\{  u,w\right\}  \right\} \\
&  =\left(  \left\{  \left\{  u,v\right\}  ,\left\{  v,w\right\}  \right\}
\cup\left\{  \left\{  u,w\right\}  \right\}  \right)  \setminus\left\{
\left\{  u,w\right\}  \right\} \\
&  =\left\{  \left\{  u,v\right\}  ,\left\{  v,w\right\}  \right\}
\ \ \ \ \ \ \ \ \ \ \left(  \text{since }\left\{  u,w\right\}  \notin\left\{
\left\{  u,v\right\}  ,\left\{  v,w\right\}  \right\}  \right)  .
\end{align*}
This completes the proof of Lemma \ref{lem.digraph.2pf-chrom.3cyc}.
\end{proof}

\begin{lemma}
\label{lem.digraph.2pf-chrom.K-free}Let $D=\left(  V,A\right)  $ be a finite
transitive loopless digraph. Let $E=\operatorname*{set}A$. Every $a\in A$
satisfies $\operatorname*{set}\underbrace{a}_{\in A}\in\operatorname*{set}%
A\subseteq E$. Hence, we can define a map $\pi:A\rightarrow E$ by%
\[
\left(  \pi\left(  a\right)  =\operatorname*{set}%
a\ \ \ \ \ \ \ \ \ \ \text{for every }a\in A\right)  .
\]
Consider this map $\pi$.

\textbf{(a)} The map $\pi:A\rightarrow E$ is bijective. Thus, an inverse map
$\pi^{-1}:E\rightarrow A$ is well-defined.

\textbf{(b)} Let $Z$ be the set $\left\{  \left(  i,k,j\right)  \in
V^{3}\ \mid\ \left(  i,k\right)  \in A\text{ and }\left(  k,j\right)  \in
A\right\}  $. Define a set $\mathfrak{K}$ by
\[
\mathfrak{K}=\left\{  \left\{  \left\{  i,k\right\}  ,\left\{  k,j\right\}
\right\}  \ \mid\ \left(  i,k,j\right)  \in Z\right\}  .
\]
Then, $\mathfrak{K}\subseteq\mathcal{P}\left(  E\right)  $.

\textbf{(c)} Let $G$ be the graph $\left(  V,E\right)  $. For every $\left(
i,j\right)  \in V^{2}$, define a subset $\operatorname*{inter}\left(
i,j\right)  $ of $V$ by%
\[
\operatorname*{inter}\left(  i,j\right)  =\left\{  k\in V\ \mid\ \left(
i,k\right)  \in A\text{ and }\left(  k,j\right)  \in A\right\}  .
\]
Define a map $\ell^{\prime}:A\rightarrow\mathbb{N}$ by setting%
\[
\left(  \ell^{\prime}\left(  i,j\right)  =\left\vert \operatorname*{inter}%
\left(  i,j\right)  \right\vert \ \ \ \ \ \ \ \ \ \ \text{for every }\left(
i,j\right)  \in A\right)  .
\]
Recall that a map $\pi^{-1}:E\rightarrow A$ is well-defined (by Lemma
\ref{lem.digraph.2pf-chrom.K-free}). Define a map $\ell:E\rightarrow
\mathbb{N}$ by $\ell=\ell^{\prime}\circ\pi^{-1}$. We shall refer to $\ell$ as
the \textit{labeling function}. For every edge $e$ of $G$, we shall refer to
$\ell\left(  e\right)  $ as the \textit{label} of $e$. If $u$, $v$ and $w$ are
three elements of $V$ satisfying $\left(  u,v\right)  \in A$ and $\left(
v,w\right)  \in A$, then%
\begin{align}
\left\{  u,w\right\}   &  \in E,\ \ \ \ \ \ \ \ \ \ \left\{  u,v\right\}  \in
E,\ \ \ \ \ \ \ \ \ \ \left\{  u,w\right\}  \in
E,\label{eq.lem.digraph.2pf-chrom.K-free.c.1}\\
\ell\left(  \left\{  u,w\right\}  \right)   &  >\ell\left(  \left\{
u,v\right\}  \right)  \ \ \ \ \ \ \ \ \ \ \text{and}\ \ \ \ \ \ \ \ \ \ \ell
\left(  \left\{  u,w\right\}  \right)  >\ell\left(  \left\{  v,w\right\}
\right)  . \label{eq.lem.digraph.2pf-chrom.K-free.c.2}%
\end{align}


\textbf{(d)} Consider the set $\mathfrak{K}$ defined in Lemma
\ref{lem.digraph.2pf-chrom.K-free} \textbf{(b)}. Consider the labeling
function $\ell$ defined in Lemma \ref{lem.digraph.2pf-chrom.K-free}
\textbf{(c)}. Definition \ref{def.BC} (applied to $X=\mathbb{N}$) shows that
the notion of a broken circuit of $G$ is well-defined (since a function
$\ell:E\rightarrow\mathbb{N}$ is given).

Every element of $\mathfrak{K}$ is a broken circuit of $G$.

\textbf{(e)} Consider the set $\mathfrak{K}$ defined in Lemma
\ref{lem.digraph.2pf-chrom.K-free} \textbf{(b)}. Let $F$ be a subset of $A$.
Then, we have the following logical equivalence:%
\[
\ \left(  \text{the digraph }\left(  V,F\right)  \text{ is }2\text{-path-free}%
\right)  \ \Longleftrightarrow\ \left(  \text{the set }\pi\left(  F\right)
\text{ is }\mathfrak{K}\text{-free}\right)  .
\]

\end{lemma}

\begin{proof}
[Proof of Lemma \ref{lem.digraph.2pf-chrom.K-free}.]We have $A\subseteq V^{2}$
(since $\left(  V,A\right)  $ is a digraph). The set $V$ is finite (since the
digraph $\left(  V,A\right)  $ is finite).

Recall that the digraph $\left(  V,A\right)  $ is transitive if and only if,
for any $u\in V$, $v\in V$ and $w\in V$ satisfying $\left(  u,v\right)  \in A$
and $\left(  v,w\right)  \in A$, we have $\left(  u,w\right)  \in A$ (by the
definition of \textquotedblleft transitive\textquotedblright). Thus, for any
$u\in V$, $v\in V$ and $w\in V$ satisfying $\left(  u,v\right)  \in A$ and
$\left(  v,w\right)  \in A$, we have
\begin{equation}
\left(  u,w\right)  \in A \label{pf.lem.digraph.2pf-chrom.K-free.transitive}%
\end{equation}
(because the digraph $\left(  V,A\right)  $ is transitive).

Recall that the digraph $\left(  V,A\right)  $ is loopless if and only if
every $v\in V$ satisfies $\left(  v,v\right)  \notin A$. Thus, every $v\in V$
satisfies%
\begin{equation}
\left(  v,v\right)  \notin A \label{pf.lem.digraph.2pf-chrom.K-free.loopless}%
\end{equation}
(since the digraph $\left(  V,A\right)  $ is loopless).

\textbf{(a)} We have
\[
\pi\left(  A\right)  =\left\{  \underbrace{\pi\left(  a\right)  }%
_{\substack{=\operatorname*{set}a\\\text{(by the definition of }\pi\text{)}%
}}\ \mid\ a\in A\right\}  =\left\{  \operatorname*{set}a\ \mid\ a\in
A\right\}  =\operatorname*{set}A=E.
\]
Thus, the map $\pi$ is surjective.

Furthermore, if $a$ and $b$ are two elements of $A$ such that $\pi\left(
a\right)  =\pi\left(  b\right)  $, then $a=b$\ \ \ \ \footnote{\textit{Proof.}
Let $a$ and $b$ be two elements of $A$ such that $\pi\left(  a\right)
=\pi\left(  b\right)  $. We must show that $a=b$.
\par
Assume the contrary. Thus, $a\neq b$.
\par
We have $a\in A\subseteq V^{2}$. Thus, we can write $a$ in the form $a=\left(
i,j\right)  $ with $i\in V$ and $j\in V$. Consider these $i$ and $j$.
\par
We have $b\in A\subseteq V^{2}$. Thus, we can write $b$ in the form $b=\left(
i^{\prime},j^{\prime}\right)  $ with $i^{\prime}\in V$ and $j^{\prime}\in V$.
Consider these $i^{\prime}$ and $j^{\prime}$.
\par
Applying (\ref{pf.lem.digraph.2pf-chrom.K-free.loopless}) to $v=i$, we obtain
$\left(  i,i\right)  \notin A$. If we had $i=j$, then we would have $\left(
i,\underbrace{i}_{=j}\right)  =\left(  i,j\right)  \in A$, which would
contradict $\left(  i,i\right)  \notin A$. Thus, we cannot have $i=j$. Hence,
we have $i\neq j$, so that $j\neq i$.
\par
Applying the map $\pi$ to both sides of the equality $a=\left(  i,j\right)  $,
we obtain
\begin{align*}
\pi\left(  a\right)   &  =\pi\left(  i,j\right)  =\operatorname*{set}\left(
i,j\right)  \ \ \ \ \ \ \ \ \ \ \left(  \text{by the definition of }\pi\right)
\\
&  =\left\{  i,j\right\}  \ \ \ \ \ \ \ \ \ \ \left(  \text{by the definition
of the map }\operatorname*{set}\right)  .
\end{align*}
The same argument (but applied to $b$, $i^{\prime}$ and $j^{\prime}$ instead
of $a$, $i$ and $j$) shows that $\pi\left(  b\right)  =\left\{  i^{\prime
},j^{\prime}\right\}  $. Hence, $\left\{  i,j\right\}  =\pi\left(  a\right)
=\pi\left(  b\right)  =\left\{  i^{\prime},j^{\prime}\right\}  $.
\par
We have $i\in\left\{  i,j\right\}  =\left\{  i^{\prime},j^{\prime}\right\}  $.
In other words, either $i=i^{\prime}$ or $i=j^{\prime}$. In other words, we
are in one of the following two cases:
\par
\textit{Case 1:} We have $i=i^{\prime}$.
\par
\textit{Case 2:} We have $i=j^{\prime}$.
\par
Let us first consider Case 1. In this case, we have $i=i^{\prime}$. Now,
$j\in\left\{  i,j\right\}  =\left\{  i^{\prime},j^{\prime}\right\}  $.
Combining this with $j\neq i=i^{\prime}$, we obtain $j\in\left\{  i^{\prime
},j^{\prime}\right\}  \setminus\left\{  i^{\prime}\right\}  \subseteq\left\{
j^{\prime}\right\}  $. Thus, $j=j^{\prime}$. Now, $a=\left(  \underbrace{i}%
_{=i^{\prime}},\underbrace{j}_{=j^{\prime}}\right)  =\left(  i^{\prime
},j^{\prime}\right)  =b$; this contradicts $a\neq b$. Hence, we have found a
contradiction in Case 1.
\par
Let us now consider Case 2. In this case, we have $i=j^{\prime}$. Now,
$j\in\left\{  i,j\right\}  =\left\{  i^{\prime},j^{\prime}\right\}  $.
Combining this with $j\neq i=j^{\prime}$, we obtain $j\in\left\{  i^{\prime
},j^{\prime}\right\}  \setminus\left\{  j^{\prime}\right\}  \subseteq\left\{
i^{\prime}\right\}  $. Thus, $j=i^{\prime}$. Hence, $\left(  \underbrace{j}%
_{=i^{\prime}},\underbrace{i}_{=j^{\prime}}\right)  =\left(  i^{\prime
},j^{\prime}\right)  =b\in A$. Also, $\left(  i,j\right)  =a\in A$. Hence,
(\ref{pf.lem.digraph.2pf-chrom.K-free.transitive}) (applied to $u=i$, $v=j$
and $w=i$) yields $\left(  i,i\right)  \in A$. This contradicts $\left(
i,i\right)  \notin A$. Thus, we have found a contradiction in Case 2.
\par
We have now found a contradiction in each of the two Cases 1 and 2. Hence, we
always have a contradiction (since Cases 1 and 2 cover all possibilities).
Thus, our assumption was wrong. Hence, $a=b$ is proven. Qed.}. In other words,
the map $\pi$ is injective. Now, the map $\pi$ is bijective (since $\pi$ is
surjective and injective). Thus, an inverse map $\pi^{-1}:E\rightarrow A$ is
well-defined. This proves Lemma \ref{lem.digraph.2pf-chrom.K-free}
\textbf{(a)}.

\textbf{(b)} The definition of $Z$ yields%
\begin{align*}
Z  &  =\left\{  \left(  i,k,j\right)  \in V^{3}\ \mid\ \left(  i,k\right)  \in
A\text{ and }\left(  k,j\right)  \in A\right\} \\
&  =\left\{  \left(  u,v,w\right)  \in V^{3}\ \mid\ \left(  u,v\right)  \in
A\text{ and }\left(  v,w\right)  \in A\right\}
\end{align*}
(here, we have renamed the index $\left(  i,k,j\right)  $ as $\left(
u,v,w\right)  $). Now, let $K\in\mathfrak{K}$. We shall prove that
$K\in\mathcal{P}\left(  E\right)  $.

We have $K\in\mathfrak{K}=\left\{  \left\{  \left\{  i,k\right\}  ,\left\{
k,j\right\}  \right\}  \ \mid\ \left(  i,k,j\right)  \in Z\right\}  $. Hence,
$K=\left\{  \left\{  i,k\right\}  ,\left\{  k,j\right\}  \right\}  $ for some
$\left(  i,k,j\right)  \in Z$. Consider this $\left(  i,k,j\right)  $. We
have
\[
\left(  i,k,j\right)  \in Z=\left\{  \left(  u,v,w\right)  \in V^{3}%
\ \mid\ \left(  u,v\right)  \in A\text{ and }\left(  v,w\right)  \in
A\right\}  .
\]
In other words, $\left(  i,k,j\right)  $ is an $\left(  u,v,w\right)  \in
V^{3}$ satisfying $\left(  u,v\right)  \in A$ and $\left(  v,w\right)  \in A$.
In other words, $\left(  i,k,j\right)  $ is an element of $V^{3}$ and
satisfies $\left(  i,k\right)  \in A$ and $\left(  k,j\right)  \in A$.

Now, $\left(  i,k\right)  \in A$. The definition of $\pi$ therefore yields
$\pi\left(  i,k\right)  =\operatorname*{set}\left(  i,k\right)  =\left\{
i,k\right\}  $ (by the definition of the map $\operatorname*{set}$). Hence,
$\left\{  i,k\right\}  =\pi\underbrace{\left(  i,k\right)  }_{\in A}\in
\pi\left(  A\right)  \subseteq E$.

Also, $\left(  k,j\right)  \in A$. The definition of $\pi$ therefore yields
$\pi\left(  k,j\right)  =\operatorname*{set}\left(  k,j\right)  =\left\{
k,j\right\}  $ (by the definition of the map $\operatorname*{set}$). Hence,
$\left\{  k,j\right\}  =\pi\underbrace{\left(  k,j\right)  }_{\in A}\in
\pi\left(  A\right)  \subseteq E$.

Now, we know that both $\left\{  i,k\right\}  $ and $\left\{  k,j\right\}  $
belong to $E$. Hence, $\left\{  \left\{  i,k\right\}  ,\left\{  k,j\right\}
\right\}  \subseteq E$. Hence, $K=\left\{  \left\{  i,k\right\}  ,\left\{
k,j\right\}  \right\}  \subseteq E$, so that $K\in\mathcal{P}\left(  E\right)
$.

Now, forget that we fixed $K$. We thus have shown that $K\in\mathcal{P}\left(
E\right)  $ for every $K\in\mathfrak{K}$. In other words, $\mathfrak{K}%
\subseteq\mathcal{P}\left(  E\right)  $. This proves Lemma
\ref{lem.digraph.2pf-chrom.K-free} \textbf{(b)}.

\textbf{(c)} We only need to prove that if $u$, $v$ and $w$ are three elements
of $V$ satisfying $\left(  u,v\right)  \in A$ and $\left(  v,w\right)  \in A$,
then (\ref{eq.lem.digraph.2pf-chrom.K-free.c.1}) and
(\ref{eq.lem.digraph.2pf-chrom.K-free.c.2}) hold.

So let $u$, $v$ and $w$ be three elements of $V$ satisfying $\left(
u,v\right)  \in A$ and $\left(  v,w\right)  \in A$. We must prove
(\ref{eq.lem.digraph.2pf-chrom.K-free.c.1}) and
(\ref{eq.lem.digraph.2pf-chrom.K-free.c.2}).

The definition of $\operatorname*{inter}\left(  u,v\right)  $ yields%
\[
\operatorname*{inter}\left(  u,v\right)  =\left\{  k\in V\ \mid\ \left(
u,k\right)  \in A\text{ and }\left(  k,v\right)  \in A\right\}  .
\]
The definition of $\operatorname*{inter}\left(  v,w\right)  $ yields%
\[
\operatorname*{inter}\left(  v,w\right)  =\left\{  k\in V\ \mid\ \left(
v,k\right)  \in A\text{ and }\left(  k,w\right)  \in A\right\}  .
\]
The definition of $\operatorname*{inter}\left(  u,w\right)  $ yields%
\[
\operatorname*{inter}\left(  u,w\right)  =\left\{  k\in V\ \mid\ \left(
u,k\right)  \in A\text{ and }\left(  k,w\right)  \in A\right\}  .
\]


We have $v\in\operatorname*{inter}\left(  u,w\right)  $%
\ \ \ \ \footnote{\textit{Proof.} We know that $v$ is an element of $V$ and
satisfies $\left(  u,v\right)  \in A$ and $\left(  v,w\right)  \in A$. In
other words, $v$ is an element $k$ of $V$ satisfying $\left(  u,k\right)  \in
A$ and $\left(  k,v\right)  \in A$. Hence,%
\[
v\in\left\{  k\in V\ \mid\ \left(  u,k\right)  \in A\text{ and }\left(
k,w\right)  \in A\right\}  =\operatorname*{inter}\left(  u,w\right)  .
\]
Qed.}.

All three sets $\operatorname*{inter}\left(  u,v\right)  $,
$\operatorname*{inter}\left(  v,w\right)  $ and $\operatorname*{inter}\left(
u,w\right)  $ are subsets of the finite set $V$, and thus are finite.

Now, $\operatorname*{inter}\left(  u,v\right)  $ is a proper subset of
$\operatorname*{inter}\left(  u,w\right)  $\ \ \ \ \footnote{\textit{Proof.}
Let $g\in\operatorname*{inter}\left(  u,v\right)  $. Thus,%
\[
g\in\operatorname*{inter}\left(  u,v\right)  =\left\{  k\in V\ \mid\ \left(
u,k\right)  \in A\text{ and }\left(  k,v\right)  \in A\right\}  .
\]
In other words, $g$ is an element $k$ of $V$ satisfying $\left(  u,k\right)
\in A$ and $\left(  k,v\right)  \in A$. In other words, $g$ is an element of
$V$ and satisfies $\left(  u,g\right)  \in A$ and $\left(  g,v\right)  \in A$.
\par
Now, (\ref{pf.lem.digraph.2pf-chrom.K-free.transitive}) (applied to $g$
instead of $u$) yields $\left(  g,w\right)  \in A$ (since $\left(  g,v\right)
\in A$ and $\left(  v,w\right)  \in A$). Hence, we now know that $g$ is an
element of $V$ and satisfies $\left(  u,g\right)  \in A$ and $\left(
g,w\right)  \in A$. In other words, $g$ is an element $k$ of $V$ satisfying
$\left(  u,k\right)  \in A$ and $\left(  k,w\right)  \in A$. Hence,%
\[
g\in\left\{  k\in V\ \mid\ \left(  u,k\right)  \in A\text{ and }\left(
k,w\right)  \in A\right\}  =\operatorname*{inter}\left(  u,w\right)  .
\]
\par
Now, forget that we fixed $g$. We therefore have proven that $g\in
\operatorname*{inter}\left(  u,w\right)  $ for every $g\in
\operatorname*{inter}\left(  u,v\right)  $. In other words,
$\operatorname*{inter}\left(  u,v\right)  \subseteq\operatorname*{inter}%
\left(  u,w\right)  $.
\par
Next, we shall prove that $\operatorname*{inter}\left(  u,v\right)
\neq\operatorname*{inter}\left(  u,w\right)  $. Indeed, assume the contrary.
Thus, $\operatorname*{inter}\left(  u,v\right)  =\operatorname*{inter}\left(
u,w\right)  $. Now,%
\[
v\in\operatorname*{inter}\left(  u,w\right)  =\operatorname*{inter}\left(
u,v\right)  =\left\{  k\in V\ \mid\ \left(  u,k\right)  \in A\text{ and
}\left(  k,v\right)  \in A\right\}  .
\]
In other words, $v$ is an element $k$ of $V$ satisfying $\left(  u,k\right)
\in A$ and $\left(  k,v\right)  \in A$. In other words, $v$ is an element of
$V$ and satisfies $\left(  u,v\right)  \in A$ and $\left(  v,v\right)  \in A$.
But (\ref{pf.lem.digraph.2pf-chrom.K-free.loopless}) yields $\left(
v,v\right)  \notin A$. This contradicts $\left(  v,v\right)  \in A$. This
contradiction proves that our assumption was false. Hence,
$\operatorname*{inter}\left(  u,v\right)  \neq\operatorname*{inter}\left(
u,w\right)  $ is proven. Combining this with $\operatorname*{inter}\left(
u,v\right)  \subseteq\operatorname*{inter}\left(  u,w\right)  $, we conclude
that $\operatorname*{inter}\left(  u,v\right)  $ is a proper subset of
$\operatorname*{inter}\left(  u,w\right)  $. Qed.}. Hence,
\begin{equation}
\left\vert \operatorname*{inter}\left(  u,v\right)  \right\vert <\left\vert
\operatorname*{inter}\left(  u,w\right)  \right\vert
\label{pf.lem.digraph.2pf-chrom.K-free.c.ieq1}%
\end{equation}
(since $\operatorname*{inter}\left(  u,v\right)  $ and $\operatorname*{inter}%
\left(  u,w\right)  $ are finite sets).

Furthermore. $\operatorname*{inter}\left(  v,w\right)  $ is a proper subset of
$\operatorname*{inter}\left(  u,w\right)  $\ \ \ \ \footnote{\textit{Proof.}
Let $g\in\operatorname*{inter}\left(  v,w\right)  $. Thus,%
\[
g\in\operatorname*{inter}\left(  v,w\right)  =\left\{  k\in V\ \mid\ \left(
v,k\right)  \in A\text{ and }\left(  k,w\right)  \in A\right\}  .
\]
In other words, $g$ is an element $k$ of $V$ satisfying $\left(  v,k\right)
\in A$ and $\left(  k,w\right)  \in A$. In other words, $g$ is an element of
$V$ and satisfies $\left(  v,g\right)  \in A$ and $\left(  g,w\right)  \in A$.
\par
Now, (\ref{pf.lem.digraph.2pf-chrom.K-free.transitive}) (applied to $g$
instead of $w$) yields $\left(  u,g\right)  \in A$ (since $\left(  u,v\right)
\in A$ and $\left(  v,g\right)  \in A$). Hence, we now know that $g$ is an
element of $V$ and satisfies $\left(  u,g\right)  \in A$ and $\left(
g,w\right)  \in A$. In other words, $g$ is an element $k$ of $V$ satisfying
$\left(  u,k\right)  \in A$ and $\left(  k,w\right)  \in A$. Hence,%
\[
g\in\left\{  k\in V\ \mid\ \left(  u,k\right)  \in A\text{ and }\left(
k,w\right)  \in A\right\}  =\operatorname*{inter}\left(  u,w\right)  .
\]
\par
Now, forget that we fixed $g$. We therefore have proven that $g\in
\operatorname*{inter}\left(  u,w\right)  $ for every $g\in
\operatorname*{inter}\left(  v,w\right)  $. In other words,
$\operatorname*{inter}\left(  v,w\right)  \subseteq\operatorname*{inter}%
\left(  u,w\right)  $.
\par
Next, we shall prove that $\operatorname*{inter}\left(  v,w\right)
\neq\operatorname*{inter}\left(  u,w\right)  $. Indeed, assume the contrary.
Thus, $\operatorname*{inter}\left(  v,w\right)  =\operatorname*{inter}\left(
u,w\right)  $. Now,%
\[
v\in\operatorname*{inter}\left(  u,w\right)  =\operatorname*{inter}\left(
v,w\right)  =\left\{  k\in V\ \mid\ \left(  v,k\right)  \in A\text{ and
}\left(  k,w\right)  \in A\right\}  .
\]
In other words, $v$ is an element $k$ of $V$ satisfying $\left(  v,k\right)
\in A$ and $\left(  k,w\right)  \in A$. In other words, $v$ is an element of
$V$ and satisfies $\left(  v,v\right)  \in A$ and $\left(  v,w\right)  \in A$.
But (\ref{pf.lem.digraph.2pf-chrom.K-free.loopless}) yields $\left(
v,v\right)  \notin A$. This contradicts $\left(  v,v\right)  \in A$. This
contradiction proves that our assumption was false. Hence,
$\operatorname*{inter}\left(  v,w\right)  \neq\operatorname*{inter}\left(
u,w\right)  $ is proven. Combining this with $\operatorname*{inter}\left(
v,w\right)  \subseteq\operatorname*{inter}\left(  u,w\right)  $, we conclude
that $\operatorname*{inter}\left(  v,w\right)  $ is a proper subset of
$\operatorname*{inter}\left(  v,w\right)  $. Qed.}. Hence,
\begin{equation}
\left\vert \operatorname*{inter}\left(  v,w\right)  \right\vert <\left\vert
\operatorname*{inter}\left(  u,w\right)  \right\vert
\label{pf.lem.digraph.2pf-chrom.K-free.c.ieq2}%
\end{equation}
(since $\operatorname*{inter}\left(  v,w\right)  $ and $\operatorname*{inter}%
\left(  u,w\right)  $ are finite sets).

From (\ref{pf.lem.digraph.2pf-chrom.K-free.transitive}), we obtain $\left(
u,w\right)  \in A$. The definition of $\pi$ thus shows that%
\[
\pi\left(  u,w\right)  =\operatorname*{set}\left(  u,w\right)  =\left\{
u,w\right\}  \ \ \ \ \ \ \ \ \ \ \left(  \text{by the definition of the map
}\operatorname*{set}\right)  .
\]
Hence, $\left\{  u,w\right\}  =\pi\underbrace{\left(  u,w\right)  }_{\in A}%
\in\pi\left(  A\right)  \subseteq E$. Thus, $\ell\left(  \left\{  u,w\right\}
\right)  $ is well-defined.

Also, recall that $\left(  u,v\right)  \in A$. The definition of $\pi$ thus
shows that%
\[
\pi\left(  u,v\right)  =\operatorname*{set}\left(  u,v\right)  =\left\{
u,v\right\}  \ \ \ \ \ \ \ \ \ \ \left(  \text{by the definition of the map
}\operatorname*{set}\right)  .
\]
Hence, $\left\{  u,v\right\}  =\pi\underbrace{\left(  u,v\right)  }_{\in A}%
\in\pi\left(  A\right)  \subseteq E$. Thus, $\ell\left(  \left\{  u,v\right\}
\right)  $ is well-defined.

Also, recall that $\left(  v,w\right)  \in A$. The definition of $\pi$ thus
shows that%
\[
\pi\left(  v,w\right)  =\operatorname*{set}\left(  v,w\right)  =\left\{
v,w\right\}  \ \ \ \ \ \ \ \ \ \ \left(  \text{by the definition of the map
}\operatorname*{set}\right)  .
\]
Hence, $\left\{  v,w\right\}  =\pi\underbrace{\left(  v,w\right)  }_{\in A}%
\in\pi\left(  A\right)  \subseteq E$. Thus, $\ell\left(  \left\{  v,w\right\}
\right)  $ is well-defined. We have also proven
(\ref{eq.lem.digraph.2pf-chrom.K-free.c.1}) by now.

Now,%
\begin{align*}
\underbrace{\ell}_{=\ell^{\prime}\circ\pi^{-1}}\left(  \left\{  u,w\right\}
\right)   &  =\left(  \ell^{\prime}\circ\pi^{-1}\right)  \left(  \left\{
u,w\right\}  \right)  =\ell^{\prime}\left(  \underbrace{\pi^{-1}\left(
\left\{  u,w\right\}  \right)  }_{\substack{=\left(  u,w\right)
\\\text{(since }\left\{  u,w\right\}  =\pi\left(  u,w\right)  \text{)}%
}}\right)  =\ell^{\prime}\left(  u,w\right) \\
&  =\left\vert \operatorname*{inter}\left(  u,w\right)  \right\vert
\ \ \ \ \ \ \ \ \ \ \left(  \text{by the definition of }\ell^{\prime}\right)
.
\end{align*}
Also,%
\begin{align*}
\underbrace{\ell}_{=\ell^{\prime}\circ\pi^{-1}}\left(  \left\{  u,v\right\}
\right)   &  =\left(  \ell^{\prime}\circ\pi^{-1}\right)  \left(  \left\{
u,v\right\}  \right)  =\ell^{\prime}\left(  \underbrace{\pi^{-1}\left(
\left\{  u,v\right\}  \right)  }_{\substack{=\left(  u,v\right)
\\\text{(since }\left\{  u,v\right\}  =\pi\left(  u,v\right)  \text{)}%
}}\right)  =\ell^{\prime}\left(  u,v\right) \\
&  =\left\vert \operatorname*{inter}\left(  u,v\right)  \right\vert
\ \ \ \ \ \ \ \ \ \ \left(  \text{by the definition of }\ell^{\prime}\right)
,
\end{align*}
and%
\begin{align*}
\underbrace{\ell}_{=\ell^{\prime}\circ\pi^{-1}}\left(  \left\{  v,w\right\}
\right)   &  =\left(  \ell^{\prime}\circ\pi^{-1}\right)  \left(  \left\{
v,w\right\}  \right)  =\ell^{\prime}\left(  \underbrace{\pi^{-1}\left(
\left\{  v,w\right\}  \right)  }_{\substack{=\left(  v,w\right)
\\\text{(since }\left\{  v,w\right\}  =\pi\left(  v,w\right)  \text{)}%
}}\right)  =\ell^{\prime}\left(  v,w\right) \\
&  =\left\vert \operatorname*{inter}\left(  v,w\right)  \right\vert
\ \ \ \ \ \ \ \ \ \ \left(  \text{by the definition of }\ell^{\prime}\right)
.
\end{align*}
Now,%
\begin{align*}
\ell\left(  \left\{  u,w\right\}  \right)   &  =\left\vert
\operatorname*{inter}\left(  u,w\right)  \right\vert >\left\vert
\operatorname*{inter}\left(  u,v\right)  \right\vert
\ \ \ \ \ \ \ \ \ \ \left(  \text{by
(\ref{pf.lem.digraph.2pf-chrom.K-free.c.ieq1})}\right) \\
&  =\ell\left(  \left\{  u,v\right\}  \right)
\end{align*}
and%
\begin{align*}
\ell\left(  \left\{  u,w\right\}  \right)   &  =\left\vert
\operatorname*{inter}\left(  u,w\right)  \right\vert >\left\vert
\operatorname*{inter}\left(  v,w\right)  \right\vert
\ \ \ \ \ \ \ \ \ \ \left(  \text{by
(\ref{pf.lem.digraph.2pf-chrom.K-free.c.ieq2})}\right) \\
&  =\ell\left(  \left\{  v,w\right\}  \right)  .
\end{align*}
Thus, (\ref{eq.lem.digraph.2pf-chrom.K-free.c.2}) holds. This proves Lemma
\ref{lem.digraph.2pf-chrom.K-free} \textbf{(c)}.

\textbf{(d)} Let $K\in\mathfrak{K}$. We shall show that $K$ is a broken
circuit of $G$.

We have
\begin{align*}
K  &  \in\mathfrak{K}=\left\{  \left\{  \left\{  i,k\right\}  ,\left\{
k,j\right\}  \right\}  \ \mid\ \left(  i,k,j\right)  \in Z\right\} \\
&  =\left\{  \left\{  \left\{  u,v\right\}  ,\left\{  v,w\right\}  \right\}
\ \mid\ \left(  u,v,w\right)  \in Z\right\}
\end{align*}
(here, we renamed the index $\left(  i,k,j\right)  $ as $\left(  u,v,w\right)
\in Z$). Hence, $K=\left\{  \left\{  u,v\right\}  ,\left\{  v,w\right\}
\right\}  $ for some $\left(  u,v,w\right)  \in Z$. Consider this $\left(
u,v,w\right)  $. We have%
\[
\left(  u,v,w\right)  \in Z=\left\{  \left(  i,k,j\right)  \in V^{3}%
\ \mid\ \left(  i,k\right)  \in A\text{ and }\left(  k,j\right)  \in
A\right\}  .
\]
In other words, $\left(  u,v,w\right)  $ is an $\left(  i,k,j\right)  \in
V^{3}$ satisfying $\left(  i,k\right)  \in A$ and $\left(  k,j\right)  \in A$.
In other words, $\left(  u,v,w\right)  $ is an element of $V^{3}$ and
satisfies $\left(  u,v\right)  \in A$ and $\left(  v,w\right)  \in A$. Hence,
Lemma \ref{lem.digraph.2pf-chrom.K-free} \textbf{(c)} shows that we have%
\begin{align*}
\left\{  u,w\right\}   &  \in E,\ \ \ \ \ \ \ \ \ \ \left\{  u,v\right\}  \in
E,\ \ \ \ \ \ \ \ \ \ \left\{  u,w\right\}  \in E,\\
\ell\left(  \left\{  u,w\right\}  \right)   &  >\ell\left(  \left\{
u,v\right\}  \right)  \ \ \ \ \ \ \ \ \ \ \text{and}\ \ \ \ \ \ \ \ \ \ \ell
\left(  \left\{  u,w\right\}  \right)  >\ell\left(  \left\{  v,w\right\}
\right)  .
\end{align*}
Let $D$ be the set $\left\{  \left\{  u,v\right\}  ,\left\{  v,w\right\}
,\left\{  u,w\right\}  \right\}  $. Now, Lemma
\ref{lem.digraph.2pf-chrom.3cyc} (applied to $C=D$) shows that the set $D$ is
a circuit of $G$ and satisfies $D\setminus\left\{  \left\{  u,w\right\}
\right\}  =\left\{  \left\{  u,v\right\}  ,\left\{  v,w\right\}  \right\}  $.
Hence, $K=\left\{  \left\{  u,v\right\}  ,\left\{  v,w\right\}  \right\}
=D\setminus\left\{  \left\{  u,w\right\}  \right\}  $. Moreover, the edge
$\left\{  u,w\right\}  $ is the unique edge in $D$ having maximum label (among
the edges in $D$)\ \ \ \ \footnote{\textit{Proof.} We have $D=\left\{
\left\{  u,v\right\}  ,\left\{  v,w\right\}  ,\left\{  u,w\right\}  \right\}
$. Now, $\left\{  u,w\right\}  \in\left\{  \left\{  u,v\right\}  ,\left\{
v,w\right\}  ,\left\{  u,w\right\}  \right\}  =D$. Thus, $\left\{
u,w\right\}  $ is an edge in $D$.
\par
Now, we need to show that the edge $\left\{  u,w\right\}  $ is the unique edge
in $D$ having maximum label (among the edges in $D$). Indeed, assume the
contrary (for the sake of contradiction). Thus, $\left\{  u,w\right\}  $ is
\textbf{not} the unique edge in $D$ having maximum label (among the edges in
$D$). In other words, there exists an edge $e\in D$ distinct from $\left\{
u,w\right\}  $ such that the label of $e$ is greater or equal to the label of
$\left\{  u,w\right\}  $ (because we already know that $\left\{  u,w\right\}
$ is an edge in $D$). Consider such an $e$.
\par
The label of $e$ is greater or equal to the label of $\left\{  u,w\right\}  $.
In other words, $\ell\left(  e\right)  $ is greater or equal to $\ell\left(
\left\{  u,w\right\}  \right)  $ (since the label of $e$ is $\ell\left(
e\right)  $ (by the definition of the \textquotedblleft
label\textquotedblright) and since the label of $\left\{  u,w\right\}  $ is
$\ell\left(  \left\{  u,w\right\}  \right)  $ (by the definition of the
\textquotedblleft label\textquotedblright)). In other words, $\ell\left(
e\right)  \geq\ell\left(  \left\{  u,w\right\}  \right)  $.
\par
We have $e\in D$ and $e\notin\left\{  u,w\right\}  $ (since $e$ is distinct
from $\left\{  u,w\right\}  $). Thus, $e\in D\setminus\left\{  \left\{
u,w\right\}  \right\}  =\left\{  \left\{  u,v\right\}  ,\left\{  v,w\right\}
\right\}  $.
\par
But $\ell\left(  e\right)  \geq\ell\left(  \left\{  u,w\right\}  \right)
>\ell\left(  \left\{  u,v\right\}  \right)  $, so that $\ell\left(  e\right)
\neq\ell\left(  \left\{  u,v\right\}  \right)  $ and thus $e\neq\left\{
u,v\right\}  $. Also, $\ell\left(  e\right)  \geq\ell\left(  \left\{
u,w\right\}  \right)  >\ell\left(  \left\{  v,w\right\}  \right)  $, so that
$\ell\left(  e\right)  \neq\ell\left(  \left\{  v,w\right\}  \right)  $ and
thus $e\neq\left\{  v,w\right\}  $. Thus, $e$ equals neither of the two edges
$\left\{  u,v\right\}  $ and $\left\{  v,w\right\}  $ (since $e\neq\left\{
u,v\right\}  $ and $e\neq\left\{  v,w\right\}  $). In other words,
$e\notin\left\{  \left\{  u,v\right\}  ,\left\{  v,w\right\}  \right\}  $.
This contradicts $e\in\left\{  \left\{  u,v\right\}  ,\left\{  v,w\right\}
\right\}  $. This contradiction shows that our assumption was wrong. Hence, we
have shown that the edge $\left\{  u,w\right\}  $ is the unique edge in $D$
having maximum label (among the edges in $D$). Qed.}.

Also, $K=\left\{  \left\{  u,v\right\}  ,\left\{  v,w\right\}  \right\}
=\underbrace{\left\{  \left\{  u,v\right\}  \right\}  }_{\substack{\subseteq
E\\\text{(since }\left\{  u,v\right\}  \in E\text{)}}}\cup\underbrace{\left\{
\left\{  v,w\right\}  \right\}  }_{\substack{\subseteq E\\\text{(since
}\left\{  v,w\right\}  \in E\text{)}}}\subseteq E\cup E=E$. Thus, $K$ is a
subset of $E$.

Altogether, we have now shown that $K$ is a subset of $E$ and satisfies
$K=D\setminus\left\{  \left\{  u,w\right\}  \right\}  $; we also know that $D$
is a circuit of $G$, and that $\left\{  u,w\right\}  $ is the unique edge in
$D$ having maximum label (among the edges in $D$). Hence, $K$ is a subset of
$E$ having the form $C\setminus\left\{  e\right\}  $, where $C$ is a circuit
of $G$, and where $e$ is the unique edge in $C$ having maximum label (among
the edges in $C$)\ \ \ \ \footnote{Namely, $K$ has this form for $C=D$ and
$e=\left\{  u,w\right\}  $.}. In other words, $K$ is a broken circuit of $G$
(since $K$ is a broken circuit of $G$ if and only if $K$ is a subset of $E$
having the form $C\setminus\left\{  e\right\}  $, where $C$ is a circuit of
$G$, and where $e$ is the unique edge in $C$ having maximum label (among the
edges in $C$)\ \ \ \ \footnote{by the definition of a \textquotedblleft broken
circuit\textquotedblright})). Now, forget that we fixed $K$. We thus have
shown that every $K\in\mathfrak{K}$ is a broken circuit of $G$. In other
words, every element of $\mathfrak{K}$ is a broken circuit of $G$. This proves
Lemma \ref{lem.digraph.2pf-chrom.K-free} \textbf{(d)}.

\textbf{(e)} The map $\pi:A\rightarrow E$ is bijective (by Lemma
\ref{lem.digraph.2pf-chrom.K-free} \textbf{(a)}). In particular, the map $\pi$
is therefore injective.

We first observe a simple fact: If $u$ and $v$ are two elements of $V$ such
that $\left(  u,v\right)  \in A$ and $\left\{  u,v\right\}  \in\pi\left(
F\right)  $, then%
\begin{equation}
\left(  u,v\right)  \in F \label{pf.lem.digraph.2pf-chrom.K-free.e.lem1}%
\end{equation}
\footnote{\textit{Proof of (\ref{pf.lem.digraph.2pf-chrom.K-free.e.lem1}):}
Let $u$ and $v$ be two elements of $V$ such that $\left(  u,v\right)  \in A$
and $\left\{  u,v\right\}  \in\pi\left(  F\right)  $. We must show that
$\left(  u,v\right)  \in F$.
\par
We have $\left\{  u,v\right\}  \in\pi\left(  F\right)  $. In other words,
there exists some $f\in F$ such that $\left\{  u,v\right\}  =\pi\left(
f\right)  $. Consider this $f$.
\par
But $\left(  u,v\right)  \in A$; therefore, the definition of $\pi$ yields%
\begin{align*}
\pi\left(  u,v\right)   &  =\operatorname*{set}\left(  u,v\right)  =\left\{
u,v\right\}  \ \ \ \ \ \ \ \ \ \ \left(  \text{by the definition of the map
}\operatorname*{set}\right) \\
&  =\pi\left(  f\right)  .
\end{align*}
Since $\pi$ is injective, this entails that $\left(  u,v\right)  =f\in F$.
This proves (\ref{pf.lem.digraph.2pf-chrom.K-free.e.lem1}).}.

Also, $\pi\left(  \underbrace{F}_{\subseteq A}\right)  \subseteq\pi\left(
A\right)  \subseteq E$. Thus, $\pi\left(  F\right)  $ is a subset of $E$.

Let us now prove the implication%
\begin{equation}
\ \left(  \text{the digraph }\left(  V,F\right)  \text{ is }2\text{-path-free}%
\right)  \ \Longrightarrow\ \left(  \text{the set }\pi\left(  F\right)  \text{
is }\mathfrak{K}\text{-free}\right)  .
\label{pf.lem.digraph.2pf-chrom.K-free.e.imp1}%
\end{equation}


\textit{Proof of (\ref{pf.lem.digraph.2pf-chrom.K-free.e.imp1}):} Assume that
the digraph $\left(  V,F\right)  $ is $2$-path-free. We shall show that the
set $\pi\left(  F\right)  $ is $\mathfrak{K}$-free.

Indeed, let $K\in\mathfrak{K}$ be such that $K\subseteq\pi\left(  F\right)  $.
We shall obtain a contradiction.

The digraph $\left(  V,F\right)  $ is $2$-path-free if and only if there exist
no three elements $u$, $v$ and $w$ of $V$ satisfying $\left(  u,v\right)  \in
F$ and $\left(  v,w\right)  \in F$ (by the definition of \textquotedblleft%
$2$-path-free\textquotedblright). Therefore, there exist no three elements
$u$, $v$ and $w$ of $V$ satisfying $\left(  u,v\right)  \in F$ and $\left(
v,w\right)  \in F$ (since we know that the digraph $\left(  V,F\right)  $ is
$2$-path-free). In other words, if $u$, $v$ and $w$ are three elements of $V$,
then%
\begin{equation}
\left(  \left(  u,v\right)  \in F\text{ and }\left(  v,w\right)  \in F\right)
\text{ is false.} \label{pf.lem.digraph.2pf-chrom.K-free.e.imp1.pf.2}%
\end{equation}


But%
\begin{align*}
K  &  \in\mathfrak{K}=\left\{  \left\{  \left\{  i,k\right\}  ,\left\{
k,j\right\}  \right\}  \ \mid\ \left(  i,k,j\right)  \in Z\right\} \\
&  =\left\{  \left\{  \left\{  u,v\right\}  ,\left\{  v,w\right\}  \right\}
\ \mid\ \left(  u,v,w\right)  \in Z\right\}
\end{align*}
(here, we renamed the index $\left(  i,k,j\right)  $ as $\left(  u,v,w\right)
\in Z$). Hence, $K=\left\{  \left\{  u,v\right\}  ,\left\{  v,w\right\}
\right\}  $ for some $\left(  u,v,w\right)  \in Z$. Consider this $\left(
u,v,w\right)  $. We have%
\[
\left(  u,v,w\right)  \in Z=\left\{  \left(  i,k,j\right)  \in V^{3}%
\ \mid\ \left(  i,k\right)  \in A\text{ and }\left(  k,j\right)  \in
A\right\}  .
\]
In other words, $\left(  u,v,w\right)  $ is an $\left(  i,k,j\right)  \in
V^{3}$ satisfying $\left(  i,k\right)  \in A$ and $\left(  k,j\right)  \in A$.
In other words, $\left(  u,v,w\right)  $ is an element of $V^{3}$ and
satisfies $\left(  u,v\right)  \in A$ and $\left(  v,w\right)  \in A$.

Now, $\left(  u,v\right)  \in A$ and $\left\{  u,v\right\}  \in\left\{
\left\{  u,v\right\}  ,\left\{  v,w\right\}  \right\}  =K\subseteq\pi\left(
F\right)  $. Hence, (\ref{pf.lem.digraph.2pf-chrom.K-free.e.lem1}) shows that
$\left(  u,v\right)  \in F$.

Also, $\left(  v,w\right)  \in A$ and $\left\{  v,w\right\}  \in\left\{
\left\{  u,v\right\}  ,\left\{  v,w\right\}  \right\}  =K\subseteq\pi\left(
F\right)  $. Hence, (\ref{pf.lem.digraph.2pf-chrom.K-free.e.lem1}) (applied to
$v$ and $w$ instead of $u$ and $v$) shows that $\left(  v,w\right)  \in F$.
Thus, we have $\left(  \left(  u,v\right)  \in F\text{ and }\left(
v,w\right)  \in F\right)  $.

But (\ref{pf.lem.digraph.2pf-chrom.K-free.e.imp1.pf.2}) shows that $\left(
\left(  u,v\right)  \in F\text{ and }\left(  v,w\right)  \in F\right)  $ is
false. This contradicts the fact that we have $\left(  \left(  u,v\right)  \in
F\text{ and }\left(  v,w\right)  \in F\right)  $.

Now, let us forget that we fixed $K$. We thus have obtained a contradiction
for each $K\in\mathfrak{K}$ satisfying $K\subseteq\pi\left(  F\right)  $.
Hence, there exists no $K\in\mathfrak{K}$ satisfying $K\subseteq\pi\left(
F\right)  $. In other words, the set $\pi\left(  F\right)  $ contains no
$K\in\mathfrak{K}$ as a subset.

The subset $\pi\left(  F\right)  $ of $E$ is $\mathfrak{K}$-free if and only
if $\pi\left(  F\right)  $ contains no $K\in\mathfrak{K}$ as a subset (by the
definition of \textquotedblleft$\mathfrak{K}$-free\textquotedblright). Thus,
the subset $\pi\left(  F\right)  $ of $E$ is $\mathfrak{K}$-free (since
$\pi\left(  F\right)  $ contains no $K\in\mathfrak{K}$ as a subset).

Now, let us forget that we assumed that the digraph $\left(  V,F\right)  $ is
$2$-path-free. We thus have shown that the set $\pi\left(  F\right)  $ is
$\mathfrak{K}$-free if the digraph $\left(  V,F\right)  $ is $2$-path-free. In
other words, we have proven the implication
(\ref{pf.lem.digraph.2pf-chrom.K-free.e.imp1}).

Next, we shall prove the following implication:%
\begin{equation}
\ \left(  \text{the set }\pi\left(  F\right)  \text{ is }\mathfrak{K}%
\text{-free}\right)  \ \Longrightarrow\ \left(  \text{the digraph }\left(
V,F\right)  \text{ is }2\text{-path-free}\right)  .
\label{pf.lem.digraph.2pf-chrom.K-free.e.imp2}%
\end{equation}


\textit{Proof of (\ref{pf.lem.digraph.2pf-chrom.K-free.e.imp2}):} Assume that
the set $\pi\left(  F\right)  $ is $\mathfrak{K}$-free. We shall show that the
digraph $\left(  V,F\right)  $ is $2$-path-free.

Let $u$, $v$ and $w$ be three elements of $V$ satisfying $\left(  u,v\right)
\in F$ and $\left(  v,w\right)  \in F$. We shall derive a contradiction.

Clearly, $\left(  u,v\right)  \in F\subseteq A$ and $\left(  v,w\right)  \in
F\subseteq A$.

The subset $\pi\left(  F\right)  $ of $E$ is $\mathfrak{K}$-free if and only
if $\pi\left(  F\right)  $ contains no $K\in\mathfrak{K}$ as a subset (by the
definition of \textquotedblleft$\mathfrak{K}$-free\textquotedblright). Thus,
$\pi\left(  F\right)  $ contains no $K\in\mathfrak{K}$ as a subset (since the
subset $\pi\left(  F\right)  $ of $E$ is $\mathfrak{K}$-free). In other words,
there exists no $K\in\mathfrak{K}$ such that $K\subseteq\pi\left(  F\right)  $.

But $\left(  u,v,w\right)  $ is an element of $V^{3}$ and satisfies $\left(
u,v\right)  \in A$ and $\left(  v,w\right)  \in A$. In other words, $\left(
u,v,w\right)  $ is an $\left(  i,k,j\right)  \in V^{3}$ satisfying $\left(
i,k\right)  \in A$ and $\left(  k,j\right)  \in A$. Hence,%
\[
\left(  u,v,w\right)  \in\left\{  \left(  i,k,j\right)  \in V^{3}%
\ \mid\ \left(  i,k\right)  \in A\text{ and }\left(  k,j\right)  \in
A\right\}  =Z.
\]
Thus, $\left\{  \left\{  u,v\right\}  ,\left\{  v,w\right\}  \right\}  $ is a
set of the form $\left\{  \left\{  i,k\right\}  ,\left\{  k,j\right\}
\right\}  $ for some $\left(  i,k,j\right)  \in Z$ (namely, for $\left(
i,k,j\right)  =\left(  u,v,w\right)  $). Hence,%
\begin{equation}
\left\{  \left\{  u,v\right\}  ,\left\{  v,w\right\}  \right\}  \in\left\{
\left\{  \left\{  i,k\right\}  ,\left\{  k,j\right\}  \right\}  \ \mid
\ \left(  i,k,j\right)  \in Z\right\}  =\mathfrak{K}.
\label{pf.lem.digraph.2pf-chrom.K-free.e.imp2.pf.2}%
\end{equation}


But $\left\{  u,v\right\}  \in\pi\left(  F\right)  $%
\ \ \ \ \footnote{\textit{Proof.} We have $\left(  u,v\right)  \in A$;
therefore, the definition of $\pi$ yields%
\[
\pi\left(  u,v\right)  =\operatorname*{set}\left(  u,v\right)  =\left\{
u,v\right\}  \ \ \ \ \ \ \ \ \ \ \left(  \text{by the definition of the map
}\operatorname*{set}\right)  .
\]
Hence, $\left\{  u,v\right\}  =\pi\underbrace{\left(  u,v\right)  }_{\in F}%
\in\pi\left(  F\right)  $, qed.} and $\left\{  v,w\right\}  \in\pi\left(
F\right)  $\ \ \ \ \footnote{\textit{Proof.} We have $\left(  v,w\right)  \in
A$; therefore, the definition of $\pi$ yields%
\[
\pi\left(  v,w\right)  =\operatorname*{set}\left(  v,w\right)  =\left\{
v,w\right\}  \ \ \ \ \ \ \ \ \ \ \left(  \text{by the definition of the map
}\operatorname*{set}\right)  .
\]
Hence, $\left\{  v,w\right\}  =\pi\underbrace{\left(  v,w\right)  }_{\in F}%
\in\pi\left(  F\right)  $, qed.}. Thus, both $\left\{  u,v\right\}  $ and
$\left\{  v,w\right\}  $ belong to the set $\pi\left(  F\right)  $. Therefore,%
\[
\left\{  \left\{  u,v\right\}  ,\left\{  v,w\right\}  \right\}  \subseteq
\pi\left(  F\right)  .
\]
Combining this with (\ref{pf.lem.digraph.2pf-chrom.K-free.e.imp2.pf.2}), we
see that the set $\left\{  \left\{  u,v\right\}  ,\left\{  v,w\right\}
\right\}  \in\mathfrak{K}$ satisfies $\left\{  \left\{  u,v\right\}  ,\left\{
v,w\right\}  \right\}  \subseteq\pi\left(  F\right)  $. Thus, there exists an
$K\in\mathfrak{K}$ such that $K\subseteq\pi\left(  F\right)  $ (namely,
$K=\left\{  \left\{  u,v\right\}  ,\left\{  v,w\right\}  \right\}  $). This
contradicts the fact that there exists no $K\in\mathfrak{K}$ such that
$K\subseteq\pi\left(  F\right)  $.

Now, forget that we fixed $u$, $v$ and $w$. We thus have obtained a
contradiction for every three elements $u$, $v$ and $w$ of $V$ satisfying
$\left(  u,v\right)  \in F$ and $\left(  v,w\right)  \in F$. Hence, there
exist no three elements $u$, $v$ and $w$ of $V$ satisfying $\left(
u,v\right)  \in F$ and $\left(  v,w\right)  \in F$.

But the digraph $\left(  V,F\right)  $ is $2$-path-free if and only if there
exist no three elements $u$, $v$ and $w$ of $V$ satisfying $\left(
u,v\right)  \in F$ and $\left(  v,w\right)  \in F$ (by the definition of
\textquotedblleft$2$-path-free\textquotedblright). Thus, the digraph $\left(
V,F\right)  $ is $2$-path-free (since there exist no three elements $u$, $v$
and $w$ of $V$ satisfying $\left(  u,v\right)  \in F$ and $\left(  v,w\right)
\in F$).

Now, let us forget that we have assumed that the set $\pi\left(  F\right)  $
is $\mathfrak{K}$-free. We thus have shown that if the set $\pi\left(
F\right)  $ is $\mathfrak{K}$-free, then the digraph $\left(  V,F\right)  $ is
$2$-path-free. In other words, we have proven the implication
(\ref{pf.lem.digraph.2pf-chrom.K-free.e.imp2}).

Now, by combining the two implications
(\ref{pf.lem.digraph.2pf-chrom.K-free.e.imp1}) and
(\ref{pf.lem.digraph.2pf-chrom.K-free.e.imp2}), we obtain the logical
equivalence%
\[
\ \left(  \text{the digraph }\left(  V,F\right)  \text{ is }2\text{-path-free}%
\right)  \ \Longleftrightarrow\ \left(  \text{the set }\pi\left(  F\right)
\text{ is }\mathfrak{K}\text{-free}\right)  .
\]
This proves Lemma \ref{lem.digraph.2pf-chrom.K-free} \textbf{(e)}.
\end{proof}

\begin{proof}
[Proof of Proposition \ref{prop.digraph.2pf-chrom}.]First, let us introduce a
general notation: If $X$ and $Y$ are two sets, and if $f:X\rightarrow Y$ is
any map, then we define a map $\mathcal{P}\left(  f\right)  :\mathcal{P}%
\left(  X\right)  \rightarrow\mathcal{P}\left(  Y\right)  $ by%
\[
\left(  \left(  \mathcal{P}\left(  f\right)  \right)  \left(  Z\right)
=f\left(  Z\right)  \ \ \ \ \ \ \ \ \ \ \text{for every }Z\in\mathcal{P}%
\left(  X\right)  \right)  .
\]
If the map $f:X\rightarrow Y$ is bijective, then%
\begin{equation}
\text{the map }\mathcal{P}\left(  f\right)  :\mathcal{P}\left(  X\right)
\rightarrow\mathcal{P}\left(  Y\right)  \text{ is bijective as well}
\label{pf.prop.digraph.2pf-chrom.P(f)bij}%
\end{equation}
(and, in fact, its inverse is $\mathcal{P}\left(  f^{-1}\right)  $).

Let $E=\operatorname*{set}A$ and $G=\underline{D}$. The definition of
$\underline{D}$ shows that $\underline{D}=\left(
V,\underbrace{\operatorname*{set}A}_{=E}\right)  =\left(  V,E\right)  $. Thus,
$G=\underline{D}=\left(  V,E\right)  $.

Define the map $\pi:A\rightarrow E$ as in Lemma
\ref{lem.digraph.2pf-chrom.K-free}. Lemma \ref{lem.digraph.2pf-chrom.K-free}
\textbf{(a)} shows that this map $\pi:A\rightarrow E$ is bijective. Hence, the
map $\mathcal{P}\left(  \pi\right)  :\mathcal{P}\left(  A\right)
\rightarrow\mathcal{P}\left(  E\right)  $ is bijective (according to
(\ref{pf.prop.digraph.2pf-chrom.P(f)bij}), applied to $X=A$, $Y=E$ and $f=\pi
$). We notice that every $F\in\mathcal{P}\left(  A\right)  $ satisfies%
\begin{align}
\left(  \mathcal{P}\left(  \pi\right)  \right)  \left(  F\right)   &
=\pi\left(  F\right)  \ \ \ \ \ \ \ \ \ \ \left(  \text{by the definition of
the map }\mathcal{P}\left(  \pi\right)  \right)
\label{pf.prop.digraph.2pf-chrom.1}\\
&  =\left\{  \underbrace{\pi\left(  a\right)  }%
_{\substack{=\operatorname*{set}a\\\text{(by the definition of }\pi\text{)}%
}}\ \mid\ a\in F\right\}  =\left\{  \operatorname*{set}a\ \mid\ a\in F\right\}
\nonumber\\
&  =\operatorname*{set}F. \label{pf.prop.digraph.2pf-chrom.2}%
\end{align}
and%
\begin{equation}
\left\vert \underbrace{\left(  \mathcal{P}\left(  \pi\right)  \right)  \left(
F\right)  }_{\substack{=\pi\left(  F\right)  \\\text{(by
(\ref{pf.prop.digraph.2pf-chrom.1}))}}}\right\vert =\left\vert \pi\left(
F\right)  \right\vert =\left\vert F\right\vert
\label{pf.prop.digraph.2pf-chrom.4}%
\end{equation}
(since the map $\pi$ is injective (since the map $\pi$ is bijective)).

Define two sets $Z$ and $\mathfrak{K}$ as in Lemma
\ref{lem.digraph.2pf-chrom.K-free} \textbf{(b)}. Define a map $\ell
:E\rightarrow\mathbb{N}$ as in Lemma \ref{lem.digraph.2pf-chrom.K-free}
\textbf{(c)}. Definition \ref{def.BC} (applied to $X=\mathbb{N}$) shows that
the notion of a broken circuit of $G$ is well-defined (since a function
$\ell:E\rightarrow\mathbb{N}$ is given).

Lemma \ref{lem.digraph.2pf-chrom.K-free} \textbf{(d)} shows that every element
of $\mathfrak{K}$ is a broken circuit of $G$. Thus, $\mathfrak{K}$ is a set of
broken circuits of $G$ (not necessarily containing all of them). Hence,
Corollary \ref{cor.chrompol.K-free} (applied to $X=\mathbb{N}$) shows that%
\begin{align*}
\chi_{G}  &  =\underbrace{\sum_{\substack{F\subseteq E;\\F\text{ is
}\mathfrak{K}\text{-free}}}}_{=\sum_{\substack{F\in\mathcal{P}\left(
E\right)  ;\\F\text{ is }\mathfrak{K}\text{-free}}}}\left(  -1\right)
^{\left\vert F\right\vert }x^{\operatorname*{conn}\left(  V,F\right)  }%
=\sum_{\substack{F\in\mathcal{P}\left(  E\right)  ;\\F\text{ is }%
\mathfrak{K}\text{-free}}}\left(  -1\right)  ^{\left\vert F\right\vert
}x^{\operatorname*{conn}\left(  V,F\right)  }\\
&  =\underbrace{\sum_{\substack{F\in\mathcal{P}\left(  A\right)  ;\\\left(
\mathcal{P}\left(  \pi\right)  \right)  \left(  F\right)  \text{ is
}\mathfrak{K}\text{-free}}}}_{\substack{=\sum_{\substack{F\in\mathcal{P}%
\left(  A\right)  ;\\\pi\left(  F\right)  \text{ is }\mathfrak{K}\text{-free}%
}}\\\text{(since every }F\in\mathcal{P}\left(  A\right)  \text{ satisfies}%
\\\left(  \mathcal{P}\left(  \pi\right)  \right)  \left(  F\right)
=\pi\left(  F\right)  \\\text{(by (\ref{pf.prop.digraph.2pf-chrom.1})))}%
}}\underbrace{\left(  -1\right)  ^{\left\vert \left(  \mathcal{P}\left(
\pi\right)  \right)  \left(  F\right)  \right\vert }}_{\substack{=\left(
-1\right)  ^{\left\vert F\right\vert }\\\text{(since }\left\vert \left(
\mathcal{P}\left(  \pi\right)  \right)  \left(  F\right)  \right\vert
=\left\vert F\right\vert \\\text{(by (\ref{pf.prop.digraph.2pf-chrom.4})))}%
}}\underbrace{x^{\operatorname*{conn}\left(  V,\left(  \mathcal{P}\left(
\pi\right)  \right)  \left(  F\right)  \right)  }}%
_{\substack{=x^{\operatorname*{conn}\left(  V,\operatorname*{set}F\right)
}\\\text{(since }\left(  \mathcal{P}\left(  \pi\right)  \right)  \left(
F\right)  =\operatorname*{set}F\\\text{(by (\ref{pf.prop.digraph.2pf-chrom.2}%
)))}}}\\
&  \ \ \ \ \ \ \ \ \ \ \left(
\begin{array}
[c]{c}%
\text{here, we have substituted }\left(  \mathcal{P}\left(  \pi\right)
\right)  \left(  F\right)  \text{ for }F\text{ in the sum,}\\
\text{since the map }\mathcal{P}\left(  \pi\right)  :\mathcal{P}\left(
A\right)  \rightarrow\mathcal{P}\left(  E\right)  \text{ is bijective}%
\end{array}
\right) \\
&  =\underbrace{\sum_{\substack{F\in\mathcal{P}\left(  A\right)  ;\\\pi\left(
F\right)  \text{ is }\mathfrak{K}\text{-free}}}}_{=\sum_{\substack{F\subseteq
A;\\\pi\left(  F\right)  \text{ is }\mathfrak{K}\text{-free}}}=\sum
_{\substack{F\subseteq A;\\\text{the set }\pi\left(  F\right)  \text{ is
}\mathfrak{K}\text{-free}}}}\left(  -1\right)  ^{\left\vert F\right\vert
}x^{\operatorname*{conn}\left(  V,\operatorname*{set}F\right)  }\\
&  =\underbrace{\sum_{\substack{F\subseteq A;\\\text{the set }\pi\left(
F\right)  \text{ is }\mathfrak{K}\text{-free}}}}_{\substack{=\sum
_{\substack{F\subseteq A;\\\text{the digraph }\left(  V,F\right)  \text{ is
}2\text{-path-free}}}\\\text{(because for any subset }F\text{ of }A\text{, the
condition}\\\left(  \text{the set }\pi\left(  F\right)  \text{ is
}\mathfrak{K}\text{-free}\right)  \text{ is equivalent to the}%
\\\text{condition }\left(  \text{the digraph }\left(  V,F\right)  \text{ is
}2\text{-path-free}\right)  \\\text{(by Lemma
\ref{lem.digraph.2pf-chrom.K-free} \textbf{(e)}))}}}\left(  -1\right)
^{\left\vert F\right\vert }x^{\operatorname*{conn}\left(
V,\operatorname*{set}F\right)  }\\
&  =\sum_{\substack{F\subseteq A;\\\text{the digraph }\left(  V,F\right)
\text{ is }2\text{-path-free}}}\left(  -1\right)  ^{\left\vert F\right\vert
}x^{\operatorname*{conn}\left(  V,\operatorname*{set}F\right)  }.
\end{align*}
Since $\underline{D}=G$, this rewrites as
\[
\chi_{\underline{D}}=\sum_{\substack{F\subseteq A;\\\text{the digraph }\left(
V,F\right)  \text{ is }2\text{-path-free}}}\left(  -1\right)  ^{\left\vert
F\right\vert }x^{\operatorname*{conn}\left(  V,\operatorname*{set}F\right)
}.
\]
This proves Proposition \ref{prop.digraph.2pf-chrom}.
\end{proof}
\end{verlong}

\section{A matroidal generalization}

\subsection{An introduction to matroids}

We shall now present a result that can be considered as a generalization of
Theorem \ref{thm.chrompol.varis} in a different direction than Theorem
\ref{thm.chromsym.varis}: namely, a formula for the characteristic polynomial
of a matroid. Let us first recall the basic notions from the theory of
matroids that will be needed to state it.

\begin{verlong}
[TODO: Make the following proofs more detailed.]
\end{verlong}

First, we introduce some basic poset-related terminology:

\begin{definition}
\label{def.poset.max}Let $P$ be a poset.

\textbf{(a)} An element $v$ of $P$ is said to be \textit{maximal} (with
respect to $P$) if and only if every $w\in P$ satisfying $w\geq v$ must
satisfy $w=v$.

\textbf{(b)} An element $v$ of $P$ is said to be \textit{minimal} (with
respect to $P$) if and only if every $w\in P$ satisfying $w\leq v$ must
satisfy $w=v$.
\end{definition}

\begin{definition}
For any set $E$, we shall regard the powerset $\mathcal{P}\left(  E\right)  $
as a poset (with respect to inclusion). Thus, any subset $\mathcal{S}$ of
$\mathcal{P}\left(  E\right)  $ also becomes a poset, and therefore the
notions of \textquotedblleft minimal\textquotedblright\ and \textquotedblleft
maximal\textquotedblright\ elements in $\mathcal{S}$ make sense. Beware that
these notions are not related to size; i.e., a maximal element of
$\mathcal{S}$ might not be a maximum-size element of $\mathcal{S}$.
\end{definition}

Now, let us define the notion of \textquotedblleft matroid\textquotedblright%
\ that we will use:

\begin{definition}
\label{def.matroid}\textbf{(a)} A \textit{matroid} means a pair $\left(
E,\mathcal{I}\right)  $ consisting of a finite set $E$ and a set
$\mathcal{I}\subseteq\mathcal{P}\left(  E\right)  $ satisfying the following axioms:

\begin{itemize}
\item \textit{Matroid axiom 1:} We have $\varnothing\in\mathcal{I}$.

\item \textit{Matroid axiom 2:} If $Y\in\mathcal{I}$ and $Z\in\mathcal{P}%
\left(  E\right)  $ are such that $Z\subseteq Y$, then $Z\in\mathcal{I}$.

\item \textit{Matroid axiom 3:} If $Y\in\mathcal{I}$ and $Z\in\mathcal{I}$ are
such that $\left\vert Y\right\vert <\left\vert Z\right\vert $, then there
exists some $x\in Z\setminus Y$ such that $Y\cup\left\{  x\right\}
\in\mathcal{I}$.
\end{itemize}

\textbf{(b)} Let $\left(  E,\mathcal{I}\right)  $ be a matroid. A subset $S$
of $E$ is said to be \textit{independent} (for this matroid) if and only if
$S\in\mathcal{I}$. The set $E$ is called the \textit{ground set} of the
matroid $\left(  E,\mathcal{I}\right)  $.
\end{definition}

Different texts give different definitions of a matroid; these definitions are
(mostly) equivalent, but not always in the obvious way\footnote{I.e., it
sometimes happens that two different texts both define a matroid as a pair
$\left(  E,U\right)  $ of a finite set $E$ and a subset $U\subseteq
\mathcal{P}\left(  E\right)  $, but they require these pairs $\left(
E,U\right)  $ to satisfy non-equivalent axioms, and the equivalence between
their definitions is more complicated than just \textquotedblleft a pair
$\left(  E,U\right)  $ is a matroid for one definition if and only if it is a
matroid for the other\textquotedblright.}. Definition \ref{def.matroid} is how
a matroid is defined in \cite[\S 10.1]{Schrij13} and in \cite[Definition
3.15]{Martin15} (where it is called a \textquotedblleft(matroid) independence
system\textquotedblright). There exist other definitions of a matroid, which
turn out to be equivalent. The definition of a matroid given in Stanley's
\cite[Definition 3.8]{Stanley} is directly equivalent to Definition
\ref{def.matroid}, with the only differences that

\begin{itemize}
\item Stanley replaces Matroid axiom 1 by the requirement that $\mathcal{I}%
\neq\varnothing$ (which is, of course, equivalent to Matroid axiom 1 as long
as Matroid axiom 2 is assumed), and

\item Stanley replaces Matroid axiom 3 by the requirement that for every
$T\in\mathcal{P}\left(  E\right)  $, all maximal elements of $\mathcal{I}%
\cap\mathcal{P}\left(  T\right)  $ have the same cardinality\footnote{Here, as
we have already explained, we regard $\mathcal{I}\cap\mathcal{P}\left(
T\right)  $ as a poset with respect to inclusion. Thus, an element $Y$ of this
poset is maximal if and only if there exists no $Z\in\mathcal{I}%
\cap\mathcal{P}\left(  T\right)  $ such that $Y$ is a proper subset of $Z$.}
(this requirement is equivalent to Matroid axiom 3 as long as Matroid axiom 2
is assumed).
\end{itemize}

We now introduce some terminology related to matroids:

\begin{definition}
\label{def.matroid.notions}Let $M=\left(  E,\mathcal{I}\right)  $ be a matroid.

\textbf{(a)} We define a function $r_{M}:\mathcal{P}\left(  E\right)
\rightarrow\mathbb{N}$ by setting%
\begin{equation}
r_{M}\left(  S\right)  =\max\left\{  \left\vert Z\right\vert \ \mid
\ Z\in\mathcal{I}\text{ and }Z\subseteq S\right\}
\ \ \ \ \ \ \ \ \ \ \text{for every }S\subseteq E. \label{eq.def.rank.def}%
\end{equation}
(Note that the right hand side of (\ref{eq.def.rank.def}) is well-defined,
because there exists at least one $Z\in\mathcal{I}$ satisfying $Z\subseteq S$
(namely, $Z=\varnothing$).) If $S$ is a subset of $E$, then the nonnegative
integer $r_{M}\left(  S\right)  $ is called the \textit{rank} of $S$ (with
respect to $M$). It is clear that $r_{M}$ is a weakly increasing function from
the poset $\mathcal{P}\left(  E\right)  $ to $\mathbb{N}$.

\textbf{(b)} If $k\in\mathbb{N}$, then a $k$\textit{-flat} of $M$ means a
subset of $E$ which has rank $k$ and is maximal among all such subsets (i.e.,
it is not a proper subset of any other subset having rank $k$). (Beware: Not
all $k$-flats have the same size.) A \textit{flat} of $M$ is a subset of $E$
which is a $k$-flat for some $k\in\mathbb{N}$. We let $\operatorname*{Flats}M$
denote the set of all flats of $M$; thus, $\operatorname*{Flats}M$ is a
subposet of $\mathcal{P}\left(  E\right)  $.

\textbf{(c)} A \textit{circuit} of $M$ means a minimal element of
$\mathcal{P}\left(  E\right)  \setminus\mathcal{I}$. (That is, a circuit of
$M$ means a subset of $E$ which is not independent (for $M$) and which is
minimal among such subsets.)

\textbf{(d)} An element $e$ of $E$ is said to be a \textit{loop} (of $M$) if
$\left\{  e\right\}  \notin\mathcal{I}$. The matroid $M$ is said to be
\textit{loopless} if no loops (of $M$) exist.
\end{definition}

Notice that the function that we called $r_{M}$ in Definition
\ref{def.matroid.notions} \textbf{(a)} is denoted by $\operatorname*{rk}$ in
Stanley's \cite[Lecture 3]{Stanley}.

One of the most classical examples of a matroid is the \textit{graphical
matroid} of a graph:

\begin{example}
\label{exam.matroid.graphical}Let $G=\left(  V,E\right)  $ be a finite graph.
Define a subset $\mathcal{I}$ of $\mathcal{P}\left(  E\right)  $ by%
\[
\mathcal{I}=\left\{  T\in\mathcal{P}\left(  E\right)  \ \mid\ T\text{ contains
no circuit of }G\text{ as a subset}\right\}  .
\]
Then, $\left(  E,\mathcal{I}\right)  $ is a matroid; it is called the
\textit{graphical matroid} (or the \textit{cycle matroid}) of $G$. It has the
following properties:

\begin{itemize}
\item The matroid $\left(  E,\mathcal{I}\right)  $ is loopless.

\item For each $T\in\mathcal{P}\left(  E\right)  $, we have%
\[
r_{\left(  E,\mathcal{I}\right)  }\left(  T\right)  =\left\vert V\right\vert
-\operatorname*{conn}\left(  V,T\right)
\]
(where $\operatorname*{conn}\left(  V,T\right)  $ is defined as in Definition
\ref{def.conn}).

\item The circuits of the matroid $\left(  E,\mathcal{I}\right)  $ are
precisely the circuits of the graph $G$.

\item The flats of the matroid $\left(  E,\mathcal{I}\right)  $ are related to
colorings of $G$. More precisely: For each set $X$ and each $X$-coloring $f$
of $G$, the set $E\cap\operatorname*{Eqs}f$ is a flat of $\left(
E,\mathcal{I}\right)  $. Every flat of $\left(  E,\mathcal{I}\right)  $ can be
obtained in this way when $X$ is chosen large enough; but often, several
distinct $X$-colorings $f$ lead to one and the same flat $E\cap
\operatorname*{Eqs}f$.
\end{itemize}
\end{example}

We recall three basic facts that are used countless times in arguing about matroids:

\begin{lemma}
\label{lem.matroid.rI}Let $M=\left(  E,\mathcal{I}\right)  $ be a matroid. Let
$T\in\mathcal{I}$. Then, $r_{M}\left(  T\right)  =\left\vert T\right\vert $.
\end{lemma}

\begin{proof}
[Proof of Lemma \ref{lem.matroid.rI}.]We have $T\in\mathcal{I}$ and
$T\subseteq T$. Thus, $T$ is a $Z\in\mathcal{I}$ satisfying $Z\subseteq T$.
Therefore, $\left\vert T\right\vert \in\left\{  \left\vert Z\right\vert
\ \mid\ Z\in\mathcal{I}\text{ and }Z\subseteq T\right\}  $, so that%
\begin{equation}
\left\vert T\right\vert \leq\max\left\{  \left\vert Z\right\vert \ \mid
\ Z\in\mathcal{I}\text{ and }Z\subseteq T\right\}  \label{pf.lem.matroid.rI.1}%
\end{equation}
(since any element of a set of integers is smaller or equal to the maximum of
this set).

On the other hand, the definition of $r_{M}$ yields%
\[
r_{M}\left(  T\right)  =\max\left\{  \left\vert Z\right\vert \ \mid
\ Z\in\mathcal{I}\text{ and }Z\subseteq T\right\}  .
\]
Hence, (\ref{pf.lem.matroid.rI.1}) rewrites as follows:%
\[
\left\vert T\right\vert \leq r_{M}\left(  T\right)  .
\]


Also,
\begin{align*}
r_{M}\left(  T\right)   &  =\max\left\{  \left\vert Z\right\vert \ \mid
\ Z\in\mathcal{I}\text{ and }Z\subseteq T\right\}  \ \ \ \ \ \ \ \ \ \ \left(
\text{by the definition of }r_{M}\right) \\
&  \in\left\{  \left\vert Z\right\vert \ \mid\ Z\in\mathcal{I}\text{ and
}Z\subseteq T\right\}
\end{align*}
(since the maximum of any set belongs to this set). Thus, there exists a
$Z\in\mathcal{I}$ satisfying $Z\subseteq T$ and $r_{M}\left(  T\right)
=\left\vert Z\right\vert $. Consider this $Z$. From $Z\subseteq T$, we obtain
$\left\vert Z\right\vert \leq\left\vert T\right\vert $, so that $r_{M}\left(
T\right)  =\left\vert Z\right\vert \leq\left\vert T\right\vert $. Combining
this with $\left\vert T\right\vert \leq r_{M}\left(  T\right)  $, we obtain
$r_{M}\left(  T\right)  =\left\vert T\right\vert $. This proves Lemma
\ref{lem.matroid.rI}.
\end{proof}

\begin{lemma}
\label{lem.matroid.Cin}Let $M=\left(  E,\mathcal{I}\right)  $ be a matroid.
Let $Q\in\mathcal{P}\left(  E\right)  \setminus\mathcal{I}$. Then, there
exists a circuit $C$ of $M$ such that $C\subseteq Q$.
\end{lemma}

\begin{proof}
[Proof of Lemma \ref{lem.matroid.Cin}.]We have $Q\in\mathcal{P}\left(
E\right)  \setminus\mathcal{I}$. Thus, there exists at least one
$C\in\mathcal{P}\left(  E\right)  \setminus\mathcal{I}$ such that $C\subseteq
Q$ (namely, $C=Q$). Thus, there also exists a \textbf{minimal} such $C$.
Consider this minimal $C$. We know that $C$ is a minimal element of
$\mathcal{P}\left(  E\right)  \setminus\mathcal{I}$ such that $C\subseteq Q$.
In other words, $C$ is an element of $\mathcal{P}\left(  E\right)
\setminus\mathcal{I}$ satisfying $C\subseteq Q$, and moreover,%
\begin{equation}
\text{every }D\in\mathcal{P}\left(  E\right)  \setminus\mathcal{I}\text{
satisfying }D\subseteq Q\text{ and }D\subseteq C\text{ must satisfy }D=C.
\label{pf.lem.matroid.Cin.1}%
\end{equation}
Thus, $C$ is a minimal element of $\mathcal{P}\left(  E\right)  \setminus
\mathcal{I}$\ \ \ \ \footnote{\textit{Proof.} We need to show that every
$D\in\mathcal{P}\left(  E\right)  \setminus\mathcal{I}$ satisfying $D\subseteq
C$ must satisfy $D=C$ (since we already know that $C\in\mathcal{P}\left(
E\right)  \setminus\mathcal{I}$).
\par
So let $D\in\mathcal{P}\left(  E\right)  \setminus\mathcal{I}$ be such that
$D\subseteq C$. Then, $D\subseteq C\subseteq Q$. Hence,
(\ref{pf.lem.matroid.Cin.1}) shows that $D=C$. This completes our proof.}. In
other words, $C$ is a circuit of $M$ (by the definition of a \textquotedblleft
circuit\textquotedblright). This circuit $C$ satisfies $C\subseteq Q$. Thus,
we have constructed a circuit $C$ of $M$ satisfying $C\subseteq Q$. Lemma
\ref{lem.matroid.Cin} is thus proven.
\end{proof}

\begin{lemma}
\label{lem.matroid.basext}Let $M=\left(  E,\mathcal{I}\right)  $ be a matroid.
Let $T$ be a subset of $E$. Let $S\in\mathcal{I}$ be such that $S\subseteq T$.
Then, there exists an $S^{\prime}\in\mathcal{I}$ satisfying $S\subseteq
S^{\prime}\subseteq T$ and $\left\vert S^{\prime}\right\vert =r_{M}\left(
T\right)  $.
\end{lemma}

\begin{proof}
[Proof of Lemma \ref{lem.matroid.basext}.]Clearly, there exists at least one
$S^{\prime}\in\mathcal{I}$ satisfying $S\subseteq S^{\prime}\subseteq T$
(namely, $S^{\prime}=S$). Hence, there exists a \textbf{maximal} such
$S^{\prime}$. Let $Q$ be such a maximal $S^{\prime}$. Thus, $Q$ is an element
of $\mathcal{I}$ satisfying $S\subseteq Q\subseteq T$.

Recall that
\begin{align*}
r_{M}\left(  T\right)   &  =\max\left\{  \left\vert Z\right\vert \ \mid
\ Z\in\mathcal{I}\text{ and }Z\subseteq T\right\}  \ \ \ \ \ \ \ \ \ \ \left(
\text{by the definition of }r_{M}\right) \\
&  \in\left\{  \left\vert Z\right\vert \ \mid\ Z\in\mathcal{I}\text{ and
}Z\subseteq T\right\}
\end{align*}
(since the maximum of any set must belong to this set). Hence, there exists
some $Z\in\mathcal{I}$ satisfying $Z\subseteq T$ and $r_{M}\left(  T\right)
=\left\vert Z\right\vert $. Denote such a $Z$ by $W$. Thus, $W$ is an element
of $\mathcal{I}$ satisfying $W\subseteq T$ and $r_{M}\left(  T\right)
=\left\vert W\right\vert $.

We have $\left\vert Q\right\vert \in\left\{  \left\vert Z\right\vert
\ \mid\ Z\in\mathcal{I}\text{ and }Z\subseteq T\right\}  $ (since
$Q\in\mathcal{I}$ and $Q\subseteq T$). Since any element of a set is smaller
or equal to the maximum of this set, this entails that $\left\vert
Q\right\vert \leq\max\left\{  \left\vert Z\right\vert \ \mid\ Z\in
\mathcal{I}\text{ and }Z\subseteq T\right\}  =r_{M}\left(  T\right)
=\left\vert W\right\vert $.

Now, assume (for the sake of contradiction) that $\left\vert Q\right\vert
\neq\left\vert W\right\vert $. Thus, $\left\vert Q\right\vert <\left\vert
W\right\vert $ (since $\left\vert Q\right\vert \leq\left\vert W\right\vert $).
Hence, Matroid axiom 3 (applied to $Y=Q$ and $Z=W$) shows that there exists
some $x\in W\setminus Q$ such that $Q\cup\left\{  x\right\}  \in\mathcal{I}$.
Consider this $x$. We have $x\in W\setminus Q\subseteq W\subseteq T$, so that
$Q\cup\left\{  x\right\}  \subseteq T$ (since $Q\subseteq T$). Also, $x\notin
Q$ (since $x\in W\setminus Q$).

Recall that $Q$ is a \textbf{maximal} $S^{\prime}\in\mathcal{I}$ satisfying
$S\subseteq S^{\prime}\subseteq T$. Thus, if some $S^{\prime}\in\mathcal{I}$
satisfies $S\subseteq S^{\prime}\subseteq T$ and $S^{\prime}\supseteq Q$, then
$S^{\prime}=Q$. Applying this to $S^{\prime}=Q\cup\left\{  x\right\}  $, we
obtain $Q\cup\left\{  x\right\}  =Q$ (since $S\subseteq Q\subseteq
Q\cup\left\{  x\right\}  \subseteq T$ and $Q\cup\left\{  x\right\}  \supseteq
Q$). Thus, $x\in Q$. But this contradicts $x\notin Q$. This contradiction
shows that our assumption (that $\left\vert Q\right\vert \neq\left\vert
W\right\vert $) was wrong. Hence, $\left\vert Q\right\vert =\left\vert
W\right\vert =r_{M}\left(  T\right)  $. Thus, there exists an $S^{\prime}%
\in\mathcal{I}$ satisfying $S\subseteq S^{\prime}\subseteq T$ and $\left\vert
S^{\prime}\right\vert =r_{M}\left(  T\right)  $ (namely, $S^{\prime}=Q$). This
proves Lemma \ref{lem.matroid.basext}.
\end{proof}

\subsection{The lattice of flats}

We shall now show a lemma that can be regarded as an alternative criterion for
a subset of $E$ to be a flat:

\begin{lemma}
\label{lem.matroid.flat-crit}Let $M=\left(  E,\mathcal{I}\right)  $ be a
matroid. Let $T$ be a subset of $E$. Then, the following statements are equivalent:

\textit{Statement }$\mathfrak{F}_{1}$\textit{:} The set $T$ is a flat of $M$.

\textit{Statement }$\mathfrak{F}_{2}$\textit{:} If $C$ is a circuit of $M$,
and if $e\in C$ is such that $C\setminus\left\{  e\right\}  \subseteq T$, then
$C\subseteq T$.
\end{lemma}

\begin{proof}
[Proof of Lemma \ref{lem.matroid.flat-crit}.]\textit{Proof of the implication
$\mathfrak{F}$}$_{1}\Longrightarrow\mathfrak{F}_{2}$\textit{:} Assume that
Statement $\mathfrak{F}_{1}$ holds. We must prove that Statement
$\mathfrak{F}_{2}$ holds.

Let $C$ be a circuit of $M$. Let $e\in C$ be such that $C\setminus\left\{
e\right\}  \subseteq T$. We must prove that $C\subseteq T$.

Assume the contrary. Thus, $C\not \subseteq T$. Combining this with
$C\setminus\left\{  e\right\}  \subseteq T$, we obtain $e\notin T$. Hence, $T$
is a proper subset of $T\cup\left\{  e\right\}  $.

We have assumed that Statement $\mathfrak{F}_{1}$ holds. In other words, the
set $T$ is a flat of $M$. In other words, there exists some $k\in\mathbb{N}$
such that $T$ is a $k$-flat of $M$. Consider this $k$.

The set $T$ is a $k$-flat of $M$, thus a subset of $E$ which has rank $k$ and
is maximal among all such subsets. In other words, $r_{M}\left(  T\right)
=k$, but every subset $S$ of $E$ for which $T$ is a proper subset of $S$ must
satisfy%
\begin{equation}
r_{M}\left(  S\right)  \neq k. \label{pf.lem.matroid.flat-crit.1.1}%
\end{equation}
Applying (\ref{pf.lem.matroid.flat-crit.1.1}) to $S=T\cup\left\{  e\right\}
$, we obtain $r_{M}\left(  T\cup\left\{  e\right\}  \right)  \neq k$. Since
$T\cup\left\{  e\right\}  \supseteq T$ (and since the function $r_{M}%
:\mathcal{P}\left(  E\right)  \rightarrow\mathbb{N}$ is weakly increasing), we
have $r_{M}\left(  T\cup\left\{  e\right\}  \right)  \geq r_{M}\left(
T\right)  =k$. Combined with $r_{M}\left(  T\cup\left\{  e\right\}  \right)
\neq k$, this yields $r_{M}\left(  T\cup\left\{  e\right\}  \right)  >k$.
Thus, $r_{M}\left(  T\cup\left\{  e\right\}  \right)  \geq k+1$.

Notice that $C\setminus\left\{  e\right\}  $ is a proper subset of $C$ (since
$e\in C$). The set $C$ is a circuit of $M$, thus a minimal element of
$\mathcal{P}\left(  E\right)  \setminus\mathcal{I}$ (by the definition of a
\textquotedblleft circuit\textquotedblright). Hence, no proper subset of $C$
belongs to $\mathcal{P}\left(  E\right)  \setminus\mathcal{I}$ (because $C$ is
minimal). In other words, every proper subset of $C$ belongs to $\mathcal{I}$.
Applying this to the proper subset $C\setminus\left\{  e\right\}  $ of $C$, we
conclude that $C\setminus\left\{  e\right\}  $ belongs to $\mathcal{I}$.
Hence, Lemma \ref{lem.matroid.basext} (applied to $S=C\setminus\left\{
e\right\}  $) shows that there exists an $S^{\prime}\in\mathcal{I}$ satisfying
$C\setminus\left\{  e\right\}  \subseteq S^{\prime}\subseteq T$ and
$\left\vert S^{\prime}\right\vert =r_{M}\left(  T\right)  $. Denote this
$S^{\prime}$ by $S$. Thus, $S$ is an element of $\mathcal{I}$ satisfying
$C\setminus\left\{  e\right\}  \subseteq S\subseteq T$ and $\left\vert
S\right\vert =r_{M}\left(  T\right)  $.

Furthermore, $S\subseteq T\subseteq T\cup\left\{  e\right\}  $. Thus, Lemma
\ref{lem.matroid.basext} (applied to $T\cup\left\{  e\right\}  $ instead of
$T$) shows that there exists an $S^{\prime}\in\mathcal{I}$ satisfying
$S\subseteq S^{\prime}\subseteq T\cup\left\{  e\right\}  $ and $\left\vert
S^{\prime}\right\vert =r_{M}\left(  T\cup\left\{  e\right\}  \right)  $.
Consider this $S^{\prime}$.

We have $\left\vert S^{\prime}\right\vert =r_{M}\left(  T\cup\left\{
e\right\}  \right)  >r_{M}\left(  T\right)  $. Hence, $S^{\prime
}\not \subseteq T$\ \ \ \ \footnote{\textit{Proof.} Assume the contrary. Thus,
$S^{\prime}\subseteq T$. Hence, $S^{\prime}$ is an element of $\mathcal{I}$
and satisfies $S^{\prime}\subseteq T$. Thus, $\left\vert S^{\prime}\right\vert
\in\left\{  \left\vert Z\right\vert \ \mid\ Z\in\mathcal{I}\text{ and
}Z\subseteq T\right\}  $.
\par
Now, the definition of $r_{M}$ yields%
\[
r_{M}\left(  T\right)  =\max\left\{  \left\vert Z\right\vert \ \mid
\ Z\in\mathcal{I}\text{ and }Z\subseteq T\right\}  \geq\left\vert S^{\prime
}\right\vert
\]
(since $\left\vert S^{\prime}\right\vert \in\left\{  \left\vert Z\right\vert
\ \mid\ Z\in\mathcal{I}\text{ and }Z\subseteq T\right\}  $). This contradicts
$\left\vert S^{\prime}\right\vert >r_{M}\left(  T\right)  $. This
contradiction proves that our assumption was wrong, qed.}. Combining this with
$S^{\prime}\subseteq T\cup\left\{  e\right\}  $, we obtain $e\in S^{\prime}$.
Combining this with $C\setminus\left\{  e\right\}  \subseteq S^{\prime}$, we
find that $\left(  C\setminus\left\{  e\right\}  \right)  \cup\left\{
e\right\}  \subseteq S^{\prime}$. Thus, $C=\left(  C\setminus\left\{
e\right\}  \right)  \cup\left\{  e\right\}  \subseteq S^{\prime}$. Since
$S^{\prime}\in\mathcal{I}$, this entails that $C\in\mathcal{I}$ (by Matroid
axiom 2). But $C\in\mathcal{P}\left(  E\right)  \setminus\mathcal{I}$ (since
$C$ is a minimal element of $\mathcal{P}\left(  E\right)  \setminus
\mathcal{I}$), so that $C\notin\mathcal{I}$. This contradicts $C\in
\mathcal{I}$. This contradiction shows that our assumption was wrong. Hence,
$C\subseteq T$ is proven. Therefore, Statement $\mathfrak{F}_{2}$ holds. Thus,
the implication $\mathfrak{F}_{1}\Longrightarrow\mathfrak{F}_{2}$ is proven.

\textit{Proof of the implication $\mathfrak{F}$}$_{2}\Longrightarrow
\mathfrak{F}_{1}$\textit{:} Assume that Statement $\mathfrak{F}_{2}$ holds. We
must prove that Statement $\mathfrak{F}_{1}$ holds.

Let $k=r_{M}\left(  T\right)  $. We shall show that $T$ is a $k$-flat of $M$.

Let $W$ be a subset of $E$ which has rank $k$ and satisfies $T\subseteq W$. We
shall show that $T=W$.

Indeed, assume the contrary. Thus, $T\neq W$. Combined with $T\subseteq W$,
this shows that $T$ is a proper subset of $W$. Thus, there exists an $e\in
W\setminus T$. Consider this $e$. We have $e\notin T$ (since $e\in W\setminus
T$).

We have
\begin{align*}
k  &  =r_{M}\left(  T\right)  =\max\left\{  \left\vert Z\right\vert
\ \mid\ Z\in\mathcal{I}\text{ and }Z\subseteq T\right\}
\ \ \ \ \ \ \ \ \ \ \left(  \text{by the definition of }r_{M}\right) \\
&  \in\left\{  \left\vert Z\right\vert \ \mid\ Z\in\mathcal{I}\text{ and
}Z\subseteq T\right\}
\end{align*}
(since the maximum of a set must belong to that set). Hence, there exists some
$Z\in\mathcal{I}$ satisfying $Z\subseteq T$ and $k=\left\vert Z\right\vert $.
Denote this $Z$ by $K$. Thus, $K$ is an element of $\mathcal{I}$ and satisfies
$K\subseteq T$ and $k=\left\vert K\right\vert $. Notice that $e\notin T$, so
that $e\notin K$ (since $K\subseteq T$).

We have $r_{M}\left(  W\right)  =k$ (since $W$ has rank $k$). Hence,
$K\cup\left\{  e\right\}  \notin\mathcal{I}$\ \ \ \ \footnote{\textit{Proof.}
Assume the contrary. Thus, $K\cup\left\{  e\right\}  \in\mathcal{I}$. Thus,
$r_{M}\left(  K\cup\left\{  e\right\}  \right)  =\left\vert K\cup\left\{
e\right\}  \right\vert $ (by Lemma \ref{lem.matroid.rI}). Thus, $r_{M}\left(
K\cup\left\{  e\right\}  \right)  =\left\vert K\cup\left\{  e\right\}
\right\vert >\left\vert K\right\vert $ (since $e\notin K$).
\par
But $K\cup\left\{  e\right\}  \subseteq W$ (since $K\subseteq T\subseteq W$
and $e\in W\setminus T\subseteq W$). Since the function $r_{M}$ is weakly
increasing, this yields $r_{M}\left(  K\cup\left\{  e\right\}  \right)  \leq
r_{M}\left(  W\right)  =k=\left\vert K\right\vert $. This contradicts
$r_{M}\left(  K\cup\left\{  e\right\}  \right)  >\left\vert K\right\vert $.
This contradiction proves that our assumption was wrong, qed.}. In other
words, $K\cup\left\{  e\right\}  \in\mathcal{P}\left(  E\right)
\setminus\mathcal{I}$. Hence, Lemma \ref{lem.matroid.Cin} (applied to
$Q=K\cup\left\{  e\right\}  $) shows that there exists a circuit $C$ of $M$
such that $C\subseteq K\cup\left\{  e\right\}  $. Consider this $C$. From
$C\subseteq K\cup\left\{  e\right\}  $, we obtain $C\setminus\left\{
e\right\}  \subseteq K\subseteq T$.

From $C\setminus\left\{  e\right\}  \subseteq K$, we conclude (using Matroid
axiom 2) that $C\setminus\left\{  e\right\}  \in\mathcal{I}$ (since
$K\in\mathcal{I}$). On the other hand, $C$ is a circuit of $M$. In other
words, $C$ is a minimal element of $\mathcal{P}\left(  E\right)
\setminus\mathcal{I}$ (by the definition of a \textquotedblleft
circuit\textquotedblright). Hence, $C\in\mathcal{P}\left(  E\right)
\setminus\mathcal{I}$, so that $C\notin\mathcal{I}$. Hence, $e\in C$ (since
otherwise, we would have $C\setminus\left\{  e\right\}  =C\notin\mathcal{I}$,
which would contradict $C\setminus\left\{  e\right\}  \in\mathcal{I}$). Now,
Statement $\mathfrak{F}_{2}$ shows that $C\subseteq T$. Hence, $e\in
C\subseteq T$, which contradicts $e\notin T$.

This contradiction shows that our assumption was wrong. Hence, $T=W$ is proven.

Now, forget that we fixed $W$. Thus, we have shown that if $W$ is a subset of
$E$ which has rank $k$ and satisfies $T\subseteq W$, then $T=W$. In other
words, $T$ is a subset of $E$ which has rank $k$ and is maximal among all such
subsets (because we already know that $T$ has rank $r_{M}\left(  T\right)
=k$). In other words, $T$ is a $k$-flat of $M$ (by the definition of a
\textquotedblleft$k$-flat\textquotedblright). Thus, $T$ is a flat of $M$. In
other words, Statement $\mathfrak{F}_{1}$ holds. This proves the implication
$\mathfrak{F}_{2}\Longrightarrow\mathfrak{F}_{1}$.

We have now proven the implications $\mathfrak{F}_{1}\Longrightarrow
\mathfrak{F}_{2}$ and $\mathfrak{F}_{2}\Longrightarrow\mathfrak{F}_{1}$.
Together, these implications show that Statements $\mathfrak{F}_{1}$ and
$\mathfrak{F}_{2}$ are equivalent. This proves Lemma
\ref{lem.matroid.flat-crit}.
\end{proof}

\begin{corollary}
\label{cor.matroid.flat-cap}Let $M=\left(  E,\mathcal{I}\right)  $ be a
matroid. Let $F_{1},F_{2},\ldots,F_{k}$ be flats of $M$. Then, $F_{1}\cap
F_{2}\cap\cdots\cap F_{k}$ is a flat of $M$. (Notice that $k$ is allowed to be
$0$ here; in this case, the empty intersection $F_{1}\cap F_{2}\cap\cdots\cap
F_{k}$ is to be interpreted as $E$.)
\end{corollary}

\begin{proof}
[Proof of Corollary \ref{cor.matroid.flat-cap}.]Lemma
\ref{lem.matroid.flat-crit} gives a necessary and sufficient criterion for a
subset $T$ of $E$ to be a flat of $M$. It is easy to see that if this
criterion is satisfied for $T=F_{1}$, for $T=F_{2}$, etc., and for $T=F_{k}$,
then it is satisfied for $T=F_{1}\cap F_{2}\cap\cdots\cap F_{k}$. In other
words, if $F_{1},F_{2},\ldots,F_{k}$ are flats of $M$, then $F_{1}\cap
F_{2}\cap\cdots\cap F_{k}$ is a flat of $M$.\ \ \ \ \footnote{Here is this
argument in slightly more detail:
\par
For every $i\in\left\{  1,2,\ldots,k\right\}  $, the following statement
holds: If $C$ is a circuit of $M$, and if $e\in C$ is such that $C\setminus
\left\{  e\right\}  \subseteq F_{i}$, then%
\begin{equation}
C\subseteq F_{i}. \label{pf.cor.matroid.flat-cap.fn1.1}%
\end{equation}
\par
\textit{Proof of (\ref{pf.cor.matroid.flat-cap.fn1.1}):} Let $i\in\left\{
1,2,\ldots,k\right\}  $. Then, the set $F_{i}$ is a flat of $M$. In other
words, Statement $\mathfrak{F}_{1}$ of Lemma \ref{lem.matroid.flat-crit} is
satisfied for $T=F_{i}$. Therefore, Statement $\mathfrak{F}_{2}$ of Lemma
\ref{lem.matroid.flat-crit} must also be satisfied for $T=F_{i}$ (since Lemma
\ref{lem.matroid.flat-crit} shows that the Statements $\mathfrak{F}_{1}$ and
$\mathfrak{F}_{2}$ are equivalent). In other words, if $C$ is a circuit of
$M$, and if $e\in C$ is such that $C\setminus\left\{  e\right\}  \subseteq
F_{i}$, then $C\subseteq F_{i}$. This proves
(\ref{pf.cor.matroid.flat-cap.fn1.1}).
\par
Now, let $C$ be a circuit of $M$, and let $e\in C$ be such that $C\setminus
\left\{  e\right\}  \subseteq F_{1}\cap F_{2}\cap\cdots\cap F_{k}$. For every
$i\in\left\{  1,2,\ldots,k\right\}  $, we have $C\setminus\left\{  e\right\}
\subseteq F_{1}\cap F_{2}\cap\cdots\cap F_{k}\subseteq F_{i}$, and therefore
$C\subseteq F_{i}$ (by (\ref{pf.cor.matroid.flat-cap.fn1.1})). So we have
shown the inclusion $C\subseteq F_{i}$ for each $i\in\left\{  1,2,\ldots
,k\right\}  $. Combining these $k$ inclusions, we obtain $C\subseteq F_{1}\cap
F_{2}\cap\cdots\cap F_{k}$.
\par
Now, forget that we fixed $C$. We thus have shown that if $C$ is a circuit of
$M$, and if $e\in C$ is such that $C\setminus\left\{  e\right\}  \subseteq
F_{1}\cap F_{2}\cap\cdots\cap F_{k}$, then $C\subseteq F_{1}\cap F_{2}%
\cap\cdots\cap F_{k}$. In other words, Statement $\mathfrak{F}_{2}$ of Lemma
\ref{lem.matroid.flat-crit} is satisfied for $T=F_{1}\cap F_{2}\cap\cdots\cap
F_{k}$. Therefore, Statement $\mathfrak{F}_{1}$ of Lemma
\ref{lem.matroid.flat-crit} must also be satisfied for $T=F_{1}\cap F_{2}%
\cap\cdots\cap F_{k}$ (since Lemma \ref{lem.matroid.flat-crit} shows that the
Statements $\mathfrak{F}_{1}$ and $\mathfrak{F}_{2}$ are equivalent). In other
words, the set $F_{1}\cap F_{2}\cap\cdots\cap F_{k}$ is a flat of $M$. Qed.}
This proves Corollary \ref{cor.matroid.flat-cap}.
\end{proof}

Corollary \ref{cor.matroid.flat-cap} (a well-known fact, which is left to the
reader to prove in \cite[\S 3.1]{Stanley}) allows us to define the
\textit{closure} of a set in a matroid:

\begin{definition}
\label{def.matroid.closure}Let $M=\left(  E,\mathcal{I}\right)  $ be a
matroid. Let $T$ be a subset of $E$. The \textit{closure} of $T$ is defined to
be the intersection of all flats of $M$ which contain $T$ as a subset. In
other words, the closure of $T$ is defined to be $\bigcap_{\substack{F\in
\operatorname*{Flats}M;\\T\subseteq F}}F$. The closure of $T$ is denoted by
$\overline{T}$.
\end{definition}

The following proposition gathers some simple properties of closures in matroids:

\begin{proposition}
\label{prop.matroid.closure.props}Let $M=\left(  E,\mathcal{I}\right)  $ be a matroid.

\textbf{(a)} If $T$ is a subset of $E$, then $\overline{T}$ is a flat of $M$
satisfying $T\subseteq\overline{T}$.

\textbf{(b)} If $G$ is a flat of $M$, then $\overline{G}=G$.

\textbf{(c)} If $T$ is a subset of $E$ and if $G$ is a flat of $M$ satisfying
$T\subseteq G$, then $\overline{T}\subseteq G$.

\textbf{(d)} If $S$ and $T$ are two subsets of $E$ satisfying $S\subseteq T$,
then $\overline{S}\subseteq\overline{T}$.

\textbf{(e)} If the matroid $M$ is loopless, then $\overline{\varnothing
}=\varnothing$.

\textbf{(f)} Every subset $T$ of $E$ satisfies $r_{M}\left(  T\right)
=r_{M}\left(  \overline{T}\right)  $.

\textbf{(g)} If $T$ is a subset of $E$ and if $G$ is a flat of $M$, then the
conditions $\left(  \overline{T}\subseteq G\right)  $ and $\left(  T\subseteq
G\right)  $ are equivalent.
\end{proposition}

\begin{proof}
[Proof of Proposition \ref{prop.matroid.closure.props}.]\textbf{(a)} The set
$\operatorname*{Flats}M$ is a subset of the finite set $\mathcal{P}\left(
E\right)  $, and thus itself finite.

Let $T$ be a subset of $E$. The closure $\overline{T}$ of $T$ is defined as
$\bigcap_{\substack{F\in\operatorname*{Flats}M;\\T\subseteq F}}F$. Now,
Corollary \ref{cor.matroid.flat-cap} shows that any intersection of finitely
many flats of $M$ is a flat of $M$. Hence, $\bigcap_{\substack{F\in
\operatorname*{Flats}M;\\T\subseteq F}}F$ (being an intersection of finitely
many flats of $M$\ \ \ \ \footnote{\textquotedblleft Finitely
many\textquotedblright\ since the set $\operatorname*{Flats}M$ is finite.}) is
a flat of $M$. In other words, $\overline{T}$ is a flat of $M$ (since
$\overline{T}=\bigcap_{\substack{F\in\operatorname*{Flats}M;\\T\subseteq F}}F$).

Also, $T\subseteq F$ for every $F\in\operatorname*{Flats}M$ satisfying
$T\subseteq F$. Hence, $T\subseteq\bigcap_{\substack{F\in\operatorname*{Flats}%
M;\\T\subseteq F}}F=\overline{T}$. This completes the proof of Proposition
\ref{prop.matroid.closure.props} \textbf{(a)}.

\textbf{(c)} Let $T$ be a subset of $E$, and let $G$ be a flat of $M$
satisfying $T\subseteq G$. Then, $G$ is an element of $\operatorname*{Flats}G$
satisfying $T\subseteq G$. Hence, $G$ is one term in the intersection
$\bigcap_{\substack{F\in\operatorname*{Flats}M;\\T\subseteq F}}F$. Thus,
$\bigcap_{\substack{F\in\operatorname*{Flats}M;\\T\subseteq F}}F\subseteq G$.
But the definition of $\overline{T}$ yields $\overline{T}=\bigcap
_{\substack{F\in\operatorname*{Flats}M;\\T\subseteq F}}F\subseteq G$. This
proves Proposition \ref{prop.matroid.closure.props} \textbf{(c)}.

\textbf{(b)} Let $G$ be a flat of $M$. Proposition
\ref{prop.matroid.closure.props} \textbf{(b)} (applied to $T=G$) yields
$\overline{G}\subseteq G$. But Proposition \ref{prop.matroid.closure.props}
\textbf{(a)} (applied to $T=G$) shows that $\overline{G}$ is a flat of $M$
satisfying $G\subseteq\overline{G}$. Combining $G\subseteq\overline{G}$ with
$\overline{G}\subseteq G$, we obtain $\overline{G}=G$. This proves Proposition
\ref{prop.matroid.closure.props} \textbf{(b)}.

\textbf{(d)} Let $S$ and $T$ be two subsets of $E$ satisfying $S\subseteq T$.
Proposition \ref{prop.matroid.closure.props} \textbf{(a)} shows that
$\overline{T}$ is a flat of $M$ satisfying $T\subseteq\overline{T}$. Now,
$S\subseteq T\subseteq\overline{T}$. Hence, Proposition
\ref{prop.matroid.closure.props} \textbf{(b)} (applied to $S$ and
$\overline{T}$ instead of $T$ and $G$) shows $\overline{S}\subseteq
\overline{T}$. This proves Proposition \ref{prop.matroid.closure.props}
\textbf{(d)}.

\textbf{(e)} Assume that the matroid $M$ is loopless. In other words, no loops
(of $M$) exist.

The definition of $r_{M}$ quickly yields $r_{M}\left(  \varnothing\right)
=0$. In other words, the set $\varnothing$ has rank $0$. We shall now show
that $\varnothing$ is a $0$-flat of $M$.

Indeed, let $W$ be a subset of $E$ which has rank $0$ and satisfies
$\varnothing\subseteq W$. We shall show that $\varnothing=W$.

Assume the contrary. Thus, $\varnothing\neq W$. Hence, $W$ has an element $w$.
Consider this $w$. The element $w$ of $E$ is not a loop (since no loops
exist). In other words, we cannot have $\left\{  w\right\}  \notin\mathcal{I}$
(since $w$ is a loop if and only if $\left\{  w\right\}  \notin\mathcal{I}$
(by the definition of a loop)). In other words, we must have $\left\{
w\right\}  \in\mathcal{I}$. Clearly, $\left\{  w\right\}  \subseteq W$ (since
$w\in W$). Thus, $\left\{  w\right\}  $ is a $Z\in\mathcal{I}$ satisfying
$Z\subseteq W$. Thus, $\left\vert \left\{  w\right\}  \right\vert \in\left\{
\left\vert Z\right\vert \ \mid\ Z\in\mathcal{I}\text{ and }Z\subseteq
W\right\}  $.

But $W$ has rank $0$. In other words,%
\begin{align*}
0  &  =r_{M}\left(  W\right)  =\max\left\{  \left\vert Z\right\vert
\ \mid\ Z\in\mathcal{I}\text{ and }Z\subseteq W\right\}
\ \ \ \ \ \ \ \ \ \ \left(  \text{by the definition of }r_{M}\right) \\
&  \geq\left\vert \left\{  w\right\}  \right\vert \ \ \ \ \ \ \ \ \ \ \left(
\text{since }\left\vert \left\{  w\right\}  \right\vert \in\left\{  \left\vert
Z\right\vert \ \mid\ Z\in\mathcal{I}\text{ and }Z\subseteq W\right\}  \right)
\\
&  =1,
\end{align*}
which is absurd. This contradiction shows that our assumption was wrong.
Hence, $\varnothing=W$ is proven.

Let us now forget that we fixed $W$. We thus have proven that if $W$ is any
subset of $E$ which has rank $0$ and satisfies $\varnothing\subseteq W$, then
$\varnothing=W$. Thus, $\varnothing$ is a subset of $E$ which has rank $0$ and
is maximal among all such subsets (because we already know that $\varnothing$
has rank $0$). In other words, $\varnothing$ is a $0$-flat of $M$ (by the
definition of a \textquotedblleft$0$-flat\textquotedblright). Thus,
$\varnothing$ is a flat of $M$. Thus, Proposition
\ref{prop.matroid.closure.props} \textbf{(b)} (applied to $G=\varnothing$)
yields $\overline{\varnothing}=\varnothing$. This proves Proposition
\ref{prop.matroid.closure.props} \textbf{(e)}.

\textbf{(f)} Let $T$ be a subset of $E$. We have $T\subseteq\overline{T}$ (by
Proposition \ref{prop.matroid.closure.props} \textbf{(a)}), and thus
$r_{M}\left(  T\right)  \leq r_{M}\left(  \overline{T}\right)  $ (since the
function $r_{M}$ is weakly increasing).

Let $k=r_{M}\left(  T\right)  $. Thus, there exists a $Q\in\mathcal{P}\left(
E\right)  $ satisfying $T\subseteq Q$ and $k=r_{M}\left(  Q\right)  $ (namely,
$Q=T$). Hence, there exists a \textbf{maximal} such $Q$. Denote this $Q$ by
$R$. Thus, $R$ is a maximal $Q\in\mathcal{P}\left(  E\right)  $ satisfying
$T\subseteq Q$ and $k=r_{M}\left(  Q\right)  $. In particular, $R$ is an
element of $\mathcal{P}\left(  E\right)  $ and satisfies $T\subseteq R$ and
$k=r_{M}\left(  R\right)  $.

Now, $R$ is a subset of $E$ (since $R\in\mathcal{P}\left(  E\right)  $) and
has rank $r_{M}\left(  R\right)  =k$. Thus, $R$ is a subset of $E$ which has
rank $k$. Furthermore, $R$ is maximal among all such
subsets\footnote{\textit{Proof.} Let $W$ be any subset of $E$ which has rank
$k$ and satisfies $W\supseteq R$. We must prove that $W=R$.
\par
We have $W\in\mathcal{P}\left(  E\right)  $, $T\subseteq R\subseteq W$ and
$k=r_{M}\left(  W\right)  $ (since $W$ has rank $k$). Thus, $W$ is a
$Q\in\mathcal{P}\left(  E\right)  $ satisfying $T\subseteq Q$ and
$k=r_{M}\left(  Q\right)  $. But recall that $R$ is a \textbf{maximal} such
$Q$. Hence, if $W\supseteq R$, then $W=R$. Therefore, $W=R$ (since we know
that $W\supseteq R$). Qed.}. Thus, $R$ is a $k$-flat of $M$ (by the definition
of a \textquotedblleft$k$-flat\textquotedblright), and therefore a flat of
$M$. Now, Proposition \ref{prop.matroid.closure.props} \textbf{(c)} (applied
to $G=R$) shows that $\overline{T}\subseteq R$. Since the function $r_{M}$ is
weakly increasing, this yields $r_{M}\left(  \overline{T}\right)  \leq
r_{M}\left(  R\right)  =k$. Combining this with $k=r_{M}\left(  T\right)  \leq
r_{M}\left(  \overline{T}\right)  $, we obtain $r_{M}\left(  \overline
{T}\right)  =k=r_{M}\left(  T\right)  $. This proves Proposition
\ref{prop.matroid.closure.props} \textbf{(f)}.

\textbf{(g)} Let $T$ be a subset of $E$. Let $G$ be a flat of $M$. Proposition
\ref{prop.matroid.closure.props} \textbf{(a)} shows that $T\subseteq
\overline{T}$. Hence, if $\overline{T}\subseteq G$, then $T\subseteq
\overline{T}\subseteq G$. Thus, we have proven the implication $\left(
\overline{T}\subseteq G\right)  \Longrightarrow\left(  T\subseteq G\right)  $.
The reverse implication (i.e., the implication $\left(  T\subseteq G\right)
\Longrightarrow\left(  \overline{T}\subseteq G\right)  $) follows from
Proposition \ref{prop.matroid.closure.props} \textbf{(c)}. Combining these two
implications, we obtain the equivalence $\left(  \overline{T}\subseteq
G\right)  \Longleftrightarrow\left(  T\subseteq G\right)  $. This proves
Proposition \ref{prop.matroid.closure.props} \textbf{(g)}.
\end{proof}

We shall now recall a few more classical notions related to posets:

\begin{definition}
\label{def.lattice}Let $P$ be a poset.

\textbf{(a)} An element $p\in P$ is said to be a \textit{global minimum} of
$P$ if every $q\in P$ satisfies $p\leq q$. Clearly, a global minimum of $P$ is
unique if it exists.

\textbf{(b)} An element $p\in P$ is said to be a \textit{global maximum} of
$P$ if every $q\in P$ satisfies $p\geq q$. Clearly, a global maximum of $P$ is
unique if it exists.

\textbf{(c)} Let $x$ and $y$ be two elements of $P$. An \textit{upper bound}
of $x$ and $y$ (in $P$) means an element $z\in P$ satisfying $z\geq x$ and
$z\geq y$. A \textit{join} (or \textit{least upper bound}) of $x$ and $y$ (in
$P$) means an upper bound $z$ of $x$ and $y$ such that every upper bound
$z^{\prime}$ of $x$ and $y$ satisfies $z^{\prime}\geq z$. In other words, a
join of $x$ and $y$ is a global minimum of the subposet $\left\{  w\in
P\ \mid\ w\geq x\text{ and }w\geq y\right\}  $ of $P$. Thus, a join of $x$ and
$y$ is unique if it exists.

\textbf{(d)} Let $x$ and $y$ be two elements of $P$. A \textit{lower bound} of
$x$ and $y$ (in $P$) means an element $z\in P$ satisfying $z\leq x$ and $z\leq
y$. A \textit{meet} (or \textit{greatest lower bound}) of $x$ and $y$ (in $P$)
means a lower bound $z$ of $x$ and $y$ such that every lower bound $z^{\prime
}$ of $x$ and $y$ satisfies $z^{\prime}\leq z$. In other words, a meet of $x$
and $y$ is a global maximum of the subposet $\left\{  w\in P\ \mid\ w\leq
x\text{ and }w\leq y\right\}  $ of $P$. Thus, a meet of $x$ and $y$ is unique
if it exists.

\textbf{(e)} The poset $P$ is said to be a \textit{lattice} if and only if it
has a global minimum and a global maximum, and every two elements of $P$ have
a meet and a join.
\end{definition}

\begin{proposition}
\label{prop.matroid.flat-lat}Let $M=\left(  E,\mathcal{I}\right)  $ be a
matroid. The subposet $\operatorname*{Flats}M$ of the poset $\mathcal{P}%
\left(  E\right)  $ is a lattice.
\end{proposition}

\begin{proof}
[Proof of Proposition \ref{prop.matroid.flat-lat}.]By the definition of a
lattice, it suffices to check the following four claims:

\textit{Claim 1:} The poset $\operatorname*{Flats}M$ has a global minimum.

\textit{Claim 2:} The poset $\operatorname*{Flats}M$ has a global maximum.

\textit{Claim 3:} Every two elements of $\operatorname*{Flats}M$ have a meet
(in $\operatorname*{Flats}M$).

\textit{Claim 4:} Every two elements of $\operatorname*{Flats}M$ have a join
(in $\operatorname*{Flats}M$).

\textit{Proof of Claim 1:} Applying Proposition
\ref{prop.matroid.closure.props} \textbf{(a)} to $T=\varnothing$, we see that
$\overline{\varnothing}$ is a flat of $M$ satisfying $\varnothing
\subseteq\overline{\varnothing}$. In particular, $\overline{\varnothing}$ is a
flat of $M$, so that $\overline{\varnothing}\in\operatorname*{Flats}M$. If $G$
is a flat of $M$, then $\overline{\varnothing}\subseteq G$ (by Proposition
\ref{prop.matroid.closure.props} \textbf{(c)}, applied to $T=\varnothing$).
Hence, $\overline{\varnothing}$ is a global minimum of the poset
$\operatorname*{Flats}M$. Thus, the poset $\operatorname*{Flats}M$ has a
global minimum. This proves Claim 1.

\textit{Proof of Claim 2:} Applying Proposition
\ref{prop.matroid.closure.props} \textbf{(a)} to $T=E$, we see that
$\overline{E}$ is a flat of $M$ satisfying $E\subseteq\overline{E}$. From
$E\subseteq\overline{E}$, we conclude that $\overline{E}=E$. Thus, $E$ is a
flat of $M$ (since $\overline{E}$ is a flat of $M$). In other words,
$E\in\operatorname*{Flats}M$. If $G$ is a flat of $M$, then $E\supseteq G$
(obviously). Hence, $E$ is a global maximum of the poset
$\operatorname*{Flats}M$. Thus, the poset $\operatorname*{Flats}M$ has a
global maximum. This proves Claim 2.

\textit{Proof of Claim 3:} Let $F$ and $G$ be two elements of
$\operatorname*{Flats}M$. We have to prove that $F$ and $G$ have a meet.

We know that $F$ and $G$ are elements of $\operatorname*{Flats}M$, thus flats
of $M$. Hence, Corollary \ref{cor.matroid.flat-cap} shows that $F\cap G$ is a
flat of $M$. In other words, $F\cap G\in\operatorname*{Flats}M$. Clearly,
$F\cap G\subseteq F$ and $F\cap G\subseteq G$; thus, $F\cap G$ is a lower
bound of $F$ and $G$ in $\operatorname*{Flats}M$. Also, every lower bound $H$
of $F$ and $G$ in $\operatorname*{Flats}M$ satisfies $H\subseteq F\cap
G$\ \ \ \ \footnote{\textit{Proof.} Let $H$ be a lower bound of $F$ and $G$ in
$\operatorname*{Flats}M$. Thus, $H\subseteq F$ and $H\subseteq G$. Combining
these two inclusions, we obtain $H\subseteq F\cap G$, qed.}. Hence, $F\cap G$
is a meet of $F$ and $G$. Thus, $F$ and $G$ have a meet. This proves Claim 3.

\textit{Proof of Claim 4:} Let $F$ and $G$ be two elements of
$\operatorname*{Flats}M$. We have to prove that $F$ and $G$ have a join.

We know that $F$ and $G$ are elements of $\operatorname*{Flats}M$, thus flats
of $M$. Proposition \ref{prop.matroid.closure.props} \textbf{(a)} (applied to
$T=F\cup G$) shows that $\overline{F\cup G}$ is a flat of $M$ satisfying
$F\cup G\subseteq\overline{F\cup G}$. Now, $\overline{F\cup G}\in
\operatorname*{Flats}M$ (since $\overline{F\cup G}$ is a flat of $M$).
Clearly, $F\subseteq F\cup G\subseteq\overline{F\cup G}$ and $G\subseteq F\cup
G\subseteq\overline{F\cup G}$; thus, $\overline{F\cup G}$ is an upper bound of
$F$ and $G$ in $\operatorname*{Flats}M$. Also, every upper bound $H$ of $F$
and $G$ in $\operatorname*{Flats}M$ satisfies $H\supseteq\overline{F\cup G}%
$\ \ \ \ \footnote{\textit{Proof.} Let $H$ be an upper bound of $F$ and $G$ in
$\operatorname*{Flats}M$. Thus, $H\supseteq F$ and $H\supseteq G$. Combining
these two inclusions, we obtain $H\supseteq F\cup G$. But $H\in
\operatorname*{Flats}M$; thus, $H$ is a flat of $M$. Since $H$ satisfies
$F\cup G\subseteq H$, we therefore obtain $\overline{F\cup G}\subseteq H$ (by
Proposition \ref{prop.matroid.closure.props} \textbf{(c)}, applied to $F\cup
G$ and $H$ instead of $T$ and $G$). In other words, $H\supseteq\overline{F\cup
G}$, qed.}. Hence, $\overline{F\cup G}$ is a join of $F$ and $G$. Thus, $F$
and $G$ have a join. This proves Claim 4.

We have now proven all four Claims 1, 2, 3, and 4. Thus, Proposition
\ref{prop.matroid.flat-lat} is proven.
\end{proof}

\begin{definition}
\label{def.matroid.flat-lat}Let $M=\left(  E,\mathcal{I}\right)  $ be a
matroid. Proposition \ref{prop.matroid.flat-lat} shows that the subposet
$\operatorname*{Flats}M$ of the poset $\mathcal{P}\left(  E\right)  $ is a
lattice. This subposet $\operatorname*{Flats}M$ is called the \textit{lattice
of flats} of $M$. (Beware: It is a subposet, but not a sublattice of
$\mathcal{P}\left(  E\right)  $, since its join is not a restriction of the
join of $\mathcal{P}\left(  E\right)  $.)
\end{definition}

The lattice of flats $\operatorname*{Flats}M$ of a matroid $M$ is denoted by
$L\left(  M\right)  $ in \cite[\S 3.2]{Stanley}.

Next, we recall the definition of the M\"{o}bius function of a poset (see,
e.g., \cite[Definition 1.2]{Stanley} or \cite[\S 5.2]{Martin15}):

\begin{definition}
\label{def.moebius}Let $P$ be a poset.

\textbf{(a)} If $x$ and $y$ are two elements of $P$ satisfying $x\leq y$, then
the set $\left\{  z\in P\ \mid\ x\leq z\leq y\right\}  $ is denoted by
$\left[  x,y\right]  $.

\textbf{(b)} A subset of $P$ is called a \textit{closed interval} of $P$ if it
has the form $\left[  x,y\right]  $ for two elements $x$ and $y$ of $P$
satisfying $x\leq y$.

\textbf{(c)} We denote by $\operatorname*{Int}P$ the set of all closed
intervals of $P$.

\textbf{(d)} If $f:\operatorname*{Int}P\rightarrow\mathbb{Z}$ is any map, then
the image $f\left(  \left[  x,y\right]  \right)  $ of a closed interval
$\left[  x,y\right]  \in\operatorname*{Int}P$ under $f$ will be abbreviated by
$f\left(  x,y\right)  $.

\textbf{(e)} Assume that every closed interval of $P$ is finite. The
\textit{M\"{o}bius function} of the poset $P$ is defined to be the unique
function $\mu:\operatorname*{Int}P\rightarrow\mathbb{Z}$ having the following
two properties:

\begin{itemize}
\item We have
\begin{equation}
\mu\left(  x,x\right)  =1\ \ \ \ \ \ \ \ \ \ \text{for every }x\in P.
\label{eq.def.moebius.rec1}%
\end{equation}


\item We have%
\begin{equation}
\mu\left(  x,y\right)  =-\sum_{\substack{z\in P;\\x\leq z<y}}\mu\left(
x,z\right)  \ \ \ \ \ \ \ \ \ \ \text{for all }x,y\in P\text{ satisfying }x<y.
\label{eq.def.moebius.rec2}%
\end{equation}

\end{itemize}

(It is easy to see that these two properties indeed determine $\mu$ uniquely.)
This M\"{o}bius function is denoted by $\mu$.
\end{definition}

We can now define the characteristic polynomial of a matroid $M$, following
\cite[(22)]{Stanley}\footnote{Our notation slightly differs from that in
\cite[(22)]{Stanley}. Namely, we use $x$ as the indeterminate, while Stanley
instead uses $t$. Stanley also denotes the global minimum $\overline
{\varnothing}$ of $\operatorname*{Flats}M$ by $\widehat{0}$.}:

\begin{definition}
\label{def.matroid.charpol}Let $M=\left(  E,\mathcal{I}\right)  $ be a
matroid. Let $m=r_{M}\left(  E\right)  $. The \textit{characteristic
polynomial} $\chi_{M}$ of the matroid $M$ is defined to be the polynomial%
\[
\sum_{F\in\operatorname*{Flats}M}\mu\left(  \overline{\varnothing},F\right)
x^{m-r_{M}\left(  F\right)  }\in\mathbb{Z}\left[  x\right]
\]
(where $\mu$ is the M\"{o}bius function of the lattice $\operatorname*{Flats}%
M$). We further define a polynomial $\widetilde{\chi}_{M}\in\mathbb{Z}\left[
x\right]  $ by $\widetilde{\chi}_{M}=\left[  \overline{\varnothing
}=\varnothing\right]  \chi_{M}$. Here, we are using the Iverson bracket
notation (as in Definition \ref{def.iverson}). If the matroid $M$ is loopless,
then%
\[
\widetilde{\chi}_{M}=\underbrace{\left[  \overline{\varnothing}=\varnothing
\right]  }_{\substack{=1\\\text{(by Proposition
\ref{prop.matroid.closure.props} \textbf{(e)})}}}\chi_{M}=\chi_{M}.
\]

\end{definition}

\begin{example}
\label{exam.matroid.charpol.graph}Let $G=\left(  V,E\right)  $ be a finite
graph. Consider the graphical matroid $\left(  E,\mathcal{I}\right)  $ defined
as in Example \ref{exam.matroid.graphical}. Then, the characteristic
polynomial $\chi_{\left(  E,\mathcal{I}\right)  }$ of this matroid is
connected to the chromatic polynomial $\chi_{G}$ of the graph $G$ as follows:%
\[
x^{\operatorname*{conn}G}\cdot\chi_{\left(  E,\mathcal{I}\right)  }\left(
x\right)  =\chi_{G}\left(  x\right)  .
\]

\end{example}

\subsection{Generalized formulas}

Let us next define broken circuits of a matroid $M=\left(  E,\mathcal{I}%
\right)  $. Stanley, in \cite[\S 4.1]{Stanley}, defines them in terms of a
total ordering $\mathcal{O}$ on the set $E$, whereas we shall use a
\textquotedblleft labeling function\textquotedblright\ $\ell:E\rightarrow X$
instead (as in the case of graphs); our setting is slightly more general than Stanley's.

\begin{definition}
\label{def.matroid.BC}Let $M=\left(  E,\mathcal{I}\right)  $ be a matroid. Let
$X$ be a totally ordered set. Let $\ell:E\rightarrow X$ be a function. We
shall refer to $\ell$ as the \textit{labeling function}. For every $e\in E$,
we shall refer to $\ell\left(  e\right)  $ as the \textit{label} of $e$.

A \textit{broken circuit} of $M$ means a subset of $E$ having the form
$C\setminus\left\{  e\right\}  $, where $C$ is a circuit of $M$, and where $e$
is the unique element of $C$ having maximum label (among the elements of $C$).
Of course, the notion of a broken circuit of $M$ depends on the function
$\ell$; however, we suppress the mention of $\ell$ in our notation, since we
will not consider situations where two different $\ell$'s coexist.
\end{definition}

We shall now state analogues (and, in light of Example
\ref{exam.matroid.charpol.graph}, generalizations, although we shall not
elaborate on the few minor technicalities of seeing them as such) of Theorem
\ref{thm.chrompol.varis}, Theorem \ref{thm.chrompol.empty}, Corollary
\ref{cor.chrompol.K-free}, Corollary \ref{cor.chrompol.NBC} and Corollary
\ref{cor.chrompol.NBCfor}:

\begin{theorem}
\label{thm.matroid.charpol.varis}Let $M=\left(  E,\mathcal{I}\right)  $ be a
matroid. Let $m=r_{M}\left(  E\right)  $. Let $X$ be a totally ordered set.
Let $\ell:E\rightarrow X$ be a function. Let $\mathfrak{K}$ be some set of
broken circuits of $M$ (not necessarily containing all of them). Let $a_{K}$
be an element of $\mathbf{k}$ for every $K\in\mathfrak{K}$. Then,%
\[
\widetilde{\chi}_{M}=\sum_{F\subseteq E}\left(  -1\right)  ^{\left\vert
F\right\vert }\left(  \prod_{\substack{K\in\mathfrak{K};\\K\subseteq F}%
}a_{K}\right)  x^{m-r_{M}\left(  F\right)  }.
\]

\end{theorem}

\begin{theorem}
\label{thm.matroid.charpol.empty}Let $M=\left(  E,\mathcal{I}\right)  $ be a
matroid. Let $m=r_{M}\left(  E\right)  $. Then,%
\[
\widetilde{\chi}_{M}=\sum_{F\subseteq E}\left(  -1\right)  ^{\left\vert
F\right\vert }x^{m-r_{M}\left(  F\right)  }.
\]

\end{theorem}

\begin{corollary}
\label{cor.matroid.charpol.K-free}Let $M=\left(  E,\mathcal{I}\right)  $ be a
matroid. Let $m=r_{M}\left(  E\right)  $. Let $X$ be a totally ordered set.
Let $\ell:E\rightarrow X$ be a function. Let $\mathfrak{K}$ be some set of
broken circuits of $M$ (not necessarily containing all of them). Then,%
\[
\widetilde{\chi}_{M}=\sum_{\substack{F\subseteq E;\\F\text{ is }%
\mathfrak{K}\text{-free}}}\left(  -1\right)  ^{\left\vert F\right\vert
}x^{m-r_{M}\left(  F\right)  }.
\]

\end{corollary}

\begin{corollary}
\label{cor.matroid.charpol.NBC}Let $M=\left(  E,\mathcal{I}\right)  $ be a
matroid. Let $m=r_{M}\left(  E\right)  $. Let $X$ be a totally ordered set.
Let $\ell:E\rightarrow X$ be a function. Then,%
\[
\widetilde{\chi}_{M}=\sum_{\substack{F\subseteq E;\\F\text{ contains no
broken}\\\text{circuit of }M\text{ as a subset}}}\left(  -1\right)
^{\left\vert F\right\vert }x^{m-r_{M}\left(  F\right)  }.
\]

\end{corollary}

\begin{corollary}
\label{cor.matroid.charpol.NBCfor}Let $M=\left(  E,\mathcal{I}\right)  $ be a
matroid. Let $m=r_{M}\left(  E\right)  $. Let $X$ be a totally ordered set.
Let $\ell:E\rightarrow X$ be an injective function. Then,%
\[
\widetilde{\chi}_{M}=\sum_{\substack{F\subseteq E;\\F\text{ contains no
broken}\\\text{circuit of }M\text{ as a subset}}}\left(  -1\right)
^{\left\vert F\right\vert }x^{m-\left\vert F\right\vert }.
\]

\end{corollary}

We notice that Corollary \ref{cor.matroid.charpol.NBCfor} is equivalent to
\cite[Theorem 4.12]{Stanley} (at least when $M$ is loopless).

Before we prove these results, let us state a lemma which will serve as an
analogue of Lemma \ref{lem.NBCm.moeb}:

\begin{lemma}
\label{lem.matroid.NBCm.moeb}Let $M=\left(  E,\mathcal{I}\right)  $ be a
matroid. Let $X$ be a totally ordered set. Let $\ell:E\rightarrow X$ be a
function. Let $\mathfrak{K}$ be some set of broken circuits of $M$ (not
necessarily containing all of them). Let $a_{K}$ be an element of $\mathbf{k}$
for every $K\in\mathfrak{K}$.

Let $F$ be any flat of $M$. Then,%
\begin{equation}
\sum_{B\subseteq F}\left(  -1\right)  ^{\left\vert B\right\vert }%
\prod_{\substack{K\in\mathfrak{K};\\K\subseteq B}}a_{K}=\left[  F=\varnothing
\right]  . \label{eq.lem.matroid.NBCm.moeb.1}%
\end{equation}
(Again, we are using the Iverson bracket notation as in Definition
\ref{def.iverson}.)
\end{lemma}

\begin{proof}
[Proof of Lemma \ref{lem.matroid.NBCm.moeb}.]Our proof will imitate the proof
of Lemma \ref{lem.NBCm.moeb} much of the time (with $E\cap\operatorname*{Eqs}%
f$ replaced by $F$); thus, we will allow ourselves some more brevity.

We WLOG assume that $F\neq\varnothing$ (since otherwise, the claim is
obvious\footnote{\textit{Proof.} Assume that $F=\varnothing$. We must show
that the claim is obvious.
\par
Let us first show that no $K\in\mathfrak{K}$ satisfies $K=\varnothing$.
Indeed, assume the contrary. Thus, there exists a $K\in\mathfrak{K}$
satisfying $K=\varnothing$. In other words, $\varnothing\in\mathfrak{K}$.
Thus, $\varnothing$ is a broken circuit of $M$ (since $\mathfrak{K}$ is a set
of broken circuits of $M$). Therefore, $\varnothing$ is obtained from a
circuit of $M$ by removing one element (by the definition of a broken
circuit). This latter circuit must therefore be a one-element set, i.e., it
has the form $\left\{  e\right\}  $ for some $e\in E$. Consider this $e$.
Thus, $\left\{  e\right\}  $ is a circuit of $M$.
\par
But $F$ is a flat of $M$. In other words, Statement $\mathfrak{F}_{1}$ (of
Lemma \ref{lem.matroid.flat-crit}) holds for $T=F$. Hence, Statement
$\mathfrak{F}_{2}$ (of Lemma \ref{lem.matroid.flat-crit}) also holds for $T=F$
(since Lemma \ref{lem.matroid.flat-crit} shows that these two statements are
equivalent). Applying Statement $\mathfrak{F}_{2}$ to $T=F$ and $C=\left\{
e\right\}  $, we thus obtain $\left\{  e\right\}  \subseteq F$ (because
$\left\{  e\right\}  \setminus\left\{  e\right\}  =\varnothing\subseteq F$).
Thus, $e\in\left\{  e\right\}  \subseteq F=\varnothing$, which is absurd. This
contradiction proves that our assumption was wrong.
\par
Hence, we have shown that no $K\in\mathfrak{K}$ satisfies $K=\varnothing$. But
from $F=\varnothing$, we see that the sum $\sum_{B\subseteq F}\left(
-1\right)  ^{\left\vert B\right\vert }\prod_{\substack{K\in\mathfrak{K}%
;\\K\subseteq B}}a_{K}$ has only one addend (namely, the addend for
$B=\varnothing$), and thus simplifies to
\begin{align*}
\underbrace{\left(  -1\right)  ^{\left\vert \varnothing\right\vert }%
}_{=\left(  -1\right)  ^{0}=1}\underbrace{\prod_{\substack{K\in\mathfrak{K}%
;\\K\subseteq\varnothing}}}_{=\prod_{\substack{K\in\mathfrak{K}%
;\\K=\varnothing}}}a_{K}  &  =\prod_{\substack{K\in\mathfrak{K}%
;\\K=\varnothing}}a_{K}=\left(  \text{empty product}\right)
\ \ \ \ \ \ \ \ \ \ \left(  \text{since no }K\in\mathfrak{K}\text{ satisfies
}K=\varnothing\right) \\
&  =1=\left[  F=\varnothing\right]  \ \ \ \ \ \ \ \ \ \ \left(  \text{since
}F=\varnothing\right)  .
\end{align*}
Thus, Lemma \ref{lem.matroid.NBCm.moeb} is proven.}). Thus, $\left[
F=\varnothing\right]  =0$.

Pick any $d\in F$ with maximum $\ell\left(  d\right)  $ (among all $d\in F$).
(This is clearly possible, since $F\neq\varnothing$.) Define two subsets
$\mathcal{U}$ and $\mathcal{V}$ of $\mathcal{P}\left(  F\right)  $ as follows:%
\begin{align*}
\mathcal{U}  &  =\left\{  T\in\mathcal{P}\left(  F\right)  \ \mid\ d\notin
T\right\}  ;\\
\mathcal{V}  &  =\left\{  T\in\mathcal{P}\left(  F\right)  \ \mid\ d\in
T\right\}  .
\end{align*}
Thus, we have $\mathcal{P}\left(  F\right)  =\mathcal{U}\cup\mathcal{V}$, and
the sets $\mathcal{U}$ and $\mathcal{V}$ are disjoint. Now, we define a map
$\Phi:\mathcal{U}\rightarrow\mathcal{V}$ by%
\[
\left(  \Phi\left(  B\right)  =B\cup\left\{  d\right\}
\ \ \ \ \ \ \ \ \ \ \text{for every }B\in\mathcal{U}\right)  .
\]
This map $\Phi$ is well-defined (because for every $B\in\mathcal{U}$, we have
$B\cup\left\{  d\right\}  \in\mathcal{V}$\ \ \ \ \footnote{This follows from
the fact that $d\in F$.}) and a bijection\footnote{Its inverse is the map
$\Psi:\mathcal{V}\rightarrow\mathcal{U}$ defined by $\left(  \Psi\left(
B\right)  =B\setminus\left\{  d\right\}  \ \ \ \ \ \ \ \ \ \ \text{for every
}B\in\mathcal{V}\right)  $.}. Moreover, every $B\in\mathcal{U}$ satisfies%
\begin{equation}
\left(  -1\right)  ^{\left\vert \Phi\left(  B\right)  \right\vert }=-\left(
-1\right)  ^{\left\vert B\right\vert }
\label{pf.lem.matroid.NBCm.moeb.short.Phi.-1}%
\end{equation}
\footnote{\textit{Proof.} This is proven exactly like we proved
(\ref{pf.lem.NBCm.moeb.short.Phi.-1}).}.

Now, we claim that, for every $B\in\mathcal{U}$ and every $K\in\mathfrak{K}$,
we have the following logical equivalence:%
\begin{equation}
\left(  K\subseteq B\right)  \ \Longleftrightarrow\ \left(  K\subseteq
\Phi\left(  B\right)  \right)  .
\label{pf.lem.matroid.NBCm.moeb.short.Phi.equiv}%
\end{equation}


\textit{Proof of (\ref{pf.lem.matroid.NBCm.moeb.short.Phi.equiv}):} Let
$B\in\mathcal{U}$ and $K\in\mathfrak{K}$. We must prove the equivalence
(\ref{pf.lem.matroid.NBCm.moeb.short.Phi.equiv}). The definition of $\Phi$
yields $\Phi\left(  B\right)  =B\cup\left\{  d\right\}  \supseteq B$, so that
$B\subseteq\Phi\left(  B\right)  $. Hence, if $K\subseteq B$, then $K\subseteq
B\subseteq\Phi\left(  B\right)  $. Therefore, the forward implication of the
equivalence (\ref{pf.lem.matroid.NBCm.moeb.short.Phi.equiv}) is proven. It
thus remains to prove the backward implication of this equivalence. In other
words, it remains to prove that if $K\subseteq\Phi\left(  B\right)  $, then
$K\subseteq B$. So let us assume that $K\subseteq\Phi\left(  B\right)  $.

We want to prove that $K\subseteq B$. Assume the contrary. Thus,
$K\not \subseteq B$. We have $K\in\mathfrak{K}$. Thus, $K$ is a broken circuit
of $M$ (since $\mathfrak{K}$ is a set of broken circuits of $M$). In other
words, $K$ is a subset of $E$ having the form $C\setminus\left\{  e\right\}
$, where $C$ is a circuit of $M$, and where $e$ is the unique element of $C$
having maximum label (among the elements of $C$) (because this is how a broken
circuit is defined). Consider these $C$ and $e$. Thus, $K=C\setminus\left\{
e\right\}  $.

The element $e$ is the unique element of $C$ having maximum label
(among the elements of $C$). Thus, if $e^{\prime}$ is any element of $C$
satisfying $\ell\left(  e^{\prime}\right)  \geq\ell\left(  e\right)  $, then%
\begin{equation}
e^{\prime}=e. \label{pf.lem.matroid.NBCm.moeb.short.Phi.equiv.pf.e'}%
\end{equation}


But $\underbrace{K}_{\subseteq\Phi\left(  B\right)  =B\cup\left\{  d\right\}
}\setminus\left\{  d\right\}  \subseteq\left(  B\cup\left\{  d\right\}
\right)  \setminus\left\{  d\right\}  \subseteq B$.

If we had $d\notin K$, then we would have $K\setminus\left\{  d\right\}  =K$
and therefore $K=K\setminus\left\{  d\right\}  \subseteq B$; this would
contradict $K\not \subseteq B$. Hence, we cannot have $d\notin K$. We thus
must have $d\in K$. Hence, $d\in K=C\setminus\left\{  e\right\}  $. Hence,
$d\in C$ and $d\neq e$.

But $C\setminus\left\{  e\right\}  =K\subseteq\Phi\left(  B\right)  \subseteq
F$ (since $\Phi\left(  B\right)  \in\mathcal{P}\left(  F\right)  $). On the
other hand, Statement $\mathfrak{F}_{1}$ (of Lemma \ref{lem.matroid.flat-crit}%
) holds for $T=F$ (since $F$ is a flat of $M$). Hence, Statement
$\mathfrak{F}_{2}$ (of Lemma \ref{lem.matroid.flat-crit}) also holds for $T=F$
(since Lemma \ref{lem.matroid.flat-crit} shows that these two statements are
equivalent). Thus, from $C\setminus\left\{  e\right\}  \subseteq F$, we obtain
$C\subseteq F$. Thus, $e\in C\subseteq F$. Consequently, $\ell\left(
d\right)  \geq\ell\left(  e\right)  $ (since $d$ was defined to be an element
of $F$ with maximum $\ell\left(  d\right)  $ among all $d\in F$).

Also, $d\in C$. Since $\ell\left(  d\right)  \geq\ell\left(  e\right)  $, we
can therefore apply (\ref{pf.lem.matroid.NBCm.moeb.short.Phi.equiv.pf.e'}) to
$e^{\prime}=d$. We thus obtain $d=e$. This contradicts $d\neq e$. This
contradiction proves that our assumption was wrong. Hence, $K\subseteq B$ is
proven. Thus, we have proven the backward implication of the equivalence
(\ref{pf.lem.matroid.NBCm.moeb.short.Phi.equiv}); this completes the proof of
(\ref{pf.lem.matroid.NBCm.moeb.short.Phi.equiv}).

Now, proceeding as in the proof of (\ref{pf.lem.NBCm.moeb.there}), we can show
that%
\[
\sum_{B\subseteq F}\left(  -1\right)  ^{\left\vert B\right\vert }%
\prod_{\substack{K\in\mathfrak{K};\\K\subseteq B}}a_{K}=\left[  F=\varnothing
\right]  .
\]
This proves Lemma \ref{lem.matroid.NBCm.moeb}.
\end{proof}

We shall furthermore use a classical and fundamental result on the M\"{o}bius
function of any finite poset:

\begin{proposition}
\label{prop.moebius.double0}Let $P$ be a finite poset. Let $\mu$ denote the
M\"{o}bius function of $P$.

\textbf{(a)} For any $x\in P$ and $y\in P$, we have%
\begin{equation}
\sum_{\substack{z\in P;\\x\leq z\leq y}}\mu\left(  x,z\right)  =\left[
x=y\right]  . \label{eq.prop.moebius.double0.a}%
\end{equation}


\textbf{(b)} For any $x\in P$ and $y\in P$, we have%
\begin{equation}
\sum_{\substack{z\in P;\\x\leq z\leq y}}\mu\left(  z,y\right)  =\left[
x=y\right]  . \label{eq.prop.moebius.double0.b}%
\end{equation}


\textbf{(c)} Let $\mathbf{k}$ be a $\mathbb{Z}$-module. Let $\left(  \beta
_{x}\right)  _{x\in P}$ be a family of elements of $\mathbf{k}$. Then, every
$z\in P$ satisfies%
\[
\beta_{z}=\sum_{\substack{y\in P;\\y\leq z}}\mu\left(  y,z\right)
\sum_{\substack{x\in P;\\x\leq y}}\beta_{x}.
\]

\end{proposition}

For the sake of completeness, let us give a self-contained proof of this
proposition (slicker arguments appear in the literature\footnote{For example,
Proposition \ref{prop.moebius.double0} \textbf{(c)} is equivalent to the
$\Longrightarrow$ implication of \cite[(5.1a)]{Martin15}.}):

\begin{vershort}
\begin{proof}
[Proof of Proposition \ref{prop.moebius.double0}.]\textbf{(a)} Let $x\in P$
and $y\in P$. We must prove the equality (\ref{eq.prop.moebius.double0.a}). We
are in one of the following three cases:

\textit{Case 1:} We have $x=y$.

\textit{Case 2:} We have $x<y$.

\textit{Case 3:} We have neither $x=y$ nor $x<y$.

Let us first consider Case 1. In this case, we have $x=y$. Hence, the sum
$\sum_{\substack{z\in P;\\x\leq z\leq y}}\mu\left(  x,z\right)  $ contains
only one addend -- namely, the addend for $z=x$. Thus,%
\begin{align*}
\sum_{\substack{z\in P;\\x\leq z\leq y}}\mu\left(  x,z\right)   &  =\mu\left(
x,x\right)  =1\ \ \ \ \ \ \ \ \ \ \left(  \text{by the definition of the
M\"{o}bius function}\right) \\
&  =\left[  x=y\right]  \ \ \ \ \ \ \ \ \ \ \left(  \text{since }x=y\right)  .
\end{align*}
Thus, (\ref{eq.prop.moebius.double0.a}) is proven in Case 1.

Let us now consider Case 2. In this case, we have $x<y$. Hence, $x\neq y$, so
that $\left[  x=y\right]  =0$. Now, $y$ is an element of $P$ satisfying $x\leq
y\leq y$. Thus, the sum $\sum_{\substack{z\in P;\\x\leq z\leq y}}\mu\left(
x,z\right)  $ contains an addend for $z=y$. Splitting off this addend, we
obtain%
\begin{align*}
\sum_{\substack{z\in P;\\x\leq z\leq y}}\mu\left(  x,z\right)   &
=\underbrace{\sum_{\substack{z\in P;\\x\leq z\leq y;\ z\neq y}}}%
_{=\sum_{\substack{z\in P;\\x\leq z<y}}}\mu\left(  x,z\right)
+\underbrace{\mu\left(  x,y\right)  }_{\substack{=-\sum_{\substack{z\in
P;\\x\leq z<y}}\mu\left(  x,z\right)  \\\text{(by (\ref{eq.def.moebius.rec2}%
))}}}\\
&  =\sum_{\substack{z\in P;\\x\leq z<y}}\mu\left(  x,z\right)  +\left(
-\sum_{\substack{z\in P;\\x\leq z<y}}\mu\left(  x,z\right)  \right)
=0=\left[  x=y\right]  .
\end{align*}
Hence, (\ref{eq.prop.moebius.double0.a}) is proven in Case 2.

Finally, let us consider Case 3. In this case, we have neither $x=y$ nor
$x<y$. Thus, we do not have $x\leq y$. Hence, there exists no $z\in P$
satisfying $x\leq z\leq y$. Thus,
\[
\sum_{\substack{z\in P;\\x\leq z\leq y}}\mu\left(  x,z\right)  =\left(
\text{empty sum}\right)  =0=\left[  x=y\right]
\]
(since we do not have $x=y$). Thus, (\ref{eq.prop.moebius.double0.a}) is
proven in Case 3.

Hence, (\ref{eq.prop.moebius.double0.a}) is proven in all three cases. This
proves Proposition \ref{prop.moebius.double0} \textbf{(a)}.

\textbf{(b)} For any two elements $u$ and $v$ of $P$, we define a subset
$\left[  u,v\right]  $ of $P$ by%
\[
\left[  u,v\right]  =\left\{  w\in P\ \mid\ u\leq w\leq v\right\}  .
\]
Thus subset $\left[  u,v\right]  $ is finite (since $P$ is finite), and thus
its size $\left\vert \left[  u,v\right]  \right\vert $ is a nonnegative integer.

We shall now prove Proposition \ref{prop.moebius.double0} \textbf{(b)} by
strong induction on $\left\vert \left[  x,y\right]  \right\vert $:

\textit{Induction step:} Let $N\in\mathbb{N}$. Assume that Proposition
\ref{prop.moebius.double0} \textbf{(b)} holds whenever $\left\vert \left[
x,y\right]  \right\vert <N$. We must now prove that Proposition
\ref{prop.moebius.double0} \textbf{(b)} holds whenever $\left\vert \left[
x,y\right]  \right\vert =N$.

We have assumed that Proposition \ref{prop.moebius.double0} \textbf{(b)} holds
whenever $\left\vert \left[  x,y\right]  \right\vert <N$. In other words, we
have assumed the following claim:

\begin{statement}
\textit{Claim 1:} For any $x\in P$ and $y\in P$ satisfying $\left\vert \left[
x,y\right]  \right\vert <N$, we have%
\[
\sum_{\substack{z\in P;\\x\leq z\leq y}}\mu\left(  z,y\right)  =\left[
x=y\right]  .
\]

\end{statement}

Now, let $x$ and $y$ be two elements of $P$ satisfying $\left\vert \left[
x,y\right]  \right\vert =N$. We are going to prove that%
\begin{equation}
\sum_{\substack{z\in P;\\x\leq z\leq y}}\mu\left(  z,y\right)  =\left[
x=y\right]  . \label{pf.prop.moebius.double0.short.b.goal}%
\end{equation}


We are in one of the following three cases:

\textit{Case 1:} We have $x=y$.

\textit{Case 2:} We have $x<y$.

\textit{Case 3:} We have neither $x=y$ nor $x<y$.

In Case 1 and in Case 3, we can prove
(\ref{pf.prop.moebius.double0.short.b.goal}) in exactly the same way as (in
our above proof of Proposition \ref{prop.moebius.double0} \textbf{(a)}) we
have proven (\ref{eq.prop.moebius.double0.a}). Thus, it remains only to prove
(\ref{pf.prop.moebius.double0.short.b.goal}) in Case 2. In other words, we can
WLOG assume that we are in Case 2.

Assume this. Hence, $x<y$, so that $\left[  x=y\right]  =0$.

For every $t\in P$ satisfying $x\leq t<y$, we have%
\begin{equation}
\left\vert \left[  x,t\right]  \right\vert <N
\label{pf.prop.moebius.double0.short.b.smaller}%
\end{equation}
\footnote{\textit{Proof of (\ref{pf.prop.moebius.double0.short.b.smaller}):}
Let $t\in P$ be such that $x\leq t<y$. We shall proceed in several steps:
\par
\begin{itemize}
\item We have
\begin{align*}
\left[  x,t\right]   &  =\left\{  w\in P\ \mid\ x\leq w\leq t\right\}
\ \ \ \ \ \ \ \ \ \ \left(  \text{by the definition of }\left[  x,t\right]
\right) \\
&  \subseteq\left\{  w\in P\ \mid\ x\leq w\leq y\right\}
\ \ \ \ \ \ \ \ \ \ \left(
\begin{array}
[c]{c}%
\text{because every }w\in P\text{ satisfying }w\leq t\\
\text{must also satisfy }w\leq y\text{ (since }t<y\text{)}%
\end{array}
\right) \\
&  =\left[  x,y\right]  \ \ \ \ \ \ \ \ \ \ \left(  \text{by the definition of
}\left[  x,y\right]  \right)  .
\end{align*}
\par
\item We have $t<y$. Thus, we do not have $y\leq t$. Hence, we do not have
$x\leq y\leq t$. Hence, $y\notin\left[  x,t\right]  $. But $y\in\left[
x,y\right]  $ (since $x\leq y\leq y$). Hence, the sets $\left[  x,t\right]  $
and $\left[  x,y\right]  $ are distinct (since the latter contains $y$ but the
former does not). Combining this with $\left[  x,t\right]  \subseteq\left[
x,y\right]  $, we conclude that $\left[  x,t\right]  $ is a proper subset of
$\left[  x,y\right]  $. Hence, $\left\vert \left[  x,t\right]  \right\vert
<\left\vert \left[  x,y\right]  \right\vert =N$. This proves
(\ref{pf.prop.moebius.double0.short.b.smaller}).
\end{itemize}
}. Therefore, for every $t\in P$ satisfying $x\leq t<y$, we have%
\begin{equation}
\sum_{\substack{z\in P;\\x\leq z\leq t}}\mu\left(  z,t\right)  =\left[
x=t\right]  \label{pf.prop.moebius.double0.short.b.smaller2}%
\end{equation}
(by Claim 1, applied to $t$ instead of $y$). Also, for every $u\in P$ and
$v\in P$, we have%
\begin{equation}
\sum_{\substack{t\in P;\\u\leq t\leq v}}\mu\left(  u,t\right)  =\left[
u=v\right]  \label{pf.prop.moebius.double0.short.b.smaller3}%
\end{equation}
\footnote{\textit{Proof of (\ref{pf.prop.moebius.double0.short.b.smaller3}):}
Let $u\in P$ and $v\in P$. Proposition \ref{prop.moebius.double0} \textbf{(a)}
(applied to $x=u$ and $y=v$) shows that $\sum_{\substack{z\in P;\\u\leq z\leq
v}}\mu\left(  u,t\right)  =\left[  u=v\right]  $. Now,%
\begin{align*}
\sum_{\substack{t\in P;\\u\leq t\leq v}}\mu\left(  u,t\right)   &
=\sum_{\substack{z\in P;\\u\leq z\leq v}}\mu\left(  u,t\right)
\ \ \ \ \ \ \ \ \ \ \left(  \text{here, we have substituted }z\text{ for
}t\text{ in the sum}\right) \\
&  =\left[  u=v\right]  .
\end{align*}
This proves (\ref{pf.prop.moebius.double0.short.b.smaller3}).}.

Now,%
\begin{align*}
&  \underbrace{\sum_{\substack{\left(  z,t\right)  \in P^{2};\\x\leq z\leq
t\leq y}}}_{=\sum_{\substack{z\in P;\\x\leq z\leq y}}\sum_{\substack{t\in
P;\\z\leq t\leq y}}}\mu\left(  z,t\right) \\
&  =\sum_{\substack{z\in P;\\x\leq z\leq y}}\underbrace{\sum_{\substack{t\in
P;\\z\leq t\leq y}}\mu\left(  z,t\right)  }_{\substack{=\left[  z=y\right]
\\\text{(by (\ref{pf.prop.moebius.double0.short.b.smaller3})}\\\text{(applied
to }u=z\text{ and }v=y\text{))}}}=\sum_{\substack{z\in P;\\x\leq z\leq
y}}\left[  z=y\right] \\
&  =\sum_{\substack{z\in P;\\x\leq z\leq y\text{ and }z=y}}\underbrace{\left[
z=y\right]  }_{\substack{=1\\\text{(since }z=y\text{)}}}+\sum_{\substack{z\in
P;\\x\leq z\leq y\text{ and }z\neq y}}\underbrace{\left[  z=y\right]
}_{\substack{=0\\\text{(since }z\neq y\text{)}}}\\
&  \ \ \ \ \ \ \ \ \ \ \left(  \text{since every }z\in P\text{ satisfies
either }z=y\text{ or }z\neq y\text{ (but not both)}\right) \\
&  =\underbrace{\sum_{\substack{z\in P;\\x\leq z\leq y\text{ and }z=y}%
}}_{=\sum_{z\in\left\{  w\in P\ \mid\ x\leq w\leq y\text{ and }w=y\right\}  }%
}1+\underbrace{\sum_{\substack{z\in P;\\x\leq z\leq y\text{ and }z\neq y}%
}0}_{=0}=\sum_{z\in\left\{  w\in P\ \mid\ x\leq w\leq y\text{ and
}w=y\right\}  }1\\
&  =\left\vert \underbrace{\left\{  w\in P\ \mid\ x\leq w\leq y\text{ and
}w=y\right\}  }_{=\left\{  y\right\}  }\right\vert =\left\vert \left\{
y\right\}  \right\vert =1.
\end{align*}
Hence,%
\begin{align*}
1  &  =\underbrace{\sum_{\substack{\left(  z,t\right)  \in P^{2};\\x\leq z\leq
t\leq y}}}_{=\sum_{\substack{t\in P;\\x\leq t\leq y}}\sum_{\substack{z\in
P;\\x\leq z\leq t}}}\mu\left(  z,t\right)  =\sum_{\substack{t\in P;\\x\leq
t\leq y}}\sum_{\substack{z\in P;\\x\leq z\leq t}}\mu\left(  z,t\right) \\
&  =\underbrace{\sum_{\substack{t\in P;\\x\leq t\leq y\text{ and }t=y}%
}}_{\substack{=\sum_{t\in\left\{  w\in P\ \mid\ x\leq w\leq y\text{ and
}w=y\right\}  }=\sum_{t\in\left\{  y\right\}  }\\\text{(since }\left\{  w\in
P\ \mid\ x\leq w\leq y\text{ and }w=y\right\}  =\left\{  y\right\}  \text{)}%
}}\sum_{\substack{z\in P;\\x\leq z\leq t}}\mu\left(  z,t\right) \\
&  \ \ \ \ \ \ \ \ \ \ +\sum_{\substack{t\in P;\\x\leq t\leq y\text{ and
}t\neq y}}\underbrace{\sum_{\substack{z\in P;\\x\leq z\leq t}}\mu\left(
z,t\right)  }_{\substack{=\left[  x=t\right]  \\\text{(by
(\ref{pf.prop.moebius.double0.short.b.smaller2})}\\\text{(since }t<y\text{
(because }t\leq y\text{ and }t\neq y\text{)) and }x\leq t\text{)}}}\\
&  \ \ \ \ \ \ \ \ \ \ \left(  \text{since every }t\in P\text{ satisfies
either }t=y\text{ or }t\neq y\text{ (but not both)}\right) \\
&  =\underbrace{\sum_{t\in\left\{  y\right\}  }\sum_{\substack{z\in P;\\x\leq
z\leq t}}\mu\left(  z,t\right)  }_{=\sum_{\substack{z\in P;\\x\leq z\leq
y}}\mu\left(  z,y\right)  }+\sum_{\substack{t\in P;\\x\leq t\leq y\text{ and
}t\neq y}}\left[  x=t\right] \\
&  =\sum_{\substack{z\in P;\\x\leq z\leq y}}\mu\left(  z,y\right)
+\sum_{\substack{t\in P;\\x\leq t\leq y\text{ and }t\neq y}}\left[
x=t\right]  .
\end{align*}
Subtracting $\sum_{\substack{z\in P;\\x\leq z\leq y}}\mu\left(  z,y\right)  $
from both sides of this equality, we obtain%
\begin{align*}
&  1-\sum_{\substack{z\in P;\\x\leq z\leq y}}\mu\left(  z,y\right) \\
&  =\sum_{\substack{t\in P;\\x\leq t\leq y\text{ and }t\neq y}}\left[
x=t\right] \\
&  =\underbrace{\sum_{\substack{t\in P;\\x\leq t\leq y\text{ and }t=x\text{
and }t\neq y}}}_{\substack{=\sum_{t\in\left\{  z\in P\ \mid\ x\leq z\leq
y\text{ and }z=x\text{ and }z\neq y\right\}  }=\sum_{t\in\left\{  x\right\}
}\\\text{(since }\left\{  z\in P\ \mid\ x\leq z\leq y\text{ and }z=x\text{ and
}z\neq y\right\}  =\left\{  x\right\}  \text{)}}}\underbrace{\left[
x=t\right]  }_{\substack{=1\\\text{(since }x=t\text{)}}}\\
&  \ \ \ \ \ \ \ \ \ \ +\sum_{\substack{t\in P;\\x\leq t\leq y\text{ and
}t\neq x\text{ and }t\neq y}}\underbrace{\left[  x=t\right]  }%
_{\substack{=0\\\text{(since }x\neq t\text{)}}}\\
&  \ \ \ \ \ \ \ \ \ \ \left(  \text{since every }t\in P\text{ satisfies
either }t=x\text{ or }t\neq x\text{ (but not both)}\right) \\
&  =\sum_{t\in\left\{  x\right\}  }1+\underbrace{\sum_{\substack{t\in
P;\\x\leq t\leq y\text{ and }t\neq x\text{ and }t\neq y}}0}_{=0}=\sum
_{t\in\left\{  x\right\}  }1=1.
\end{align*}
Solving this equality for $\sum_{\substack{z\in P;\\x\leq z\leq y}}\mu\left(
z,y\right)  $, we obtain%
\[
\sum_{\substack{z\in P;\\x\leq z\leq y}}\mu\left(  z,y\right)  =1-1=0=\left[
x=y\right]
\]
(since $x<y$). Thus, (\ref{pf.prop.moebius.double0.short.b.goal}) is proven.

Let us now forget that we fixed $x$ and $y$. We thus have proven that for any
$x\in P$ and $y\in P$ satisfying $\left\vert \left[  x,y\right]  \right\vert
=N$, we have%
\[
\sum_{\substack{z\in P;\\x\leq z\leq y}}\mu\left(  z,y\right)  =\left[
x=y\right]  .
\]
In other words, Proposition \ref{prop.moebius.double0} \textbf{(b)} holds
whenever $\left\vert \left[  x,y\right]  \right\vert =N$. This completes the
induction step. Thus, Proposition \ref{prop.moebius.double0} \textbf{(b)} is
proven by induction.

\textbf{(c)} For every $v\in P$, we have%
\begin{align*}
&  \sum_{\substack{y\in P;\\y\leq v}}\mu\left(  y,v\right)  \sum
_{\substack{x\in P;\\x\leq y}}\beta_{x}\\
&  =\sum_{\substack{z\in P;\\z\leq v}}\mu\left(  z,v\right)  \sum
_{\substack{x\in P;\\x\leq z}}\beta_{x}\ \ \ \ \ \ \ \ \ \ \left(
\begin{array}
[c]{c}%
\text{here, we have renamed the summation}\\
\text{index }y\text{ as }z\text{ in the outer sum}%
\end{array}
\right) \\
&  =\underbrace{\sum_{\substack{z\in P;\\z\leq v}}\sum_{\substack{x\in
P;\\x\leq z}}}_{=\sum_{x\in P}\sum_{\substack{z\in P;\\x\leq z\leq v}}}%
\mu\left(  z,v\right)  \beta_{x}=\sum_{x\in P}\underbrace{\sum_{\substack{z\in
P;\\x\leq z\leq v}}\mu\left(  z,v\right)  }_{\substack{=\left[  x=v\right]
\\\text{(by Proposition \ref{prop.moebius.double0} \textbf{(b)}}%
\\\text{(applied to }y=v\text{))}}}\beta_{x}\\
&  =\sum_{x\in P}\left[  x=v\right]  \beta_{x}=\sum_{\substack{x\in
P;\\x=v}}\underbrace{\left[  x=v\right]  }_{\substack{=1\\\text{(since
}x=v\text{)}}}\beta_{x}+\sum_{\substack{x\in P;\\x\neq v}}\underbrace{\left[
x=v\right]  }_{\substack{=0\\\text{(since }x\neq v\text{)}}}\beta_{x}\\
&  \ \ \ \ \ \ \ \ \ \ \left(  \text{since every }x\in P\text{ satisfies
either }x=v\text{ or }x\neq v\text{ (but not both)}\right) \\
&  =\sum_{\substack{x\in P;\\x=v}}\beta_{x}+\underbrace{\sum_{\substack{x\in
P;\\x\neq v}}0\beta_{x}}_{=0}=\sum_{\substack{x\in P;\\x=v}}\beta_{x}%
=\beta_{v}\ \ \ \ \ \ \ \ \ \ \left(  \text{since }v\in P\right)  .
\end{align*}
Renaming $v$ as $z$ in this result, we obtain precisely Proposition
\ref{prop.moebius.double0} \textbf{(c)}.
\end{proof}
\end{vershort}

\begin{verlong}
\begin{proof}
[Proof of Proposition \ref{prop.moebius.double0}.]We first notice that%
\begin{equation}
\left\{  w\in P\ \mid\ x\leq w\leq x\right\}  =\left\{  x\right\}
\label{pf.prop.moebius.double0.xx}%
\end{equation}
for every $x\in P$\ \ \ \ \footnote{\textit{Proof of
(\ref{pf.prop.moebius.double0.xx}):} Let $x\in P$.
\par
Clearly, $x$ is an element of $P$ and satisfies $x\leq x\leq x$. Thus, $x$ is
an element $w$ of $P$ satisfying $x\leq w\leq x$. In other words,
$x\in\left\{  w\in P\ \mid\ x\leq w\leq x\right\}  $. Hence, $\left\{
x\right\}  \subseteq\left\{  w\in P\ \mid\ x\leq w\leq x\right\}  $.
\par
On the other hand, let $z\in\left\{  w\in P\ \mid\ x\leq w\leq x\right\}  $.
Thus, $z$ is an element $w$ of $P$ satisfying $x\leq w\leq x$. In other words,
$z$ is an element of $P$ and satisfies $x\leq z\leq x$. Combining $x\leq z$
with $z\leq x$, we obtain $x=z$ (since the partial order on $P$ is
antisymmetric). Hence, $z=x\in\left\{  x\right\}  $.
\par
Now, forget that we fixed $z$. We thus have proven that $z\in\left\{
x\right\}  $ for every $z\in\left\{  w\in P\ \mid\ x\leq w\leq x\right\}  $.
In other words, $\left\{  w\in P\ \mid\ x\leq w\leq x\right\}  \subseteq
\left\{  x\right\}  $. Combining this with $\left\{  x\right\}  \subseteq
\left\{  w\in P\ \mid\ x\leq w\leq x\right\}  $, we obtain $\left\{  w\in
P\ \mid\ x\leq w\leq x\right\}  =\left\{  x\right\}  $. This proves
(\ref{pf.prop.moebius.double0.xx}).}.

\textbf{(a)} Let $x\in P$ and $y\in P$. We must prove the equality
(\ref{eq.prop.moebius.double0.a}). If we do not have $x\leq y$, then
(\ref{eq.prop.moebius.double0.a}) holds for obvious
reasons\footnote{\textit{Proof.} Assume that we do not have $x\leq y$. Then,
there exists no $z\in P$ satisfying $x\leq z\leq y$ (because if such a $z$
would exist, then it would satisfy $x\leq z\leq y$, which would contradict the
fact that we do not have $x\leq y$). Therefore, the sum $\sum_{\substack{z\in
P;\\x\leq z\leq y}}\mu\left(  x,z\right)  $ is an empty sum. Hence,%
\[
\sum_{\substack{z\in P;\\x\leq z\leq y}}\mu\left(  x,z\right)  =\left(
\text{empty sum}\right)  =0.
\]
But we do not have $x=y$ (because if we had $x=y$, then we would have $x\leq
y$, which would contradict the fact that we do not have $x\leq y$). Thus,
$\left[  x=y\right]  =0$. Now, $\sum_{\substack{z\in P;\\x\leq z\leq y}%
}\mu\left(  x,z\right)  =0=\left[  x=y\right]  $. Hence,
(\ref{eq.prop.moebius.double0.a}) is proven, qed.}. Hence, for the rest of our
proof of (\ref{eq.prop.moebius.double0.a}), we can WLOG assume that $x\leq y$.
Assume this.

We have $x\leq y$. Thus, either $x=y$ or $x<y$. In other words, we are in one
of the following two cases:

\textit{Case 1:} We have $x=y$.

\textit{Case 2:} We have $x<y$.

Let us first consider Case 1. In this case, we have $x=y$, so that $\left[
x=y\right]  =1$. On the other hand,
\[
\left\{  w\in P\ \mid\ x\leq w\leq\underbrace{y}_{=x}\right\}  =\left\{  w\in
P\ \mid\ x\leq w\leq x\right\}  =\left\{  x\right\}
\]
(by (\ref{pf.prop.moebius.double0.xx})). Now,%
\[
\underbrace{\sum_{\substack{z\in P;\\x\leq z\leq y}}}_{\substack{=\sum
_{z\in\left\{  w\in P\ \mid\ x\leq w\leq y\right\}  }=\sum_{z\in\left\{
x\right\}  }\\\text{(since }\left\{  w\in P\ \mid\ x\leq w\leq y\right\}
=\left\{  x\right\}  \text{)}}}\mu\left(  x,z\right)  =\sum_{z\in\left\{
x\right\}  }\mu\left(  x,z\right)  =\mu\left(  x,x\right)  =1
\]
(by (\ref{eq.def.moebius.rec1})). Comparing this with $\left[  x=y\right]
=1$, we obtain $\sum_{\substack{z\in P;\\x\leq z\leq y}}\mu\left(  x,z\right)
=\left[  x=y\right]  $. Hence, (\ref{eq.prop.moebius.double0.a}) is proven in
Case 1.

Let us now consider Case 2. In this case, we have $x<y$. Thus,
(\ref{eq.def.moebius.rec2}) yields
\begin{align*}
\mu\left(  x,y\right)   &  =-\underbrace{\sum_{\substack{z\in P;\\x\leq z<y}%
}}_{\substack{=\sum_{\substack{z\in P;\\x\leq z\text{ and }z<y}}=\sum
_{\substack{z\in P;\\x\leq z\text{ and }z\leq y\text{ and }z\neq
y}}\\\text{(since the condition }\left(  z<y\right)  \text{ is equivalent
to}\\\text{the condition }\left(  z\leq y\text{ and }z\neq y\right)  \text{)}%
}}\mu\left(  x,z\right)  =-\underbrace{\sum_{\substack{z\in P;\\x\leq z\text{
and }z\leq y\text{ and }z\neq y}}}_{=\sum_{\substack{z\in P;\\x\leq z\leq
y\text{ and }z\neq y}}}\mu\left(  x,z\right) \\
&  =-\sum_{\substack{z\in P;\\x\leq z\leq y\text{ and }z\neq y}}\mu\left(
x,z\right)  .
\end{align*}
Hence, $\sum_{\substack{z\in P;\\x\leq z\leq y\text{ and }z\neq y}}\mu\left(
x,z\right)  =-\mu\left(  x,y\right)  $.

On the other hand,
\[
\left\{  w\in P\ \mid\ x\leq w\leq y\text{ and }w=y\right\}  =\left\{
y\right\}
\]
\footnote{\textit{Proof.} Let $z\in\left\{  w\in P\ \mid\ x\leq w\leq y\text{
and }w=y\right\}  $. Thus, $z$ is an element $w$ of $P$ satisfying $x\leq
w\leq y$ and $w=y$. In other words, $z$ is an element of $P$ and satisfies
$x\leq z\leq y$ and $z=y$. Now, $z=y\in\left\{  y\right\}  $. Let us now
forget that we fixed $z$. We thus have proven that $z\in\left\{  y\right\}  $
for every $z\in\left\{  w\in P\ \mid\ x\leq w\leq y\text{ and }w=y\right\}  $.
In other words, $\left\{  w\in P\ \mid\ x\leq w\leq y\text{ and }w=y\right\}
\subseteq\left\{  y\right\}  $.
\par
On the other hand, $y$ is an element of $P$ and satisfies $x\leq y\leq y$ and
$y=y$. In other words, $y$ is an element $w$ of $P$ satisfying $x\leq w\leq y$
and $w=y$. Hence, $y\in\left\{  w\in P\ \mid\ x\leq w\leq y\text{ and
}w=y\right\}  $. Thus, $\left\{  y\right\}  \subseteq\left\{  w\in
P\ \mid\ x\leq w\leq y\text{ and }w=y\right\}  $. Combining this with
$\left\{  w\in P\ \mid\ x\leq w\leq y\text{ and }w=y\right\}  \subseteq
\left\{  y\right\}  $, we obtain $\left\{  w\in P\ \mid\ x\leq w\leq y\text{
and }w=y\right\}  =\left\{  y\right\}  $. Qed.}.

But every $z\in P$ satisfies either $z=y$ or $z\neq y$ (but not both). Thus,%
\begin{align*}
\sum_{\substack{z\in P;\\x\leq z\leq y}}\mu\left(  x,z\right)   &
=\underbrace{\sum_{\substack{z\in P;\\x\leq z\leq y\text{ and }z=y}%
}}_{\substack{=\sum_{z\in\left\{  w\in P\ \mid\ x\leq w\leq y\text{ and
}w=y\right\}  }=\sum_{z\in\left\{  y\right\}  }\\\text{(since }\left\{  w\in
P\ \mid\ x\leq w\leq y\text{ and }w=y\right\}  =\left\{  y\right\}  \text{)}%
}}\mu\left(  x,z\right)  +\underbrace{\sum_{\substack{z\in P;\\x\leq z\leq
y\text{ and }z\neq y}}\mu\left(  x,z\right)  }_{=-\mu\left(  x,y\right)  }\\
&  =\underbrace{\sum_{z\in\left\{  y\right\}  }\mu\left(  x,z\right)  }%
_{=\mu\left(  x,y\right)  }+\left(  -\mu\left(  x,y\right)  \right)
=\mu\left(  x,y\right)  +\left(  -\mu\left(  x,y\right)  \right)  =0.
\end{align*}
Thus, (\ref{eq.prop.moebius.double0.a}) is proven in Case 2.

We have now proven (\ref{eq.prop.moebius.double0.a}) in each of the two Cases
1 and 2. Thus, (\ref{eq.prop.moebius.double0.a}) always holds (since Cases 1
and 2 cover all possibilities). Proposition \ref{prop.moebius.double0}
\textbf{(a)} is thus proven.

\textbf{(b)} For any two elements $u$ and $v$ of $P$, we define a subset
$\left[  u,v\right]  $ of $P$ by%
\[
\left[  u,v\right]  =\left\{  w\in P\ \mid\ u\leq w\leq v\right\}  .
\]
Thus subset $\left[  u,v\right]  $ is finite (since $P$ is finite), and thus
its size $\left\vert \left[  u,v\right]  \right\vert $ is a nonnegative integer.

We shall now prove Proposition \ref{prop.moebius.double0} \textbf{(b)} by
strong induction on $\left\vert \left[  x,y\right]  \right\vert $:

\textit{Induction step:} Let $N\in\mathbb{N}$. Assume that Proposition
\ref{prop.moebius.double0} \textbf{(b)} holds whenever $\left\vert \left[
x,y\right]  \right\vert <N$. We must now prove that Proposition
\ref{prop.moebius.double0} \textbf{(b)} holds whenever $\left\vert \left[
x,y\right]  \right\vert =N$.

We have assumed that Proposition \ref{prop.moebius.double0} \textbf{(b)} holds
whenever $\left\vert \left[  x,y\right]  \right\vert <N$. In other words, we
have assumed the following claim:

\begin{statement}
\textit{Claim 1:} For any $x\in P$ and $y\in P$ satisfying $\left\vert \left[
x,y\right]  \right\vert <N$, we have%
\[
\sum_{\substack{z\in P;\\x\leq z\leq y}}\mu\left(  z,y\right)  =\left[
x=y\right]  .
\]

\end{statement}

Now, let $x$ and $y$ be two elements of $P$ satisfying $\left\vert \left[
x,y\right]  \right\vert =N$. We are going to prove that%
\begin{equation}
\sum_{\substack{z\in P;\\x\leq z\leq y}}\mu\left(  z,y\right)  =\left[
x=y\right]  . \label{pf.prop.moebius.double0.b.goal}%
\end{equation}


If we do not have $x\leq y$, then (\ref{pf.prop.moebius.double0.b.goal}) holds
for obvious reasons\footnote{\textit{Proof.} Assume that we do not have $x\leq
y$. Then, there exists no $z\in P$ satisfying $x\leq z\leq y$ (because if such
a $z$ would exist, then it would satisfy $x\leq z\leq y$, which would
contradict with the fact that we do not have $x\leq y$). Therefore, the sum
$\sum_{\substack{z\in P;\\x\leq z\leq y}}\mu\left(  z,y\right)  $ is an empty
sum. Hence,%
\[
\sum_{\substack{z\in P;\\x\leq z\leq y}}\mu\left(  z,y\right)  =\left(
\text{empty sum}\right)  =0.
\]
But we do not have $x=y$ (because if we had $x=y$, then we would have $x\leq
y$, which would contradict with the fact that we do not have $x\leq y$). Thus,
$\left[  x=y\right]  =0$. Now, $\sum_{\substack{z\in P;\\x\leq z\leq y}%
}\mu\left(  z,y\right)  =0=\left[  x=y\right]  $. Hence,
(\ref{pf.prop.moebius.double0.b.goal}) is proven, qed.}. Hence, for the rest
of our proof of (\ref{pf.prop.moebius.double0.b.goal}), we can WLOG assume
that $x\leq y$. Assume this.

If $x=y$, then (\ref{pf.prop.moebius.double0.b.goal}) holds for obvious
reasons\footnote{\textit{Proof.} Assume that $x=y$. Then,%
\[
\left\{  w\in P\ \mid\ x\leq w\leq\underbrace{y}_{=x}\right\}  =\left\{  w\in
P\ \mid\ x\leq w\leq x\right\}  =\left\{  x\right\}
\]
(by (\ref{pf.prop.moebius.double0.xx})). Now,
\[
\underbrace{\sum_{\substack{z\in P;\\x\leq z\leq y}}}_{\substack{=\sum
_{z\in\left\{  w\in P\ \mid\ x\leq w\leq y\right\}  }=\sum_{z\in\left\{
x\right\}  }\\\text{(since }\left\{  w\in P\ \mid\ x\leq w\leq y\right\}
=\left\{  x\right\}  \text{)}}}\mu\left(  z,y\right)  =\sum_{z\in\left\{
x\right\}  }\mu\left(  z,y\right)  =\mu\left(  x,\underbrace{y}_{=x}\right)
=\mu\left(  x,x\right)  =1
\]
(by (\ref{eq.def.moebius.rec1})). Comparing this with
\[
\left[  x=y\right]  =1\ \ \ \ \ \ \ \ \ \ \left(  \text{since }x=y\right)  ,
\]
we obtain $\sum_{\substack{z\in P;\\x\leq z\leq y}}\mu\left(  z,y\right)
=\left[  x=y\right]  $. Hence, (\ref{pf.prop.moebius.double0.b.goal}) is
proven, qed.}. Hence, for the rest of our proof of
(\ref{pf.prop.moebius.double0.b.goal}), we can WLOG assume that we don't have
$x=y$. Assume this. Thus, we don't have $x=y$. In other words, we have $x\neq
y$. Combining this with $x\leq y$, we obtain $x<y$.

Notice that $\left[  x=y\right]  =0$ (since we don't have $x=y$).

For any pair $\left(  z,t\right)  \in P^{2}$, we have the following logical
equivalence:%
\begin{equation}
\left(  x\leq z\leq y\text{ and }z\leq t\leq y\right)  \ \Longleftrightarrow
\ \left(  x\leq t\leq y\text{ and }x\leq z\leq t\right)
\label{pf.prop.moebius.double0.b.equiv}%
\end{equation}
\footnote{\textit{Proof of (\ref{pf.prop.moebius.double0.b.equiv}):} Let
$\left(  z,t\right)  \in P^{2}$. We shall first prove the logical implication%
\begin{equation}
\left(  x\leq z\leq y\text{ and }z\leq t\leq y\right)  \ \Longrightarrow
\ \left(  x\leq t\leq y\text{ and }x\leq z\leq t\right)  .
\label{pf.prop.moebius.double0.b.equiv.pf.1}%
\end{equation}
\par
\textit{Proof of (\ref{pf.prop.moebius.double0.b.equiv.pf.1}):} Assume that
$\left(  x\leq z\leq y\text{ and }z\leq t\leq y\right)  $ holds. We must prove
that $\left(  x\leq t\leq y\text{ and }x\leq z\leq t\right)  $ holds. We have
$x\leq z\leq t$, so that $x\leq t\leq y$. Also, $x\leq z\leq t$. Thus,
$\left(  x\leq t\leq y\text{ and }x\leq z\leq t\right)  $ holds. This proves
the implication (\ref{pf.prop.moebius.double0.b.equiv.pf.1}).
\par
Next, we shall prove the logical implication%
\begin{equation}
\left(  x\leq t\leq y\text{ and }x\leq z\leq t\right)  \ \Longrightarrow
\ \left(  x\leq z\leq y\text{ and }z\leq t\leq y\right)  .
\label{pf.prop.moebius.double0.b.equiv.pf.2}%
\end{equation}
\par
\textit{Proof of (\ref{pf.prop.moebius.double0.b.equiv.pf.2}):} Assume that
$\left(  x\leq t\leq y\text{ and }x\leq z\leq t\right)  $ holds. We must prove
that $\left(  x\leq z\leq y\text{ and }z\leq t\leq y\right)  $ holds. We have
$z\leq t\leq y$, thus $x\leq z\leq y$. Also, $z\leq t\leq y$. Thus, $\left(
x\leq z\leq y\text{ and }z\leq t\leq y\right)  $ holds. This proves the
implication (\ref{pf.prop.moebius.double0.b.equiv.pf.2}).
\par
Combining the two implications (\ref{pf.prop.moebius.double0.b.equiv.pf.1})
and (\ref{pf.prop.moebius.double0.b.equiv.pf.2}), we obtain the logical
equivalence%
\[
\left(  x\leq z\leq y\text{ and }z\leq t\leq y\right)  \ \Longleftrightarrow
\ \left(  x\leq t\leq y\text{ and }x\leq z\leq t\right)  .
\]
This proves (\ref{pf.prop.moebius.double0.b.equiv}).}.

For every $t\in P$ satisfying $x\leq t<y$, we have%
\begin{equation}
\left\vert \left[  x,t\right]  \right\vert <N
\label{pf.prop.moebius.double0.b.smaller}%
\end{equation}
\footnote{\textit{Proof of (\ref{pf.prop.moebius.double0.b.smaller}):} Let
$t\in P$ be such that $x\leq t<y$. We shall proceed in several steps:
\par
\begin{itemize}
\item We have $\left[  x,t\right]  =\left\{  w\in P\ \mid\ x\leq w\leq
t\right\}  $ (by the definition of $\left[  x,t\right]  $) and $\left[
x,y\right]  =\left\{  w\in P\ \mid\ x\leq w\leq y\right\}  $ (by the
definition of $\left[  x,y\right]  $).
\par
\item Let us first prove that $\left[  x,t\right]  \subseteq\left[
x,y\right]  $.
\par
Indeed, let $g\in\left[  x,t\right]  $. Then, $g\in\left[  x,t\right]
=\left\{  w\in P\ \mid\ x\leq w\leq t\right\}  $. In other words, $g$ is an
element $w$ of $P$ satisfying $x\leq w\leq t$. In other words, $g$ is an
element of $P$ and satisfies $x\leq g\leq t$. We have $g\leq t<y$, so that
$g\leq y$. Combining this with $x\leq g$, we obtain $x\leq g\leq y$. Thus, $g$
is an element of $P$ and satisfies $x\leq g\leq y$. In other words, $g$ is an
element $w$ of $P$ satisfying $x\leq w\leq y$. In other words, $g\in\left\{
w\in P\ \mid\ x\leq w\leq y\right\}  $. Since $\left[  x,y\right]  =\left\{
w\in P\ \mid\ x\leq w\leq y\right\}  $, this rewrites as $g\in\left[
x,y\right]  $.
\par
Now, forget that we fixed $g$. Thus, we have shown that $g\in\left[
x,y\right]  $ for each $g\in\left[  x,t\right]  $. In other words, $\left[
x,t\right]  \subseteq\left[  x,y\right]  $.
\par
\item Now, let us prove that $y\notin\left[  x,t\right]  $.
\par
Indeed, assume the contrary. Thus, $y\in\left[  x,t\right]  $. Hence,
$y\in\left[  x,t\right]  =\left\{  w\in P\ \mid\ x\leq w\leq t\right\}  $. In
other words, $y$ is an element $w$ of $P$ satisfying $x\leq w\leq t$. In other
words, $g$ is an element of $P$ and satisfies $x\leq y\leq t$. But $y\leq t$
contradicts $t<y$ (since $P$ is a poset). This contradiction proves that our
assumption was wrong. Hence, $y\notin\left[  x,t\right]  $ is proven.
\par
\item Next, let us prove that $y\in\left[  x,y\right]  $.
\par
Indeed, $y$ is an element of $P$ and satisfies $x\leq y\leq y$. In other
words, $y$ is an element $w$ of $P$ satisfying $x\leq w\leq y$. In other
words, $y\in\left\{  w\in P\ \mid\ x\leq w\leq y\right\}  $. Since $\left[
x,y\right]  =\left\{  w\in P\ \mid\ x\leq w\leq y\right\}  $, this rewrites as
$y\in\left[  x,y\right]  $.
\par
\item Let us now prove that $\left[  x,t\right]  \neq\left[  x,y\right]  $.
\par
Indeed, assume the contrary. Thus, $\left[  x,t\right]  =\left[  x,y\right]
$. Now, $y\notin\left[  x,t\right]  =\left[  x,y\right]  $ contradicts
$y\notin\left[  x,y\right]  $. This contradiction proves that our assumption
was wrong. Hence, $\left[  x,t\right]  \neq\left[  x,y\right]  $ is proven.
\par
\item Combining $\left[  x,t\right]  \subseteq\left[  x,y\right]  $ with
$\left[  x,t\right]  \neq\left[  x,y\right]  $, we conclude that $\left[
x,t\right]  $ is a proper subset of $\left[  x,y\right]  $. Thus, $\left\vert
\left[  x,t\right]  \right\vert <\left\vert \left[  x,y\right]  \right\vert $
(since both $\left[  x,t\right]  $ and $\left[  x,y\right]  $ are finite
sets). Hence, $\left\vert \left[  x,t\right]  \right\vert <\left\vert \left[
x,y\right]  \right\vert =N$. This proves
(\ref{pf.prop.moebius.double0.b.smaller}).
\end{itemize}
}. Therefore, for every $t\in P$ satisfying $x\leq t<y$, we have%
\begin{equation}
\sum_{\substack{z\in P;\\x\leq z\leq t}}\mu\left(  z,t\right)  =\left[
x=t\right]  \label{pf.prop.moebius.double0.b.smaller2}%
\end{equation}
\footnote{\textit{Proof of (\ref{pf.prop.moebius.double0.b.smaller2}):} Let
$t\in P$ be such that $x\leq t<y$. Then, $\left\vert \left[  x,t\right]
\right\vert <N$ (by (\ref{pf.prop.moebius.double0.b.smaller})). Hence, Claim 1
(applied to $t$ instead of $y$) shows that $\sum_{\substack{z\in P;\\x\leq
z\leq t}}\mu\left(  z,t\right)  =\left[  x=t\right]  $. This proves
(\ref{pf.prop.moebius.double0.b.smaller2}).}. Also, for every $u\in P$ and
$v\in P$, we have%
\begin{equation}
\sum_{\substack{t\in P;\\u\leq t\leq v}}\mu\left(  u,t\right)  =\left[
u=v\right]  \label{pf.prop.moebius.double0.b.smaller3}%
\end{equation}
\footnote{\textit{Proof of (\ref{pf.prop.moebius.double0.b.smaller3}):} Let
$u\in P$ and $v\in P$. Proposition \ref{prop.moebius.double0} \textbf{(a)}
(applied to $x=u$ and $y=v$) shows that $\sum_{\substack{z\in P;\\u\leq z\leq
v}}\mu\left(  u,t\right)  =\left[  u=v\right]  $. Now,%
\begin{align*}
\sum_{\substack{t\in P;\\u\leq t\leq v}}\mu\left(  u,t\right)   &
=\sum_{\substack{z\in P;\\u\leq z\leq v}}\mu\left(  u,t\right)
\ \ \ \ \ \ \ \ \ \ \left(  \text{here, we have substituted }z\text{ for
}t\text{ in the sum}\right) \\
&  =\left[  u=v\right]  .
\end{align*}
This proves (\ref{pf.prop.moebius.double0.b.smaller3}).}.

Furthermore,
\begin{equation}
\left\{  z\in P\ \mid\ x\leq z\leq y\text{ and }z=y\right\}  =\left\{
y\right\}  \label{pf.prop.moebius.double0.b.set1}%
\end{equation}
\footnote{\textit{Proof of (\ref{pf.prop.moebius.double0.b.set1}):} We know
that $y$ is an element of $P$ and satisfies $x\leq y\leq y$ and $y=y$. In
other words, $y$ is an element $z$ of $P$ satisfying $x\leq z\leq y$ and
$z=y$. In other words, $y\in\left\{  z\in P\ \mid\ x\leq z\leq y\text{ and
}z=y\right\}  $. Thus, $\left\{  y\right\}  \subseteq\left\{  z\in
P\ \mid\ x\leq z\leq y\text{ and }z=y\right\}  $.
\par
On the other hand, let $g\in\left\{  z\in P\ \mid\ x\leq z\leq y\text{ and
}z=y\right\}  $ be arbitrary. Thus, $g$ is an element $z$ of $P$ satisfying
$x\leq z\leq y$ and $z=y$. In other words, $g$ is an element of $P$ and
satisfies $x\leq g\leq y$ and $g=y$. Thus, $g=y\in\left\{  y\right\}  $. Now,
forget that we fixed $g$. We thus have shown that $g\in\left\{  y\right\}  $
for every $g\in\left\{  z\in P\ \mid\ x\leq z\leq y\text{ and }z=y\right\}  $.
In other words, $\left\{  z\in P\ \mid\ x\leq z\leq y\text{ and }z=y\right\}
\subseteq\left\{  y\right\}  $. Combining this with $\left\{  y\right\}
\subseteq\left\{  z\in P\ \mid\ x\leq z\leq y\text{ and }z=y\right\}  $, we
obtain%
\[
\left\{  z\in P\ \mid\ x\leq z\leq y\text{ and }z=y\right\}  =\left\{
y\right\}  .
\]
This proves (\ref{pf.prop.moebius.double0.b.set1}).}. Thus,%
\begin{align}
&  \left\{  w\in P\ \mid\ x\leq w\leq y\text{ and }w=y\right\} \nonumber\\
&  =\left\{  z\in P\ \mid\ x\leq z\leq y\text{ and }z=y\right\} \nonumber\\
&  \ \ \ \ \ \ \ \ \ \ \left(  \text{here, we have renamed the index }w\text{
as }z\right) \nonumber\\
&  =\left\{  y\right\}  . \label{pf.prop.moebius.double0.b.set2}%
\end{align}
Also,%
\begin{equation}
\left\{  z\in P\ \mid\ x\leq z\leq y\text{ and }z=x\text{ and }z\neq
y\right\}  =\left\{  x\right\}  \label{pf.prop.moebius.double0.b.set3}%
\end{equation}
\footnote{\textit{Proof of (\ref{pf.prop.moebius.double0.b.set3}):} We know
that $x$ is an element of $P$ and satisfies $x\leq x\leq y$ and $x=x$ and
$x\neq y$. In other words, $x$ is an element $z$ of $P$ satisfying $x\leq
z\leq y$ and $z=x$ and $z\neq y$. In other words, $x\in\left\{  z\in
P\ \mid\ x\leq z\leq y\text{ and }z=x\text{ and }z\neq y\right\}  $. Thus,
$\left\{  x\right\}  \subseteq\left\{  z\in P\ \mid\ x\leq z\leq y\text{ and
}z=x\text{ and }z\neq y\right\}  $.
\par
On the other hand, let $g\in\left\{  z\in P\ \mid\ x\leq z\leq y\text{ and
}z=x\text{ and }z\neq y\right\}  $ be arbitrary. Thus, $g$ is an element $z$
of $P$ satisfying $x\leq z\leq y$ and $z=x$ and $z\neq y$. In other words, $g$
is an element of $P$ and satisfies $x\leq g\leq y$ and $g=x$ and $g\neq y$.
Thus, $g=x\in\left\{  x\right\}  $. Now, forget that we fixed $g$. We thus
have shown that $g\in\left\{  x\right\}  $ for every $g\in\left\{  z\in
P\ \mid\ x\leq z\leq y\text{ and }z=x\text{ and }z\neq y\right\}  $. In other
words, $\left\{  z\in P\ \mid\ x\leq z\leq y\text{ and }z=x\text{ and }z\neq
y\right\}  \subseteq\left\{  x\right\}  $. Combining this with $\left\{
x\right\}  \subseteq\left\{  z\in P\ \mid\ x\leq z\leq y\text{ and }z=x\text{
and }z\neq y\right\}  $, we obtain%
\[
\left\{  z\in P\ \mid\ x\leq z\leq y\text{ and }z=x\text{ and }z\neq
y\right\}  =\left\{  x\right\}  .
\]
This proves (\ref{pf.prop.moebius.double0.b.set3}).}.

Now,%
\begin{align*}
&  \underbrace{\sum_{\substack{\left(  z,t\right)  \in P^{2};\\x\leq z\leq
y\text{ and }z\leq t\leq y}}}_{=\sum_{\substack{z\in P;\\x\leq z\leq y}%
}\sum_{\substack{t\in P;\\z\leq t\leq y}}}\mu\left(  z,t\right) \\
&  =\sum_{\substack{z\in P;\\x\leq z\leq y}}\underbrace{\sum_{\substack{t\in
P;\\z\leq t\leq y}}\mu\left(  z,t\right)  }_{\substack{=\left[  z=y\right]
\\\text{(by (\ref{pf.prop.moebius.double0.b.smaller3})}\\\text{(applied to
}u=z\text{ and }v=y\text{))}}}=\sum_{\substack{z\in P;\\x\leq z\leq y}}\left[
z=y\right] \\
&  =\sum_{\substack{z\in P;\\x\leq z\leq y\text{ and }z=y}}\underbrace{\left[
z=y\right]  }_{\substack{=1\\\text{(since }z=y\text{)}}}+\sum_{\substack{z\in
P;\\x\leq z\leq y\text{ and }z\neq y}}\underbrace{\left[  z=y\right]
}_{\substack{=0\\\text{(since }z=y\text{ is false}\\\text{(since }z\neq
y\text{))}}}\\
&  \ \ \ \ \ \ \ \ \ \ \left(  \text{since every }z\in P\text{ satisfies
either }z=y\text{ or }z\neq y\text{ (but not both)}\right) \\
&  =\sum_{\substack{z\in P;\\x\leq z\leq y\text{ and }z=y}}1+\underbrace{\sum
_{\substack{z\in P;\\x\leq z\leq y\text{ and }z\neq y}}0}_{=0}=\sum
_{\substack{z\in P;\\x\leq z\leq y\text{ and }z=y}}1\\
&  =\left\vert \left\{  z\in P\ \mid\ x\leq z\leq y\text{ and }z=y\right\}
\right\vert \cdot1\\
&  =\left\vert \underbrace{\left\{  z\in P\ \mid\ x\leq z\leq y\text{ and
}z=y\right\}  }_{\substack{=\left\{  y\right\}  \\\text{(by
(\ref{pf.prop.moebius.double0.b.set1}))}}}\right\vert =\left\vert \left\{
y\right\}  \right\vert =1.
\end{align*}
Hence,%
\begin{align*}
1  &  =\underbrace{\sum_{\substack{\left(  z,t\right)  \in P^{2};\\x\leq z\leq
y\text{ and }z\leq t\leq y}}}_{\substack{=\sum_{\substack{\left(  z,t\right)
\in P^{2};\\x\leq t\leq y\text{ and }x\leq z\leq t}}\\\text{(because for any
}\left(  z,t\right)  \in P^{2}\text{, the}\\\text{condition }\left(  x\leq
z\leq y\text{ and }z\leq t\leq y\right)  \\\text{is equivalent to }\left(
x\leq t\leq y\text{ and }x\leq z\leq t\right)  \\\text{(by
(\ref{pf.prop.moebius.double0.b.equiv})))}}}\mu\left(  z,t\right)
=\underbrace{\sum_{\substack{\left(  z,t\right)  \in P^{2};\\x\leq t\leq
y\text{ and }x\leq z\leq t}}}_{=\sum_{\substack{t\in P;\\x\leq t\leq y}%
}\sum_{\substack{z\in P;\\x\leq z\leq t}}}\mu\left(  z,t\right) \\
&  =\sum_{\substack{t\in P;\\x\leq t\leq y}}\sum_{\substack{z\in P;\\x\leq
z\leq t}}\mu\left(  z,t\right) \\
&  =\underbrace{\sum_{\substack{t\in P;\\x\leq t\leq y\text{ and }t=y}%
}}_{\substack{=\sum_{t\in\left\{  w\in P\ \mid\ x\leq w\leq y\text{ and
}w=y\right\}  }=\sum_{t\in\left\{  y\right\}  }\\\text{(since }\left\{  w\in
P\ \mid\ x\leq w\leq y\text{ and }w=y\right\}  =\left\{  y\right\}  \text{)}%
}}\sum_{\substack{z\in P;\\x\leq z\leq t}}\mu\left(  z,t\right) \\
&  \ \ \ \ \ \ \ \ \ \ +\sum_{\substack{t\in P;\\x\leq t\leq y\text{ and
}t\neq y}}\underbrace{\sum_{\substack{z\in P;\\x\leq z\leq t}}\mu\left(
z,t\right)  }_{\substack{=\left[  x=t\right]  \\\text{(by
(\ref{pf.prop.moebius.double0.b.smaller2})}\\\text{(since }t<y\text{ (because
}t\leq y\text{ and }t\neq y\text{)) and }x\leq t\text{)}}}\\
&  \ \ \ \ \ \ \ \ \ \ \left(  \text{since every }t\in P\text{ satisfies
either }t=y\text{ or }t\neq y\text{ (but not both)}\right) \\
&  =\underbrace{\sum_{t\in\left\{  y\right\}  }\sum_{\substack{z\in P;\\x\leq
z\leq t}}\mu\left(  z,t\right)  }_{=\sum_{\substack{z\in P;\\x\leq z\leq
y}}\mu\left(  z,y\right)  }+\sum_{\substack{t\in P;\\x\leq t\leq y\text{ and
}t\neq y}}\left[  x=t\right] \\
&  =\sum_{\substack{z\in P;\\x\leq z\leq y}}\mu\left(  z,y\right)
+\sum_{\substack{t\in P;\\x\leq t\leq y\text{ and }t\neq y}}\left[
x=t\right]  .
\end{align*}
Subtracting $\sum_{\substack{z\in P;\\x\leq z\leq y}}\mu\left(  z,y\right)  $
from both sides of this equality, we obtain%
\begin{align*}
1-\sum_{\substack{z\in P;\\x\leq z\leq y}}\mu\left(  z,y\right)   &
=\sum_{\substack{t\in P;\\x\leq t\leq y\text{ and }t\neq y}}\left[  x=t\right]
\\
&  =\underbrace{\sum_{\substack{t\in P;\\x\leq t\leq y\text{ and }t=x\text{
and }t\neq y}}}_{\substack{=\sum_{t\in\left\{  z\in P\ \mid\ x\leq z\leq
y\text{ and }z=x\text{ and }z\neq y\right\}  }=\sum_{t\in\left\{  x\right\}
}\\\text{(since }\left\{  z\in P\ \mid\ x\leq z\leq y\text{ and }z=x\text{ and
}z\neq y\right\}  =\left\{  x\right\}  \text{)}}}\underbrace{\left[
x=t\right]  }_{\substack{=1\\\text{(since }x=t\\\text{(since }t=x\text{))}}}\\
&  \ \ \ \ \ \ \ \ \ \ +\sum_{\substack{t\in P;\\x\leq t\leq y\text{ and
}t\neq x\text{ and }t\neq y}}\underbrace{\left[  x=t\right]  }%
_{\substack{=0\\\text{(since }x=t\text{ is false}\\\text{(since }x\neq t\text{
(since }t\neq x\text{)))}}}\\
&  \ \ \ \ \ \ \ \ \ \ \left(  \text{since every }t\in P\text{ satisfies
either }t=x\text{ or }t\neq x\text{ (but not both)}\right) \\
&  =\sum_{t\in\left\{  x\right\}  }1+\underbrace{\sum_{\substack{t\in
P;\\x\leq t\leq y\text{ and }t\neq x\text{ and }t\neq y}}0}_{=0}=\sum
_{t\in\left\{  x\right\}  }1=1.
\end{align*}
Solving this equality for $\sum_{\substack{z\in P;\\x\leq z\leq y}}\mu\left(
z,y\right)  $, we obtain%
\[
\sum_{\substack{z\in P;\\x\leq z\leq y}}\mu\left(  z,y\right)  =1-1=0=\left[
x=y\right]  .
\]
Thus, (\ref{pf.prop.moebius.double0.b.goal}) is proven.

Let us now forget that we fixed $x$ and $y$. We thus have proven that for any
$x\in P$ and $y\in P$ satisfying $\left\vert \left[  x,y\right]  \right\vert
=N$, we have%
\[
\sum_{\substack{z\in P;\\x\leq z\leq y}}\mu\left(  z,y\right)  =\left[
x=y\right]  .
\]
In other words, Proposition \ref{prop.moebius.double0} \textbf{(b)} holds
whenever $\left\vert \left[  x,y\right]  \right\vert =N$. This completes the
induction step. Thus, Proposition \ref{prop.moebius.double0} \textbf{(b)} is
proven by induction.

\textbf{(c)} For every $v\in P$, we have%
\begin{align*}
\sum_{\substack{y\in P;\\y\leq v}}\mu\left(  y,v\right)  \sum_{\substack{x\in
P;\\x\leq y}}\beta_{x}  &  =\sum_{\substack{z\in P;\\z\leq v}}\mu\left(
z,v\right)  \sum_{\substack{x\in P;\\x\leq z}}\beta_{x}%
\ \ \ \ \ \ \ \ \ \ \left(
\begin{array}
[c]{c}%
\text{here, we have renamed the summation}\\
\text{index }y\text{ as }z\text{ in the outer sum}%
\end{array}
\right) \\
&  =\underbrace{\sum_{\substack{z\in P;\\z\leq v}}\sum_{\substack{x\in
P;\\x\leq z}}}_{=\sum_{\substack{\left(  x,z\right)  \in P^{2};\\z\leq v\text{
and }x\leq z}}=\sum_{x\in P}\sum_{\substack{z\in P;\\x\leq z\text{ and }z\leq
v}}}\mu\left(  z,v\right)  \beta_{x}\\
&  =\sum_{x\in P}\underbrace{\sum_{\substack{z\in P;\\x\leq z\text{ and }z\leq
v}}}_{=\sum_{\substack{z\in P;\\x\leq z\leq v}}}\mu\left(  z,v\right)
\beta_{x}=\sum_{x\in P}\sum_{\substack{z\in P;\\x\leq z\leq v}}\mu\left(
z,v\right)  \beta_{x}\\
&  =\sum_{x\in P}\underbrace{\left(  \sum_{\substack{z\in P;\\x\leq z\leq
v}}\mu\left(  z,v\right)  \right)  }_{\substack{=\left[  x=v\right]
\\\text{(by Proposition \ref{prop.moebius.double0} \textbf{(b)}}%
\\\text{(applied to }y=v\text{))}}}\beta_{x}=\sum_{x\in P}\left[  x=v\right]
\beta_{x}\\
&  =\sum_{\substack{x\in P;\\x=v}}\underbrace{\left[  x=v\right]
}_{\substack{=1\\\text{(since }x=v\text{)}}}\beta_{x}+\sum_{\substack{x\in
P;\\x\neq v}}\underbrace{\left[  x=v\right]  }_{\substack{=0\\\text{(since
}x=v\text{ is false}\\\text{(since }x\neq v\text{))}}}\beta_{x}\\
&  \ \ \ \ \ \ \ \ \ \ \left(  \text{since every }x\in P\text{ satisfies
either }x=v\text{ or }x\neq v\text{ (but not both)}\right) \\
&  =\sum_{\substack{x\in P;\\x=v}}\beta_{x}+\underbrace{\sum_{\substack{x\in
P;\\x\neq v}}0\beta_{x}}_{=0}=\sum_{\substack{x\in P;\\x=v}}\beta_{x}%
=\beta_{v}\ \ \ \ \ \ \ \ \ \ \left(  \text{since }v\in P\right)  .
\end{align*}
Renaming $v$ as $z$ in this statement, we obtain the following: For every
$z\in P$, we have%
\[
\sum_{\substack{y\in P;\\y\leq z}}\mu\left(  y,z\right)  \sum_{\substack{x\in
P;\\x\leq y}}\beta_{x}=\beta_{z}.
\]
This proves Proposition \ref{prop.moebius.double0} \textbf{(c)}.
\end{proof}
\end{verlong}

\begin{proof}
[Proof of Theorem \ref{thm.matroid.charpol.varis}.]If $T$ is a subset of $E$,
then $\overline{T}$ is a flat of $M$ (by Proposition
\ref{prop.matroid.closure.props} \textbf{(a)}). In other words, if $T$ is a
subset of $E$, then $\overline{T}\in\operatorname*{Flats}M$. Renaming $T$ as
$B$ in this statement, we conclude that if $B$ is a subset of $E$, then
$\overline{B}\in\operatorname*{Flats}M$.

For every $F\in\operatorname*{Flats}M$, define an element $\beta_{F}%
\in\mathbf{k}$ by%
\[
\beta_{F}=\sum_{\substack{B\subseteq E;\\\overline{B}=F}}\left(  -1\right)
^{\left\vert B\right\vert }\left(  \prod_{\substack{K\in\mathfrak{K}%
;\\K\subseteq B}}a_{K}\right)  .
\]
Now, using Lemma \ref{lem.matroid.NBCm.moeb}, we can easily see that%
\begin{equation}
\sum_{\substack{G\in\operatorname*{Flats}M;\\G\subseteq F}}\beta_{G}=\left[
F=\varnothing\right]  \ \ \ \ \ \ \ \ \ \ \text{for every }F\in
\operatorname*{Flats}M \label{pf.thm.matroid.charpol.varis.b-via-a}%
\end{equation}
\footnote{\textit{Proof of (\ref{pf.thm.matroid.charpol.varis.b-via-a}):} Let
$F\in\operatorname*{Flats}M$. Thus, $F$ is a flat of $M$.
\par
If $B$ is a subset of $E$, then the statements $\left(  \overline{B}\subseteq
F\right)  $ and $\left(  B\subseteq F\right)  $ are equivalent. (This follows
from Proposition \ref{prop.matroid.closure.props} \textbf{(g)}, applied to
$T=B$ and $G=F$.)
\par
Now,%
\begin{align*}
&  \sum_{\substack{G\in\operatorname*{Flats}M;\\G\subseteq F}%
}\underbrace{\beta_{G}}_{\substack{=\sum_{\substack{B\subseteq E;\\\overline
{B}=G}}\left(  -1\right)  ^{\left\vert B\right\vert }\left(  \prod
_{\substack{K\in\mathfrak{K};\\K\subseteq B}}a_{K}\right)  \\\text{(by the
definition of }\beta_{G}\text{)}}}\\
&  =\underbrace{\sum_{\substack{G\in\operatorname*{Flats}M;\\G\subseteq
F}}\sum_{\substack{B\subseteq E;\\\overline{B}=G}}}_{\substack{=\sum
_{\substack{B\subseteq E;\\\overline{B}\subseteq F}}\\\text{(because if
}B\text{ is a subset of }E\text{,}\\\text{then }\overline{B}\in
\operatorname*{Flats}M\text{)}}}\left(  -1\right)  ^{\left\vert B\right\vert
}\left(  \prod_{\substack{K\in\mathfrak{K};\\K\subseteq B}}a_{K}\right) \\
&  =\underbrace{\sum_{\substack{B\subseteq E;\\\overline{B}\subseteq F}%
}}_{\substack{=\sum_{\substack{B\subseteq E;\\B\subseteq F}}\\\text{(because
if }B\text{ is a subset of }E\text{, then}\\\text{the statements }\left(
\overline{B}\subseteq F\right)  \text{ and }\left(  B\subseteq F\right)
\text{ are}\\\text{equivalent)}}}\left(  -1\right)  ^{\left\vert B\right\vert
}\left(  \prod_{\substack{K\in\mathfrak{K};\\K\subseteq B}}a_{K}\right)
=\underbrace{\sum_{\substack{B\subseteq E;\\B\subseteq F}}}_{=\sum_{B\subseteq
F}}\left(  -1\right)  ^{\left\vert B\right\vert }\left(  \prod_{\substack{K\in
\mathfrak{K};\\K\subseteq B}}a_{K}\right) \\
&  =\sum_{B\subseteq F}\left(  -1\right)  ^{\left\vert B\right\vert }\left(
\prod_{\substack{K\in\mathfrak{K};\\K\subseteq B}}a_{K}\right)  =\left[
F=\varnothing\right]  \ \ \ \ \ \ \ \ \ \ \left(  \text{by
(\ref{eq.lem.matroid.NBCm.moeb.1})}\right)  .
\end{align*}
This proves (\ref{pf.thm.matroid.charpol.varis.b-via-a}).}.

Let $\mu$ be the M\"{o}bius function of the lattice $\operatorname*{Flats}M$.
The element $\overline{\varnothing}$ is the global minimum of the poset
$\operatorname*{Flats}M$.\ \ \ \ \footnote{This was proven during our proof of
Proposition \ref{prop.matroid.flat-lat}.} In particular, $\overline
{\varnothing}\in\operatorname*{Flats}M$ and $\overline{\varnothing}\subseteq
F$. Hence, $\mu\left(  \overline{\varnothing},F\right)  $ is well-defined.

Now, fix $F\in\operatorname*{Flats}M$. Proposition \ref{prop.moebius.double0}
\textbf{(c)} (applied to $P=\operatorname*{Flats}M$ and $z=F$) shows that%
\begin{align}
\beta_{F}  &  =\sum_{\substack{y\in\operatorname*{Flats}M;\\y\subseteq F}%
}\mu\left(  y,F\right)  \sum_{\substack{x\in\operatorname*{Flats}%
M;\\x\subseteq y}}\beta_{x}\nonumber\\
&  \ \ \ \ \ \ \ \ \ \ \left(  \text{since the relation }\leq\text{ of the
poset }\operatorname*{Flats}M\text{ is }\subseteq\right) \nonumber\\
&  =\sum_{\substack{H\in\operatorname*{Flats}M;\\H\subseteq F}}\mu\left(
H,F\right)  \underbrace{\sum_{\substack{G\in\operatorname*{Flats}%
M;\\G\subseteq H}}\beta_{G}}_{\substack{=\left[  H=\varnothing\right]
\\\text{(by (\ref{pf.thm.matroid.charpol.varis.b-via-a}), applied to}\\H\text{
instead of }F\text{)}}}\nonumber\\
&  \ \ \ \ \ \ \ \ \ \ \left(  \text{here, we renamed the summation indices
}y\text{ and }x\text{ as }H\text{ and }G\right) \nonumber\\
&  =\sum_{\substack{H\in\operatorname*{Flats}M;\\H\subseteq F}}\mu\left(
H,F\right)  \left[  H=\varnothing\right] \nonumber\\
&  =\sum_{\substack{H\in\operatorname*{Flats}M;\\H\subseteq F;\\H=\varnothing
}}\mu\left(  H,F\right)  \underbrace{\left[  H=\varnothing\right]
}_{\substack{=1\\\text{(since }H=\varnothing\text{)}}}+\sum_{\substack{H\in
\operatorname*{Flats}M;\\H\subseteq F;\\H\neq\varnothing}}\mu\left(
H,F\right)  \underbrace{\left[  H=\varnothing\right]  }%
_{\substack{=0\\\text{(since }H\neq\varnothing\text{)}}}\nonumber\\
&  =\underbrace{\sum_{\substack{H\in\operatorname*{Flats}M;\\H\subseteq
F;\\H=\varnothing}}}_{\substack{=\sum_{\substack{H\in\operatorname*{Flats}%
M;\\H=\varnothing}}\\\text{(since the condition }H\subseteq F\\\text{is
automatically implied by}\\\text{the condition }H=\varnothing\text{)}}%
}\mu\left(  H,F\right) \nonumber\\
&  =\sum_{\substack{H\in\operatorname*{Flats}M;\\H=\varnothing}}\mu\left(
H,F\right)  . \label{pf.thm.matroid.charpol.varis.3}%
\end{align}
Now, we shall prove that
\begin{equation}
\beta_{F}=\left[  \overline{\varnothing}=\varnothing\right]  \mu\left(
\overline{\varnothing},F\right)  . \label{pf.thm.matroid.charpol.varis.5}%
\end{equation}


\textit{Proof of (\ref{pf.thm.matroid.charpol.varis.5}):} We are in one of the
following two cases:

\textit{Case 1:} We have $\overline{\varnothing}=\varnothing$.

\textit{Case 2:} We have $\overline{\varnothing}\neq\varnothing$.

Let us consider Case 1 first. In this case, we have $\overline{\varnothing
}=\varnothing$. Hence, $\varnothing=\overline{\varnothing}\in
\operatorname*{Flats}M$. Thus, the sum $\sum_{\substack{H\in
\operatorname*{Flats}M;\\H=\varnothing}}\mu\left(  H,F\right)  $ has exactly
one addend: namely, the addend for $H=\varnothing$. Thus, $\sum
_{\substack{H\in\operatorname*{Flats}M;\\H=\varnothing}}\mu\left(  H,F\right)
=\mu\left(  \underbrace{\varnothing}_{=\overline{\varnothing}},F\right)
=\mu\left(  \overline{\varnothing},F\right)  $. Thus,
(\ref{pf.thm.matroid.charpol.varis.3}) becomes $\beta_{F}=\sum_{\substack{H\in
\operatorname*{Flats}M;\\H=\varnothing}}\mu\left(  H,F\right)  =\mu\left(
\overline{\varnothing},F\right)  $. Comparing this with $\underbrace{\left[
\overline{\varnothing}=\varnothing\right]  }_{\substack{=1\\\text{(since
}\overline{\varnothing}=\varnothing\text{)}}}\mu\left(  \overline{\varnothing
},F\right)  =\mu\left(  \overline{\varnothing},F\right)  $, we obtain
$\beta_{F}=\left[  \overline{\varnothing}=\varnothing\right]  \mu\left(
\overline{\varnothing},F\right)  $. Thus,
(\ref{pf.thm.matroid.charpol.varis.5}) is proven in Case 1.

Let us now consider Case 2. In this case, we have $\overline{\varnothing}%
\neq\varnothing$. Thus, there exists no $H\in\operatorname*{Flats}M$ such that
$H=\varnothing$\ \ \ \ \footnote{\textit{Proof.} Assume the contrary. Thus,
there exists some $H\in\operatorname*{Flats}M$ such that $H=\varnothing$. In
other words, $\varnothing\in\operatorname*{Flats}M$. Hence, $\varnothing$ is a
flat of $M$. Proposition \ref{prop.matroid.closure.props} \textbf{(b)}
(applied to $G=\varnothing$) thus shows that $\overline{\varnothing
}=\varnothing$. This contradicts $\overline{\varnothing}\neq\varnothing$. This
contradiction proves that our assumption was wrong, qed.}. Hence, the sum
$\sum_{\substack{H\in\operatorname*{Flats}M;\\H=\varnothing}}\mu\left(
H,F\right)  $ is empty. Thus, $\sum_{\substack{H\in\operatorname*{Flats}%
M;\\H=\varnothing}}\mu\left(  H,F\right)  =\left(  \text{empty sum}\right)
=0$, so that (\ref{pf.thm.matroid.charpol.varis.3}) becomes $\beta_{F}%
=\sum_{\substack{H\in\operatorname*{Flats}M;\\H=\varnothing}}\mu\left(
H,F\right)  =0$. Comparing this with $\underbrace{\left[  \overline
{\varnothing}=\varnothing\right]  }_{\substack{=0\\\text{(since }%
\overline{\varnothing}\neq\varnothing\text{)}}}\mu\left(  \overline
{\varnothing},F\right)  =0$, we obtain $\beta_{F}=\left[  \overline
{\varnothing}=\varnothing\right]  \mu\left(  \overline{\varnothing},F\right)
$. Thus, (\ref{pf.thm.matroid.charpol.varis.5}) is proven in Case 2.

Now, we have proven (\ref{pf.thm.matroid.charpol.varis.5}) in both possible
Cases 1 and 2. Thus, (\ref{pf.thm.matroid.charpol.varis.5}) always holds.

Now, let us forget that we fixed $F$. We thus have proven
(\ref{pf.thm.matroid.charpol.varis.5}) for each $F\in\operatorname*{Flats}M$.

Now,%
\begin{align}
&  \sum_{F\subseteq E}\left(  -1\right)  ^{\left\vert F\right\vert }\left(
\prod_{\substack{K\in\mathfrak{K};\\K\subseteq F}}a_{K}\right)  x^{m-r_{M}%
\left(  F\right)  }\nonumber\\
&  =\underbrace{\sum_{B\subseteq E}}_{\substack{=\sum_{F\in
\operatorname*{Flats}M}\sum_{\substack{B\subseteq E;\\\overline{B}%
=F}}\\\text{(because if }B\text{ is a subset of }E\text{,}\\\text{then
}\overline{B}\in\operatorname*{Flats}M\text{)}}}\left(  -1\right)
^{\left\vert B\right\vert }\left(  \prod_{\substack{K\in\mathfrak{K}%
;\\K\subseteq B}}a_{K}\right)  \underbrace{x^{m-r_{M}\left(  B\right)  }%
}_{\substack{=x^{m-r_{M}\left(  \overline{B}\right)  }\\\text{(since
Proposition \ref{prop.matroid.closure.props} \textbf{(f)} (applied to
}T=B\text{)}\\\text{shows that }r_{M}\left(  B\right)  =r_{M}\left(
\overline{B}\right)  \text{)}}}\nonumber\\
&  \ \ \ \ \ \ \ \ \ \ \left(  \text{here, we have renamed the summation index
}F\text{ as }B\right) \nonumber\\
&  =\sum_{F\in\operatorname*{Flats}M}\sum_{\substack{B\subseteq E;\\\overline
{B}=F}}\left(  -1\right)  ^{\left\vert B\right\vert }\left(  \prod
_{\substack{K\in\mathfrak{K};\\K\subseteq B}}a_{K}\right)
\underbrace{x^{m-r_{M}\left(  \overline{B}\right)  }}%
_{\substack{_{\substack{=x^{m-r_{M}\left(  F\right)  }}}\\\text{(since
}\overline{B}=F\text{)}}}\nonumber\\
&  =\sum_{F\in\operatorname*{Flats}M}\underbrace{\sum_{\substack{B\subseteq
E;\\\overline{B}=F}}\left(  -1\right)  ^{\left\vert B\right\vert }\left(
\prod_{\substack{K\in\mathfrak{K};\\K\subseteq B}}a_{K}\right)  }%
_{\substack{=\beta_{F}=\left[  \overline{\varnothing}=\varnothing\right]
\mu\left(  \overline{\varnothing},F\right)  \\\text{(by
(\ref{pf.thm.matroid.charpol.varis.5}))}}}x^{m-r_{M}\left(  F\right)
}\nonumber\\
&  =\sum_{F\in\operatorname*{Flats}M}\left[  \overline{\varnothing
}=\varnothing\right]  \mu\left(  \overline{\varnothing},F\right)
x^{m-r_{M}\left(  F\right)  }\nonumber\\
&  =\left[  \overline{\varnothing}=\varnothing\right]  \sum_{F\in
\operatorname*{Flats}M}\mu\left(  \overline{\varnothing},F\right)
x^{m-r_{M}\left(  F\right)  }. \label{pf.thm.matroid.charpol.varis.9}%
\end{align}


But the definition of $\chi_{M}$ yields $\chi_{M}=\sum_{F\in
\operatorname*{Flats}M}\mu\left(  \overline{\varnothing},F\right)
x^{m-r_{M}\left(  F\right)  }$. The definition of $\widetilde{\chi}_{M}$
yields%
\begin{align*}
\widetilde{\chi}_{M}  &  =\left[  \overline{\varnothing}=\varnothing\right]
\underbrace{\chi_{M}}_{=\sum_{F\in\operatorname*{Flats}M}\mu\left(
\overline{\varnothing},F\right)  x^{m-r_{M}\left(  F\right)  }}=\left[
\overline{\varnothing}=\varnothing\right]  \sum_{F\in\operatorname*{Flats}%
M}\mu\left(  \overline{\varnothing},F\right)  x^{m-r_{M}\left(  F\right)  }\\
&  =\sum_{F\subseteq E}\left(  -1\right)  ^{\left\vert F\right\vert }\left(
\prod_{\substack{K\in\mathfrak{K};\\K\subseteq F}}a_{K}\right)  x^{m-r_{M}%
\left(  F\right)  }\ \ \ \ \ \ \ \ \ \ \left(  \text{by
(\ref{pf.thm.matroid.charpol.varis.9})}\right)  .
\end{align*}
This proves Theorem \ref{thm.matroid.charpol.varis}.
\end{proof}

\begin{proof}
[Proof of Corollary \ref{cor.matroid.charpol.K-free}.]Corollary
\ref{cor.matroid.charpol.K-free} can be derived from Theorem
\ref{thm.matroid.charpol.varis} in the same way as Corollary
\ref{cor.chromsym.K-free} was derived from Theorem \ref{thm.chromsym.varis}.
\end{proof}

\begin{proof}
[Proof of Theorem \ref{thm.matroid.charpol.empty}.]Theorem
\ref{thm.matroid.charpol.empty} can be derived from Theorem
\ref{thm.matroid.charpol.varis} in the same way as Theorem
\ref{thm.chromsym.empty} was derived from Theorem \ref{thm.chromsym.varis}.
\end{proof}

\begin{proof}
[Proof of Corollary \ref{cor.matroid.charpol.NBC}.]Corollary
\ref{cor.matroid.charpol.NBC} follows from Corollary
\ref{cor.matroid.charpol.K-free} when $\mathfrak{K}$ is set to be the set of
\textbf{all} broken circuits of $M$.
\end{proof}

\begin{proof}
[Proof of Corollary \ref{cor.matroid.charpol.NBCfor}.]If $F$ is a subset of
$E$ such that $F$ contains no broken circuit of $M$ as a subset, then%
\begin{equation}
r_{M}\left(  F\right)  =\left\vert F\right\vert
\label{pf.cor.matroid.charpol.NBCfor.1}%
\end{equation}
\footnote{\textit{Proof of (\ref{pf.cor.matroid.charpol.NBCfor.1}):} Let $F$
be a subset of $E$ such that $F$ contains no broken circuit of $M$ as a
subset.
\par
We shall show that $F\in\mathcal{I}$. Indeed, assume the contrary. Thus,
$F\notin\mathcal{I}$, so that $F\in\mathcal{P}\left(  E\right)  \setminus
\mathcal{I}$. Hence, there exists a circuit $C$ of $M$ such that $C\subseteq
F$ (according to Lemma \ref{lem.matroid.Cin}, applied to $Q=F$). Consider this
$C$. The set $C$ is a circuit, and thus nonempty (because the empty set is in
$\mathcal{I}$). Let $e$ be the unique element of $C$ having maximum label.
(This is clearly well-defined, since the labeling function $\ell$ is
injective). Then, $C\setminus\left\{  e\right\}  $ is a broken circuit of $M$
(by the definition of a broken circuit). Thus, $F$ contains a broken circuit
of $M$ as a subset (since $C\setminus\left\{  e\right\}  \subseteq C\subseteq
F$). This contradicts the fact that $F$ contains no broken circuit of $M$ as a
subset. This contradiction shows that our assumption was wrong. Hence,
$F\in\mathcal{I}$ is proven.
\par
Thus, Lemma \ref{lem.matroid.rI} (applied to $T=F$) shows that $r_{M}\left(
F\right)  =\left\vert F\right\vert $, qed.}. Now, Corollary
\ref{cor.matroid.charpol.NBC} yields%
\[
\widetilde{\chi}_{M}=\sum_{\substack{F\subseteq E;\\F\text{ contains no
broken}\\\text{circuit of }M\text{ as a subset}}}\left(  -1\right)
^{\left\vert F\right\vert }\underbrace{x^{m-r_{M}\left(  F\right)  }%
}_{\substack{=x^{m-\left\vert F\right\vert }\\\text{(by
(\ref{pf.cor.matroid.charpol.NBCfor.1}))}}}=\sum_{\substack{F\subseteq
E;\\F\text{ contains no broken}\\\text{circuit of }M\text{ as a subset}%
}}\left(  -1\right)  ^{\left\vert F\right\vert }x^{m-\left\vert F\right\vert
}.
\]
This proves Corollary \ref{cor.matroid.charpol.NBCfor}.
\end{proof}

\begin{thebibliography}{999999999}                                                                                        %


\bibitem[BlaSag86]{BlaSag86}%
\href{http://users.math.msu.edu/users/sagan/Papers/Old/bpt.pdf}{Andreas Blass,
Bruce Eli Sagan, \textit{Bijective Proofs of Two Broken Circuit Theorems},
Journal of Graph Theory, Vol. 10 (1986), pp. 15--21.} (See also
\href{https://www.math.ucdavis.edu/~deloera/MISC/BIBLIOTECA/trunk/Sagan2.pdf}{another
scan}.)

\bibitem[Bollob79]{Bollobas}B\'{e}la Bollob\'{a}s, \textit{Graph Theory: An
Introductory Course}, Springer, Softcover reprint of the hardcover 1st edition 1979.

\bibitem[Bona11]{Bona11}Mikl\'{o}s B\'{o}na, \textit{A Walk Through
Combinatorics: An Introduction to Enumeration and Graph Theory}, 3rd edition,
World Scientific 2011.

\bibitem[Grinbe16]{Grinbe16}Darij Grinberg, \textit{Double posets and the
antipode of }$\operatorname*{QSym}$, version 1.4.\newline\url{http://web.mit.edu/~darij/www/algebra/dp-abstr.pdf}

\bibitem[GriRei14]{Reiner}Darij Grinberg, Victor Reiner, \textit{Hopf algebras
in Combinatorics}, August 25, 2015. arXiv preprint
\href{http://www.arxiv.org/abs/1409.8356v3}{\texttt{arXiv:1409.8356v3}}.
\newline A version which is more often updated can be found at
\url{http://www.math.umn.edu/~reiner/Classes/HopfComb.pdf} .

\bibitem[Martin15]{Martin15}%
\href{https://www.math.ku.edu/~jmartin/LectureNotes.pdf}{Jeremy L. Martin,
\textit{Lecture Notes on Algebraic Combinatorics}, February 18, 2016.}

\bibitem[Stanle95]{Stanley-chsym}%
\href{http://www.sciencedirect.com/science/article/pii/S0001870885710201}{Richard
P. Stanley, \textit{A Symmetric Function Generalization of the Chromatic
Polynomial of a Graph}, Advances in Mathematics, Volume 111, Issue 1, March
1995, pp. 166--194.}

\bibitem[Stanle99]{Stanley-EC2}Richard P. Stanley, \textit{Enumerative
Combinatorics, volume 2}, Cambridge University Press 1999.

\bibitem[Stanley06]{Stanley}%
\href{http://www-math.mit.edu/~rstan/arrangements/arr.html}{Richard Stanley,
\textit{An Introduction to Hyperplane Arrangements}, lecture series at the
Park City Mathematics Institute, July 12-19, 2004}. See also
\href{http://www-math.mit.edu/~rstan/arrangements/errata.pdf}{the errata, last
updated 2016}.

\bibitem[Schrij13]{Schrij13}%
\href{http://homepages.cwi.nl/~lex/files/dict.pdf}{Alexander Schrijver,
\textit{A Course in Combinatorial Optimization}, February 3, 2013.}
\end{thebibliography}


\end{document}